\documentclass[11pt]{article}
\usepackage[T1]{fontenc}
\usepackage[utf8]{inputenc}
\usepackage{lmodern}
\usepackage[margin=1in]{geometry}
\usepackage{amsmath,amssymb,amsthm,mathtools,mathrsfs}
\usepackage{microtype}
\usepackage{enumitem}
\usepackage{array,booktabs,longtable}
\usepackage{graphicx}
\usepackage{tikz-cd}
\usepackage{xcolor}
\usepackage{xurl}
\usepackage[hidelinks,unicode]{hyperref}
\usepackage[capitalise,nameinlink,noabbrev]{cleveref}
\AddToHook{env/thebibliography/begin}{\raggedright}
\let\hthreeOriginalThebibliography\thebibliography
\renewcommand{\thebibliography}[1]{%
  \hthreeOriginalThebibliography{#1}%
  \addcontentsline{toc}{section}{\refname}%
}
\hypersetup{pdftitle={The height-three chromatic overlap: an explicit filtration and its attachments},pdfauthor={OpenAI}}
\numberwithin{equation}{section}
\newtheorem{theorem}{Theorem}[section]
\newtheorem{proposition}[theorem]{Proposition}
\newtheorem{lemma}[theorem]{Lemma}
\newtheorem{corollary}[theorem]{Corollary}
\theoremstyle{definition}
\newtheorem{definition}[theorem]{Definition}
\theoremstyle{remark}
\newtheorem{remark}[theorem]{Remark}
\newcommand{\Fp}{\mathbb F_p}
\newcommand{\Zp}{\mathbb Z_p}
\newcommand{\Qp}{\mathbb Q_p}
\newcommand{\Z}{\mathbb Z}
\newcommand{\Q}{\mathbb Q}
\newcommand{\F}{\mathbb F}
\newcommand{\cO}{\mathcal O}
\newcommand{\cG}{\mathcal G}
\newcommand{\cD}{\mathcal D}
\newcommand{\cM}{\mathcal M}
\DeclareMathOperator{\GL}{GL}
\DeclareMathOperator{\Gal}{Gal}
\DeclareMathOperator{\Ext}{Ext}
\DeclareMathOperator{\Hom}{Hom}
\DeclareMathOperator{\Lie}{Lie}
\DeclareMathOperator{\Aut}{Aut}
\DeclareMathOperator{\Spec}{Spec}
\DeclareMathOperator{\Spf}{Spf}
\DeclareMathOperator{\Spa}{Spa}
\DeclareMathOperator{\Spd}{Spd}
\DeclareMathOperator*{\colim}{colim}
\DeclareMathOperator{\coker}{coker}
\DeclareMathOperator{\cofib}{cofib}
\DeclareMathOperator{\fib}{fib}
\DeclareMathOperator{\RHom}{RHom}
\newcommand{\RGamma}{R\Gamma}
\newcommand{\LK}[1]{L_{K(#1)}}
\newcommand{\LN}[1]{L_{#1}}
\setlist[enumerate]{itemsep=3pt,topsep=5pt}
\setlist[itemize]{itemsep=3pt,topsep=5pt}
\setlength{\emergencystretch}{2em}
\allowdisplaybreaks[2]
\ifdefined\pdfinfoomitdate\pdfinfoomitdate=1\fi
\ifdefined\pdftrailerid\pdftrailerid{}\fi
\ifdefined\pdfsuppressptexinfo\pdfsuppressptexinfo=15\fi
\DeclareMathOperator{\hthreeMap}{Map}
\DeclareMathOperator{\hthreeNrd}{Nrd}
\DeclareMathOperator{\hthreeTr}{Tr}
\DeclareMathOperator{\hthreeRes}{Res}
\DeclareMathOperator{\hthreerank}{rank}
\DeclareMathOperator{\hthreeim}{im}
\newcommand{\hthreeid}{\mathrm{id}}
\newcommand{\hthreects}{\mathrm{cts}}
\newcommand{\hthreehol}{\mathrm{hol}}
\newcommand{\hthreepk}{\mathrm{pk}}
\DeclareMathOperator{\hthreeMod}{Mod}
\DeclareMathOperator{\hthreeEnd}{End}
\DeclareMathOperator{\hthreeSL}{SL}

\newcommand{\Sp}{\mathbb S_p^\wedge}
\DeclareMathOperator{\Map}{Map}
\DeclareMathOperator{\Mod}{Mod}
\newcommand{\id}{\mathrm{id}}

\title{The height-three chromatic overlap:\\
an explicit filtration and its attachments}
\author{OpenAI}
\date{September 27, 2026}
\begin{document}
\maketitle
\begin{abstract}
For every prime \(p\geq5\), we construct an explicit eight-stage
filtration of \(L_2L_{K(3)}\mathbb S_p^\wedge\) by the local-sphere
layers in the height-three chromatic-splitting pattern. The map from
its first stage to the overlap is the canonical unit, and two signed
fracture formulas identify all attachments for the chosen local maps
and compatibility homotopy. A companion canonical-map theorem further
implies that the first height-one attachment is nonzero.
\end{abstract}
\tableofcontents
\section{Introduction}\label{sec:introduction}

Fix a prime \(p\).  Write \(L_i=L_{E(i)}\) for localization at
Johnson--Wilson theory and \(L_{K(i)}\) for localization at Morava
\(K\)-theory.  We use the derived \(p\)-completion
\(S=\mathbb S_p^\wedge\) of the sphere and study the overlap
\[
 T=L_{K(3)}S,\qquad X=L_2T.
\]
The object \(X\) is the part of the height-three local sphere visible
below height three.  Our question is how to build it from lower local
spheres while retaining the maps by which the pieces are joined.

Chromatic fracture makes the gluing problem precise.  An
\(E(2)\)-local spectrum is a homotopy pullback of its \(K(2)\)-local
and \(E(1)\)-local parts over their common localization.  The gluing
map is part of the data; the fracture square does not determine it
from the two corners alone
\cite[Section~2.4, diagram~(2.19)]{BB}.
Hopkins' chromatic splitting conjecture, as recorded by Hovey
\cite[Conjecture~4.2]{Hovey}, proposes classes, factorizations, and
summands that would split the relevant canonical cofiber sequence.
The revised strong formulation of Barthel--Beaudry
\cite[Conjecture~6.3 and Remark~6.4]{BB} has the same lower-sphere
pattern at the height and primes considered here.  We study an ordered
filtration with those cofibers and a canonical first unit, allowing
nonzero connecting maps.

The low-height results already distinguish these formulations.
Hovey's July 1993 account records the height-one case and the
height-two case at primes greater than three, attributing the latter
to Hopkins using Shimomura--Yabe \cite[pp.~17--19]{Hovey}.
Goerss--Henn--Mahowald later proved
the height-two splitting at the prime three
\cite[Theorem~1.2]{GoerssHennMahowald2014}.  Their proof identifies
the actual rational comparison map before applying chromatic fracture
\cite[proof of Theorem~5.11, p.~1288]{GoerssHennMahowald2014}.
At height two and the prime two, Beaudry disproved the original strong
form \cite[Theorem~1.4]{Beaudry2017}; Beaudry--Goerss--Henn established
a corrected decomposition with Moore-spectrum terms that retains the
weak unit splitting \cite[Theorem~1.1.6]{BGH2022}.

Two established calculations explain the pieces at height three.
At odd primes the height-two splitting gives two \(E(1)\)-local
pieces and two rational pieces
\cite[Theorem~6.7]{BB}.  At every prime and positive height, the theorem of
Barthel--Schlank--Stapleton--Weinstein computes the rational homotopy
of the \(K(n)\)-local sphere as an exterior algebra over \(\Qp\)
on generators of degrees \(1-2i\), \(1\leq i\leq n\)
\cite[Theorem~A]{BSSW}.  At height three the resulting degrees are
\[
 0,\ -1,\ -3,\ -4,\ -5,\ -6,\ -8,\ -9.
\]
This algebra calculation determines neither integral representatives
nor their images under lower-height localization.  Those map
identifications are what allow the layers to be assembled.

\subsection{The explicit filtration}

All spectra and maps in the theorem carry their natural \(S\)-module
structures.  A filtration means a sequence of homotopy cofiber
sequences: its \(i\)-th cofiber \(C_i\) has a connecting map
\(\partial_i:C_i\to\Sigma F_{i-1}\).

\begin{theorem}[Explicit height-three filtration]\label{thm:main}
For every prime \(p\geq5\), there are \(E(2)\)-local \(S\)-modules
and maps
\[
 0=F_0\longrightarrow F_1\longrightarrow\cdots\longrightarrow F_8,
 \qquad e:F_8\xrightarrow{\simeq}X,
\]
whose successive cofibers \(C_i=\cofib(F_{i-1}\to F_i)\), in order,
are
\[
\begin{array}{c|l@{\qquad\qquad}c|l}
i&C_i&i&C_i\\ \hline
1&L_2S&5&\Sigma^{-5}H\Qp\\
2&\Sigma^{-1}L_2S&6&\Sigma^{-6}H\Qp\\
3&\Sigma^{-3}L_1S&7&\Sigma^{-8}H\Qp\\
4&\Sigma^{-4}L_1S&8&\Sigma^{-9}H\Qp.
\end{array}
\]
Under \(F_1\simeq L_2S\), the composite
\(F_1\to F_8\xrightarrow e X\) is the canonical map
\(L_2(S\to L_{K(3)}S)\).
\end{theorem}

The attachment identification is a separate formal part of the
result.  Once the stages and their fracture data have been defined,
\Cref{thm:attachment-identification} collects every connecting map
and both signed roofs.  The application in \Cref{sec:application}
establishes \Cref{thm:main} and the height-three instance of
\Cref{thm:attachment-identification} from one choice of local maps,
compatibility homotopy, and coherent inverse data.  The first unit is
canonical; the attachment identification is relative to those choices.

Existence of a canonical-unit filtration with these cofibers is also
the height-three case of
\cite[Theorem~1.1]{CompanionFilteredChromaticSplitting}, whose range
\(p>n+1\) becomes \(p\geq5\) here.  That theorem supplies one family
of product bases at all lower positive heights.  The present proof
constructs the height-three maps independently of that theorem,
identifies the rational degree-three comparison, and records every
attachment relative to one compatible set of choices.  Neither
filtration statement asserts the strong wedge or splits its first unit.
The construction of \Cref{thm:main} is independent of the
rational-obstruction theorem used in \Cref{sec:attachment}.  That
separate theorem supplies additional information about the
extensions: for this explicit filtration, the first height-one
connecting map is nonzero.

\subsection{Three blocks and two gluing problems}

Put \(Q=L_0S\simeq H\Qp\), and group the layers as
\[
\begin{aligned}
 W&=L_2S\vee\Sigma^{-1}L_2S,\\
 U_{\mathrm{mid}}&=\Sigma^{-3}L_1S\vee\Sigma^{-4}L_1S,\\
 V&=\Sigma^{-5}Q\vee\Sigma^{-6}Q
       \vee\Sigma^{-8}Q\vee\Sigma^{-9}Q.
\end{aligned}
\]
The first localized basis is an actual map
\[
 w=(1,\zeta):W\longrightarrow X,
\]
where \(\zeta:\Sigma^{-1}S\to T\) is the integral determinant
class.  It becomes an equivalence after \(K(2)\)-localization.
Hence its cofiber \(Y\) is \(E(1)\)-local.  To insert the two
summands of \(U_{\mathrm{mid}}\), one needs a map to \(Y\), not merely the
knowledge that \(Y\) has the expected local homotopy groups.

There are two descriptions of that map.  At height one, an integral
orientation class \(\delta\in\pi_{-3}L_{K(1)}T\) participates in
the actual basis
\[
 (1,\zeta,\delta,\zeta\delta):
 R_1\vee\Sigma^{-1}R_1\vee\Sigma^{-3}R_1\vee\Sigma^{-4}R_1
 \xrightarrow{\simeq}L_{K(1)}T,\qquad R_1=L_{K(1)}S.
\]
The first two components are localizations of the same maps used
in \(w\); quotienting them gives the \(K(1)\)-local part of a map
from \(U_{\mathrm{mid}}\).  Rationally, the degree-three primitive
\(\alpha_3\in\pi_{-3}L_0T\) maps to \(c\delta\) for a nonzero
\(c\in\Qp\), and its product with \(\zeta\) maps to
\(c\zeta\delta\).  Rescaling only the rational primitive by
\(c^{-1}\) makes the two descriptions agree on the rationalized
\(K(1)\)-local overlap.  The height-one fracture square then
constructs
\[
 h:U_{\mathrm{mid}}\longrightarrow Y.
\]
The specified compatibility homotopy, rather than a dimension count,
is the input to this gluing step.

Pulling \(X\to Y\) back along the first summand and then all of \(U_{\mathrm{mid}}\)
gives \(F_3\) and \(F_4\).  The cofiber of \(F_4\to X\) is rational.
The four remaining exterior products give an actual equivalence
from \(V\) to that cofiber, and their ordered partial sums give
\(F_5,\ldots,F_8\) by pullback.  The terminal projection is an
equivalence because the full map from \(V\) is.  The proof in
\Cref{sec:fracture} follows the two fracture squares far enough to
identify their connecting maps, including both rotation signs.

\subsection{Why the localized bases require geometry}

The two coefficient comparisons use the loci of the height-three
deformation where the connected part has height two and height one.
Their \'etale parts have heights one and two, respectively; these are
the coheight-one and ordinary strata.  Finite marking towers on these
loci are related to extensions of vector bundles on the
Fargues--Fontaine curve.  The bundle and local-system framework
comes from Fargues--Fontaine and Scholze--Weinstein
\cite{FarguesFontaine,ScholzeWeinstein}; relative full faithfulness
of derived sections is supplied by Ansch\"utz--Le Bras
\cite[Corollary~3.11]{AnschuetzLeBrasFourier}.
Retaining the integral Tate frames and the determinant lattice makes
the group actions and the determinant product identity part of the
comparison.  These data are later used to identify the coefficient
classes of the unit and determinant maps in \(w=(1,\zeta)\).

The Kummer and translation calculations in the coheight-one work of
Barthel--Mann--Ray--Schlank--Senger--Weinstein--Zhou
\cite[Sections~3.3--3.6]{BMR} provide methodological context for
the cohomology of the completed strata.  The geometric answers must
then be carried to the algebraic coefficients in ordinary Morava
cooperations.
A cochain in their marking union must occur at one finite stage on
its compact domain; a cochain with Laurent coefficients must have
one pole bound.  Completion does not preserve those requirements by
definition, so the comparisons must establish the required finite
bounds before returning to the algebraic coefficients.

On the locus with connected height two, the completed calculation
controls cohomology at each finite stage of the connected marking tower.
The comparison in \Cref{h3:sec:coheight-one} proves finite cohomology
for a cofinal family of finite products of complete local fields, then
removes completion while preserving the group actions and the constants
map.  The descent calculation in \Cref{h3:sec:descent} then uses these
algebraic coefficients to prove that the actual unit
and determinant map \(w=(1,\zeta)\) is a \(K(2)\)-equivalence.

On the locus with connected height one, the completed calculation
provides the constants comparison determined by the Tate frame of rank
two and the connected part.
To reach the ordinary coefficients, \Cref{h3:sec:ordinary} first proves
finite cohomology for each compact lattice of power series.  An
intermediate chromatic descent calculation then supplies finite cohomology
for the union of coefficients with bounded poles.  Only after this step is
the analytic completion removed.  The resulting constants comparison
is the input to \Cref{h3:sec:integral-basis}, where finite Witt
coefficients carry the orientation through compatible length transitions
and derived completion produces \(\delta\).

The rational comparison uses these integral maps and the established
rational exterior algebra, whose height-three proof is reconstructed
in \Cref{h3:app:rational-boundary} following \cite{BSSW}.
\Cref{h3:sec:rational} transports a trace class in degree three along
the actual Tate frame of rank two on the ordinary stratum.  The finite
Witt comparisons carry this class through the coefficient tower, whose
derived limit is taken before rationalization.  The resulting comparison
proves that its image in the line \(\Qp\delta\) is nonzero.  This supplies
the scalar \(c\) in the gluing overview.
Product naturality for the same \(\zeta\) gives the same scalar on
\(\zeta\alpha_3\).

\subsection{Reading the proof}

\Cref{sec:fracture} first constructs the ordered filtration from
precisely stated local maps and a compatibility homotopy, and identifies
all its attachments for those same choices.  The remainder of the paper
constructs the required maps.

\Cref{h3:sec:framework} fixes the coefficient, group, and cochain
conventions.  \Cref{h3:sec:towers} constructs the finite Tate-coordinate
towers and their integral frame comparisons, and
\Cref{h3:sec:analytic} computes the completed strata with their actual
constants maps.  \Cref{h3:sec:coheight-one} removes completion on the
coheight-one stratum, using finite local-field stages and the full
reduced-norm quotient.  \Cref{h3:sec:descent} then compares ordinary
residue cooperations with those stages and obtains the actual
height-two basis.

\Cref{h3:sec:ordinary} proves the second coefficient comparison in the
order required by its dependencies: root and Laurent-tail control,
compact lattice finiteness, intermediate chromatic finiteness,
bounded-pole finiteness, and only then the natural ordinary comparison.
\Cref{h3:sec:integral-basis} carries the orientation through finite Witt
coefficients and derived completion to the actual height-one basis.
The trace argument of \Cref{h3:sec:rational} identifies the localization of the
particular degree-three primitive.

Finally, \Cref{sec:application} inserts these maps into the earlier
fracture construction, proving the main theorem and its attachment
identification from one common choice.  \Cref{sec:attachment} proves a
conditional criterion for the first middle attachment and applies the
companion canonical-map theorem through the local completion,
localization, and scalar comparisons proved there.
\Cref{h3:app:rational-boundary} then reconstructs the rational exterior
algebra from its arithmetic and geometric inputs.

\section{A finite fracture construction}\label{sec:fracture}

This section isolates the topological assembly.  Its inputs are a
height-two equivalence, a compatible height-one and rational pair of
maps, and a basis for the final rational quotient.  The construction
produces the stages themselves and computes every connecting map.
It does not use the coefficient calculations of
\Cref{h3:sec:ordinary}.

\subsection{Module localizations and rotation signs}

Put \(Q=L_0S\simeq H\Qp\) and \(J=K(1)\vee K(2)\).
We work in \(\Mod_S\).  Homological Bousfield localization \(L_E\)
lifts to this category with the same underlying spectrum, whether
or not it is smashing.  Indeed,
\(S^{\wedge n}\wedge M\to S^{\wedge n}\wedge L_EM\) is an
\(E\)-equivalence for every \(n\).  Mapping into the local
spectrum \(L_EM\) extends the action maps and their coherent unit
and associativity homotopies.  The free-module bar resolution
then gives, for every local \(S\)-module \(N\),
\begin{equation}\label{eq:module-localization}
 \Map_{\Mod_S}(L_EM,N)\simeq\Map_{\Mod_S}(M,N).
\end{equation}
Thus the units below are module maps.  The forgetful functor creates
limits and detects equivalences, so the spectral fracture pullbacks
are also pullbacks of \(S\)-modules.

Rationalization is extension of scalars:
\(Q\otimes_SM\simeq L_0M\).  In particular, for a rational
module \(D\),
\begin{equation}\label{eq:rational-adjunction}
 \Map_{\Mod_S}(M,D)
   \simeq\Map_{\Mod_Q}(L_0M,D).
\end{equation}
A map from a finite sum of shifts of \(Q\) is determined up to
homotopy by the images of its module generators.  A homotopy between
the rational maps in \eqref{eq:rational-adjunction} gives the
corresponding \(S\)-module homotopy after precomposition by
\(M\to L_0M\).

The height-two fracture square for an \(E(2)\)-local module \(B\)
is
\begin{equation}\label{eq:height-two-fracture}
 B\simeq L_{K(2)}B
       \mathbin{\times}_{L_1L_{K(2)}B}L_1B.
\end{equation}
For an \(E(1)\)-local module the corresponding square uses \(K(1)\)
and \(L_0\) \cite[Section~2.4, diagram~(2.19)]{BB}.
Consequently an \(E(2)\)-local, \(K(2)\)-acyclic module is
\(E(1)\)-local, and an \(E(2)\)-local module acyclic for both
\(K(1)\) and \(K(2)\) is rational.  There is also a fracture
square separating all positive heights:
\begin{equation}\label{eq:positive-fracture}
 B\simeq L_JB\mathbin{\times}_{L_0L_JB}L_0B.
\end{equation}
To see this, \(L_JB\) is \(E(2)\)-local because
\(\langle E(2)\rangle=\langle K(0)\vee J\rangle\).
The fiber \(N=\fib(B\to L_JB)\) is therefore \(E(2)\)-local
and \(J\)-acyclic, hence rational.  Exactness of \(L_0\) identifies
it with \(\fib(L_0B\to L_0L_JB)\), which proves
\eqref{eq:positive-fracture}.

For \(f:A\to B\) and \(C=\cofib(f)\), we fix the
cofiber-rotation convention
\[
 A\xrightarrow{f}B\longrightarrow C\longrightarrow\Sigma A
 \xrightarrow{-\Sigma f}\Sigma B.
\]
If \(i:N\to A\) is the fiber of \(A\to D\), the induced
identification \(\cofib(A\to D)\simeq\Sigma N\) therefore makes
the boundary to \(\Sigma A\) equal to \(-\Sigma i\).

\begin{lemma}[Boundary of a localization comparison]\label{lem:boundary}
Let \(a:A\to B\) be a map in a stable category, and let \(L\) be
an exact localization for which \(La\) is an equivalence.
Choose an inverse of \(La\) and its inverse homotopies, and put
\[
 \tau:B\longrightarrow LB\xrightarrow{(La)^{-1}}LA,\qquad
 N=\fib(A\longrightarrow LA).
\]
The inverse homotopy makes \(\tau\) a map under \(A\), and
therefore gives
\[
 \bar\tau:\cofib(a)\longrightarrow\cofib(A\to LA)
                         \simeq\Sigma N.
\]
If \(i:N\to A\) is the fiber inclusion, then the connecting map
of the cofiber sequence of \(a\) is
\begin{equation}\label{eq:boundary-lemma}
 \partial_a=(-\Sigma i)\bar\tau.
\end{equation}
If \(M\) is another exact localization and \(N\) is \(M\)-local,
the middle cofiber can equivalently be written
\(\cofib(MA\to MLA)\), using its localization identification.
\end{lemma}

\begin{proof}
The under-\(A\) homotopy supplies a map of cofiber sequences
\[
\begin{tikzcd}[column sep=large]
A \arrow[r,"a"] \arrow[d,equal]
  & B \arrow[r] \arrow[d,"\tau"]
  & \cofib(a) \arrow[r,"\partial_a"] \arrow[d,"\bar\tau"]
  & \Sigma A \arrow[d,equal]\\
A \arrow[r]
  & LA \arrow[r]
  & \Sigma N \arrow[r,"-\Sigma i"]
  & \Sigma A.
\end{tikzcd}
\]
The right square is \eqref{eq:boundary-lemma}; its sign is the
rotation sign fixed above.  If \(N\) is \(M\)-local, so is
\(\cofib(A\to LA)\simeq\Sigma N\).  Its \(M\)-localization
map is then an equivalence, and exactness identifies the localized
cofiber with \(\cofib(MA\to MLA)\).
\end{proof}

\subsection{The stages and their boundaries}

Use the following ordered blocks:
\[
\begin{aligned}
 W&=W_1\vee W_2=L_2S\vee\Sigma^{-1}L_2S,\\
 U_{\mathrm{mid}}&=U_3\vee U_4=\Sigma^{-3}L_1S\vee\Sigma^{-4}L_1S,\\
 V&=V_5\vee V_6\vee V_7\vee V_8\\
  &=\Sigma^{-5}Q\vee\Sigma^{-6}Q
       \vee\Sigma^{-8}Q\vee\Sigma^{-9}Q.
\end{aligned}
\]

\begin{proposition}[Ordered fracture filtration]\label{prop:fracture}
Let \(X\) be an \(E(2)\)-local \(S\)-module.  Suppose the
following module maps and homotopies are given.
\begin{enumerate}[label=\textup{(\roman*)}]
\item A map \(w:W\to X\) for which \(L_{K(2)}w\) is an
  equivalence.  Put \(Y=\cofib(w)\), with quotient \(q:X\to Y\).
\item An equivalence
  \(\kappa:L_{K(1)}U_{\mathrm{mid}}\xrightarrow{\simeq}L_{K(1)}Y\), a rational
  map \(r:L_0U_{\mathrm{mid}}\to L_0Y\), and a specified homotopy of \(Q\)-module
  maps between
  \begin{equation}\label{eq:compatibility}
  \begin{gathered}
  L_0U_{\mathrm{mid}}\longrightarrow L_0L_{K(1)}U_{\mathrm{mid}}
          \xrightarrow{L_0\kappa}L_0L_{K(1)}Y,\\
  L_0U_{\mathrm{mid}}\xrightarrow rL_0Y\longrightarrow L_0L_{K(1)}Y.
  \end{gathered}
  \end{equation}
\item A map \(b:V\to L_0X\) for which
  \begin{equation}\label{eq:quotient-basis}
  t_0:V\xrightarrow bL_0X\xrightarrow{L_0q}L_0Y
                 \longrightarrow\cofib(r)
  \end{equation}
  is an equivalence.
\end{enumerate}
Then there is a filtration in \(E(2)\)-local \(S\)-modules
\[
 0=F_0\longrightarrow F_1\longrightarrow\cdots\longrightarrow F_8
 \xrightarrow[\simeq]{e}X
\]
with ordered cofibers \(W_1,W_2,U_3,U_4,V_5,V_6,V_7,V_8\).
The composite \(F_2=W\to F_8\xrightarrow eX\) is \(w\).
\end{proposition}

\begin{proof}
\emph{Constructing the middle map.}
The cofiber \(Y\) is \(E(2)\)-local and \(K(2)\)-acyclic by
\textup{(i)}, hence \(E(1)\)-local.  Exactness identifies it
naturally with \(\cofib(L_1w)\).  Its height-one fracture square is
\[
 Y\simeq L_{K(1)}Y\mathbin{\times}_{L_0L_{K(1)}Y}L_0Y.
\]
The two proposed maps from \(U_{\mathrm{mid}}\) to its corners are
\[
 U_{\mathrm{mid}}\longrightarrow L_{K(1)}U_{\mathrm{mid}}\xrightarrow{\kappa}L_{K(1)}Y,
 \qquad
 U_{\mathrm{mid}}\longrightarrow L_0U_{\mathrm{mid}}\xrightarrow rL_0Y.
\]
By \eqref{eq:rational-adjunction}, the specified rational homotopy
in \eqref{eq:compatibility} identifies their composites to the
overlap as \(S\)-module maps.  The pullback therefore gives
\begin{equation}\label{eq:middle-map}
 h:U_{\mathrm{mid}}\longrightarrow Y,\qquad L_{K(1)}h=\kappa,\quad L_0h=r.
\end{equation}
Here each equality includes the homotopy obtained from the
corresponding pullback projection: localizing that homotopy and
using idempotence gives the displayed identification.  In
particular, \(h\) is a \(K(1)\)-equivalence.

\emph{The first four stages.}
Let \(\partial_w:Y\to\Sigma W\) be the boundary of \(w\), and put
\[
 d=\partial_wh:U_{\mathrm{mid}}\longrightarrow\Sigma W,\qquad
 d_i=d|_{U_i}\quad(i=3,4).
\]
Define
\begin{equation}\label{eq:first-stages}
 F_0=0,\quad F_1=W_1,\quad F_2=W,\quad
 F_3=X\mathbin{\times}_YU_3,\quad
 F_4=X\mathbin{\times}_YU_{\mathrm{mid}}.
\end{equation}
The first nonzero transition is the wedge inclusion \(W_1\to W\).
The identification \(X\times_Y0\simeq W\) makes the next transitions
the pullbacks of \(0\to U_3\to U_3\vee U_4\).
In a stable category a pullback square is also a pushout square.
Their cofibers are consequently \(U_3\) and \(U_4\), and
\[
 F_3\simeq\fib(d_3:U_3\to\Sigma W),\qquad
 F_4\simeq\fib(d:U_{\mathrm{mid}}\to\Sigma W).
\]
Naturality of the cofiber sequence for each pullback gives
\begin{equation}\label{eq:first-boundaries}
 \partial_1=0,\qquad \partial_2=0,\qquad
 \partial_3=d_3,\qquad
 \partial_4=\Sigma(F_2\to F_3)d_4.
\end{equation}
For the last formula, the quotient of the pullback square for
\(U_3\to U_3\vee U_4\) identifies its quotient with \(U_4\).
There is a map of cofiber sequences from
\(W\to F_4\to U_{\mathrm{mid}}\) to
\(F_3\to F_4\to U_4\), whose components are
\(W\to F_3\), the identity of \(F_4\), and the projection
\(\operatorname{pr}_4:U_{\mathrm{mid}}\to U_4\).
Its boundary square gives
\(\partial_4\operatorname{pr}_4=\Sigma(W\to F_3)d\).
Restricting to the summand \(U_4\) gives the displayed formula.

We next express \(d\) through the height-two fracture map.  Choose
coherent inverse data for \(L_{K(2)}w\), and form
\[
 \tau_2:L_1X\longrightarrow L_1L_{K(2)}X
       \xrightarrow{(L_1L_{K(2)}w)^{-1}}L_1L_{K(2)}W.
\]
It is a map under \(L_1W\), so it induces
\[
 \bar\tau_2:Y\simeq\cofib(L_1w)
       \longrightarrow\cofib(L_1W\to L_1L_{K(2)}W).
\]
Let \(j_W:\fib(W\to L_{K(2)}W)\to W\) be the fiber inclusion.
Then
\begin{equation}\label{eq:first-roof}
\begin{split}
 d:\ U_{\mathrm{mid}}\xrightarrow hY\xrightarrow{\bar\tau_2}
 &\cofib(L_1W\longrightarrow L_1L_{K(2)}W)\\
 &\simeq\Sigma\fib(W\longrightarrow L_{K(2)}W)
   \xrightarrow{-\Sigma j_W}\Sigma W.
\end{split}
\end{equation}
Indeed, \eqref{eq:height-two-fracture} identifies the fiber of
\(W\to L_{K(2)}W\) with the fiber of
\(L_1W\to L_1L_{K(2)}W\); it is \(E(1)\)-local.
\Cref{lem:boundary}, with \(L=L_{K(2)}\) and \(M=L_1\),
then says that the composite from \(Y\) in
\eqref{eq:first-roof} is exactly \(\partial_w\).  This proves
the formula and fixes its minus sign.

\emph{The rational quotient.}
Let \(a:F_4\to X\) be the pullback projection and set
\(Z=\cofib(a)\).  Taking the cofiber of the pullback square defining
\(F_4\) gives a natural equivalence
\begin{equation}\label{eq:quotient-identification}
 Z\simeq\cofib(h:U_{\mathrm{mid}}\longrightarrow Y).
\end{equation}
This cofiber is \(E(1)\)-local and \(K(1)\)-acyclic by
\eqref{eq:middle-map}, hence rational.  Thus \(a\) is an
equivalence after localization at \(K(1)\) or \(K(2)\), and
therefore after \(L_J\).  Rationalizing
\eqref{eq:quotient-identification} identifies
\(Z\simeq\cofib(r)\).

Write \(z:X\to Z\) for the quotient.  Since \(Z\) is rational,
\eqref{eq:rational-adjunction} gives its factorization
\(\widetilde z:L_0X\to Z\).  Under the preceding cofiber
identification, the map in \textup{(iii)} is
\begin{equation}\label{eq:terminal-quotient}
 t=\widetilde z\,b:V\xrightarrow{\simeq}Z.
\end{equation}
This records the factorization used to compare the last attachments
with fracture.

\emph{The last four stages.}
Let \(V_{\leq j}\) be the first \(j\) summands of \(V\), for
\(0\leq j\leq4\), and use the restriction of \(t\) to define
\begin{equation}\label{eq:last-stages}
 F_{4+j}=X\mathbin{\times}_ZV_{\leq j}.
\end{equation}
At \(j=0\), the cofiber sequence of \(a\) identifies this
definition with the previous \(F_4\), over \(X\).
The cofiber of \(V_{\leq j-1}\to V_{\leq j}\) is \(V_{4+j}\).
Exact pullback gives the same cofiber for
\(F_{4+j-1}\to F_{4+j}\).

Let \(\partial_a:Z\to\Sigma F_4\) be the boundary of \(a\),
put \(g=\partial_at\), and write \(g_i=g|_{V_i}\).
The quotient-of-a-pullback argument used above gives every remaining
connecting map:
\begin{equation}\label{eq:last-boundaries}
 \partial_i=\Sigma(F_4\to F_{i-1})g_i:
 V_i\longrightarrow\Sigma F_{i-1},\qquad 5\leq i\leq8.
\end{equation}

The map \(g\) is described by the positive-height fracture square.
Choose coherent inverse data for \(L_Ja\), and let
\(j_a:\fib(F_4\to L_JF_4)\to F_4\) be the fiber inclusion.
The formula is
\begin{equation}\label{eq:second-roof}
\begin{aligned}[b]
 g:\ V\xrightarrow bL_0X\longrightarrow L_0L_JX
 &\xrightarrow{(L_0L_Ja)^{-1}}L_0L_JF_4\\
 &\longrightarrow\cofib(L_0F_4\longrightarrow L_0L_JF_4)\\
 &\simeq\Sigma\fib(F_4\longrightarrow L_JF_4)
   \xrightarrow{-\Sigma j_a}\Sigma F_4.
\end{aligned}
\end{equation}
We verify in particular why this roof begins with \(b\), rather
than only with the quotient equivalence \(t\).
Put
\[
 C_a=\cofib(F_4\to L_JF_4),\qquad
 \tau_J=(L_Ja)^{-1}\circ(X\to L_JX),
\]
and let \(q_a:L_JF_4\to C_a\) be the quotient.  The chosen inverse
homotopy makes \(\tau_Ja\) the localization unit of \(F_4\).
\Cref{lem:boundary} gives \(\bar\tau_J:Z\to C_a\), with a
specified homotopy
\[
 \bar\tau_Jz\simeq q_a\tau_J.
\]
Both \(Z\) and \(C_a\) are rational: for \(C_a\) this follows
from \eqref{eq:positive-fracture}.  If
\(\rho:C_a\xrightarrow{\simeq}L_0C_a\) is its localization map,
the rational adjunction transports the last homotopy to
\[
 \rho\bar\tau_J\widetilde z\simeq (L_0q_a)(L_0\tau_J).
\]
Using \(t=\widetilde z b\), this becomes
\[
 \rho\bar\tau_Jt\simeq
 (L_0q_a)(L_0L_Ja)^{-1}\circ
 (L_0X\to L_0L_JX)\circ b.
\]
Exactness identifies \(L_0C_a\) with the cofiber in
\eqref{eq:second-roof}.  The boundary formula
\(\partial_a=(-\Sigma j_a)\bar\tau_J\) now identifies that
entire roof with \(\partial_at=g\), including the second minus sign.

Finally, \eqref{eq:terminal-quotient} makes the projection
\[
 e:F_8=X\times_ZV\longrightarrow X
\]
an equivalence.  All stages use finite limits or colimits of
\(E(2)\)-local \(S\)-modules, so they remain in that category.
Their maps to \(X\) extend \(w\); hence
\(F_2\to F_8\xrightarrow eX\) is \(w\).
\end{proof}

\begin{theorem}[Attachment identification for the ordered fracture filtration]
\label{thm:attachment-identification}
Fix the data \(w,\kappa,r,b\) and the compatibility homotopy of
\Cref{prop:fracture}.  Choose coherent inverse data for
\(L_{K(2)}w\), and carry out that proposition's construction.
Use exactly its fracture map \(h\), pullback stages \(F_i\),
projection \(a:F_4\to X\), quotient \(Z=\cofib(a)\),
equivalence \(t=\widetilde z b:V\xrightarrow{\simeq}Z\), and
terminal projection \(e:F_8\to X\), choosing the coherent inverse
data for \(L_Ja\) at the indicated step.  Under the resulting
ordered cofiber identifications
\[
 W_1,\ W_2,\ U_3,\ U_4,\ V_5,\ V_6,\ V_7,\ V_8,
\]
the connecting maps of this same filtration are
\[
\begin{gathered}
 \partial_1=0,\qquad \partial_2=0,\qquad
 \partial_3=d_3,\qquad
 \partial_4=\Sigma(F_2\to F_3)d_4,\\
 \partial_i=\Sigma(F_4\to F_{i-1})g_i
 \qquad(5\leq i\leq8),
\end{gathered}
\]
where \(d=\partial_wh\), \(g=\partial_at\),
\(d_i=d|_{U_i}\), and \(g_i=g|_{V_i}\).  The maps \(d\) and \(g\)
are the signed roofs
\[
\begin{aligned}
 d:\ U_{\mathrm{mid}}\xrightarrow hY\xrightarrow{\bar\tau_2}
 &\cofib(L_1W\longrightarrow L_1L_{K(2)}W)\\
 &\simeq\Sigma\fib(W\longrightarrow L_{K(2)}W)
   \xrightarrow{-\Sigma j_W}\Sigma W,
\end{aligned}
\]
and
\[
\begin{aligned}
 g:\ V\xrightarrow bL_0X\longrightarrow L_0L_JX
 &\xrightarrow{(L_0L_Ja)^{-1}}L_0L_JF_4\\
 &\longrightarrow\cofib(L_0F_4\longrightarrow L_0L_JF_4)\\
 &\simeq\Sigma\fib(F_4\longrightarrow L_JF_4)
   \xrightarrow{-\Sigma j_a}\Sigma F_4.
\end{aligned}
\]
Here \(\bar\tau_2,j_W,j_a\) are the maps defined from the same
construction and inverse data above.
\end{theorem}

\begin{proof}
The proof of \Cref{prop:fracture} established
\eqref{eq:first-boundaries} and \eqref{eq:last-boundaries} under
these ordered cofiber identifications.  It identified their
components by \eqref{eq:first-roof} and \eqref{eq:second-roof}
using precisely the stated inverse data.  The collected formulas
therefore apply to the same \(h,F_i,a,Z,t,e\).
\end{proof}

\begin{remark}
The choices of compatibility homotopy and coherent inverses are
part of this realization.  \Cref{thm:attachment-identification}
does not claim that the
filtration is unique or that its attachments are canonical without
those choices.  Together with \Cref{prop:fracture}, it identifies the
original map \(w\), the terminal projection, and the connecting maps
associated with the stated data.
In the height-three application, only the first component of \(w\)
is asserted to be the canonical unit.
\end{remark}

\section{Coefficients, stabilizers, and comparison inputs}\label{h3:sec:framework}

The coefficient arguments require three kinds of data: the local sphere
and its finite constant-field extensions, the stabilizer actions on the
deformation rings, and the topology used on continuous cochains.  We fix
these data here and state the external bundle and boundary results at the
precise scope used later.  The actual comparisons between algebraic and
completed coefficients will be proved in the following sections.

\subsection{Localizations and coefficient fields}

Fix \(p\geq5\). We work with \(p\)-local spectra and write
\[
 S=\mathbb S_p^\wedge,\qquad T=L_{K(3)}S,\qquad
 X=L_2T,\qquad R_i=L_{K(i)}S.
\]
We use the natural \(S\)-module structures throughout. In particular,
\(L_0S=H\Q_p\), and the rational maps used in fracture are
\(H\Q_p\)-linear. The rationalization of \(L_iS\), for \(i=1,2\), is
\(H\Q_p\); it should be distinguished from \(L_0R_i\).

Choose a finite field \(k\supseteq\F_{p^6}\). Further finite enlargement
will be specified when needed to define markings or characters.
Let \(W(k)\) be its Witt ring. The finite étale extension
\(\Z_p\to W(k)\) has a unique finite étale lift to an \(E_\infty\)
\(S\)-algebra \(S_k\), with
\[
 \pi_*S_k=\pi_*S\otimes_{\Z_p}W(k).
\]
This is the finite étale classification for ring spectra;
see \cite[Theorem~2.32]{Mathew}. We write \(R_{i,k}=L_{K(i)}S_k\).
A \(\Z_p\)-basis of \(W(k)\) gives an equivalence of the underlying
\(S\)-module with a finite wedge of copies of \(S\), so extension to
\(S_k\) is faithful. The Galois automorphisms lift coherently through
the same classification. They will be used to descend constructions
made over \(k\).

Let \(E_n(k)\) be Morava \(E\)-theory for the height-\(n\) Honda formal
group over \(k\). At height three our coefficient convention is
\[
 E_{3,*}(k)=W(k)[[x,y]][u^{\pm1}],\qquad |u|=-2.
\]
The invariant cotangent, or periodicity, line is
\(\omega=\pi_2E_3(k)\); it is generated by \(u^{-1}\) over the
degree-zero ring. Set
\[
 A=k[[x,y]],\qquad x=u_1,\quad y=u_2,\qquad B_0=A[x^{-1}].
\]
Twists by powers of \(\omega\) and finite characters are retained
whenever a finiteness argument requires them. A trivialization made
after finite extension will always be descended.

\subsection{Stabilizers and reduced norm}

Let \(D_n\) be the central division algebra over \(\Q_p\) of invariant
\(1/n\), and let
\[
 P_n=\mathcal O_{D_n}^{\times}.
\]
This is the ordinary Honda stabilizer. Its action on \(E_n(k)\) fixes
\(W(k)\). The coefficient-field Galois group acts semilinearly and gives
the extended stabilizer
\[
 G_n(k)=P_n\rtimes\Gal(k/\F_p).
\]
The field contains the endomorphism field for every \(n=1,2,3\) used here.

At height three, define
\[
 H=\ker\bigl(\hthreeNrd:P_3\to\Z_p^\times\bigr),\qquad
 H'=\hthreeNrd^{-1}(\mu_{p-1}).
\]
The reduced norm is surjective, and therefore
\begin{equation}\label{h3:eq:framework-norm-quotients}
 H'/H\simeq\mu_{p-1},\qquad
 P_3/H'\simeq1+p\Z_p\simeq\Z_p.
\end{equation}
For example, the norm on the maximal unramified subfield is surjective
on units. The logarithm on the last quotient is normalized using a
topological generator. The resulting integral determinant class is
denoted \(\zeta\in\pi_{-1}T\); its construction as an actual sphere map
is recalled in \Cref{h3:sec:descent}.

\begin{lemma}\label{h3:lem:framework-groups}
The compact groups \(P_3,H,P_2,P_1,\GL_2(\Z_p)\) and
\(\hthreeSL_2(\Z_p)\) have no elements of order \(p\). Their \(p\)-cohomological
dimensions are respectively \(9,8,4,1,4,3\). Their adjoint orientation
characters are trivial. Their trivial \(\Z_p\)-modules admit finite
resolutions by finitely generated projective completed group-ring
modules.
\end{lemma}

\begin{proof}
An element of order \(p\) in \(D_n^\times\) would generate a copy of
\(\Q_p(\zeta_p)\). Its degree \(p-1\) must divide \(n^2\), since
\(D_n\) would be a vector space over that subfield. For \(n=1,2,3\)
and \(p\geq5\), this is impossible except for the apparent numerical
possibility \(n=2,p=5\); in that case a commutative subfield of a
degree-two central division algebra has degree at most two, whereas
\(p-1=4\). Equivalently, the usual subfield-degree theorem rules out
all these cases at once. An element of order \(p\) in
\(\GL_2(\Q_p)\) would give a faithful two-dimensional representation
of the nontrivial irreducible factor of \(X^p-1\), whose degree is
\(p-1>2\), and is likewise impossible.

The Lie dimensions are the stated ones. Reduced norm has surjective
differential, so its kernel has dimension eight. Conjugation has
determinant one on a central simple algebra; its restriction to the
trace-zero summand also has determinant one. The same calculation
applies to the matrix groups. Thus their adjoint orientation
characters are trivial.

For a compact \(p\)-adic analytic group without \(p\)-torsion,
\(p\)-cohomological dimension equals its Lie dimension
\cite[Section~1, Corollary~(1)]{Serre1965CohomologicalDimension}.
Writing \(\Lambda=\Z_p[[G]]\), Venjakob
\cite[pp.~275--276]{Venjakob} records both Noetherianity of \(\Lambda\)
and \(\operatorname{pd}_{\Lambda}\Z_p=\operatorname{cd}_p(G)\),
with agreement between finitely generated compact and abstract
modules.  Successive finite free presentations therefore have finitely
generated kernels and a projective final syzygy, giving the asserted
finite projective resolution. The associated duality is used with the
continuous coefficient models specified below.
\end{proof}

No torsion-freeness claim about an arbitrary finite Galois factor is
needed. Its action is instead handled by semilinear descent.

\begin{lemma}\label{h3:lem:framework-semilinear}
Let \(k'/k\) be a finite extension of finite fields with Galois group
\(\Gamma\). On \(W(k')\)-modules carrying a semilinear \(\Gamma\)-action,
the invariants functor is exact. The same assertion holds for reductions
modulo \(p^\ell\), and the associated finite-group higher cohomology
vanishes. These assertions respect continuous equivariant maps.
\end{lemma}

\begin{proof}
The residue-field trace is surjective. Lifting a trace-one element to
\(W(k')\) and dividing by its unit trace produces \(a\in W(k')\) with
\(\sum_{\gamma\in\Gamma}\gamma(a)=1\). For a semilinear module \(M\), the
finite sum
\[
 P(m)=\sum_{\gamma\in\Gamma}\gamma(a)\,\gamma(m)
\]
is an additive projection onto \(M^\Gamma\). If an invariant element
in a quotient has a lift, applying \(P\) to that lift gives an invariant
lift. This proves exactness, and the same formula works modulo
\(p^\ell\). For an equivariant homogeneous cochain \(f\) of positive
degree, the operator
\[
 (hf)(g_0,\ldots,g_{q-1})
 =\sum_{\gamma\in\Gamma}\gamma(a)
 f(\gamma,g_0,\ldots,g_{q-1})
\]
is still equivariant, by semilinearity and reindexing the sum. The
alternating homogeneous differential satisfies \(dh+hd=\hthreeid\) in
positive degrees, since the sum of the coefficients is one. This proves
the higher vanishing. All sums and scalar operations are finite and
continuous. No division by \(|\Gamma|\) has been made.
\end{proof}

\subsection{Continuous cochains and filtered coefficients}

For a topological group \(G\) and a continuous coefficient module \(M\),
we use the inhomogeneous continuous cochain complex
\[
 C^q_{\hthreects}(G,M)=C_{\hthreects}(G^q,M)
\]
with its usual differential. Complete power-series and valuation
lattices have their linear topologies. Their discrete quotients have
the quotient topology. For a profinite parameter space \(Q\), all
constructions are also made with coefficients
\(C_{\hthreects}(Q,M)\).

A filtered localization or algebraic union is interpreted at the level
of complexes. If \(M=\colim_\lambda M_\lambda\) is such a coefficient
object, our notation means
\begin{equation}\label{h3:eq:framework-filtered-cochains}
 C^\bullet_{\hthreects}(G,M)
 :=\colim_\lambda C^\bullet_{\hthreects}(G,M_\lambda).
\end{equation}
Thus a cochain has a common finite stage or pole bound. This is the
coefficient convention produced by ordinary tensor products in the
residue tests. It is generally different from the complex of all
metric-continuous maps into a topologized union.

We use completed group-ring resolutions to compute continuous
cohomology only after comparing them with the continuous cochain model
on the coefficient category in question. For complete linearly
topologized modules the comparison, compact duality, and the required
strictness statements are established in
\Cref{h3:sec:coheight-one}. For the geometric integral complexes in
\Cref{h3:sec:rational}, profinite parameters are retained through descent
and the derived limits. The finite solid resolution argument there
justifies the later calculation on their point-valued cohomology rows.

The holomorphic coefficient ring \(B_{\hthreehol}\) consists of integral-exponent
Laurent series converging on
\[
 0<|x|<1,\qquad |y|<1
\]
over the trivially valued field \(k\). We give it the inverse-limit
topology from Gauss norms on an exhaustion by closed polyannuli.
The analogous analytic completed perfection is \(B_{\hthreepk}\).
The precise exhaustion, coefficient growth conditions and root
projections are specified in \Cref{h3:sec:ordinary}. Bounded-pole
submodules of \(B_0\) carry their compact lattice topology. We
distinguish the coheight-one Cartier operator \(\mathcal C\) from
the ordinary root-projection operator \(\mathcal P\). Their
different Frobenius identities will be proved on their respective
coefficient modules.

\subsection{Geometric coefficient theory}

We use perfectoid tilting and the Fargues--Fontaine curve over an
algebraically closed perfectoid field. The notation
\(\mathcal O^\flat\) denotes the analytic tilted structure sheaf;
an integral valuation lattice is indicated separately. In particular,
the affinoid perfectoid acyclicity used below concerns this analytic
sheaf. We do not replace an integral almost-vanishing statement by
ordinary vanishing.

Let
\[
 V_s=\mathcal O(1/s)
\]
be the stable rank-\(s\), degree-one bundle. We use the classification
of vector bundles by slopes, \(H^0(\mathcal O)=\Q_p\), and the
relative vanishing of \(H^1\) for trivial bundles and positive
basic bundles as a \(v\)-sheaf. On a fixed basic stratum, bundles admit \(v\)-local
standard forms with the indicated locally profinite automorphism
groups. Slope-zero bundles correspond to rational local systems;
their integral lattices are additional data. These are the precise
background bundle results needed from
\cite{FarguesFontaine,KedlayaLiu,ScholzeWeinstein}; the precise relative arguments appear in
\Cref{h3:geom:curve-inputs}.

The Honda universal cover identifies the section sheaf of \(V_s\)
with the perfected open Honda ball, compatibly with addition and
endomorphisms. The Banach--Colmez description is provided by
\cite{LeBras,ScholzeWeinstein}. For maps over a varying perfectoid base we use
relative full faithfulness of derived sections
\cite[Corollary~3.11]{AnschuetzLeBrasFourier}; its application to positive bundles
is made explicit in \Cref{h3:sec:towers}.
\Cref{h3:sec:towers} verifies the finite-level coordinate comparison
and the relative integral frame reductions used in this paper.
\Cref{h3:sec:analytic} proves the nonproper relative support calculations;
they are not supplied by affinoid acyclicity alone.

\subsection{The two boundary calculations}

We use two boundary calculations. The rational height-three
calculation is
\begin{equation}\label{h3:eq:framework-rational-boundary}
 \pi_*L_0T
 \simeq \Lambda_{\Q_p}(\alpha_1,\alpha_3,\alpha_5),
 \qquad |\alpha_i|=-i.
\end{equation}
For the primes \(p\geq5\) considered here, \Cref{h3:prop:rational-boundary}
gives the argument of \cite[Theorem~A]{BSSW}, including its arithmetic
and tower-comparison inputs. Its identification with the particular
constant trace class needed here is treated separately in
\Cref{h3:sec:rational}; a rational dimension calculation cannot identify
a localization map.

For context, the known odd-prime height-two splitting is
\begin{equation}\label{h3:eq:framework-height-two-boundary}
 L_1R_2\simeq
 L_1(S\vee\Sigma^{-1}S)
 \vee L_0(\Sigma^{-3}S\vee\Sigma^{-4}S),
\end{equation}
as stated in \cite[Theorem~6.7]{BB}.  Exactness also gives
\(R_2/p\simeq L_{K(2)}(S/p)\), where \(S/p\) is the finite Moore
spectrum.  For \(p\geq5\), \cite[Theorem~15.1]{HS99} therefore implies
that \(\pi_a(R_2/p)\) is finite for every integer \(a\).
The finite groups \(\pi_a(R_2/p)\) give one boundary for the
intermediate finiteness argument in \Cref{h3:sec:ordinary}.  That
argument derives the required finiteness of \(L_1(R_2/p)\) directly
from the height-two group calculation of \cite[Remark~7.8]{Behrens2012}.
The coheight-one test, the ordinary comparison, and all the attaching
maps are proved below.

Morava descent will be used with an explicit convergence statement:
the completed Amitsur error tower has a uniform nilpotence bound, so
the exact tests in this proof preserve its totalization, as derived
from descendability in \Cref{h3:prop:exact-descent}.  The cooperation
calculation of Devinatz--Hopkins
\cite[Theorem~2(ii) and Proposition~2.2]{DevinatzHopkins} is used
first for the usual coefficient ring $W(\F_{p^n})$ and extended
stabilizer.  The passage to the chosen $k$, the relative algebra over
$S_k$, and its ordinary-stabilizer action-cobar terms are local
arguments of \Cref{h3:sec:descent}, which also retains the actual unit
and cofaces.

\section{Tate coordinates and integral frames}\label{h3:sec:towers}

Our goal is to describe the two completed height strata by extension
spaces carrying their actual integral frames.  The finite torsion
coordinates first identify the completed Tate-basis towers.  Relative
sections on the Fargues--Fontaine curve then identify their quotient
frames and determinant characters.  We finish by descending the
markings and the finite root torsors to the algebraic coefficient rings
needed for ordinary descent.

Fix a finite field $k$ containing $\F_{p^6}$.  For $s=1,2,3$ let
$\Gamma_s$ be a Honda formal group over $k$, with
$[p]_{\Gamma_s}(T)=T^{p^s}$, and put
\[
 D_s=\hthreeEnd_k(\Gamma_s)[1/p],\qquad
 P_s=\Aut_k(\Gamma_s)=\mathcal O_{D_s}^{\times}.
\]
We use the convention in which the corresponding Fargues--Fontaine bundle
$V_s=\mathcal O(1/s)$ has rank $s$ and degree one.  The notation $H_s$
denotes the perfected open Honda ball: on an affinoid perfectoid
$k$-algebra $(R,R^+)$ it is $R^{\circ\circ}$, with addition defined by
$\Gamma_s$.  It is a sheaf of $\Q_p$-vector spaces, since $[p]$ is
invertible.  All assertions about a locus in these sheaves are understood
fiberwise at geometric points.

Let $\mathcal G$ be the $p$-divisible group obtained from the universal
height-three deformation, reduced modulo $p$, over $A=k[[x,y]]$.
Here and below this is the algebraic $p$-divisible group formed by its
finite kernels.  In detail, preparation applied to $[p^l](T)$ gives
finite free algebras
\[
 A[[T]]/([p^l](T))
\]
of rank $p^{3l}$.  Their coordinates are topologically nilpotent for the
maximal-ideal topology, so substitution in the formal group law defines
their algebraic group operations.  Preparation applied to
$[p](T)-Z$ proves the required finite flat transition maps.  We base
change these finite groups when we pass to a height stratum; we do not
complete the generic fiber at its identity and thereby discard its
\'etale part.

For $m=1,2$ put $r=3-m$ and
\[
 A_m=A/(u_1,\ldots,u_{m-1})=k[[u_m,\ldots,u_2]],
 \qquad u_1=x,\quad u_2=y,
\]
where the ideal is zero for $m=1$.  We use the same symbol
$\mathcal G$ for its base change to $\operatorname{Spec}(A_m)$ and
for subsequent base changes to covers.  Let $\mathcal X_m$ be the completed perfection of the analytic locus
where the displayed parameters are topologically nilpotent and $u_m$
is invertible.  The $p$-divisible group on this locus has connected
height $m$ and \'etale height $r$.  Let $\mathcal T_m\to\mathcal X_m$
be the tower of integral bases of its \'etale Tate module.  Thus its
structure group is $\GL_r(\Z_p)$.

\subsection{Finite levels and their norms}

\begin{lemma}\label{h3:geom:regular-levels}
For $l\geq1$, form the reduced closure of the locus of $r$ points of
$\mathcal G[p^l]$ whose images are a basis of the \'etale quotient at
the generic point of $A_m$.  Its coordinate algebra $B_l$ is finite
over $A_m$, and its lifted point coordinates $t_{1,l},\ldots,t_{r,l}$
give isomorphisms
\[
 B_l=k[[t_{1,l},\ldots,t_{r,l}]],\qquad
 C_l:=A_m[t_{1,l}^{p^{ml}},\ldots,t_{r,l}^{p^{ml}}]
     =k[[t_{1,l}^{p^{ml}},\ldots,t_{r,l}^{p^{ml}}]].
\]
The subalgebra on the left is already complete.  The algebra
$C_l[u_m^{-1}]$ is the finite \'etale algebra of bases of
$\mathcal G^{\mathrm{et}}[p^l]$ over $A_m[u_m^{-1}]$.
Over this basis cover, $B_l[u_m^{-1}]$ is the algebra of all
generator lifts, without taking a reduction.
The transition maps $B_j\to B_l$, $j\leq l$, are finite injective maps,
and
\[
 t_{i,j}\in
 k[[t_{1,l}^{p^{m(l-j)}},\ldots,t_{r,l}^{p^{m(l-j)}}]].
\]
\end{lemma}

\begin{proof}
One precise definition of the closure is the image of the finite
$A_m$-algebra of $r$-tuples in $\mathcal G[p^l]$ in the reduced
coordinate algebra of the indicated generic locus.  It is therefore
finite and reduced, the map $A_m\to B_l$ is injective, and every
irreducible component dominates $\Spec A_m$.  At the closed point all
torsion coordinates are nilpotent.  Consequently $B_l$ is local with
residue field $k$, and its maximal ideal is generated by the height
parameters and the $t_{i,l}$.

Write $d=p^m$ and $q=p^3$.  Put
$a_i=[p^{l-1}]_{\mathcal G}(t_{i,l})$, and for
$v=(v_1,\ldots,v_r)\in\F_p^r$ put
$a_v=\sum_{\mathcal G}[v_i]a_i$.  If $W_p(T)$ is the distinguished
polynomial for $[p]_{\mathcal G}(T)$, then
\begin{equation}\label{h3:geom:basis-divisor}
 W_p(T)=\prod_{v\in\F_p^r}(T-a_v)^d
       =\prod_{v\in\F_p^r}(T^d-a_v^d).
\end{equation}
Indeed, over each geometric generic field the connected part has
length $d$ and the $a_v$ enumerate the distinct \'etale points.
Both sides are monic of degree $q$.  Equality of their coefficients
there implies equality in the reduced generic closure.

On $A_m$, the series $[p](T)$ factors through $T^d$.
It follows that every $a_v^d$ is a series over $A_m$ in
$w_i=t_{i,l}^{d^l}$, with zero constant term as a series in the $w_i$.
The algebra $C_l=A_m[w_1,\ldots,w_r]$ is finite over the complete ring
$A_m$, hence complete.  Equation~\eqref{h3:geom:basis-divisor} is an
identity over $C_l$.  Modulo $(w_1,\ldots,w_r)$ it becomes $W_p(T)=T^q$.
The preparation unit for $[p](T)$ lies over $A_m$; thus every
coefficient below degree $q$ in $[p](T)$ vanishes in this quotient.
The successive height-coordinate congruences
\[
 [p](T)\equiv u_iT^{p^i}
 \pmod{(u_m,\ldots,u_{i-1},T^{p^i+1})},\qquad m\leq i\leq2,
\]
now imply, successively, that all $u_i$ vanish.  Hence the maximal
ideal of $C_l$ is generated by the $r$ elements $w_i$.
The ring has dimension $r$, because it is integral over $A_m$.
It is a regular local ring.  The surjection
$k[[W_1,\ldots,W_r]]\to C_l$, $W_i\mapsto w_i$, is an
isomorphism: a nonzero kernel would lower dimension in the regular
domain on the left.  Applying the identical maximal-ideal argument
with $t_{i,l}$ in place of $w_i$ proves $B_l=k[[t_{1,l},\ldots,t_{r,l}]]$.
This also proves the asserted identification of the inclusion
$C_l\subset B_l$, rather than just abstract regularity of both rings.

Over $u_m\ne0$, relative Frobenius of order $ml$ identifies the
quotient of $\mathcal G[p^l]$ by its connected part with its image
under $T\mapsto T^{d^l}$.  Here is a check that retains the
prepared finite algebra.  Write $[p^l](T)=g_l(T^{d^l})$ and
prepare $g_l(S)=U_l(S)V_l(S)$, with $U_l$ a unit and $V_l$
monic of degree $p^{rl}$.  Its linear coefficient $a_l$ is a
unit multiple of $u_m^{1+d+\cdots+d^{l-1}}$.  Invariant
differentials for this Frobenius-divided homomorphism give
\[
 g_l'(S)=a_l\frac{h_{\mathrm{source}}(S)}
                         {h_{\mathrm{target}}(g_l(S))},
\]
where both $h$ are unit series.  In the finite algebra defined
by $V_l$, this says that $U_lV_l'$ becomes a unit after
inverting $u_m$.  Thus that algebra is \'etale of rank
$p^{rl}$.  The kernel of relative Frobenius has rank $d^l$
and is the connected part.  Taking bases gives exactly
$S_l=C_l[u_m^{-1}]$.

Over $S_l$, the scheme of lifts of its universal basis is
finite locally free of rank $d^{lr}$.  Its coordinate algebra
surjects onto $B_l[u_m^{-1}]$: the latter is generated by the
same lifted coordinates.  But the identified power-series
rings show that $B_l[u_m^{-1}]$ is free over $S_l$ of this
same rank, with the monomials $\prod_i t_{i,l}^{e_i}$,
$0\leq e_i<d^l$, as a basis.  A surjection between finite
locally free modules of the same rank is an isomorphism.
This proves the assertion about the full lift scheme; in
particular the earlier generic reduction has removed no
infinitesimal part on the exact-height locus.

Every \'etale basis at level $j$ lifts to a basis at level $l$ after a
geometric field extension.  The transition map on generic coordinate
algebras is therefore injective.  Its restriction $B_j\to B_l$ is
injective by generic density and finite since $B_l$ is finite over
$A_m$.  Finally
$t_{i,j}=[p^{l-j}]_{\mathcal G}(t_{i,l})$ factors through
$t_{i,l}^{d^{l-j}}$.  Substituting the expressions for the base
parameters in $C_l$ gives the last assertion.
\end{proof}

The finite-level calculation has retained both the \'etale basis and
every infinitesimal lift.  To pass to the completed tower, we now give
explicit inverse coordinates on the universal cover.

\begin{lemma}\label{h3:geom:universal-cover}
Let $(R,R^+)$ be an affinoid perfectoid $k$-algebra and specialize the
parameters of $A_m$ to $R^{\circ\circ}$.  For the resulting formal law
$\mathcal G_a$, there is a natural additive isomorphism
\[
 H_3(R)\ \xrightarrow{\ \Phi_a\ }\
 \varprojlim_{[p]}\mathcal G_a(R^{\circ\circ}).
\]
Writing $P_a=[p]_{\mathcal G_a}$, its coordinates and inverse are
\begin{equation}\label{h3:geom:cover-formulas}
 \Phi_a(z)_j=\lim_{n\to\infty}
 P_a^{\circ n}\bigl(z^{1/p^{3(j+n)}}\bigr),\qquad
 \Phi_a^{-1}((t_j))=\lim_{j\to\infty}t_j^{p^{3j}}.
\end{equation}
These formulas commute with change of coefficients and with changes
of deformation framing.  Over an algebraically closed perfectoid
field, the kernel of the zeroth projection is the integral Tate
module of the \'etale quotient of $\mathcal G_a$.
\end{lemma}

\begin{proof}
Choose a power-multiplicative spectral norm and $\lambda<1$ bounding
the specialized parameters.  All coefficients of
$\mathcal G_a-\Gamma_3$ and of $P_a(T)-T^q$, $q=p^3$, have norm at
most $\lambda$.  Integral formal series are $1$-Lipschitz on the open
unit ball.  Since $P_a$ factors through $T^d$, $d=p^m$, we have
\[
 |P_a^{\circ n}(v)-P_a^{\circ n}(w)|\leq |v-w|^{d^n}.
\]
Consecutive approximations in the first formula differ by at most
$\lambda^{d^n}$.  Their limits exist, are topologically nilpotent,
and satisfy
\[
 |\Phi_a(z)_j-z^{1/q^j}|\leq\lambda,
 \qquad P_a(\Phi_a(z)_{j+1})=\Phi_a(z)_j.
\]
For a compatible sequence $(t_j)$, consecutive terms in the second
formula differ by at most $\lambda^{q^j}$.  Its limit $z$ satisfies
$|z-t_j^{q^j}|\leq\lambda^{q^j}$.  Substitution in the first formula
and the displayed contraction show in both directions that the
formulas are inverse.  This reasoning never requires a common
strict upper bound on all the $|t_j|$.

The coefficients of $\Gamma_3$ are fixed by the $q$th-power map.
The difference between the two addition laws is at most $\lambda$;
after $n$ applications of $P_a$ it is at most $\lambda^{d^n}$.
Taking the limit proves additivity.  If $g$ changes the deformation
framing, its evaluated coordinate change differs from its Honda
reduction by coefficients of norm at most $\lambda$, after increasing
$\lambda$ if necessary.  The same contraction proves
$\Phi_{ga}(gz)=g\Phi_a(z)$.  Uniform convergence on smaller balls
proves continuity and compatibility with maps of perfectoid
algebras.

Over an algebraically closed field the points with $t_0=0$ are exactly
the compatible $p^j$-torsion points.  A finite connected group over a
perfect field has just its identity as a geometric point, so these
points form the Tate module of the \'etale quotient.  The zeroth
projection is surjective: preparation applied to $[p](T)-w$ gives a
monic distinguished polynomial, each of whose roots is small when
$w$ is small.  Successively choosing roots gives a compatible
sequence with prescribed $t_0$.
\end{proof}

Write $(H_3^r)_{\mathrm{ind}}$ for the locus of tuples whose
images on every geometric fiber are $\Q_p$-linearly independent
for the Honda vector-space structure on $H_3$.

\begin{theorem}[Coordinate comparison]\label{h3:thm:coordinate-comparison}
There is a natural $P_3\times\GL_r(\Z_p)$-equivariant isomorphism
\[
 \mathcal T_m\ \simeq\ (H_3^r)_{\mathrm{ind}}.
\]
The map takes a compatible basis of generator lifts to
$z_i=\lim_l t_{i,l}^{p^{3l}}$.  It identifies the relevant completed
function algebras with their spectral norms on exhausted closed
smaller balls.  The assertion holds on the relative perfectoid site
and remains valid after adjoining arbitrary profinite parameters.
\end{theorem}

\begin{proof}
First retain the reduced finite-level closures, including their
higher-height boundary.  Perfection identifies a point of the
\'etale quotient with its unique connected generator lift, so their
perfected inverse tower is the tower from
\Cref{h3:geom:regular-levels}.

We check the complete algebras after extension to an algebraically
closed perfectoid field $C^\flat$; these field extensions form
v-covers, and all constructions are defined over $k$.
Put $q=p^3$, $d=p^m$ and $y_{i,l}=t_{i,l}^{q^l}$.
For $0<\rho<1$ in the value group, let $X_l(\rho)$ be the closed
polydisc $\max_i|y_{i,l}|\leq\rho$.  If $\lambda$ is the maximum
norm of the height parameters at a point of this polydisc, then
\Cref{h3:geom:regular-levels} at level one gives
\[
 \lambda\leq\|t_{\cdot,1}\|^d,\qquad
 \|t_{\cdot,1}\|\leq\max(\rho^{1/q},\lambda).
\]
The second inequality follows by finite telescoping of the normalized
division coordinates.  If $\lambda>\rho^{1/q}$, the two inequalities
would give $0<\lambda\leq\lambda^d<\lambda$.  Therefore
\begin{equation}\label{h3:geom:uniform-radius}
 \lambda\leq\rho^{d/q}.
\end{equation}
For consecutive levels it follows that
\[
 \|y_{\cdot,l+1}-y_{\cdot,l}\|
       \leq\lambda^{q^l}
       \leq\rho^{d q^{l-1}}<\rho\qquad(l\geq1).
\]
The finite transition maps take $X_{l+1}(\rho)$ to $X_l(\rho)$.
Conversely every point of $X_l(\rho)$ has a lift, by finiteness,
injectivity and lying-over, and the last strict inequality forces
every lift to remain in $X_{l+1}(\rho)$.  They are thus surjective
maps of these closed polydiscs.  Pullback is isometric for their
spectral norms, before and after completed perfection.

In the completed direct limit of the perfected Banach algebras, the
$y_{i,l}$ converge to elements $z_i$.  The homomorphism from the
perfected Gauss algebra of a polydisc of radius $\rho$ sending its
coordinates to $z_i$ is isometric.  Indeed, for any polynomial in
finitely many fractional powers, evaluation on $y_{\cdot,l}$ has
exactly its Gauss norm: $t_{\cdot,l}$ are free polydisc coordinates
and $q^l$ is a power of $p$.  These evaluations converge to evaluation
on $z$.  A nonzero limiting polynomial therefore has precisely that
norm.  Completing preserves the isometry.

Its image is all the completed direct limit.  For $l\geq j$, write
\[
 t_{i,j}=F_{i,j,l}
 (t_{1,l}^{d^{l-j}},\ldots,t_{r,l}^{d^{l-j}}),
 \qquad F_{i,j,l}\in k[[T_1,\ldots,T_r]].
\]
The formulas in \Cref{h3:geom:universal-cover} show that replacing
$t_{i,l}$ by $z_i^{1/q^l}$ changes this expression by at most
\[
 \lambda^{d^{l-j}}
 \leq \rho^{(d/q)d^{l-j}},
\]
which tends to zero.  Its substituted series converges on the
$z$-polydisc.  Thus the closed image contains every $t_{i,j}$ and,
by applying roots to the same estimates, all its $p$-power roots.
These elements generate the completed direct limit.  Varying
$\rho$ proves the comparison for the open balls.

It remains to identify the open locus.  For
$a=(a_1,\ldots,a_r)\in\Z_p^r$, form the compatible points
$t_{a,j}=\sum_{\mathcal G}[a_i]t_{i,j}$.
The normalized limit is $\sum_{\Gamma_3}[a_i]z_i$ by
\Cref{h3:geom:universal-cover}.  If it vanishes, then
$|t_{a,l}|\leq\lambda$ at every level, and hence
\[
 |t_{a,j}|=|[p^{l-j}](t_{a,l})|
       \leq\lambda^{d^{l-j}}\longrightarrow0.
\]
Thus every $t_{a,j}$ vanishes.  Conversely a relation in the Tate
module gives the same relation among the $z_i$.
At height exactly $m$, \eqref{h3:geom:basis-divisor} gives a basis of
the \'etale quotient modulo $p$, and the compatible points give an
injective map $\Z_p^r$ into its Tate module.  At greater height that
Tate module has rank less than $r$, so there is a relation.
Multiplying a rational relation by a power of $p$ reduces to the
integral assertion.  Exact height $m$ is therefore exactly the
independent locus.  This identifies the corresponding open
subsheaves; it is not only a bijection on their field-valued points.

Equivariance is part of \Cref{h3:geom:universal-cover}.  All the
estimates are spectral estimates on affinoid perfectoid algebras.
They apply to continuous functions with a profinite parameter by
taking their uniform norm, so the same isomorphisms retain those
parameters.
\end{proof}

\subsection{Sections and extension spaces on the curve}

We recall precisely the curve results we need.  For a perfectoid
test object $S$, write $X_S$ for its relative Fargues--Fontaine
curve; this notation does not posit a morphism of schemes
$X_S\to S$.  Relative sections and relative cohomology below mean
the corresponding sheaves on the perfectoid site.

\begin{lemma}\label{h3:geom:curve-inputs}
There are natural isomorphisms of additive sheaves
\[
 H_s\simeq\bigl(S\longmapsto H^0(X_S,V_s)\bigr).
\]
The automorphism sheaf of $V_s$ is the locally profinite group
$D_s^{\times}$.  A family fiberwise isomorphic to $V_s$ is locally
of that form for the pro-\'etale topology.  If $\tau$ denotes
the morphism of v-sites associated with $X_S$, then
\[
 R\tau_*\mathcal O^a=\underline{\Q_p}^{\,a},\qquad
 R\tau_*E=\tau_*E[0]
\]
for every family $E$ fiberwise isomorphic to a $V_s$.
Relative sections are fully faithful on $\Q_p$-linear sheaf
morphisms between these bundles and slope-zero bundles.
All assertions commute with base change and v-descent.
\end{lemma}

\begin{proof}
For the unramified degree-$s$ extension of $\Q_p$, the
Lubin--Tate universal cover is the section sheaf of $\mathcal O(1)$
on the curve with that coefficient field
\cite[Proposition II.2.2]{FarguesScholze}.  Unramified pushforward gives $V_s$.
Reduction of the universal cover and tilting identify it with the
Honda universal cover; one can check reduction directly by the
contracting formulas of \Cref{h3:geom:universal-cover}, using
$|[p](v)-[p](w)|\leq\max(|p|\,|v-w|,|v-w|^2)$ before reduction.
This proves the first assertion with its addition and endomorphisms.
For the local-form assertion, the vector-bundle equivalence and
pure-module theorem identify a pure family of slope $c/d$ in lowest terms,
after adjoining $\F_{p^d}$, with a $(c,d)$-$\Q_p$ local
system \cite[Theorems 8.7.8 and 8.5.12]{KedlayaLiu}.  Its semilinear
Frobenius supplies the action of the constant cyclic division
algebra \cite[Definition 8.5.7]{KedlayaLiu}.  These local systems
trivialize in the pro-\'etale topology
\cite[Lemma 9.1.11]{KedlayaLiu}; a module basis over that division
algebra gives the specified basic local form.  Its automorphism
sheaf is therefore the locally profinite $D_s^\times$.
The slope-zero case is also
\cite[Corollary 8.7.10]{KedlayaLiu}.
To retain the ground field $k$, observe that every geometric
Honda endomorphism $f(T)=\sum a_iT^i$ commutes with
$[p](T)=T^{p^s}$.  Hence $a_i^{p^s}=a_i$, so all integral
endomorphisms, and then all of $D_s$, are defined over
$\F_{p^s}\subset k$.  The actual Honda bundle over $k$
therefore defines a map from $B\underline{D_s^\times}$ to
the rank-$s$, degree-one basic locus.  After extending $k$
to its algebraic closure, the preceding local-form theorem
and its identification of automorphisms make this an
equivalence.  Equivalences of v-stacks descend along that
cover, so the frame torsors already have descent over $k$.

For the required cohomology, on the curve with unramified
coefficient field of degree $s$, the bundle $\mathcal O(1)$
is the positive geometric twist of the globally \'etale
bundle $\mathcal O$.  Thus its $H^1$ vanishes on every
affinoid perfectoid test object
\cite[Corollary 8.8.7 and Theorem 8.7.13]{KedlayaLiu}.
Unramified finite pushforward gives the same assertion for
the standard $V_s$.  A family of that basic type first
standardizes pro-\'etale locally, so its $H^1$ sheaf vanishes.
For $\mathcal O^a$, its $H^0$ sheaf is
$\underline{\Q_p}^{\,a}$.  An extension of $\mathcal O$
by $\mathcal O^a$ is pointwise of slope zero; the local-system
equivalence makes it split pro-\'etale locally.  This proves
vanishing of its $H^1$ sheaf, without asserting vanishing
over an unrefined affinoid.  Higher cohomology on affinoid
tests is zero by \cite[Theorem 8.7.13]{KedlayaLiu}.

The needed full faithfulness is relative: for any small v-stack
$S$, the functor
\[
 R\tau_*:\operatorname{Perf}(X_S)\longrightarrow D(S_v,\Q_p)
\]
is fully faithful \cite[Corollary 3.11]{AnschuetzLeBrasFourier}.  For the
basic positive and slope-zero bundles just considered, the
preceding vanishing identifies $R\tau_*E$ with its section
sheaf in degree zero.  Taking
degree-zero morphisms therefore identifies bundle maps with
$\Q_p$-linear maps of section sheaves.  This construction is
compatible with every base change on $S$, so it also applies to
relative forms of the basic bundles.
\end{proof}

Put
\[
 N_m=\bigl(S\longmapsto H^1(X_S,V_m^{\vee})\bigr),\qquad
 \mathcal E_m=(N_m^r)_{\mathrm{ind}}.
\]
The extension interpretation uses sheafified relative $H^1$ and
classifies extensions with both ends framed.  Independence in
$N_m^r$ means independence over $\Q_p$ on every geometric fiber.

\begin{lemma}\label{h3:geom:independent-extensions}
Independent sections of $V_3$ define a subbundle
$\mathcal O^r\hookrightarrow V_3$ with quotient fiberwise $V_m$.
Conversely, an extension of $V_m$ by $\mathcal O^r$ has middle
bundle fiberwise $V_3$ if and only if its class lies in
$\mathcal E_m$.
\end{lemma}

\begin{proof}
First work on an absolute curve.  Let $M$ be the saturation of the
generic span of $r$ independent sections and put $b=\hthreerank M$.
A wedge of $b$ generically independent sections gives a nonzero
section of $\det M$, so $\deg M\geq0$.  Stability of $V_3$ gives
$\deg M\leq b/3<1$.  Since degree is integral, it is zero.
The wedge has no zeros, because a nonzero effective divisor has
positive degree.  The selected sections trivialize $M$.
Since $H^0(X,\mathcal O)=\Q_p$, independence of the original
sections forces $b=r$; hence their map is a subbundle map.
All slopes of its quotient are positive: a quotient of nonpositive
slope would give a nonzero map from $V_3$ to a bundle of slope at
most zero.  A positive bundle of rank $m\leq2$ and degree one is
$V_m$, by the classification of bundles and integrality of degrees
\cite[Theorem II.2.14]{FarguesScholze}.

An extension of $V_m$ by $\mathcal O^r$ has no negative-slope
quotient, since neither end has a nonzero map to a negative bundle.
A nonnegative bundle of rank three and degree one is either $V_3$
or has a quotient $\mathcal O$.  Such a quotient restricts
nontrivially to $\mathcal O^r$, because $\Hom(V_m,\mathcal O)=0$.
The resulting nonzero linear functional on $\Q_p^r$ kills the
extension class.  Conversely, if a nonzero functional kills that
class, its pushout extension splits and gives a quotient
$\mathcal O$.  This proves the equivalence.  Fiberwise
surjectivity of a morphism of vector bundles is a relative
subbundle condition, and \Cref{h3:geom:curve-inputs} supplies local
standard forms and the relative extension sheaf, so the argument
gives the asserted relative loci.
\end{proof}

\subsection{Integral frames and determinant normalization}

The extension locus is now identified as a rational bundle problem.
The next step retains the integral Tate basis and connected marking,
because their determinant lattices determine the group actions in the
coefficient calculation.

A $P_s$-structured form of $V_s$ is a form whose $D_s^\times$-torsor
of frames has been reduced to $P_s$.  Since
\[
 P_s=\hthreeNrd^{-1}(\Z_p^\times),
\]
this is determinant-lattice data.  It is not an integral lattice
in a rank-$s$ slope-zero local system.  We choose determinant
identifications between the $V_s$ over $k$.  They exist there:
the determinant of the cyclic Honda isocrystal differs from the
degree-one rank-one isocrystal by its cyclic permutation sign
$(-1)^{s-1}$.  For even $s$, choose $a\in\F_{p^2}^\times$
with $a^p=-a$; its Teichm\"uller lift changes the Frobenius
multiplier by $-1$.  Thus $\F_{p^2}\subset k$ removes the sign.
We specify a
common valuation normalization in the following proof.

For the height-two statement below, let
$\operatorname{Surj}(V_3,V_2)$ be the relative sheaf of bundle maps
$V_3\to V_2$ that are surjective on every geometric fiber.  Write
$N_2^*=N_2\setminus\{0\}$ for the locus of extension classes that
are nonzero on every geometric fiber.  The zero class represents the
split extension, so this is exactly the locus where the extension of
$V_2$ by $\mathcal O$ is nonsplit on every geometric fiber.

\begin{proposition}[Integral frame comparison]\label{h3:prop:integral-frames}
Put $\mathcal Y_m=[\mathcal X_m/P_3]$.  There is an equivalence
of v-stacks over $k$
\begin{equation}\label{h3:geom:frame-stack}
 \mathcal Y_m\simeq
 [\mathcal E_m/(\GL_r(\Z_p)\times P_m)].
\end{equation}
It carries an actual universal sequence
\begin{equation}\label{h3:geom:universal-sequence}
 0\longrightarrow\mathcal L_r\otimes_{\Z_p}\mathcal O
 \longrightarrow\mathcal V_3
 \longrightarrow\mathcal V_m\longrightarrow0,
\end{equation}
where $\mathcal L_r$ is the integral \'etale Tate local system and
$\mathcal V_s$ are $P_s$-structured forms.  The two presentations
give the classifying maps to $BP_3$, $B\GL_r(\Z_p)$ and $BP_m$,
and their determinant characters satisfy
\begin{equation}\label{h3:geom:determinant-product}
 \hthreeNrd_3=\det_{\GL_r}\,\hthreeNrd_m.
\end{equation}
With the convention $(A,b)e=A_*((b^{-1})^*e)$ for extension classes,
the kernel scalar $u$ acts on $N_m^r$ by $u$ and the quotient scalar
$u$ by $u^{-1}$.

For $m=2$, let $\widetilde{\mathcal X}_2$ be the full
connected-marking tower.  It admits a normalized integral
\'etale generator, fixed by $H=\ker(\hthreeNrd:P_3\to\Z_p^\times)$.
There are compatible equivalences
\begin{equation}\label{h3:geom:normalized-tower}
 \widetilde{\mathcal X}_2\simeq\operatorname{Surj}(V_3,V_2),
 \qquad [\widetilde{\mathcal X}_2/H]\simeq N_2^*.
\end{equation}
The full residual norm quotient acts in the second equivalence by
scalar units, in the kernel convention just specified.  No section
of $P_3\to\Z_p^\times$ is involved.
\end{proposition}

\begin{proof}
Put $R_m=A_m[u_m^{-1}]$.  Apply
\cite[Theorem~1 and its proof, pp.~1--4]{LurieFormalGroups2010}
to the formal completion $\widehat{\mathcal G}_{R_m}$ of $\mathcal G$
over $R_m$ and to $\Gamma_m$, both of exact height $m$.  That theorem defines the
isomorphism ring using every coefficient of the full identity
\[
 F_{\Gamma_m}(f(X),f(Y))=f(F_{\mathcal G}(X,Y))
\]
and expresses it as a filtered union of injective finite \'etale
$R_m$-algebras.  Pullback to $\mathcal X_m$ therefore gives the
pro-\'etale tower of genuine connected markings; any finite set of
Taylor coefficients is defined at one finite stage.  Composition
makes it a torsor under the automorphism sheaf of $\Gamma_m$.
The Honda coefficient calculation in \Cref{h3:geom:curve-inputs}
makes the latter the constant profinite group
$P_m=\Aut_k(\Gamma_m)$: its finite \'etale coefficient stages have
all geometric coefficients in $\F_{p^m}\subset k$.

For the needed bound, pull a genuine marking
$f(T)=\sum_{n\geq1}a_nT^n:\widehat{\mathcal G}_R
\xrightarrow{\sim}\Gamma_{m,R}$ to a characteristic-$p$
affinoid perfectoid test algebra $R$ on this tower, with its
power-multiplicative spectral norm.  Put $q=p^m$.  The factorization
of the $p$-series used in \Cref{h3:geom:regular-levels} reads
\[
 [p]_{\mathcal G}(T)=P(T^q),\qquad
 P(S)=u_mS+\sum_{j\geq2}c_jS^j,
\]
where $u_m\in R^\times\cap R^{\circ\circ}$ and $c_j\in R^\circ$;
these bounds come from the integral deformation coefficients.
The necessary identity $f([p]_{\mathcal G}(T))=f(T)^q$ gives
\[
 a_1^{q-1}=u_m,\qquad
 a_n^q-u_m^na_n=\sum_{1\leq i<n}a_i[S^n]P(S)^i
 \quad(n\geq2).
\]
Here $[S^n]$ denotes coefficient extraction.  First
$\|a_1\|=\|u_m\|^{1/(q-1)}<1$.  If the earlier coefficients have
norm at most one and $A=\|a_n\|>1$, the second equation gives
$A^q\leq\max\{\|u_m\|^nA,1\}\leq A$, a contradiction.
Thus every $a_n$ belongs to $R^\circ$.  The series $f(z)$ converges
in $R^{\circ\circ}$, uniformly on each closed ball of radius less
than one.  The full formal-law identity then evaluates on small
points and proves additivity; evaluation commutes with maps of
perfectoid test algebras.  Compose it with the zeroth projection
in \Cref{h3:geom:universal-cover}.  The resulting additive sheaf map
$H_3\to H_m$ is $\Q_p$-linear: it commutes with integral
formal-group multiplications and with inverse multiplication
by $p$.  It kills the $\Q_p$-span of the Tate basis: it kills
all $p$-power torsion because $[p]$ is injective on $H_m$.
By \Cref{h3:geom:independent-extensions,h3:geom:curve-inputs}, it
therefore factors through the section sheaf of the quotient bundle.
By the relative full faithfulness in \Cref{h3:geom:curve-inputs},
this is the section map of a unique relative bundle morphism
from the quotient to $V_m$.  On every geometric fiber the
section map is nonzero: restrict near the origin, where the
formal isomorphism has nonzero linear coefficient.  After
identifying that fiber of the quotient with $V_m$, the bundle
map is a nonzero element of the division algebra $D_m$, hence
invertible.  A bundle morphism invertible on every geometric
fiber is an isomorphism.  Uniqueness under relative full
faithfulness makes these frames compatible with base change
and with the group actions.

Connected markings form a $P_m$-torsor, and the constructed frames
change by that same maximal compact subgroup.  The relative frame
theorem shows that their only possible discrepancy with the
determinant-lattice reduction of the quotient is its integer
valuation.  This integer is locally constant on the Tate-basis
space and invariant under $P_3$, $P_m$ and $\GL_r(\Z_p)$:
all their determinant changes are units.  For $m=2$, that
space is $H_3\setminus\{0\}$ by
\Cref{h3:thm:coordinate-comparison}.  It is geometrically
connected: after extending the base field it is exhausted by
intersecting connected closed annuli, and perfection preserves
their underlying topology.  The integer is therefore constant,
and we choose the fixed determinant identification to absorb it.
For $m=1$, let the locally constant valuation discrepancy be
$\nu_{\mathrm{val}}$.  Multiplication of the quotient frame by
$p^{-\nu_{\mathrm{val}}}$ corrects it.  Since
$D_1^\times=\Q_p^\times$ is central and $\nu_{\mathrm{val}}$
is invariant under all three actions, this correction is
equivariant and glues without any connectedness assumption.
These choices give the common determinant-lattice normalization.

Now mark both the connected group and the integral Tate basis.
By \Cref{h3:thm:coordinate-comparison} this tower is the sheaf of
exact sequences
\[
 0\to\mathcal O^r\to V_3\to V_m\to0
\]
whose induced determinant identification is a unit multiple of the
chosen one.  The left frame changes are exactly
$\GL_r(\Z_p)\times P_m$.  Forgetting the middle frame instead
gives $\mathcal E_m$, and its middle frames preserving that
determinant lattice form exactly a $P_3$-torsor.  This pair of
torsor presentations proves \eqref{h3:geom:frame-stack}.
Descent of the displayed sequence proves
\eqref{h3:geom:universal-sequence}; taking its determinants proves
\eqref{h3:geom:determinant-product}.  Pushout at the kernel and
pullback at the quotient give the stated action on extension
classes.  This construction uses actual frames throughout, so the
maps to the classifying stacks are part of the equivalence.

For $r=1$, the fixed determinant identification and quotient frame
identify the kernel line with $\mathcal O$.  On the integral
Tate-basis torsor its determinant multiplier is a unit.  Dividing
the generator by that multiplier gives the unique generator with
multiplier one.  Uniqueness glues this normalization and shows
that $g\in P_3$ changes it by $\hthreeNrd(g)$, with the inverse if all
frame conventions are reversed.  It is therefore fixed by $H$.
Forgetting this now determined generator identifies the connected
marking tower with all surjections $V_3\to V_2$.  Indeed, a
surjection $q$ determines its unique kernel generator $i_q$
of determinant multiplier one.  Coordinate comparison on the
entire nonzero Honda ball reconstructs a deformation and a
primitive integral Tate generator from $i_q$, at every valuation.
Scaling $i_q$ by a power of $p$ can change that deformation;
the comparison always returns a basis for the reconstructed one.
Choose a connected marking locally and let $q_h$ be its quotient
frame.  The constant discrepancy already normalized above
ensures that the determinant multiplier of $(i_q,q_h)$ is a
unit.  Write $q=dq_h$ with $d\in D_2^\times$.  Since
$(i_q,q)$ has multiplier one, $v_p(\hthreeNrd d)=0$, so $d\in P_2$.
Thus the required change of marking is an integral Honda
automorphism and yields a unique connected formal isomorphism.
Allowing the middle
bundle to vary while keeping its determinant frame gives precisely
the nonsplit extensions of $V_2$ by $\mathcal O$, namely $N_2^*$.
Changing the middle determinant scales the kernel generator
by its full reduced norm.  This norm maps $P_3$ onto
$\Z_p^\times$, as is seen on the units of the unramified
cubic subfield of $D_3$.
This proves \eqref{h3:geom:normalized-tower}, with the full norm
quotient action.
\end{proof}

\subsection{Algebraic markings and infinitesimal splitting torsors}

The stack comparison supplies normalized markings after completed
perfection.  Ordinary cooperations require those markings at finite
algebraic stages.  We first prove a uniform descent statement that also
allows the number of field factors to grow, and then identify the
resulting finite flat root torsors.

Set $K=k((y))$.  Let $L$ be the algebraic union of finite
connected-marking algebras over $K$, and write
\[
 L^{\mathrm c}=\widehat{L^{\mathrm{perf}}}.
\]
Finite products of complete fields are allowed in this notation;
all assertions below are compatible with their idempotents.  The
commuting actions of $P_3$ and $P_2$ are the actions on the
deformation and on its connected marking respectively.

\begin{lemma}\label{h3:geom:uniform-etale-descent}
Let $L=\colim_i B_i$ be an injective union of finite \'etale
$K$-algebras, with their maximum valuation norms, and put
$C=\widehat{L^{\mathrm{perf}}}$.  For every finite \'etale
$L$-algebra $E$, the map
\[
 \Hom_{L\text{-alg}}(E,L)\longrightarrow
 \Hom_{L\text{-alg}}(E,C)
\]
is bijective.  Each section on the right is defined at one
finite separable stage $B_j$, even if the numbers of field
factors in the $B_i$ are unbounded.
\end{lemma}

\begin{proof}
Every $B_i$ and every $B_i^{1/p^n}$ is a finite product of
complete valued fields.  First let $e\in C$ be an idempotent.
Choose $a\in B_i^{1/p^n}$ with $\|a-e\|<1$.  Then
$\|a\|\leq1$ and $\|a^2-a\|<1$, and $2a-1$ is a unit
whose inverse has norm at most one.  Hensel's lemma in this
one complete finite product gives an idempotent $e_i$ with
$\|e_i-a\|<1$.  Two idempotents at distance less than one
are equal: each of $e(1-e_i)$ and $e_i(1-e)$ is an idempotent
of norm less than one and hence vanishes.  Thus $e=e_i$.
Purely inseparable extensions create no idempotents, so
$e_i\in B_i$.  Any finite list of idempotents descends to
a common stage.

The algebra $E$ descends to a finite \'etale algebra at some
$B_i$.  After a finite idempotent partition it admits a
monogenic \'etale presentation: each component field of
$B_i$ is infinite, so a primitive element may be chosen for
its finite product of separable extensions.  The preceding
paragraph descends the idempotents chosen by a section to
$C$.  It remains to descend its finitely many roots.

Suppose $f(z)=0$ and $f'(z)$ is invertible in $C$, with
$f\in B_i[T]$.  Approximate both $z$ and $f'(z)^{-1}$
simultaneously by $a,c\in B_j^{1/p^n}$, for one $j\geq i$
and one $n$.  We may require $\|1-cf'(a)\|<1$.
The geometric series in this complete finite product makes
$f'(a)$ invertible with inverse norm at most $\|c\|$.
Put $C_0=\max(1,\|c\|)$, and let $M\geq1$ bound the
quadratic Taylor remainder of $f$ on the unit ball about
$a$.  By making the approximations closer, require
$\|f(a)\|C_0^2M<1$ and
$\|f(a)\|C_0<\min(1,(C_0M)^{-1})$.
Newton iteration stays in $B_j^{1/p^n}$, its inverse
derivatives stay bounded by $C_0$, and its residual norms
are bounded successively by
$M C_0^2\|f(a)\|^2$ and their iterates.
It converges there to a root $b$ with
$\|b-a\|\leq C_0\|f(a)\|$.
Taking the original approximation inside the same simple-root
ball makes $b=z$: the divided difference
$(f(b)-f(z))/(b-z)$ is a unit throughout that ball.
The element $b$ is separable over each component of $B_j$
and lies in its purely inseparable extension, so $b\in B_j$.
The finitely many roots use a common $j$.  Injectivity follows
from $L\hookrightarrow C$, proving the assertion.
\end{proof}

\begin{lemma}[Descent of the normalized marking]\label{h3:lem:marking-descent}
The normalized \'etale generator of
\Cref{h3:prop:integral-frames} descends to a compatible integral
marking over $L$.  It is $H$-invariant, and the full $P_3$ action
changes it by the constant unit $\hthreeNrd(g)$, up to the common inverse
convention.  At level $l$, the torsor of its lifts to
$\mathcal G[p^l]$ has coordinate algebra
\begin{equation}\label{h3:geom:height-two-roots}
 L^{1/p^{2l}}.
\end{equation}
The connected marking makes this a torsor for
$\Gamma_2[p^l]_L$.  The group $H$ acts trivially on the translation
group, while $P_2$ acts by constant Honda automorphisms.

For a finite connected-marking level $B=L^\Pi$, with $\Pi\subset P_2$
open, put $B^{\mathrm b}=\widehat{B^{\mathrm{perf}}}$; the superscript
here denotes completed perfection, not tilting.  There is a
pro-\'etale $\Pi$-torsor
\begin{equation}\label{h3:geom:completed-level-torsor}
 \operatorname{Spa}(L^{\mathrm c})\longrightarrow
 \operatorname{Spa}(B^{\mathrm b}),
\end{equation}
and consequently
\begin{equation}\label{h3:geom:finite-level-stack}
 [\operatorname{Spa}(B^{\mathrm b})/H]\simeq[N_2^*/\Pi].
\end{equation}
These equivalences are natural under level changes and after
extension to an algebraically closed perfectoid base.
\end{lemma}

\begin{proof}
For each $l$ the normalized generator is a point of
\[
 \operatorname{Isom}_L
 ((\Z/p^l)_L,\mathcal G^{\mathrm{et}}[p^l]_L).
\]
By \Cref{h3:geom:uniform-etale-descent}, it descends to $L$ at
one finite separable stage.  Compatibility and equivariance
descend through the injective coefficient extension, so the
descended points give the asserted integral marking.  This
argument descends a finite \'etale marking, not a splitting of
the connected--\'etale extension.

The fiber over this marked generator in $\mathcal G[p^l]$ is a
torsor for its connected kernel.  Over the \'etale generator
cover, \Cref{h3:geom:regular-levels} identifies its two fields as
\[
 k((t_l^{p^{2l}}))\subset k((t_l)).
\]
They have purely inseparable degree $p^{2l}$.  The marked
generator supplies a map from the separable field on the left
to $L$.  Root extensions commute with separable scalar
extension, so the torsor algebra after this base change is
exactly \eqref{h3:geom:height-two-roots}; no degree is lost.
The connected marking identifies its group with
$\Gamma_2[p^l]$.  The equivariance already proved fixes both
markings under $H$ and changes them by constant automorphisms
under $P_2$, giving the stated actions on translations.

To prove the torsor assertion, choose a cofinal system of
normal open subgroups $\Pi'\subset\Pi$.  The algebra
$B_{\Pi'}=L^{\Pi'}$ is a finite \'etale $\Pi/\Pi'$-torsor over
$B=L^\Pi$.
Perfection commutes with finite \'etale base change, and
completion commutes with finite modules.  Consequently
\[
 A_{\Pi'}:=B_{\Pi'}\otimes_B B^{\mathrm b}
       \simeq\widehat{B_{\Pi'}^{\mathrm{perf}}}
\]
is again a finite \'etale $\Pi/\Pi'$-torsor.  Their transition
maps are surjective.  The affinoid perfectoid inverse limit
is $\operatorname{Spa}(L^{\mathrm c})$: its algebra is
the uniform completed union of the $A_{\Pi'}$, equal to
$\widehat{L^{\mathrm{perf}}}$.
Passing the finite torsor identities to inverse limits gives
\[
 \operatorname{Spa}(L^{\mathrm c})
 \mathbin{\times}_{\operatorname{Spa}(B^{\mathrm b})}
 \operatorname{Spa}(L^{\mathrm c})
 \simeq\underline\Pi\times\operatorname{Spa}(L^{\mathrm c}).
\]
Each geometric fiber is a nonempty inverse limit of finite
nonempty sets with surjective transition maps.  Thus the
map is surjective and is the pro-\'etale $\Pi$-torsor in
\eqref{h3:geom:completed-level-torsor}.  These arguments permit
arbitrarily many factors at later finite stages.
Finally $H$ and $\Pi$ commute,
so quotienting this torsor and applying
\eqref{h3:geom:normalized-tower} gives
\eqref{h3:geom:finite-level-stack}.  No cohomological Galois-descent
assertion is used in this deduction.
\end{proof}

We fix the periodicity convention needed for the coefficient
arguments.  The invariant cotangent line is
$\omega=\pi_2E_3$; if $|u|=-2$, its generator is $U=u^{-1}$.
In particular $v_1=xU^{p-1}$ modulo $p$.

\begin{lemma}[Invariant differential]\label{h3:lem:invariant-differential}
There is an equivariant isomorphism
\begin{equation}\label{h3:geom:canonical-lines}
 \bigwedge^2\Omega^1_{A/k,\hthreects}\simeq\omega^3[\det],
 \qquad
 \Omega^1_{K/k}\simeq\omega_K^{p+2}[\det].
\end{equation}
Here $\det$ is the reduced norm character of $P_3$.
The connected marking over $L$ supplies a generator $U_0$ of
$\omega_L$ fixed by $P_3$.  Consequently
\eqref{h3:geom:canonical-lines} gives a nonzero differential
$\eta\in\Omega^1_{L/k}$ fixed by $H$, defined at a finite
separable marking level.  Changes of the connected marking
act on these frames by their constant cotangent characters.
\end{lemma}

\begin{proof}
The integral canonical-line theorem gives the first isomorphism
\cite[Theorem 20]{GHdual}.  In that theorem $\omega=E_2$, and
the determinant is the reduced norm, so its line and action
conventions are the ones just fixed.  If $g(U)=q_gU$, comparison
of the first nonzero graded $p$-series coefficient gives
$g(x)=q_g^{-(p-1)}x$ modulo $p$.  Equivalently the graded
coefficient $v_1=xU^{p-1}$ is invariant there.  Hence
$(x)/(x^2)\simeq\omega^{-(p-1)}$ on the height-two slice.
Taking determinants in its conormal sequence yields
\[
 \Omega^1_{k[[y]]/k,\hthreects}
 \simeq\omega^{3+(p-1)}[\det]
 =\omega^{p+2}[\det].
\]
Localizing at $y$ proves the second isomorphism.

There is no change of differential theory at the generic point.
Every power series in one variable over a perfect field can
be grouped uniquely by its exponent modulo $p$; thus $y$ is
a $p$-basis of $K$ and $\Omega^1_{K/k}=K\,dy$.
Separable extension gives
$\Omega^1_{L/k}=L\otimes_K\Omega^1_{K/k}$.
Transport a nonzero invariant differential on $\Gamma_2$
along the tautological connected marking.  This gives $U_0$.
The $P_3$ action transports the marking, so it fixes this
frame.  The image of $U_0^{p+2}\otimes1_{\det}$ under
\eqref{h3:geom:canonical-lines} is $\eta$.  Its remaining
character is the reduced norm, trivial on $H$.  Its
coefficient and that of $U_0$ occur at a finite separable
stage, since each belongs to the algebraic union $L$.
The action of $P_2$ is the constant change of Honda marking,
and differentiation gives the asserted constant characters.
\end{proof}

\begin{lemma}[The ordinary splitting torsor]\label{h3:lem:ordinary-torsor}
Let $B_0=A[x^{-1}]$.  The ordinary connected--\'etale sequence
of $\mathcal G$ over $B_0$ is
\[
 0\longrightarrow\mathcal C\longrightarrow\mathcal G
 \longrightarrow\mathcal E\longrightarrow0,
\]
where $\mathcal C$ has height one and is of multiplicative
type, and $\mathcal E$ is \'etale of height two.
For every $l\geq1$, the scheme of splittings of its
$p^l$-torsion sequence is canonically
\begin{equation}\label{h3:geom:ordinary-root-torsor}
 \Spec B_0^{1/p^l}\longrightarrow\Spec B_0.
\end{equation}
It is a finite flat torsor for
\[
 \operatorname{Hom}(\mathcal E[p^l],\mathcal C[p^l]),
\]
a form of $\mu_{p^l}^2$.  These identifications commute
with the $P_3$ action, with truncation in $l$, and with
Frobenius.  They therefore define a natural tower of
multiplicative-type torsors over the uncompleted ring $B_0$.
\end{lemma}

\begin{proof}
On the ordinary locus the relative Frobenius kernels define
the connected height-one subgroup; dividing by them gives
the finite \'etale quotients.  These compatible finite
groups give $\mathcal C$ and $\mathcal E$.  A height-one
one-dimensional connected $p$-divisible group becomes
$\mu_{p^\infty}$ on its connected-marking cover; at each
finite level the required cover is \'etale.  Thus
$\mathcal C[p^l]$ is of multiplicative type.

The $p^l$-torsion sequence is fpqc exact.  Locally choose
a basis of $\mathcal E[p^l]$ and lift its two generators
to $\mathcal G[p^l]$.  Every lift is killed by $p^l$, so
the two lifts define a homomorphism and a splitting.
Two splittings differ by a unique homomorphism to
$\mathcal C[p^l]$.  Consequently the splitting scheme
is a finite flat torsor for the displayed group, of rank
$p^{2l}$.

The \'etale basis cover is
$B'=C_l[x^{-1}]$, with $m=1$ in
\Cref{h3:geom:regular-levels}.  On this cover, the splitting
scheme is the scheme of generator lifts, whose algebra is
\[
 B_l[x^{-1}]
 =k[[t_{1,l},t_{2,l}]][x^{-1}]
 =(C_l[x^{-1}])^{1/p^l}=(B')^{1/p^l}.
\]
Relative Frobenius on an \'etale algebra is an isomorphism.
Therefore
\[
 (B')^{1/p^l}\simeq
 B'\otimes_{B_0}B_0^{1/p^l}.
\]
These local identifications descend canonically.  More
explicitly, in the splitting algebra the $p^l$th power
of every element lies in the base algebra, as can be
checked after the faithfully flat cover $B'$.  Send
that element to the unique $p^l$th root of its power in
$B_0^{1/p^l}$.  The resulting algebra map is an
isomorphism after the cover, hence an isomorphism.
This description shows both its uniqueness and its
independence of all basis and connected-marking choices.

The splitting construction is functorial for changes of
deformation framing.  Its canonical root identification
is therefore $P_3$-equivariant.  Restricting a splitting
at level $l+1$ to level $l$ is the quotient by the kernel
of the corresponding homomorphism of the translation
groups.  Under the unique root identifications it is
the inclusion $B_0^{1/p^l}\subset B_0^{1/p^{l+1}}$.
Frobenius compatibility follows from the same uniqueness
and the natural relative Frobenius maps of the finite
groups.  Thus all the stated compatibilities hold over
$B_0$ itself, rather than only on an auxiliary \'etale
extension.
\end{proof}

\section{Cohomology of the completed strata}\label{h3:sec:analytic}

We compute the two completed strata together with their maps from
constants.  These maps retain the class from the rank-two Tate frame
through the changes of coefficients later in the proof.
The Kummer and translation calculations of
Barthel--Mann--Ray--Schlank--Senger--Weinstein--Zhou
\cite[Sections~3.3--3.6]{BMR} provide the methodological setting.
We prove the relative estimates, support maps, and limit arguments
needed to retain the particular constants maps in this setting.

We use the relative extension sheaf $N_m$ and integral-frame stacks
of \Cref{h3:sec:towers}. In this section we abbreviate
\[
 E_m=\mathcal E_m=(N_m^{\,3-m})_{\mathrm{ind}},\qquad
 Y_m=\mathcal Y_m\simeq
 [E_m/(\GL_{3-m}(\Z_p)\times P_m)].
\]
Here $N_m$ is the sheafified relative $H^1$ of $V_m^\vee$ defined
before \Cref{h3:geom:independent-extensions}, and the stack equivalence
is \Cref{h3:prop:integral-frames}. Independence is over $\Q_p$.
In particular $E_2=N_2^*$, whereas
$E_1$ is the space of independent pairs in $N_1$.  All coefficient maps
below are maps on these stacks or on their marking torsors.  They
therefore commute with the actions and the changes of marking in
\Cref{h3:prop:integral-frames}.

\subsection{Relative coefficients and the affine calculation}

Fix an algebraically closed perfectoid untilt $C$ of $C^\flat$ and an
embedding $k\subset C^\flat$.  If $S=\operatorname{Spa}(R,R^+)$ is an
affinoid perfectoid space over $C^\flat$, a relative affine space means
the diamond of the affine space over its induced untilt.  Its
characteristic-$p$ structure sheaf is denoted by $\mathcal O^\flat$.
The notation $R(-1)$ retains the Kummer orientation line; choosing
compatible roots of unity in $C$ identifies it with $R$.  This choice
does not make the arithmetic action on this line trivial.

For the ordinary stratum over this relative base, we will prove that
the structure map
$Y_1\to B(\GL_2(\Z_p)\times P_1)$ induces an equivalence
\[
 R\Gamma_{\hthreects}(\GL_2(\Z_p)\times P_1,R)
   \longrightarrow R\Gamma(Y_1,\mathcal O^\flat).
\]
For coheight one, we first compute $N_2^*$: its cohomology consists
of constants $R$ in degree zero and a free $R$-line in degree three,
whose scalar action must be retained.  We will obtain both descriptions
by computing support cohomology in affine spaces and then taking
translation and frame-group cohomology.  Descent to $k$ will turn these
relative calculations into the completed coefficient comparisons,
including the finite connected-marking stages needed for coheight one.

Our relative assertions also include a profinite parameter space $T$.
On coefficients this replaces $R$ by $C_{\hthreects}(T,R)$.  Products have
their product topology, and countable inverse limits are derived until
their higher derived limits have been shown to vanish.  These
conventions are particularly relevant for the distribution modules
introduced below.
Finiteness means finite-dimensional cohomology over $C^\flat$ at the
geometric base $R=C^\flat$ and $T=*$; after descent, the field is $k$.
For an arbitrary relative base $R$, the corresponding answer is obtained
by scalar extension from $C^\flat$.
The corresponding parameter assertion for a computing complex
$K^\bullet$ is the natural identification
\[
 H^a(C_{\hthreects}(T,K^\bullet))
       =C_{\hthreects}(T,H^a(K^\bullet)).
\]
Thus, after descent to $k$, a finite-dimensional answer gives a
finite free module over $C_{\hthreects}(T,k)$.  It need not be
finite-dimensional over $k$ when $T$ is infinite.  The proofs
below establish this identity for the actual computing complexes.

\begin{lemma}[The relative additive presentation]\label{h3:ana:relative-bc}
Let $F/\Q_p$ be the unramified extension of degree $m$, embedded in $C$.
There is an identification over $C^\flat$
\[
 N_m\simeq \mathbb G_{a,C}^{\diamond}/\underline F.
\]
It is compatible with base change in $S$ and with scalar multiplication
by $\Q_p^\times$.  In particular
\[
 N_m^*\simeq
 [(\mathbb A^1_C\setminus\underline F)^{\diamond}/\underline F].
\]
\end{lemma}
\begin{proof}
On the Fargues--Fontaine curve with coefficient field $F$, the divisor
of the chosen untilt gives
\[
 0\longrightarrow\mathcal O_F(-1)\longrightarrow\mathcal O_F
   \longrightarrow i_*\mathcal O_{S^\sharp}\longrightarrow0.
\]
The relative section sheaf of $\mathcal O_F$ is $\underline F$,
its relative first cohomology vanishes, and
$\mathcal O_F(-1)$ has no relative sections.  These are the
slope-zero and negative-slope cases of the relative bundle
calculation in \Cref{h3:geom:curve-inputs} and \cite{FarguesFontaine,KedlayaLiu}. Pushforward along the
unramified coefficient extension identifies
$\mathcal O_F(-1)$ with $V_m^\vee$: its isocrystal has rank $m$,
degree $-1$, and the cyclic Frobenius of slope $-1/m$.
The cohomology sequence is consequently the displayed quotient.
All its arrows are relative arrows of sheaves, so the presentation
holds over $S$, not only over geometric fields.  The inverse image of
the zero section is precisely $\underline F$.
\end{proof}

Here is the estimate used in the nonproper calculation.  It is useful
to give it explicitly, since affinoid perfectoid acyclicity alone does
not compute the cohomology of an untilted affine line.

\begin{lemma}[Restriction on Kummer annuli]\label{h3:ana:annulus-contraction}
Fix a topological generator of $\Z_p(1)$, let $\epsilon\in C^\flat$
be its compatible system of roots, and put $\theta=|\epsilon-1|<1$.
Consider an untilted closed annulus
$A[r,b]=\{r\leq |z|\leq b\}$ and a smaller annulus
$A[r',b']$ satisfying
\[
 r'\geq r/\theta,\qquad b'\leq\theta b.
\]
In the Kummer cochain models, restriction is homotopic to its
constant-monomial part.  That part is
$R\oplus R(-1)[-1]$.  The homotopy on the nonconstant part is
bounded by one in these two Gauss norms.  The assertion holds
uniformly after $R$ is replaced by $C_{\hthreects}(T,R)$.
\end{lemma}
\begin{proof}
Adjoin all $p$-power roots of $z$.  The resulting annulus is
perfectoid, and its $\Z_p(1)$-torsor computes the cohomology by
\[
 [\,M\xrightarrow{\gamma-1}M\,].
\]
Here $M$ is the tilted Laurent algebra, completed for the two Gauss
norms.  Its elements have unique expansions
$\sum_{u\in\Z[1/p]}a_u z^u$; the weighted coefficients form a
$c_0$ family.  In this statement $c_0$ concerns all the fractional
exponents: for each positive bound only finitely many weighted
coefficients exceed that bound.  It is stronger than a condition
only as $|u|$ tends to infinity.

The differential on the $u$-th monomial is multiplication by
$\epsilon^u-1$.  Write $u=v p^j$, where $v\in\Z$ is prime to $p$.
For $u\ne0$, characteristic $p$ and $|v|_p=1$ give
\[
 |\epsilon^u-1|=\theta^{p^j},\qquad p^j\leq |u|_{\mathbb R}.
\]
For $u>0$, the outer-radius estimate is
\[
 \frac{(b'/b)^u}{|\epsilon^u-1|}
 \leq\theta^{u-p^j}\leq1.
\]
For $u<0$, the inner-radius estimate is
\[
 \frac{(r'/r)^u}{|\epsilon^u-1|}
 \leq\theta^{|u|-p^j}\leq1.
\]
Thus division by $\epsilon^u-1$ after restriction defines a
continuous contracting homotopy on all nonzero monomials.  The
$c_0$ condition is preserved by these inequalities.  The zero
monomial has zero differential, giving the two stated terms;
the second is $\Hom(\Z_p(1),R)$ and hence $R(-1)$.

The two-term model follows by applying continuous Hom to
\[
0\to\F_p[[\Z_p]]\xrightarrow{\gamma-1}\F_p[[\Z_p]]\to\F_p\to0.
\]
Its comparison with the completed bar resolution is continuous.
Affinoid perfectoid acyclicity
\cite[Proposition~8.8]{ScholzeDiamonds} then identifies it with the
cohomology of the annulus.  The displayed coefficientwise
inequalities are unchanged for the supremum norm on
$C_{\hthreects}(T,R)$, proving the parameter assertion.
\end{proof}

\begin{lemma}[The analytic inverse limits]\label{h3:ana:analytic-limits}
Let $M_0\leftarrow M_1\leftarrow\cdots$ be a tower of complete
nonarchimedean Banach spaces with dense transition maps of norm at
most one.  Its derived inverse limit has no higher cohomology.
The assertion remains true after applying $C_{\hthreects}(T,-)$.
In particular it applies to the tilted Laurent algebras on an
exhaustion by closed smaller annuli or polyannuli.
\end{lemma}
\begin{proof}
It suffices to show that the Roos difference map
\[
 \prod_n M_n\longrightarrow\prod_n M_n,\qquad
 (x_n)\longmapsto(x_n-f_nx_{n+1})
\]
is surjective.  Given $(y_n)$, construct a solution of the first
$j$ equations at stage $j$.  To add the next equation, choose
$x_{j+1}$ so that $f_jx_{j+1}$ approximates $x_j-y_j$ to within
$2^{-j}$.  Set the new $x_j$ equal to $y_j+f_jx_{j+1}$ and
propagate this correction backwards through the earlier equations.
The transition norm bound makes the change in each earlier
coordinate at most $2^{-j}$.  Every fixed coordinate therefore
converges, and the resulting tuple solves all the equations.
This proves the claim.

Density persists for continuous functions on $T$: approximate
on a finite clopen partition and lift approximations to the
finitely many values.  Completeness and the norm bound persist
as well, so the proof applies unchanged.  Fractional Laurent
polynomials belong to every algebra in the specified exhaustion
and are dense in each of them, verifying its hypotheses.
\end{proof}

\begin{lemma}[Affine space and zero-section purity]\label{h3:ana:affine-purity}
For every $d\geq1$ and every $S,T$ as above, the natural maps give
\[
 R\Gamma(\mathbb A^d_S,\mathcal O^\flat)=R,
 \qquad
 R\Gamma_{0}(\mathbb A^d_S,\mathcal O^\flat)
   =R(-d)[-2d].
\]
The first identification is the map of constants.  The second is
the product of the normal Kummer classes.  It is natural under
translation of the section and invertible linear changes of the
normal coordinates.  The equalities are relative equalities,
including the continuous parameters $T$.
\end{lemma}
\begin{proof}
We first treat a line.  An untilted disc is covered by the
perfectoid disc obtained by adjoining all roots of its coordinate.
A section of $\mathcal O^\flat$ on the original disc pulls back to
an invariant section on this cover.  The monomial computation in
\Cref{h3:ana:annulus-contraction}, now with nonnegative exponents,
shows that every such section is constant.  Consequently the
degree-zero cohomology of a relative disc is $R$.

Restrictions from sufficiently large discs to a fixed smaller disc
kill positive cohomology.  To see this without making an assertion
about the branched cover at the center, choose a point $a$ outside
the smaller disc and inside the larger one.  The inclusion of the
smaller disc factors through two nested annuli centered at $a$,
with the radius gaps of \Cref{h3:ana:annulus-contraction}.  The larger
disc can be chosen to contain these annuli.  Restriction between
them kills the nonconstant part and all degrees above one.
The remaining degree-one class is the mod-$p$ Kummer class of
$z-a$.  On a sufficiently small disc about zero it vanishes:
if $|z/a|<|p|^{p/(p-1)}$, the binomial series for
$(1-z/a)^{1/p}$ converges, and a $p$-th root of $-a$ supplies
a first $p$-th root of $z-a$.  All choices can be made using
constants from $C$, with bounds independent of $S$ and $T$.

Exhaust the line by such discs.  For each positive degree the
inverse system of cohomology groups is pro-zero, and degree zero
is the constant system $R$.  The countable derived-limit sequence
therefore proves the first assertion for a line.  In particular,
no assertion that density alone kills a derived limit is used.

Exhaust the punctured line by annuli whose inner radii tend to zero
and outer radii tend to infinity, taking a cofinal subsequence with
the stated radius gaps.  The homotopies of
\Cref{h3:ana:annulus-contraction} show that its cohomology is
$R\oplus R(-1)[-1]$.  The constants identify the first summand.
The fiber of restriction from the line to the punctured line is
therefore $R(-1)[-2]$.  Its generator is the boundary of the
Kummer class.  Multiplication of a coordinate by an invertible
constant changes its Kummer class by the class of that constant,
which vanishes locally after adjoining a root.  Translation simply
changes the chosen section.  This proves the asserted naturality.

For $d$ coordinates apply the relative line calculation one
coordinate at a time.  Successive cohomology with support in their
zero loci gives the product of the $d$ Kummer boundary classes,
in degree $2d$ with twist $-d$.  Invertible changes of coordinates
preserve this local fundamental class: elementary additions are
checked on the complement by the same nearby-coordinate Kummer
calculation, diagonal changes were just checked, and coordinate
permutations permute classes of even degree.  These operations
generate the invertible linear changes locally on the base.
The constructions and the limiting arguments above are uniform
on continuous profinite families, proving the relative statement.
\end{proof}

\subsection{Tubes, distributions, and translation orientation}

The next support calculation replaces the single zero section by a
locally profinite family of sections.  We surround the family by disjoint
tubes and refine those tubes to the support.  The resulting support
classes will be identified with compatible residue assignments: when
tubes merge, their residues add.  The distribution module below records
this summation rule.
For a compact totally disconnected space $K$ define
\[
 \mathcal D(K,R)=
 \varprojlim_{\mathscr P}\bigoplus_{U\in\mathscr P}R,
 \qquad
 (a_V)_{V\in\mathscr Q}\longmapsto
 \left(\sum_{V\subset U}a_V\right)_{U\in\mathscr P},
 \tag{$\ast$}\label{h3:ana:distribution-definition}
\]
where $\mathscr Q$ refines the finite clopen partition
$\mathscr P$.  For a locally compact totally disconnected space
$Z=\coprod_i K_i$, put
$\mathcal D_{\mathrm{lf}}(Z,R)=\prod_i\mathcal D(K_i,R)$.
Equivalently this is the continuous dual, in this product sense,
of compactly supported locally constant functions on $Z$.
This description makes it independent of the decomposition.
It imposes no common norm bound on the different components.
The transition maps in \eqref{h3:ana:distribution-definition} are
split surjections, by choosing one child of each partition member.
For a profinite $T$ there are natural identifications
\begin{equation}
 C_{\hthreects}(T,\mathcal D_{\mathrm{lf}}(Z,R))
 =\mathcal D_{\mathrm{lf}}(Z,C_{\hthreects}(T,R)).
 \label{h3:ana:distribution-parameters}
\end{equation}
They follow directly by commuting continuous maps with products
and inverse limits.

\begin{lemma}[Support in a family of rational lattices]\label{h3:ana:tube-purity}
The following supports in an untilted relative affine line have
cohomology
\[
 R\Gamma_Z(\mathbb A^1_S,\mathcal O^\flat)
   \simeq\mathcal D_{\mathrm{lf}}(Z,R)(-1)[-2]:
\]
\begin{enumerate}[label=\textup{(\roman*)}]
\item a constant finite extension $F\subset C$ of $\Q_p$;
\item the image of $(a,b)\mapsto az+b$, with $a,b\in\Q_p$
and $z$ a section of the relative affine line whose value in each
geometric untilt lies outside $\Q_p$; one may also restrict $a$ to a
compact open subset.
\end{enumerate}
These are relative identifications, natural for the translations
and linear changes preserving the indicated support.  They retain
all profinite parameters.  Taking products of the first support
in two separate coordinates gives
\[
 R\Gamma_{F^2}(\mathbb A^2_S,\mathcal O^\flat)
   \simeq\mathcal D_{\mathrm{lf}}(F^2,R)(-2)[-4].
\]
\end{lemma}
\begin{proof}
We describe both the neighborhoods and the limit that computes
their support cohomology.  On a geometric fiber the map in (ii)
is a continuous linear injection from the finite-dimensional
$\Q_p$-Banach space $\Q_p^2$ into the geometric untilt.  Its induced
norm is equivalent to the maximum norm, and hence is bounded above
and below on the unit sphere.  The lower bound can be made uniform on a
neighborhood of that fiber.  Indeed the parameter unit sphere
in $\Q_p^2$ is compact; for each of its points the corresponding
$az+b$ is nonzero, and finitely many neighborhoods give a common
positive lower bound.  Shrink the base by their intersection.
Homothety by powers of $p$ now gives constants $c_-,c_+>0$ such that
\[
 c_-\max(|a|,|b|)\leq |az+b|
       \leq c_+\max(|a|,|b|)
\]
uniformly on that neighborhood.  The same assertion for (i) is
the norm comparison for the fixed embedding $F\subset C$.

Write $V$ for the parameter space and $d=\dim_{\Q_p}V$;
thus $d=[F:\Q_p]$ in \textup{(i)} and $d=2$ in \textup{(ii)}.
For the compact open restriction on $a$, use $V=\Q_p^2$
and start at a level at which the restricted strip is a
union of lattice cosets; compact openness permits this.
Choose a lattice adapted to the norm, as follows.  At the
chosen geometric point, the norm of the injection is locally
constant and nowhere zero on the compact parameter unit sphere,
so it takes only finitely many values there.  Choose a radius
in a gap and let $\Lambda$ be its closed parameter ball.
This is a $\Z_p$-lattice.  After a further rational restriction
of the base there are constants $\rho_-<\rho<\rho_+$ such that
\[
 |\ell(v)|\leq\rho_-\quad(v\in\Lambda),\qquad
 |\ell(v)|\geq\rho_+\quad(v\notin\Lambda),
\]
where $\ell$ denotes the indicated injection.  Only finitely
many inequalities are needed for this restriction: use a
$\Z_p$-basis of $\Lambda$ for the upper bound and representatives
$w_j$ of the nonzero cosets of $(p^{-1}\Lambda)/\Lambda$
for the lower bound.  Indeed, for $v\notin\Lambda$, a
nonnegative power of $p$ takes $v$ to
$v'=w_j+\lambda\in p^{-1}\Lambda\setminus\Lambda$.
The strict gap gives
$|\ell(v')|=|\ell(w_j)|\geq\rho_+$, and scalar rescaling
gives the desired bound for $v$.  These are inequalities
on the rational base neighborhood and hence hold on
every perfectoid test over it.

For every $n$ take discs of radius $|p|^n\rho$ about the
images of representatives modulo $p^n\Lambda$.
Distinct centers are separated by at least $|p|^n\rho_+$,
so the discs are disjoint and nested under refinement.
Only finitely many meet any bounded relative affine disc:
the lower norm bound restricts their parameters to a bounded,
hence compact, subset of $V$, which has finitely many cosets
modulo $p^n\Lambda$.  Thus the family is locally finite.

We verify its limiting support on an arbitrary affinoid
perfectoid test $S'\to S$.  Choose two buffered radii between
$\rho_-$ and $\rho_+$.  A section $w$ belonging to every
smaller closed tube system belongs to every larger open
tube system.  The inverse images of the latter's disjoint
components are open and closed in $S'$.  They define
compatible locally constant maps
$c_n:S'\to V/p^n\Lambda$.  Quasicompactness gives finitely
many values at each level.  Since
$V=\varprojlim_n V/p^n\Lambda$ with its locally profinite
topology, these maps give a continuous parameter $c:S'\to V$.
Choose representatives of the $c_n$.  Their images under
$\ell$ converge uniformly, since the successive differences
have norm at most $|p|^n\rho_-$.  The tube condition gives
$\|w-\ell(c_n)\|\leq |p|^n\rho_{\mathrm{outer}}\to0$;
completeness therefore gives $w=\ell(c)$.  Conversely a
continuous parameter $c$ belongs to all the smaller tube
systems by the upper bound on $\Lambda$.  The buffered
neighborhoods consequently have the required intersection
as relative v-sheaves, including analytic boundary points.
Their complements exhaust the desired complement, and
support cohomology is their derived inverse limit.

The exterior of a disc is exhausted by annuli.  By
\Cref{h3:ana:analytic-limits} its cohomology is the two-term
Kummer complex of Laurent series converging on that exterior.
Its degree-zero invariants are $R$.  Taking the fiber from
the affine line consequently puts the tube support entirely
in degree two.  Constant-monomial extraction gives a split
surjective residue map in that degree to $R(-1)$.
Its kernel is the nonconstant cokernel of the difference
operator.

There is a canonical description of this residue: map to
the direct limit of support cohomology over all larger
enclosing discs.  The annulus homotopy identifies that
limit with $R(-1)$.  Recentring gives the same identification,
since sufficiently far out the ratio of the two coordinates
has a first $p$-th root.  Inclusions of discs therefore
commute with residues to this common target.  Excision
identifies disjoint children with a direct sum, so the
residue of their merger is literally the sum of their
residues.  Surjectivity is also witnessed by the class of
any section inside the tube, which has residue one.

The kernels of the residue maps form a pro-zero system with
a fixed level gap.  Choose $j$ such that $|p|^j<\theta$, with
$\theta$ as in \Cref{h3:ana:annulus-contraction}; increase $j$
if necessary for the two buffered radii.  The ratio of a
level-$(n+j)$ child's radius to its level-$n$ parent's radius
is $|p|^j$, independently of $n$ and of the parent coset.
Recenter the parent disc at this child's center.  Its exterior
is unchanged because the center difference is strictly smaller
than the parent's radius.  On restriction between these
exteriors the inverse of the difference operator is bounded
on negative monomials by the fixed radius gap.  For positive
monomials use a source outer radius at least the target outer
radius divided by $\theta$, which is always available on an
exterior.  Thus the entire child residue kernel maps to zero
in the parent's support cohomology.

Each parent has exactly $p^{dj}$ children.  Finite summation
shows that the map from all level-$(n+j)$ residue kernels to
all level-$n$ kernels is zero.  For infinitely many parents
the homotopies operate coordinatewise in the product topology;
they require no common bound on the coefficient values.
For a profinite $T$ the same inequalities apply in each
$C_{\hthreects}(T,R)$ supremum norm.  Thus one level gap annihilates
the kernels on every profinite test, and not merely on
geometric fibers.

This proves that the actual system of tube support cohomology
is pro-isomorphic to the system of finite sums $R(-1)$ under
summation.  It does not require that the classes of different
sections in a fixed tube be equal.  Those classes can differ
by a residue-zero class before passage to the limit.
The residue system has split surjective transitions, while
the discarded systems are pro-zero.  Both their $\varprojlim^1$
terms consequently vanish.  Locally finite disjoint components
give products by excision.  This proves the formula
\eqref{h3:ana:distribution-definition} with its product topology.

All norm estimates and all residue maps were made over the
base neighborhood.  Applying them to $C_{\hthreects}(T,R)$ proves
the parameter assertion, and descent glues the neighborhoods
of the base.  The product statement follows by performing the
two relative support calculations successively; its partition
system consists of rectangles and is cofinal in the finite
clopen partitions of the product.
\end{proof}

\begin{lemma}[Translation cohomology]\label{h3:ana:translations}
Let $F/\Q_p$ have degree $m$, acting by translations.  Then
\[
 R\Gamma_{\hthreects}(F,R)=R,\qquad
 R\Gamma_{\hthreects}(F,\mathcal D_{\mathrm{lf}}(F,R))
       =R\otimes_{\F_p}\operatorname{or}_F[-m].
\]
Here $\operatorname{or}_F$ is the inverse orientation of the
$m$-dimensional translation group.  With the convention that
$a$ acts on points by $z\mapsto az$, a scalar
$a\in\Z_p^\times$ acts on this line by $\bar a^{-m}$.
The assertions and their support maps hold with continuous
profinite parameters.
\end{lemma}
\begin{proof}
Choose a lattice $K\simeq\Z_p^m$ in $F$.  On its distributions
the finite Koszul complex is the Koszul cochain complex for
\[
 R[[X_1,\ldots,X_m]],\qquad X_i=\gamma_i-1.
\]
The power-series notation here means the inverse limit of finite
group algebras, as in \eqref{h3:ana:distribution-definition}, hence
a product of coefficient copies.  Multiplication by a variable
is injective and its quotient is obtained by setting that
variable to zero.  Successively applying these statements
shows that the Koszul cochain complex has just its top
cohomology, the augmentation quotient $R$, in degree $m$.
The coefficient splitting by powers of $X_i$ makes these
statements exact also with continuous parameters.

There is an identification of translation modules
\[
 \mathcal D_{\mathrm{lf}}(F,R)
  =\operatorname{Coind}_{K}^{F}\mathcal D(K,R).
\]
Indeed $F/K$ is discrete, and both sides are the product of
one copy of $\mathcal D(K,R)$ for each coset, with the same
translation rule.  Evaluation on $K$, and the bar contraction
using representatives for these open cosets, give continuous
cohomological Shapiro for this coinduced module.  This proves
the second formula.  Changing a lattice basis acts on the
top cochain generator by the inverse determinant of its
change matrix.  Scalar multiplication by $a$ has determinant
$a^m$, proving the asserted character.

For the first formula, exhaust $F$ by $p^{-j}K$.  In every cochain
degree, restriction between these open-and-closed subgroups is
surjective by extension by zero.  The inverse limit is the continuous
cochain complex on $F$, since these open subgroups cover $F$; the same
statements hold with the parameter $T$.  On trivial
characteristic-$p$ coefficients the Koszul differentials
are zero, and compact-lattice cohomology is the exterior
algebra on the continuous linear forms of the lattice.
Restriction from $p^{-j-1}K$ to $p^{-j}K$ multiplies each
degree-one form by $p$ and is therefore zero in every positive
degree.  In degree zero restriction is the identity.
The derived inverse limit is consequently $R$.
The same finite complexes and the same vanishing transition
maps compute the parameter versions.
\end{proof}

\begin{proposition}[The punctured negative space]\label{h3:prop:punctured-cohomology}
For $m=1,2$ the natural constants map fits into a fiber sequence
\[
 R\longrightarrow R\Gamma(N_m^*,\mathcal O^\flat)
   \longrightarrow
 R(-1)\otimes\operatorname{or}_F[-m-1].
\]
It is equivariant for scalar units and natural in $S$ and in
profinite parameters.  Thus the only cohomology groups are
$R$ in degree zero and an orientation line in degree $m+1$.
The scalar $a\in\Z_p^\times$ has weight $\bar a^{-m}$ on
the latter.  Reversing the action convention reverses all
these exponents simultaneously.
\end{proposition}
\begin{proof}
Use the support fiber sequence for
$\underline F\subset\mathbb A^1_C$ in
\Cref{h3:ana:relative-bc}.  Its middle term is $R$ by
\Cref{h3:ana:affine-purity}, and its first term is
$\mathcal D_{\mathrm{lf}}(F,R)(-1)[-2]$ by
\Cref{h3:ana:tube-purity}.  Take the translation cohomology
of \Cref{h3:ana:translations}.  The resulting support term
has degree $m+2$; rotation gives the stated cofiber in
degree $m+1$.  The map from $R$ throughout this construction
is the coefficient unit.  Scalar multiplication preserves
the normal Kummer class and contributes exactly the
translation orientation computed in that lemma.
\end{proof}

\subsection{Independent pairs and compact frame cohomology}

The punctured calculation has produced the two cohomology rows for
coheight one.  For the ordinary stratum, we must instead remove the
dependent pairs in $N_1^2$.  We compute the two support terms with
their frame characters before taking the compact frame-group
cohomology.

Put $N=N_1$, $E=E_1$, and $G_f=\GL_2(\Z_p)\times P_1$.
Let $D\subset N^2\setminus\{0\}$ be the locus of proportional
pairs, so that $E=(N^2\setminus\{0\})\setminus D$.
The direction of a proportional pair belongs to
$\mathbb P^1(\Q_p)$.

\begin{lemma}[The two support terms]\label{h3:ana:ordinary-support}
The support at the double origin in $N^2$ is a line in degree
six, of scalar weight $-2$.
The support in $D$ has cohomology in degrees three and five.
In each degree it is a distribution module on
$\mathbb P^1(\Q_p)$ with a one-dimensional continuous fiber.
The scalar weights of these two fibers are respectively
$-1$ and $-2$.

These descriptions are equivariant relative descriptions,
with continuous profinite parameters.  For the stabilizer
of the direction $(1,0)$, written
\[
 B=\left\{\begin{pmatrix}a&b\\0&d\end{pmatrix}
                \in\GL_2(\Z_p)\right\},
\]
the two fibers have characters $\bar d^{-1}$ and
$\bar a^{-1}\bar d^{-1}$.  An element $v\in P_1=\Z_p^\times$
has characters $\bar v$ and $\bar v^2$, respectively.
Normal Kummer twists are understood in these descriptions.
\end{lemma}
\begin{proof}
The preimage of the double origin under
$\mathbb A^2_C\to N^2$ is $\Q_p^2$.
The product support calculation gives its distributions in
degree four.  Taking the two translation directions adds
degree two and inverse determinant orientation, giving
degree six and scalar weight $-2$.

On a slope chart, lift the first coordinate to
$z_1\notin\Q_p$.  The support in the second coordinate is
\[
 z_2=a z_1+b,\qquad a,b\in\Q_p,
\]
with $a$ in a compact open slope patch when desired.
\Cref{h3:ana:tube-purity} gives distributions in $(a,b)$ in
normal support degree two.  Taking translations in the
second coordinate is the $b$-translation calculation of
\Cref{h3:ana:translations}.  It leaves distributions in $a$
and adds degree one with inverse orientation.  First-coordinate
translation replaces $b$ by $b-ac$ and leaves $a$ fixed.
It consequently preserves the resulting orientation line.
The remaining base is $N^*$; \Cref{h3:prop:punctured-cohomology}
adds its degrees zero and two.  We obtain degrees three and
five, and weights $-1$ and $-2$ with the same sign.
To justify this computation with distributions in the slope,
first use a finite clopen slope partition.  It gives a finite
sum of the relative punctured calculation.  Then take the
partition limit.  The summation maps are split and the
discarded residue kernels are uniformly pro-zero by
\Cref{h3:ana:tube-purity}; hence this limit introduces no
additional cohomology.  This also specifies the topology
of the resulting distribution rows.

For equivariance, a point mass at a direction is the support
map from that direction's tubes.  These maps commute with
linear changes of the pair.  Finite sums of point masses
are dense in the finite-partition topology: on any finite
partition every prescribed collection of values is realized
by such a sum.  Transport of the point masses therefore
determines transport on the full distribution module.
At the direction $(1,0)$ the tangential coordinate is changed
by $a$ and the normal coordinate by $d$.  The inverse
translation orientations just computed are $\bar d^{-1}$
and $\bar a^{-1}\bar d^{-1}$; the upper-right entry acts
trivially on both fibers.  The quotient frame $v$ scales
both coordinates by $v^{-1}$, giving the stated $P_1$
characters.  These are finite characters.  The calculations
of tubes and translations were relative, so this transport
argument also applies with all the indicated parameters.
\end{proof}

\begin{proposition}[The ordinary constants map]\label{h3:prop:ordinary-constants}
For $p\geq5$, the structure map $Y_1\to BG_f$ and the
coefficient unit induce an equivalence
\[
 R\Gamma_{\hthreects}(G_f,R)
   \xrightarrow{\ \simeq\ }
 R\Gamma(Y_1,\mathcal O^\flat).
\]
This is an equivalence of derived algebras and is natural
over $S$, with continuous profinite parameters.  It retains
both frame-group actions before taking their cohomology.
In particular the degree-three class from the rank-two
Tate frame is carried by this actual map from constants.
\end{proposition}
\begin{proof}
By \Cref{h3:ana:affine-purity,h3:ana:translations}, constants
compute the cohomology of $N^2$.  Remove the double origin
and then the proportional locus.  These two support fiber
sequences are equivariant for $G_f$ and their maps from
constants are the coefficient unit.

The central subgroup $\mu_{p-1}I_2\subset\GL_2(\Z_p)$
acts on their support cohomology by the characters of
weights $-1$ and $-2$ in \Cref{h3:ana:ordinary-support}.
Neither character is trivial when $p\geq5$.  Its averaging
idempotent is defined in characteristic $p$, so its derived
invariants are exact and annihilate both support complexes.
Taking the remaining frame-group cohomology therefore
turns the constants arrow into an equivalence.
The arrow was induced by the structure map and coefficient
unit, which are multiplicative.  Its equivalence on
underlying complexes proves the assertion for derived
algebras as well.
\end{proof}

For later finiteness arguments we need twists even when
the preceding averaging no longer kills the support.
The appropriate induction is induction of measures.

\begin{lemma}[Cohomology of induced measure modules]\label{h3:ana:measure-induction}
Let $G$ be a compact $p$-adic analytic group without
$p$-torsion, let $B\subset G$ be a closed analytic subgroup
without $p$-torsion, and suppose $G\to G/B$ has continuous
local sections.  If $V$ is a finite-dimensional continuous
$k$-representation of $B$, then
\[
 M=k[[G]]\widehat\otimes_{k[[B]]}V
\]
has finite-dimensional continuous $G$-cohomology in every
degree.  The same cohomology after replacing coefficient
copies of $k$ by $R$ is the corresponding finite-dimensional
cohomology tensored with $R$, including continuous
profinite parameters.
\end{lemma}
\begin{proof}
Write $g=\dim G$.  Analytic group duality identifies
cohomology with homology in complementary degree:
\[
 H^a_{\hthreects}(G,M)
   \simeq H^{\hthreects}_{g-a}
        (G,\operatorname{or}_G\otimes_k M).
\]
The orientation module is retained in this formula.
One obtains it by dualizing a finite projective completed
group-ring resolution of $k$ and reversing its degrees.
Such resolutions exist for the stated groups by
\cite[Proposition~3.3]{AW}; their dualizing module is the top inverse
adjoint orientation
\cite[Proposition~4.16, Remark~4.23, and Proposition~4.40]{BGHS}.
Thus the formula applies to compact
measure modules, not just to finite coefficients.

Completed homological Shapiro gives
\[
 H^{\hthreects}_{g-a}(G,\operatorname{or}_G\otimes M)
   \simeq H^{\hthreects}_{g-a}
       (B,\operatorname{or}_G|_B\otimes V).
\]
For completeness, local sections identify $k[[G]]$, as a
right completed $k[[B]]$-module, with a completed free
module on $G/B$.  Inducing a completed projective
$B$-resolution and then taking $G$-coinvariants is
therefore the same complex as taking $B$-coinvariants.
This proves the displayed Shapiro identification.
The last homology is finite-dimensional, since a finite
projective $B$-resolution has finite-dimensional terms
after tensoring with the finite-dimensional fiber.
This argument includes the duality shift; it makes no
unshifted cohomological Shapiro assertion for $M$.

To justify extension to $R$, start with compact $k$-modules
and apply the exact functor
\[
 A\longmapsto\Hom_k(A^\vee,R),\qquad
 A^\vee=\Hom_{\hthreects,k}(A,k).
\]
After a choice of a basis of $A^\vee$ its value is a
product of copies of $R$.  Exactness follows by splitting
sequences of the discrete vector spaces $A^\vee$.
Finite projective resolutions use only finite powers
and idempotent summands, so their comparison, duality,
and Shapiro maps commute with this functor.  Apply the
same argument to $C_{\hthreects}(T,R)$, or use
\eqref{h3:ana:distribution-parameters}.
\end{proof}

\begin{corollary}[Twisted ordinary finiteness]\label{h3:ana:twisted-finiteness}
If a coefficient line on $Y_1$ becomes a finite-character
constant line on $E_1$, its geometric cohomology is
finite-dimensional in every degree and vanishes outside
a finite range.  This applies to all powers of the
periodicity line and all finite tame determinant twists.
For a profinite parameter $T$, the corresponding cohomology
is $C_{\hthreects}(T,-)$ of this finite-dimensional answer.
\end{corollary}
\begin{proof}
The ambient constants and the origin-support line have
finite-dimensional fibers with finite-character action.
The two remaining support rows in
\Cref{h3:ana:ordinary-support} are completed measure induction
from $B\times P_1$ to $G_f$.  To check the local-section
hypothesis, a primitive vector over $\Z_p$ can, on either
standard slope chart, be completed continuously to an
invertible matrix.  These sections also show that
$\GL_2(\Z_p)$ acts transitively on $\mathbb P^1(\Q_p)$.
Its stabilizer is the displayed upper triangular group.
Both $G_f$ and $B\times P_1$ are compact analytic without
$p$-torsion when $p\geq5$: a matrix of order $p$ in
dimension two would require the degree-$p-1$ cyclotomic
polynomial, and $\Z_p^\times$ has no $p$-torsion.
\Cref{h3:ana:measure-induction} proves finiteness of their
support cohomology.  The two support sequences prove the
claim for the independent locus.

On the connected Honda marking the periodicity line is
framed by the leading derivative of the formal-group
isomorphism.  At height one its change of frame is
the reduction $\Z_p^\times\to\F_p^\times$.  The determinant
identity of \Cref{h3:prop:integral-frames} expresses a tame
middle determinant character through the two frame
characters.  The indicated coefficient lines therefore
have exactly the finite-character property used above.
\end{proof}

\subsection{Descent to the finite constant field}

The geometric calculations are over $C^\flat$.  We now return to the
finite field $k$ while retaining the stack maps and the residual group
actions.  This step uses base-field Frobenius on the $k$-defined
spaces; the chosen untilt chart remains a device for computing their
geometric cohomology.

For a $k$-defined space $X$, base-field Frobenius means the
map $\hthreeid_X\times\varphi_q$ on
$X\times_{\operatorname{Spd}k}\operatorname{Spd}C^\flat$.
In a coordinate chart defined over $k$ it raises coefficients
to the $q$-th power and fixes the variables.  It is distinct
from Frobenius of all the variables in a perfected coefficient
ring.  A chosen untilt-divisor presentation need not be
preserved by this map.

\begin{lemma}[Base-field Frobenius with parameters]\label{h3:ana:base-frobenius}
Let $q=|k|$, and let $M_k$ be the analytic function space over $k$
on a chosen perfected punctured ball or on the perfected ordinary
polyannular open.  Let $M_C$ be the corresponding function space after
extension from $k$ to $C^\flat$.  Let $\varphi$ act as the $q$-th power
on coefficients and fix the variables.  Then
\[
 0\longrightarrow M_k\longrightarrow M_C
      \xrightarrow{\varphi-1}M_C\longrightarrow0
\]
is exact, including after $C_{\hthreects}(T,-)$ for a profinite
space $T$.  The same statement applies to finite products
of the punctured-ball charts at finite marking levels.
Consequently the coefficient comparison is an equality
with derived base-Frobenius invariants, naturally for all
group actions defined over $k$.
\end{lemma}
\begin{proof}
Use the expansions in fractional monomials on exhausted
closed smaller polyannuli.  An invariant expansion has
each coefficient in $\F_q=k$, and hence is precisely an
element of $M_k$.  For a coefficient $a\in C^\flat$, choose
a root $b$ of $b^q-b=a$ of least absolute value.  The
Newton polygon, or the convergent series
$b=-\sum_{j\geq0}a^{q^j}$ when $|a|<1$, gives
\[
 |b|=|a|\quad (|a|<1),\qquad
 |b|\leq1\quad (|a|=1),\qquad
 |b|=|a|^{1/q}\quad (|a|>1).
\]
In particular $|b|\leq|a|$.  Solving for each coefficient
preserves the $c_0$ convergence conditions in every
smaller polyannulus norm.  This proves surjectivity and
also supplies solutions arbitrarily small whenever the
input is sufficiently small in any finite collection
of these norms.
The Frobenius map itself is continuous in the full analytic
topology: for a Gauss poly-radius $\rho$,
$\|\varphi f\|_\rho=\|f\|_{\rho^{1/q}}^q$.
The slightly larger poly-radius on the right belongs to
the same open domain.

For continuous parameter families, first approximate a
given function on a finite clopen partition of $T$ and
lift the finitely many values.  The remaining error is
small.  Solve for its coefficients with the small
branch just described, refine the partition, and repeat
with errors tending to zero in successive defining
norms.  The estimates give a uniformly convergent sum
of continuous corrections; its limit is a continuous
lift.  This argument proves exactness on parameters
directly.  For the spaces defined by an exhaustion,
the continuous Laurent-polynomial approximation and
the complete norm topologies meet the hypotheses of
\Cref{h3:ana:analytic-limits}.  Thus these analytic inverse
limits have no additional derived term.

The maps $M_k\to M_C$ and $\varphi-1$ themselves are
natural for changes of variables defined over $k$.
Their auxiliary coefficient choices need not be natural.
The exact sequence therefore identifies the original
coefficient complex with the fiber of $\varphi-1$,
equivariantly and with parameters.  If a finite marking
chart has residue field $k'/k$, geometric base change
gives a product indexed by the embeddings of $k'$.
The canonical base Frobenius permutes these factors.
On a cycle of length $e=[k':k]$, eliminating successive
coordinates reduces its difference equation to
coefficient Frobenius $q^e$ minus identity.  The preceding
estimates and continuous lifting apply with $q^e$.
Thus finite products and finite constant-field descent
retain the claimed exact sequence, with the actual
permutation of factors.
\end{proof}

Write $B_{\hthreepk}$ for the algebra of functions on $\mathcal X_1$,
the completed perfection of the ordinary locus, with the inverse-limit
topology of uniform convergence on the exhausted closed polyannular
charts.  \Cref{h3:sec:ordinary} gives an equivalent description using
finite root levels of Laurent series and a countable family of defining
norms.

\begin{theorem}[Completed cohomology and its constants maps]
\label{h3:thm:completed-cohomology}
The completed strata of \Cref{h3:thm:coordinate-comparison}
have the following properties.
\begin{enumerate}[label=\textup{(\roman*)}]
\item For the ordinary completed coefficient ring
$B_{\hthreepk}$, the natural constants map is an equivalence
\[
 R\Gamma_{\hthreects}(\GL_2(\Z_p)\times P_1,k)
   \xrightarrow{\ \simeq\ }
 R\Gamma_{\hthreects}(P_3,B_{\hthreepk}).
\]
It is the map defined by $Y_1\to B(\GL_2(\Z_p)\times P_1)$,
and is compatible with geometric constant extension,
products, coefficient Frobenius, and profinite parameters.
Every periodicity twist and finite tame character twist
has finite cohomology in each degree, in a bounded range.

\item Let $L$ be the connected height-two marking algebra,
$U\subset P_2$ an open subgroup, $B=L^U$, and
$B^{\mathrm b}=\widehat{B^{\mathrm{perf}}}$.  Finite products of
fields are allowed.  Then each
$H^a_{\hthreects}(H,B^{\mathrm b})$ is finite-dimensional over $k$,
and its degree-zero group is $k$ by constants.
The same finiteness holds for lines framed on a finite
connected-marking cover, including the periodicity and
tame determinant lines used below.
On taking the filtered colimit over marking levels there
is a natural fiber sequence
\[
 k\longrightarrow
 \colim_U R\Gamma_{\hthreects}(H,(L^U)^{\mathrm b})
   \longrightarrow k_{\lambda}[-3].
\]
The first arrow is the inclusion of constants.  The
full residual norm quotient $P_3/H=\Z_p^\times$ acts
on $k_\lambda$ by $a\mapsto\bar a^{-2}$ in the
point-action convention above.  The commuting $P_2$
action is retained.  It is smooth on the displayed
finite-dimensional rows.  The statement includes
restrictions between marking levels.  With a profinite
parameter $T$, each displayed cohomology group is replaced
by its $C_{\hthreects}(T,-)$, as specified at the start of the section.
\end{enumerate}
In particular, invariants of the order-$p-1$ subgroup
of the full norm quotient kill the extra line in
\textup{(ii)}.  This subgroup is in the quotient;
no section of $P_3\to\Z_p^\times$ is required.
\end{theorem}
\begin{proof}
First compute over $C^\flat$.  For the ordinary locus,
\Cref{h3:prop:ordinary-constants} is precisely the desired
map on the quotient-stack presentation in
\Cref{h3:prop:integral-frames}.  The perfected polyannular
exhaustion has affinoid perfectoid charts.  Its coefficient
complex is the function complex $B_{\hthreepk}$, since the
successive polynomial approximation in
\Cref{h3:ana:analytic-limits} gives the required vanishing
of the derived analytic limit.  Thus this is also the
ordinary continuous $P_3$-cochain comparison with the
stated topology.

Take derived invariants of base-field Frobenius using
\Cref{h3:ana:base-frobenius}.  On the constants side the
fiber of $\varphi-1$ is $k$.  This proves the untwisted
comparison over $k$ with its actual coefficient map.
For a twist, \Cref{h3:ana:twisted-finiteness} gives a bounded
complex with finite-dimensional geometric cohomology.
Frobenius acts semilinearly and bijectively on these
groups.  On a finite-dimensional $C^\flat$-space a
bijective $q$-semilinear operator $F$ has $F-1$
surjective and a finite-dimensional fixed space over
$k$.  Indeed write $F(x)=A x^{[q]}$ with $A$ invertible.
The system $A x^{[q]}-x=b$ is equivalent to monic
equations $x_i^q=(A^{-1}(x+b))_i$.  Their pairwise
coprime leading monomials show that the coordinate
algebra is free over the $b$-coordinate algebra on
the $q^n$ monomials with each exponent less than $q$.
The derivative is $-I$, so the map is finite \'etale,
surjective, and has kernel of order $q^n$.  The kernel
is a $k$-vector space and therefore has dimension $n$.
This proves the asserted semilinear calculation.
Hence Frobenius descent
preserves the finite-dimensional conclusion.  All
these operators commute with the frame actions and
with the coefficient maps.

For \textup{(ii)}, define Frobenius before choosing an
untilt-divisor presentation.  The space $N_2$ and its
$P_2$ and residual norm actions are defined over $k$.
Consequently $\hthreeid_{N_2}\times\varphi_q$ induces a
$q$-semilinear automorphism $\Phi$ of
\[
 \mathcal K_C=R\Gamma(N_{2,C}^*,\mathcal O^\flat),\qquad
 N_{2,C}=N_2\times_{\operatorname{Spd}k}
                       \operatorname{Spd}C^\flat,
\]
commuting with both actions.  The presentation
$N_{2,C}^*\simeq[(\mathbb A_C^1\setminus F)/F]$ computes
$\mathcal K_C$.  It is not required to commute with
$\Phi$: transporting $\Phi$ through it can change the
chosen untilt divisor.  Put
$\mathcal K_k=\fib(\Phi-1:\mathcal K_C\to\mathcal K_C)$.

Write $X_B=\operatorname{Spa}(B^{\mathrm b})$ and
$M_{B,C}=\Gamma(X_B\times_k\operatorname{Spd}C^\flat,
\mathcal O^\flat)$.  The perfected punctured-ball
exhaustion and \Cref{h3:ana:analytic-limits} give no
higher cohomology on this geometric chart.  On its
$k$-defined coordinates the canonical base Frobenius
is coefficient Frobenius, with any finite residue-field
factor permutation retained.  Thus
\Cref{h3:ana:base-frobenius}, applied in every bar degree
with parameter $H^a$, identifies
$R\Gamma_{\hthreects}(H,B^{\mathrm b})$ with the fiber of
$\Phi-1$ on $R\Gamma_{\hthreects}(H,M_{B,C})$.

The actual marking torsor of
\Cref{h3:geom:finite-level-stack} gives a $k$-defined
equivalence $[X_B/H]\simeq[N_2^*/U]$.
Base-changing this equivalence supplies the vertical
maps in the commutative diagram
\begin{equation}\label{h3:ana:frobenius-square}
\begin{tikzcd}[column sep=large]
 R\Gamma_{\hthreects}(H,M_{B,C})
   \arrow[r,"\Phi-1"]\arrow[d,"\simeq"']
 &R\Gamma_{\hthreects}(H,M_{B,C})\arrow[d,"\simeq"]\\
 R\Gamma_{\hthreects}(U,\mathcal K_C)\arrow[r,"\Phi-1"]
 &R\Gamma_{\hthreects}(U,\mathcal K_C).
\end{tikzcd}
\end{equation}
These maps come from the quotient stacks over $k$.
The diagram is natural for $P_2$-transport between
marking levels and is equivariant for the full
residual norm and the coherent Frobenius actions.
Taking fibers, and commuting the two derived limits,
gives the natural identity
\[
 R\Gamma_{\hthreects}(H,B^{\mathrm b})
      \simeq R\Gamma_{\hthreects}(U,\mathcal K_k).
\]

There are no additional rows in this fiber.
\Cref{h3:prop:punctured-cohomology} gives geometric
cohomology only in degrees zero and three.  In every
degree the fiber sequence gives
\[
 0\longrightarrow
 \operatorname{coker}(\Phi-1\mid H^{i-1}\mathcal K_C)
 \longrightarrow H^i\mathcal K_k
 \longrightarrow
 \ker(\Phi-1\mid H^i\mathcal K_C)
 \longrightarrow0.
\]
The semilinear calculation above makes the two
potentially nonzero cokernels vanish.  Thus
$\mathcal K_k$ has the two stated $k$-rows, with
degree zero equal to $k$ and trivial $P_2$ action.
Smoothness of the upper action can also be seen
directly: if $e$ is a basis, $\Phi(e)=Ae$, and
$u(e)=c(u)e$, commutation gives
$c(u)^q=c(u)$.  Hence $c(u)\in k^\times$.
The scalar norm character $\bar a^{-2}$ survives
this descent since it already takes values in
$\F_p^\times$.

A compact analytic $U\subset P_2$ has finite
completed resolutions: it has no $p$-torsion,
because an order-$p$ element in its degree-two
division algebra would require a subfield of
degree $p-1>2$.  The quotient spectral sequence
with the two finite $k$-rows therefore proves
finite-stage finiteness and identifies degree
zero with $k$.  The resolutions consist of
finite powers and idempotent summands.  The
continuous Artin--Schreier lifting above and
these finite resolutions identify the actual
parameter cohomology with $C_{\hthreects}(T,-)$ of
the point-valued finite groups.
For a line framed on a finite connected-marking
cover, the same argument has a finite-character
$U$-fiber.  The tautological connected frame
trivializes the periodicity line under $H$, and a
tame determinant character is trivial on $H$.
This proves the additional finite-stage assertion.

Finally choose a cofinal sequence of open uniform
subgroups $U$, small enough to act trivially on
these two finite rows.  Restriction to a sufficiently
deep $p$-power subgroup is zero in positive
cohomological degrees: on the mod-$p$ exterior
cohomology of a uniform group it is zero in degree
one and hence in every positive degree.  Filtered
colimits of the quotient spectral sequences
therefore retain precisely the two rows themselves.
The constants arrow gives the stated fiber
sequence.  Its scalar character was computed on
$N_2^*$ in \Cref{h3:prop:punctured-cohomology}.
By determinant normalization in
\Cref{h3:prop:integral-frames}, this scalar is the
\emph{full} reduced norm of an element of $P_3$.
The induced quotient action on $H$-cohomology is
well defined because inner conjugation by $H$
acts trivially there.  Exact averaging in the
prime-to-$p$ norm quotient gives the final assertion.
\end{proof}

\begin{remark}\label{h3:ana:witt-naturality}
The structure map used in
\Cref{h3:prop:ordinary-constants,h3:thm:completed-cohomology}
also defines
\[
 R\Gamma(G_f,W_\ell(k))\longrightarrow
 R\Gamma(Y_1,W_\ell\mathcal O^\flat)
\]
for every $\ell$.  These maps commute with Verschiebung,
Frobenius, truncation, and the maps from
$\Z/p^\ell$-constants.  Their equivalence follows
by induction on the finite Verschiebung filtration
from the length-one equivalence.  Thus the constants
comparison retains the actual frame classes through
the finite Witt comparisons used later.
\end{remark}

\section{Removing completion in coheight one}\label{h3:sec:coheight-one}

The calculation of the completed height-two stratum must be compared with
the algebraic coefficient algebra occurring in ordinary Morava cooperations.
This section proves that comparison.  The two main points are finiteness at
each complete local-field stage and the use of the \emph{full} reduced-norm
quotient.  In particular, averaging over central tame units fails at
$p=7$.

Put
\[
 K=k((y)),\qquad J=P_2,\qquad
 H=\ker(\hthreeNrd:P_3\longrightarrow\Z_p^\times),\qquad
 H'=\hthreeNrd^{-1}(\mu_{p-1}).
\]
Let $L$ be the connected-marking algebra of
\Cref{h3:prop:integral-frames,h3:lem:marking-descent}.  Its commuting $P_3$- and
$J$-actions make it an ind-\'etale $J$-torsor over $K$.
For a normal open subgroup $U\subset J$, write $B_U=L^U$.
It is a finite Galois \'etale $K$-algebra with group $J/U$; in particular,
it is a finite product of complete discretely valued fields.  All field
valuations are normalized by $v(y)=1$.  Statements about valuation ideals
in a product are imposed on every factor.  Both $P_3$ and $J$ preserve
these valuations, allowing permutations of the factors.

We write $C(T,M)=C_{\hthreects}(T,M)$ for continuous functions from a
profinite space $T$ to a topological module $M$, and use this notation
termwise for cochain complexes. We also put
\[
 R=L^{\mathrm{perf}}=\colim_{U,e}B_U^{1/p^e},\qquad
 B_U^b=\widehat{B_U^{\mathrm{perf}}}.
\]
As throughout the paper, a cochain complex with coefficients in an
algebraic union means the filtered colimit of the cochain complexes of
the displayed complete stages.  Thus, for example,
\begin{equation}\label{h3:co:filtered-convention}
 C^\bullet_{\hthreects}(H,R)
 :=\colim_{U,e}C^\bullet_{\hthreects}(H,B_U^{1/p^e}).
\end{equation}
We do not identify this complex with continuous functions into a union
endowed with its subspace metric.  Group cohomology in negative degrees
is understood to be zero.

The input from \Cref{h3:thm:completed-cohomology} has the following precise
form.  Each $H^a_{\hthreects}(H,B_U^b)$ is finite.  Moreover,
\begin{equation}\label{h3:co:completed-input}
 H^a\!\left(\colim_U C^\bullet_{\hthreects}(H,B_U^b)\right)
 =
 \begin{cases}
 k,&a=0,\\
 k(2\epsilon),&a=3,\\
 0,&a\ne0,3,
 \end{cases}
 \qquad \epsilon\in\{1,-1\}.
\end{equation}
The degree-zero map is the constants map.  In the notation $k(2\epsilon)$,
an element of the full quotient $H'/H=\mu_{p-1}$ represented by a norm
unit $u$ acts as $u^{2\epsilon}$.  The comparison, the commuting $J$-action,
and the filtered transition maps are retained with every profinite
parameter space.  No statement about Galois descent of a single completed
infinite field is required below.

\begin{theorem}[Coheight-one comparison]\label{h3:thm:coheight-one-comparison}
For every $p\geq5$, the natural constants map is a quasi-isomorphism
\begin{equation}\label{h3:co:main-constants}
 k\xrightarrow{\ \sim\ }
 \colim_U C^\bullet_{\hthreects}(H',B_U).
\end{equation}
It is equivariant for the commuting connected-frame action and for the
remaining height-three action.  For every profinite space $T$, its
parameterized version
\[
 C(T,k)\xrightarrow{\ \sim\ }
 \colim_U C\bigl(T,C^\bullet_{\hthreects}(H',B_U)\bigr)
\]
is also a quasi-isomorphism.  All maps are the maps induced by constants
and inclusions of coefficient algebras.
\end{theorem}

The proof has three stages.  First, the distribution algebra of the
connected Honda marking torsor and its Cartier residue transpose prove
finiteness at every sufficiently deep complete local-field stage.
Second, Frobenius contracts the positive valuation lattices and compares
the algebraic perfection with its completion.  Finally, the full norm
quotient removes the remaining nonconstant row and the distribution
target.  Every comparison in the conclusion is the original constants
map, even though the proof uses auxiliary operators.

\subsection{Continuous cohomology and finite resolutions}

We first establish the topological facts needed for both the present
comparison and the ordinary-stratum argument of \Cref{h3:sec:ordinary}.
A linearly topologized $\F_p$-vector space below is complete and separated;
an action has invariant neighborhoods if its open linear subspaces have
a cofinal subsystem preserved by the group.

\begin{lemma}[Finite resolutions and continuous duality]
\label{h3:co:compact-resolution}
Let $G$ be a compact $p$-adic analytic group with no element of order
$p$, and put $d=\dim G$.  The trivial module $\F_p$ has a resolution of
length $d$ by finitely generated projective $\F_p[[G]]$-modules.
For every complete linearly topologized representation $M$ with invariant
neighborhoods, applying continuous $\F_p[[G]]$-linear Hom to this
resolution computes $C^\bullet_{\hthreects}(G,M)$ up to a continuous chain
homotopy equivalence.  Its terms are continuous direct summands of finite
powers of $M$.

If $M$ is compact, $H^a_{\hthreects}(G,M)$ is compact Hausdorff.
For $G=H$ or $G=P_3$, put $d_H=8$ and $d_{P_3}=9$.  Their
adjoint orientation characters are trivial, and there are natural
identifications
\begin{equation}\label{h3:co:compact-duality}
 H^a_{\hthreects}(G,M)^\vee
 \cong H^{d_G-a}_{\hthreects}(G,M^\vee),
 \qquad
 M^\vee=\Hom_{\hthreects,\F_p}(M,\F_p),
\end{equation}
where $M^\vee$ has its discrete topology and contragredient action.
The transposed coefficient map induces the dual cohomology map.
The same formula, with compact and discrete sides interchanged, holds
for a discrete continuous $\F_p$-representation and its compact dual.
In particular, finite coefficient modules have finite cohomology.
\end{lemma}

\begin{proof}
The groups needed here satisfy the hypotheses by
\Cref{h3:lem:framework-groups}; their open subgroups inherit the absence
of $p$-torsion and have the same Lie dimension.
For the general assertion, use the Noetherian finite-global-dimension
theorem for the completed
group ring \cite[Proposition~3.3]{AW} and the fact that its integral
dualizing cohomology is a free
rank-one module in degree $d$ and vanishes in the other degrees, with
the adjoint orientation action
\cite[Proposition~4.16 and Remark~4.23]{BGHS}.
These statements give a finite projective resolution over
$\Z_p[[G]]$.  Noetherianity makes its terms finitely generated.  Give each
term its compact topology as a summand of a finite power of the group
ring.  The maps and their images are continuous and closed.  Its syzygies
are $p$-torsion-free, so reduction modulo $p$ is exact.
The integral completed bar resolution is projective in compact modules:
the completed free $\Z_p$-module on a profinite set is compact and
torsion-free, hence projective, since its Pontryagin dual is divisible
and therefore injective.  Induction to $\Z_p[[G]]$ preserves that
projectivity.  The continuous projective comparison therefore shows
that the dual of the finite integral resolution computes the quoted
continuous dualizing cohomology.  Its reduction modulo $p$ does so as
well, because both its terms and its unique integral cohomology module
are $p$-torsion-free.  The reduced dual complex
has cohomology only in degree $d$.  If the reduced resolution extends
beyond $d$, its last dual differential surjects onto a projective module;
splitting and dualizing removes a contractible pair.  Repetition gives
the stated length-$d$ resolution.

Here is why the resolution computes the asserted topological cochains.
For a profinite space $Z$, continuous maps $Z\to M$ identify with
continuous linear maps $\F_p[[Z]]\to M$.  Modulo an invariant open
subspace of $M$, the map has finite image and its linear extension
factors through a finite clopen partition of $Z$.  Taking the inverse
limit over these subspaces gives the claimed extension into $M$.
The completed bar resolution is projective in compact modules: a
surjection of compact $\F_p$-vector spaces admits a continuous linear
section, obtained by dualizing and splitting the corresponding injection
of discrete vector spaces.  Such a section lifts the profinite set of
bar generators; the lift extends group-ring linearly.  The projective
comparison construction therefore gives continuous chain maps and
homotopies between the completed bar resolution and the finite
resolution.  Applying continuous Hom into $M$ proves the assertion.
Continuity for the uniform cochain topologies can be checked after
quotienting by each invariant open subspace.

For compact coefficients, the finite complex has compact terms and
closed boundaries, which proves the compact Hausdorff assertion.
Now take $G=H$ or $G=P_3$.  The adjoint orientation of each is trivial
by \Cref{h3:lem:framework-groups}: conjugation has determinant one on
$D_3$, and also on its trace-zero summand.  Thus the dualizing module
has the trivial right action, in degree $d_G=8$ or $9$ respectively.

Let $Q_\bullet$ be the finite resolution for this $G$, and set
$Q_i^*=\Hom_{\F_p[[G]]}(Q_i,\F_p[[G]])$.
The reversed group-ring dual resolves the trivial right module, by
the dualizing calculation above.  Regard $Q_i^*$ as a left module
through the involution $g\mapsto g^{-1}$.  Finite-projective evaluation
gives
\[
 \Hom_{\hthreects,\F_p[[G]]}(Q_i,M)^\vee
 \cong
 \Hom_{\hthreects,\F_p[[G]]}(Q_i^*,M^\vee).
\]
For a finite free module this is evaluation coordinate by coordinate;
passing to a projective summand preserves the identification.  It
commutes with both differentials and coefficient maps.  Continuous
$\F_p$-duality is exact on compact vector spaces, as follows by extending
a functional on a finite-dimensional quotient.  Reversing degrees about
$d_G$ and taking cohomology proves \eqref{h3:co:compact-duality} and its
naturality for both groups.  For a discrete continuous representation,
each element has a finite $G$-orbit; finite sets therefore lie in finite
$G$-stable submodules.  Their annihilators give invariant open
neighborhoods in the compact dual.  Dualizing the compact formula again
proves the discrete version.  Finally, the finite-resolution complex
has finite terms when $M$ is finite.
\end{proof}

Here the coefficient actions have the required invariant neighborhoods.
For the local-field and root stages, valuation balls are invariant.
For the ordinary spaces $B_{\hthreehol}$ and $B_{\hthreepk}$ of
\Cref{h3:sec:ordinary}, write $A=k[[x,y]]$ and use the valuations
$v_{a,b}$ with $v_{a,b}(x)=a$ and $v_{a,b}(y)=b$.  Since
$g(x)=xu_g$ with $u_g\in A^\times$ and $g(y)\in(x,y)$,
\begin{equation}\label{h3:co:uniform-action-bound}
 v_{a,b}(gf)\geq v_{a,\min(a,b)}(f)
 \qquad(g\in P_3,\ a,b>0).
\end{equation}
Indeed units and their inverses have valuation zero, so the estimate
holds termwise on Laurent polynomials.  Density extends it to
$B_{\hthreehol}$, then to its finite root levels and $B_{\hthreepk}$.
The coefficient lines used in \Cref{h3:sec:ordinary} have unit action
multipliers in their free $A$-lattices, so the same estimate holds
in those frames.  Given a basic open linear neighborhood $V$,
\eqref{h3:co:uniform-action-bound} supplies a neighborhood contained in
$\bigcap_{g\in P_3}g^{-1}V$.  This intersection is therefore open,
invariant, and contained in $V$.  The bound also holds uniformly for
continuous functions on a compact cochain domain, with any profinite
parameter space.  Thus these spaces satisfy the hypothesis of
\Cref{h3:co:compact-resolution} for the actual uniform cochain topologies.

An exact coefficient sequence induces an exact sequence of ordinary
continuous cochain complexes whenever its surjection has a continuous
section as a map of spaces.  Equivariance and additivity of the section
are unnecessary: compose a cochain with the section.  We will use this
for discrete valuation quotients and for finite-dimensional linear maps
over complete local fields, including their Frobenius transports.  In a
filtered union it suffices that every target stage have a continuous
lift at one larger source stage, and that each stage's kernel lie in a
kernel stage with its subspace topology.

\begin{lemma}[Strictness]\label{h3:co:strictness}
Let $A^\bullet$ be a complex of Polish abelian groups with continuous
differentials.  If $H^a(A^\bullet)$ is finite, its quotient topology is
discrete.  More precisely, the boundary subgroup $B^a$ is open and closed
in the cycle group $Z^a$, and
$\partial:A^{a-1}\twoheadrightarrow B^a$ is open.  Consequently, for
every neighborhood $V$ of zero in $A^{a-1}$ there is a neighborhood $W$
of zero in $Z^a$ such that every member of $W$ has a primitive in $V$.
\end{lemma}

\begin{proof}
The cycles form a closed Polish group.  The boundary subgroup is an
analytic subset of it, being the continuous image of a Polish space,
and hence has the Baire property \cite[Theorem~21.6]{Kechris}.
Its finite index implies that it is nonmeagre.  A nonmeagre set with the
Baire property is comeagre in some nonempty open set; intersecting two
sufficiently small translates of that open set shows that its difference
set contains a neighborhood of zero.  Applied to the subgroup $B^a$,
this proves openness, and an open subgroup is also closed.

We also recall the open-mapping step.  If $q:E\twoheadrightarrow F$ is
a continuous surjective homomorphism of Polish groups and $V$ is a
neighborhood of zero in $E$, choose open $V'$ with $V'-V'\subset V$.
Separability gives a countable covering of $E$ by translates of $V'$.
Thus $q(V')$ is an analytic nonmeagre subset of $F$, and its difference
set contains a neighborhood of zero.  This difference set lies in
$q(V)$.  Hence $q$ is open.  Apply this to the differential with target
$B^a$, which is already open in $Z^a$.
\end{proof}

\begin{lemma}[Profinite parameters]\label{h3:co:parameters}
Let $T$ be a profinite space.
\begin{enumerate}
\item A continuous map $T\to F$ lifts through every continuous open
surjective homomorphism $E\twoheadrightarrow F$ of Polish abelian groups.
\item If a Polish complex $A^\bullet$ has finite cohomology in every
degree, then
\[
 H^a C(T,A^\bullet)\cong C(T,H^a(A^\bullet)),
\]
where the finite cohomology group on the right is discrete.
\item For a filtered system of such complexes, the corresponding
identity commutes with the cochain colimit, with the colimit of the
cohomology groups given its discrete topology.
\end{enumerate}
\end{lemma}

\begin{proof}
For the first assertion, choose complete translation-invariant metrics
on the groups and successively smaller neighborhoods of zero in $E$
whose diameters tend to zero and whose sums converge.  Openness provides
corresponding neighborhoods in $F$.  Compactness and zero-dimensionality
give a finite clopen partition on which the given map is approximated
by finitely many chosen images.  Lift those values to $E$.
Repeat on the error, refining the partition and choosing the correction
inside the next prescribed neighborhood of zero in $E$.  The successive
locally constant lifts converge uniformly to a continuous lift; their
images converge to the original map.  This construction uses no
metrizability assumption on $T$.

A continuous family of cycles gives a continuous family of cohomology
classes.  If that family is zero, the first assertion and
\Cref{h3:co:strictness} lift it continuously through
$A^{a-1}\twoheadrightarrow B^a$.  Conversely, a continuous map from
$T$ to the finite discrete cohomology group has finite image; choose a
cycle for each value to obtain a continuous family.  This proves the
second assertion.  For the third, a continuous function to a discrete
filtered colimit has finite image.  Its finitely many values and equality
relations are represented at one stage.  Filtered colimits of vector
spaces are exact, which proves the claimed compatibility.
\end{proof}

All coefficient spaces to which \Cref{h3:co:strictness} is applied are
Polish.  A finite product of local fields with finite residue fields is
Polish.  Its perfected completion has a countable dense subset obtained
from the finite root levels.  For a compact metrizable profinite space
$Z$, $C(Z,M)$ is Polish when $M$ is Polish: uniform limits give
completeness, and finite clopen partitions with a countable dense subset
of $M$ give separability.  The groups $H^a$ serving as cochain domains
are compact metrizable.  The parameterized complex $C(T,A^\bullet)$
itself is not asserted to be Polish for arbitrary $T$.

We shall also use the following elementary countability fact.  If $G$ is
compact and second-countable and $N$ is countable discrete, then every
$C^a_{\hthreects}(G,N)$ is countable.  Indeed $G^a$ has countably many clopen
sets, and a cochain is specified by a finite clopen partition and
finitely many values.  For discrete coefficients the same finite-image
description shows that continuous cochains commute with filtered
colimits, even when the transition maps are not injective.

\subsection{Distribution operators and the invariant differential}

The normalized generator-lift torsors of
\Cref{h3:lem:marking-descent} are torsors over $L$
under the constant group schemes $\Gamma_2[p^l]$.  Their coordinate
algebras and transition embeddings are
\begin{equation}\label{h3:co:root-torsors}
 R_l=L^{1/p^{2l}},\qquad R_l\longrightarrow R_{l+1}.
\end{equation}
Here the transition on the torsors is multiplication by $p$.
The group $H$ commutes with translations, and $J$ conjugates them by
continuous constant automorphisms.  These assertions concern the finite
flat torsors, not just their geometric points.  In particular, their
radicial rank $p^{2l}$ is retained under the separable connected-marking
extension.  Before perfection, \Cref{h3:geom:regular-levels} gives the
regular finite algebras of this rank, and \Cref{h3:lem:marking-descent}
identifies their generator-lift algebra after the separable base change.

\begin{lemma}[Distribution parameter]\label{h3:co:distribution}
After a finite enlargement of $k$, the inverse-limit distribution
algebra of \eqref{h3:co:root-torsors} is $k[[D]]$, with a parameter $D$
having the following properties.
\begin{enumerate}
\item It acts $L$-linearly and locally nilpotently on $R$.  For every
$i\geq0$,
\[
 \ker(D^{p^i}:R\to R)=L^{1/p^i}.
\]
It is surjective, and $D^{p^i-1}:L^{1/p^i}\to L$ is nonzero on every
factor.
\item It commutes with $H$.  Its residual $\mu_{p-1}$-character is
$\chi(u)=u^{s_{\mathrm{dist}}}$, where $s_{\mathrm{dist}}\in\{1,-1\}$.  Thus
\begin{equation}\label{h3:co:distribution-sequence}
 0\longrightarrow L\longrightarrow R
 \xrightarrow{D} R(\chi^{-1})\longrightarrow0
\end{equation}
is an exact sequence of $H'$-modules.
\item If $q=p^i$, where $i$ is even and divisible by $[k:\F_p]$, then
\begin{equation}\label{h3:co:frobenius-distribution}
 D(f^q)=(D^qf)^q,\qquad
 D^q(z^{1/q})=(Dz)^{1/q}.
\end{equation}
In particular $D^q$ is $L^{1/q}$-linear.
\item On any complete finite separable and root stage these operators
are continuous after passage to a larger such stage.  Each target stage
admits a continuous additive right inverse of $D$ with values in one
larger stage.  Therefore \eqref{h3:co:distribution-sequence} is exact on
the filtered continuous cochain complexes, with profinite parameters.
\end{enumerate}
\end{lemma}

\begin{proof}
The Cartier dual of the height-two, dimension-one Honda group is again
connected of height two and dimension one. Indeed, on the Honda model
$\Phi^2=[p]$, so duality gives $V^2=[p]$ on its dual. Together with
$\Phi V=[p]$, cancellation on the divisible group gives $\Phi=V$;
thus each finite dual torsion kernel is killed by a Frobenius iterate
and is connected. The dimension formula and formal-group construction
in \cite[Section~2.2, Proposition~1, and Section~2.3, Proposition~3]{Tate1967PDivisible}
give dimension one and a power-series coordinate, already over $\F_p$.
No identification with a chosen coordinate law on the original Honda
group is needed. A parameter on the dual identifies
the finite distribution algebra with
\[
 \mathcal O(\Gamma_2[p^l])^*
 =k[D]/(D^{p^{2l}}).
\]
The duals of the torsor transition maps are the inclusions of the dual
torsion kernels.  Their restriction maps give the inverse limit
$k[[D]]$.

After a faithfully flat trivialization of a torsor, its functions form
the dual regular module of this truncated polynomial ring.  The pairing
given by the coefficient of $D^{p^{2l}-1}$ in a product is perfect, so
that module is one regular nilpotent block.  Ranks descend under faithful
flatness.  Consequently the kernel of $D^e$ on $R_l$ has rank $e$ over
$L$ for $1\leq e\leq p^{2l}$.

For $q=p^i$, the formal-group coproduct satisfies
\[
 \Delta(D^q)\in(D^q\otimes1,1\otimes D^q).
\]
The module-algebra identity therefore makes $\ker D^q$ a subalgebra.
On a field factor it is an intermediate field of degree $q$.
The element $y$ is a $p$-basis of $K$ and remains a $p$-basis after
separable extension: if $E/K$ is finite separable, then
$[E:E^p]=[K:K^p]=p$, and $y\notin E^p$ because a separable extension
contains no nontrivial purely inseparable subextension of $K$.
An intermediate field of degree $q$ in $E^{1/p^{2l}}/E$ lies in
$E^{1/q}$, since every element has purely inseparable degree at most
$q$; equality follows from the degrees.  Applying this on every factor
and then passing to $L$ proves the kernel formula.  The block calculation
also proves nonvanishing of the top operator.  In a larger block the
old kernel is contained in the image of $D$, which proves surjectivity
on $R$.

By \Cref{h3:lem:marking-descent}, a full reduced-norm unit $u$ rescales the
normalized \'etale generator by $u$ or $u^{-1}$.  On the fixed connected
framing cover this conjugates translations by the corresponding scalar
automorphism.  Cartier duality preserves scalar multiplication, so the
dual tangent character is $\chi(u)=u^{s_{\mathrm{dist}}}$ with $s_{\mathrm{dist}}=\pm1$.
Since $H$ acts trivially on translations, this is an action of the
quotient $H'/H$, whether or not that quotient has a subgroup section.
Starting with a parameter $D_0$, replace it by
\[
 D=\frac1{p-1}\sum_{u\in\mu_{p-1}}
       \chi(u)^{-1}u(D_0).
\]
Its linear term is the linear term of $D_0$, and it is an eigenparameter.
This proves the second assertion, including the inverse twist on the
target of \eqref{h3:co:distribution-sequence}.

For the Frobenius assertion, the Honda model over $\F_p$ satisfies
$\Phi^2=[p]$.  On the full divisible group,
$V\Phi=[p]=\Phi^2$; cancellation of the faithfully flat epimorphism
$\Phi$ gives $V=\Phi$.  Cartier duality exchanges the two operators.
For the stated even $i$, their $i$th iterates are $[p^{i/2}]$.
To make the scalar and duality conventions explicit, put
$A_l=\mathcal O(\Gamma_2[p^l])$ and pair it with its finite dual by
$\langle b,\xi\rangle=\xi(b)$.  Since $q$ fixes $k$, Cartier-dual
naturality for the $i$th relative Frobenius reads
\begin{equation}\label{h3:co:frobenius-dual-pairing}
 \langle b^q,D\rangle
   =\langle b,(V^i)^*D\rangle
   =\langle b,D^q\rangle.
\end{equation}
For the second equality, choose an $\F_p$-coordinate $t$ on the Cartier dual,
which descends to $\F_p$.  Tame averaging leaves $D=\sum_n c_nt^n$ defined over
$k$, and hence
$(V^i)^*D=\sum_n c_nt^{nq}=D^q$ because $c_n^q=c_n$.
Here multiplication in the distribution algebra corresponds to
composition of operators.  If a finite torsor coaction is written
$\rho(f)=\sum_j f_j\otimes b_j$, \eqref{h3:co:frobenius-dual-pairing}
therefore gives
\[
 D(f^q)=\sum_j f_j^qD(b_j^q)
       =\sum_j f_j^qD^q(b_j)
       =\left(\sum_j f_jD^q(b_j)\right)^q.
\]
The last equality uses $q$-power invariance of the constants.
The coefficients $f_j$ retain their $q$th powers throughout; this finite
Hopf-coaction calculation requires no splitting of the torsor.
This proves \eqref{h3:co:frobenius-distribution}.  If $a\in L^{1/q}$,
apply it to $af$ and use $a^q\in L$ and $L$-linearity of $D$ to obtain
$D^q(af)=aD^qf$.  The cancellation above was on a divisible group;
no cancellation on a finite torsion kernel has been used.

On a fixed root level only finitely many powers of $D$ act.  In the
basis of successive fractional powers of the $p$-basis $y$, their
matrices have finitely many coefficients in $L$, and these lie in one
finite separable stage.  The resulting matrices between finite-dimensional
spaces over complete local fields are continuous.  To construct a
right inverse on a stage, choose a $D$-preimage of each member of its
finite field basis and put those preimages in one larger complete
separable/root stage.  Linear extension gives a continuous additive
section.  The kernel at a finite stage has the required subspace
topology; intersection with the separable union is that stage's
separable part, by uniqueness of purely inseparable roots.  The
cochain-exactness criterion preceding \Cref{h3:co:strictness} proves the
last assertion.  All of these arguments are factorwise for a finite
product, with finitely many choices at each stage.
\end{proof}

The injection of $L$ into the filtered completed coefficient algebra,
together with degree zero of \eqref{h3:co:completed-input}, gives
\begin{equation}\label{h3:co:fixed-constants}
 L^H=k.
\end{equation}
Indeed an invariant element belongs to an $H$-stable finite stage;
completion injects on its invariants, and filtered degree-zero cohomology
of the target is exactly $k$.

The $p$-basis calculation also identifies ordinary and continuous
differentials: on a factor $\kappa((s))$ both are freely generated by
$ds$, since derivations kill $p$th powers and the constant field is
perfect.  Thus Cartier below is the intrinsic field operator, with
its usual continuous Laurent-series formula.

Let $\eta$ be the nowhere vanishing differential of
\Cref{h3:lem:invariant-differential}.  It is defined at a finite separable
stage and is $H$-invariant.  The canonical-line and conormal
identifications used there give, on this slice,
\[
 \Omega^1_{K/k}\cong\omega_K^{\otimes(p+2)}[\det].
\]
The connected marking trivializes $\omega_L$, and the full norm kernel
fixes the determinant factor.  We write $\mathrm{Car}$ for Cartier on
differentials and set
\begin{equation}\label{h3:co:cartier-definition}
 \mathcal C(f)=\frac{\mathrm{Car}(f\eta)}{\eta}.
\end{equation}
This is the coheight-one Cartier operator.  It is distinct from the
ordinary root projection $\mathcal P$ of \Cref{h3:sec:ordinary}.

\begin{lemma}[Cartier and distribution traces]\label{h3:co:cartier-trace}
For every $i\geq1$, put $q=p^i$.  There exists $c_i\in k^\times$ such that
\begin{equation}\label{h3:co:cartier-distribution-trace}
 D^{q-1}(f^{1/q})=c_i\mathcal C^i(f)\qquad(f\in L).
\end{equation}
In particular,
\begin{equation}\label{h3:co:cartier-frobenius-zero}
 \mathcal C(f^p)=0.
\end{equation}
At every finite separable stage $B$ containing $\eta$ there is an exact
sequence of continuous $\F_p[H]$-modules
\begin{equation}\label{h3:co:cartier-four-term}
 0\longrightarrow B\xrightarrow{\mathrm{Fr}}B
 \xrightarrow{d/\eta}B\xrightarrow{\mathcal C}B\longrightarrow0.
\end{equation}
Here and throughout this section, \(\mathrm{Fr}(f)=f^p\) is coefficient
\(p\)-Frobenius.  The ordinary-stratum argument in
\Cref{h3:sec:ordinary} separately fixes a \(q\)-power operator.
The two short exact factors of \eqref{h3:co:cartier-four-term} admit
continuous additive sections onto their images.  Under the residue pairing
\begin{equation}\label{h3:co:residue-pairing}
 \langle f,g\rangle_B
 =\sum_j\hthreeTr_{\kappa_j/\F_p}\hthreeRes_{s_j}(fg\eta),
 \qquad B=\prod_j\kappa_j((s_j)),
\end{equation}
the $H$-action is isometric and
$\langle\mathcal C f,g\rangle_B=\langle f,g^p\rangle_B$.
\end{lemma}

\begin{proof}
The functional $I_i(f)=D^{q-1}(f^{1/q})$ takes values in $L$, is nonzero
on every factor, is $H$-equivariant, and satisfies
$I_i(a^qf)=aI_i(f)$.  The same semilinearity and equivariance hold for
$\mathcal C^i$.  The Cartier pairing gives an isomorphism
\begin{equation}\label{h3:co:perfect-cartier-pairing}
 \mathrm{Fr}^i_*L\longrightarrow
 \Hom_L(\mathrm{Fr}^i_*L,L),\qquad
 b\longmapsto\bigl[f\longmapsto\mathcal C^i(bf)\bigr].
\end{equation}
Here $a$ acts on $\mathrm{Fr}^i_*L$ as multiplication by $a^q$.
For a field factor, both sides have dimension $q$; the map is injective,
since for $b\ne0$ one can choose $z$ with $\mathcal C^i(z)\ne0$ and
evaluate at $b^{-1}z$.  This proves the isomorphism on products and
separable filtered unions as well; equivalently it is the perfect
coefficient-extraction pairing in the common $p$-basis.
Thus there is a unique multiplier $b_i$ with
$I_i(f)=\mathcal C^i(b_if)$.  It is a unit because both functionals
are nonzero on every field factor.  Equivariance and uniqueness give
$b_i\in L^H=k$, by \eqref{h3:co:fixed-constants}.  Semilinearity gives
\eqref{h3:co:cartier-distribution-trace}, with $c_i=b_i^{1/q}$.

This comparison applies in particular when $i=1$; it does not require
the even-iterate identity \eqref{h3:co:frobenius-distribution}.  Since $D$
kills $L$, it yields
$c_1\mathcal C(f^p)=D^{p-1}f=0$, proving
\eqref{h3:co:cartier-frobenius-zero}.  The latter is therefore a consequence
of the invariant torsor and differential, not an identity for an
arbitrary normalization of Cartier.

On a factor $\kappa((s))$ the logarithmic Laurent formula is
\begin{equation}\label{h3:co:cartier-laurent}
 \mathrm{Car}\!\left(\sum_n a_ns^n\frac{ds}{s}\right)
 =\sum_n a_{pn}^{1/p}s^n\frac{ds}{s}.
\end{equation}
It proves continuity and surjectivity.  The kernel consists of the forms
whose coefficients in indices divisible by $p$ vanish; a continuous
additive primitive is $\sum_{p\nmid n}(a_n/n)s^n$.
The kernel of $d$ is $B^p$, and its root map is continuous.
A continuous section of Cartier sends
$\sum_n a_ns^n ds/s$ to $\sum_n a_n^ps^{pn}ds/s$.
Multiplication or division by the fixed nonzero $\eta$ transports these
sections to \eqref{h3:co:cartier-four-term}.  The maps themselves are
$H$-equivariant because $\eta$ is invariant and Cartier is intrinsic.

Residue is invariant under change of uniformizer.  The sum of finite
field traces is preserved by residue-field automorphisms and factor
permutations.  This proves $H$-invariance of
\eqref{h3:co:residue-pairing}.  Finally,
$\hthreeRes\mathrm{Car}(\alpha)=(\hthreeRes\alpha)^{1/p}$ and
$g\mathrm{Car}(\alpha)=\mathrm{Car}(g^p\alpha)$.
The finite-field trace is unchanged by $p$th roots.  Applying these
identities to $\alpha=f\eta$ proves the transposition formula.
\end{proof}

\begin{lemma}[Good finite stages]\label{h3:co:good-stages}
There is a cofinal family of normal open subgroups $U\subset J$ such
that, for $B=B_U$, the differential $\eta$ is defined over $B$, there is
$\theta\in B^{1/p}$ with $D\theta=1$, and
\begin{equation}\label{h3:co:good-jet}
 \sigma(D)-D\in D^{p+2}k[[D]]\qquad(\sigma\in U).
\end{equation}
The operator $D$ restricts continuously to $B^{1/p}$.
Writing $\eta=h_j(s_j)ds_j/s_j$ on each field factor, the lattice
\begin{equation}\label{h3:co:positive-differential-lattice}
 P_B=\prod_j h_j^{-1}s_j\kappa_j[[s_j]]
\end{equation}
is compact, open, and preserved by $H$ and $\mathcal C$.
Its annihilator for \eqref{h3:co:residue-pairing} is $\mathcal O_B$, and
the annihilator of $\mathcal O_B$ is $P_B$.
\end{lemma}

\begin{proof}
The nilpotent block on $L^{1/p}$ puts $1$ in the image of $D$, so choose
$\theta$ there.  The finite data $\theta^p$ and $\eta$ lie at a finite
separable stage.  The action of $J$ on the finite set of jets
$k[[D]]/(D^{p+2})$ is continuous.  Intersect their open stabilizers with
any prescribed open subgroup and pass to a normal core.  This proves
cofinality and \eqref{h3:co:good-jet}.  For $f\in B^{1/p}$, $D^pf=0$ and
\[
 \sigma(Df)-Df=(\sigma(D)-D)f=0.
\]
The kernel filtration and uniqueness of roots therefore put $Df$ in
$(L^{1/p})^U=B^{1/p}$.  Its finite matrix over the complete stage proves
continuity.

The condition defining $P_B$ says that $f\eta$ is an integral ordinary
differential on each factor.  That condition is independent of the
uniformizer and is preserved by $H$.  Formula
\eqref{h3:co:cartier-laurent} preserves strictly positive logarithmic
order, proving Cartier stability.  The annihilator statements follow
coefficient by coefficient: an integral $g$ pairs trivially with all
strictly positive logarithmic coefficients, while every pole of $g$
is detected by a suitable positive coefficient of $f\eta$.  Conversely
every nonpositive coefficient of $f\eta$ is detected by a monomial in
$\mathcal O_B$.  Nondegeneracy of the finite-field trace makes these
detections nonzero when needed.
\end{proof}

\subsection{Finiteness at the local-field stages}

We now use the distribution and Cartier operators to prove the missing
finite-stage bound.  The only finiteness assumed at this point is for
the completed perfected fields.  Compact residue duality first controls
the stable Cartier image of a lattice; a separate Frobenius-kernel
argument then forces the uncompleted field cohomology to be finite.

An order lattice in a finite product of local fields means a product of
fractional powers of their maximal ideals.

\begin{proposition}[Finite-stage finiteness]\label{h3:prop:finite-stage-finiteness}
Let $B=B_U$ be a good stage as in \Cref{h3:co:good-stages}.
For every $a\geq0$, each of the following cohomology groups is finite:
\[
 H^a_{\hthreects}(H,B),\qquad
 H^a_{\hthreects}(H,\Lambda),\qquad
 H^a_{\hthreects}(H,B/\Lambda),
\]
where $\Lambda$ is any $H$-stable order lattice in $B$.
The assertion holds at each finite root stage and for a coefficient line
which has an $H$-invariant frame at a finite separable stage, with its
order lattices and their quotients.  In particular it holds for the
periodicity and determinant lines used in this paper after choosing a
sufficiently deep stage.
\end{proposition}

We prove the proposition without assuming finite cohomology of the
uncompleted fields.

\begin{lemma}[Positive completed coefficients]\label{h3:co:completed-positive}
Let $B$ be a good stage.  All $H^a_{\hthreects}(H,B^b)$ are finite, and
\[
 R\Gamma_{\hthreects}(H,(B^b)^{++})=0,
 \qquad (B^b)^{++}=\{f:v(f)>0\}.
\]
The cohomology groups of $\mathcal O_{B^b}$ and
$B^b/\mathcal O_{B^b}$ are finite as well.
\end{lemma}

\begin{proof}
The first assertion is the finite-stage statement of
\Cref{h3:thm:completed-cohomology}.  Let $z$ be a continuous positive-valued
$a$-cocycle.  Compactness of $H^a$ supplies $\varepsilon>0$ with
$v(z)\geq\varepsilon$ on the entire cochain domain.  Hence
$z^{p^n}\to0$ uniformly.  For $a>0$, apply \Cref{h3:co:strictness} to the
complete coefficient complex, prescribing the open neighborhood of
positive-valued $(a-1)$-cochains.  A sufficiently high Frobenius power
has a primitive $b$ in that neighborhood.  Inverse Frobenius is a
continuous additive automorphism of the completed perfect algebra and
preserves positivity, so $b^{1/p^n}$ is a positive primitive of $z$.
For $a=0$, the finite invariant group is discrete; the invariant elements
$z^{p^n}$ tend to zero and are eventually zero.  Injectivity of Frobenius
then gives $z=0$.

The quotient of $\mathcal O_{B^b}$ by its strictly positive ideal is the
finite product of the residue fields of $B$.  The two valuation-cutoff
sequences are exact on cochains because their quotients are discrete.
Their long exact sequences, the vanishing just proved, and finite
cohomology of finite coefficients establish the last assertions.
\end{proof}

Fix a good $B$ and put
\[
 P=P_B,\qquad M_a=H^a_{\hthreects}(H,P),\qquad
 X_a=H^a_{\hthreects}(H,B),\qquad f_a:M_a\longrightarrow X_a.
\]
The group $M_a$ is compact Hausdorff by
\Cref{h3:co:compact-resolution}.  Since $B/P$ is countable discrete, the
countability observation above and the cochain-exact lattice sequence
give
\begin{equation}\label{h3:co:countable-kernel-cokernel}
 \ker f_a\ \text{and}\ \operatorname{coker}f_a
 \quad\text{are countable.}
\end{equation}

We first prove that every $X_a$ is countable.  If $a$ were the largest
degree where countability failed, the Cartier exact sequences would make
$M_a/\mathcal C M_a$ finite.  The next two lemmas combine this bound with
a finite stable Cartier image to control the Cartier kernel in $X_a$.
Since $\mathcal C\mathrm{Fr}=0$, the Frobenius image would then be
countable; the following kernel estimate makes its kernel countable as
well, a contradiction.  Once every $X_a$ is countable, the comparison
$f_a$ makes each compact $M_a$ finite.  A second use of residue duality
gives finite cohomology of the lattice quotient, and the lattice sequence
then makes $X_a$ finite.

\begin{lemma}[Finite stable Cartier image]\label{h3:co:stable-image}
For every $a$, the subgroup
\[
 I_a=\bigcap_{n\geq0}\mathcal C^nM_a
\]
is finite.
\end{lemma}

\begin{proof}
The residue pairing and \Cref{h3:co:good-stages} give topological
identifications
\begin{equation}\label{h3:co:residue-duals}
 P^\vee=B/\mathcal O_B,\qquad (B/P)^\vee=\mathcal O_B.
\end{equation}
A continuous functional on $P$ uses only finitely many Laurent
coefficients and is represented by a finite principal part in the first
quotient.  An arbitrary functional on the discrete $B/P$ is represented
by a power series, which proves the second identification and its
topology.  The trace transposition in \Cref{h3:co:cartier-trace} and
\eqref{h3:co:compact-duality} identify the dual of the inverse system
$(M_a,\mathcal C)$ with
\[
 \bigl(H^{8-a}_{\hthreects}(H,B/\mathcal O_B),\mathrm{Fr}\bigr).
\]
Its coefficient colimit is
\begin{equation}\label{h3:co:discrete-perfection-quotient}
 \colim_{\mathrm{Fr}}B/\mathcal O_B
 =B^{\mathrm{perf}}/\mathcal O_{B^{\mathrm{perf}}}
 =B^b/\mathcal O_{B^b}.
\end{equation}
At index $n$ the first map sends $f$ to $f^{1/p^n}$.
The second equality follows by approximating any completed element with
integral error.  All these quotients are discrete, so their continuous
cochains commute with the colimit.  By
\Cref{h3:co:completed-positive}, the direct limit of the displayed
cohomology system is finite.  Compact--discrete duality shows that
$\varprojlim_{\mathcal C}M_a$ is finite.

Its zeroth projection has image $I_a$.  To prove surjectivity onto that
intersection, fix $x\in I_a$.  For each $n$ there is a compatible
length-$n$ chain of preimages ending at $x$.  These conditions define
nested nonempty closed subsets of the compact product
$\prod_{n\geq0}M_a$.  Their intersection is an infinite compatible
chain.  Thus $I_a$ is an image of the finite inverse limit.
\end{proof}

\begin{lemma}[A compact Cartier module]\label{h3:co:compact-cartier}
Let $M$ be a compact Hausdorff $\F_p$-vector space with a continuous
endomorphism $C$.  Suppose that $I=\bigcap_n C^nM$ and $M/CM$ are finite.
Then $C(I)=I$, the quotient $Q=M/I$ is a finitely generated
$\F_p[[t]]$-module with $t$ acting as $C$, and both $Q[C^\infty]$ and
$\ker(C:M\to M)$ are finite.
If a $C$-stable subgroup $N\subset M$ has finite forward $C$-orbits
elementwise, then $N$ is finite and
\begin{equation}\label{h3:co:stable-quotient}
 \bigcap_n\hthreeim(C^nM\longrightarrow M/N)
 =\hthreeim(I\longrightarrow M/N).
\end{equation}
\end{lemma}

\begin{proof}
For $x\in I$, the sets $C^{-1}(x)\cap C^nM$ are nested nonempty compact
sets, so they meet.  This proves $C(I)=I$.
The compact images $C^nQ$ have intersection zero and shrink uniformly
to zero: otherwise their intersections with the complement of a fixed
open neighborhood would be nested nonempty compact sets.  Therefore
$\sum_{n\geq0}a_n C^n x$ converges for every formal power series
$\sum a_nt^n$, uniformly in the tails, and defines an
$\F_p[[t]]$-action on $Q$.

Lift a finite basis of $Q/CQ$ to $e_1,\dots,e_s$.  Repeatedly expressing
an element modulo $CQ$ gives
\[
 x=\sum_{i=1}^s\sum_{j=0}^{N-1}a_{ij}C^je_i+C^Nx_N.
\]
The remainder tends uniformly to zero.  Taking the limit expresses $x$
as a sum of $s$ power-series multiples of the $e_i$.  The structure
theorem over the discrete valuation ring $\F_p[[t]]$ now writes $Q$ as
a finite free module plus finitely many modules $\F_p[[t]]/(t^e)$.
Its $C$-power torsion and $C$-kernel are finite.  The kernel on $M$ has
kernel in $I$ and image in the kernel on $Q$, so it is finite too.

The image in $Q$ of an element with finite forward orbit both tends
to zero and belongs to a finite set; it is eventually zero.  Thus the
image of $N$ is in the finite group $Q[C^\infty]$, and its kernel lies
in finite $I$.  Hence $N$ is finite and closed.  A coset meeting every
$C^nM$ meets their intersection, by compactness of its nested
intersections with those images.  This proves
\eqref{h3:co:stable-quotient}.
\end{proof}

We record one algebraic consequence used in applying this lemma.
Suppose $f:M\to X$ commutes with $C$ and has countable kernel and
cokernel.  If $X/CX$ is countable, so is $M/CM$.  Indeed, with
$N=\ker f$, $A=\hthreeim f$, and $E=X/A$, the exact portions
\[
 N/CN\longrightarrow M/CM\longrightarrow A/CA\longrightarrow0,
 \qquad
 \ker(C|E)\longrightarrow A/CA\longrightarrow X/CX
\]
prove the assertion.  If $N$ and $\ker(C|M)$ are finite, then
$\ker(C|A)$ is finite, by its exact sequence with $N/CN$, and
$\ker(C|X)$ is countable since its remaining quotient embeds in $E$.

\begin{lemma}[Frobenius-kernel bound]\label{h3:co:frobenius-kernel}
If $\bigcap_n\hthreeim(\mathcal C^nM_a\to X_a)$ is finite, then
$\ker(\mathrm{Fr}:X_a\to X_a)$ is countable.
\end{lemma}

\begin{proof}
For $a=0$, Frobenius is injective.  For $a>0$, its kernel is the kernel
of the inclusion into the root-stage cohomology under the
$H$-equivariant homeomorphism $B^{1/p}\to B$.  Consequently each kernel
class has a presentation $[\partial b]$, where
\[
 b\in\mathcal B_a
 :=\{b\in C^{a-1}_{\hthreects}(H,B^{1/p}):
               \partial b\in C^a_{\hthreects}(H,B)\}.
\]
Put $w=Db$.  It is a $B^{1/p}$-valued cocycle because $D$ kills $B$.
The quotient map
\[
 \mathcal B_a\longrightarrow
 C^{a-1}_{\hthreects}\bigl(H,B^{1/p}/(B^{1/p})^{++}\bigr),\qquad
 b\longmapsto Db\bmod ++,
\]
has countable target.  Its kernel $\mathcal B_a^+$ consists of
presentations with positive-valued $w$ and has countable index.
It suffices to put the classes presented by this kernel into the finite
intersection in the hypothesis.

Fix such a $b$, and let $i$ range through arbitrarily large positive
integers divisible by two and by $[k:\F_p]$; put $q=p^i$.
For $\theta$ in \Cref{h3:co:good-stages}, define
\[
 b'=\theta^{1/q}w\in C^{a-1}_{\hthreects}(H,B^{1/(pq)}).
\]
By \eqref{h3:co:frobenius-distribution},
$D^q\theta^{1/q}=1$, and $D^q$ is $L^{1/q}$-linear.
As $w\in B^{1/p}\subset L^{1/q}$, this proves $D^qb'=w$.
Taking differentials and using the kernel formula gives
\begin{equation}\label{h3:co:primitive-root-level}
 \partial b'\in C^a_{\hthreects}(H,B^{1/q}).
\end{equation}
Here $b'$ is fixed by $U$, and
$(L^{1/q})^U=B^{1/q}$ by uniqueness of roots.

We must also descend $D^{q-1}b'$.  For $\sigma\in U$ write
$\sigma(D)=D(1+h_\sigma)$, with
$h_\sigma\in D^{p+1}k[[D]]$.  Then
\[
 \sigma(D)^{q-1}-D^{q-1}
 \in D^{q+p}k[[D]].
\]
The whole error kills $b'$, since
$D^{q+p}b'=D^pw=D^{p+1}b=0$ and $D^pb=0$.
Local nilpotence makes each error series finite on the relevant stage.
Thus $D^{q-1}b'$ is $U$-fixed, and
$D(D^{q-1}b'-b)=0$ proves that $D^{q-1}b'-b$ is a $B$-valued
cochain.  For each $q$ all the operators act continuously after one
finite-stage enlargement.  The descended subspaces carry their
valuation subspace topologies, so these descended cochains are
continuous as well.

Let $\alpha_i=(\partial b')^q$.  By
\eqref{h3:co:primitive-root-level}, it is a $B$-valued cocycle.
The descended difference and \Cref{h3:co:cartier-trace} give
\begin{equation}\label{h3:co:kernel-deep-image}
 [\partial b]=[D^{q-1}\partial b']
             =c_i\mathcal C^i[\alpha_i]\quad\text{in }X_a.
\end{equation}
Compactness gives a common $\varepsilon>0$ with $v(w)\geq\varepsilon$.
Because $H$ preserves valuations,
\[
 v(\alpha_i)\geq q\varepsilon+v(\theta).
\]
This eventually exceeds the finitely many thresholds defining $P$.
Hence $\alpha_i$ is a $P$-valued cocycle for all sufficiently large
admissible $i$.  The scalar $c_i$ can be absorbed into the lattice:
$c_i\mathcal C^i(\alpha_i)=\mathcal C^i(c_i^q\alpha_i)$ and
$c_i^qP=P$.  Equation \eqref{h3:co:kernel-deep-image} puts our class in a
cofinal set of the decreasing images $f_a(\mathcal C^iM_a)$, hence
in their intersection.  The subgroup $\mathcal B_a^+$ has finite image
and countable index, so the entire Frobenius kernel is countable.
\end{proof}

\begin{proof}[Proof of \Cref{h3:prop:finite-stage-finiteness}]
Suppose some $X_a$ were uncountable and choose the largest such $a$.
It exists by the degree-eight bound in
\Cref{h3:co:compact-resolution}.  Split \eqref{h3:co:cartier-four-term} at
$V=\hthreeim(d/\eta)=\ker\mathcal C$:
\[
 0\longrightarrow B\xrightarrow{\mathrm{Fr}}B\longrightarrow V
 \longrightarrow0,
 \qquad
 0\longrightarrow V\longrightarrow B\xrightarrow{\mathcal C}B
 \longrightarrow0.
\]
The first sequence places $H^{a+1}_{\hthreects}(H,V)$ between a quotient of
$X_{a+1}$ and a subgroup of $X_{a+2}$, so this group is countable.
The second embeds $X_a/\mathcal C X_a$ into it.  By
\eqref{h3:co:countable-kernel-cokernel} and the observation after
\Cref{h3:co:compact-cartier}, $M_a/\mathcal C M_a$ is countable.
It is compact Hausdorff since $\mathcal C M_a$ is compact and closed.
A countable compact Hausdorff group is finite: Baire applied to its
singleton cover gives an isolated point, and a discrete compact group
is finite.

Apply \Cref{h3:co:compact-cartier} using the finite stable image from
\Cref{h3:co:stable-image}.  To treat $N_a=\ker f_a$, use its connecting
presentation from $H^{a-1}_{\hthreects}(H,B/P)$.  A representing cochain has
finite image, hence a uniform lower bound on the logarithmic orders of
its differential coefficients.  The identity
\[
 \mathcal C^n(f)\eta=\mathrm{Car}^n(f\eta)
\]
and \eqref{h3:co:cartier-laurent} show that for sufficiently large $n$ all
negative indices in that finite image disappear modulo $P$.
Its iterate is therefore a cocycle in the finite, $H$- and
$\mathcal C$-stable module
\[
 P(0)/P,\qquad
 P(0)=\prod_jh_j^{-1}\kappa_j[[s_j]].
\]
Finite cohomology of this module and naturality of the connecting map
show that every element of $N_a$ has finite forward Cartier orbit.
For $a=0$ the kernel is already zero.  The compact-module lemma makes
$N_a$ finite and gives
\[
 \ker(\mathcal C:X_a\to X_a)\text{ countable},\qquad
 \bigcap_n f_a(\mathcal C^nM_a)=f_a(I_a)\text{ finite}.
\]
By \eqref{h3:co:cartier-frobenius-zero}, the Frobenius image in $X_a$ is
countable.  An uncountable group with countable image has uncountable
kernel, contradicting \Cref{h3:co:frobenius-kernel}.

All $X_a$ are therefore countable.  Equation
\eqref{h3:co:countable-kernel-cokernel} makes the compact $M_a$ countable,
hence finite.  Any $H$-stable order lattice is commensurable with $P$;
intersecting the two lattices gives finite quotients.  The associated
long exact sequences prove finite cohomology for every such lattice,
including $\mathcal O_B$.  By the second identification in
\eqref{h3:co:residue-duals} and \eqref{h3:co:compact-duality},
\[
 H^a_{\hthreects}(H,B/P)^\vee
 \cong H^{8-a}_{\hthreects}(H,\mathcal O_B)
\]
is finite.  Thus the quotient cohomology is finite, and the lattice
sequence now proves finiteness of $X_a$ itself.  Other lattice quotients
follow by commensurability.

The homeomorphism $B^{1/p^e}\to B$, $f\mapsto f^{p^e}$, is
$H$-equivariant and $\F_p$-linear and carries order lattices to order
lattices.  It proves the assertion for finite root stages.  Finally,
an $H$-invariant frame identifies a coefficient line at a sufficiently
deep finite stage with $B$ as an $H$-module, and its lattices with order
lattices.  The connected marking supplies such a frame for every
periodicity power; the full norm kernel fixes every determinant twist.
Enlarging the good stage to contain finitely many frame coefficients
proves the stated twisted assertion.
\end{proof}

\subsection{The natural completion comparison and the norm quotient}

Finite cohomology is now known at every cofinal good field and root
stage.  This is the input needed to contract positive valuation
cochains before comparing the algebraic and completed coefficient
complexes.  We then apply the residual norm characters to the resulting
natural comparison.

\begin{lemma}[Positive-order decompletion]\label{h3:co:decompletion}
The inclusions induce a natural quasi-isomorphism
\begin{equation}\label{h3:co:decompletion-map}
 \colim_{U,e}C^\bullet_{\hthreects}(H,B_U^{1/p^e})
 \xrightarrow{\ \sim\ }
 \colim_U C^\bullet_{\hthreects}(H,B_U^b).
\end{equation}
It retains the commuting actions and remains a quasi-isomorphism after
forming continuous families parametrized by any profinite space.
\end{lemma}

\begin{proof}
Fix a good $B$.  The positive order lattice $B^{++}$ has finite
cohomology by \Cref{h3:prop:finite-stage-finiteness}.  A positive-valued
cocycle has a uniform positive valuation bound, so its Frobenius powers
tend uniformly to zero.  \Cref{h3:co:strictness} makes each cohomology class
eventually zero under Frobenius.  Since the cohomology is finite, a
single power in each degree kills every class.

The homeomorphisms $B^{1/p^e}\to B$ identify inclusion of one positive
root stage into the next with Frobenius.  Their cochain colimit is
therefore acyclic.  The positive ideal in $B^b$ is acyclic by
\Cref{h3:co:completed-positive}.  Density gives an isomorphism of discrete
coefficient modules
\[
 B^{\mathrm{perf}}/(B^{\mathrm{perf}})^{++}
 \cong B^b/(B^b)^{++}.
\]
A continuous cochain to either quotient has finite image, whose values
can be lifted at one root stage.  Thus this is also an isomorphism of
the prescribed quotient cochain complexes.  Comparing the two
valuation-cutoff exact sequences proves the quasi-isomorphism for this
$B$.  Passing through the cofinal good stages proves
\eqref{h3:co:decompletion-map}.

Each finite root stage has finite cohomology by
\Cref{h3:prop:finite-stage-finiteness}, and each completed stage has finite
cohomology by \Cref{h3:co:completed-positive}.  Apply
\Cref{h3:co:parameters} to their Polish cochain complexes and pass to the
filtered colimits.  It proves the assertion with any profinite
parameter space.  All comparisons used inclusions and valuation
quotients, so their naturality retains all the indicated actions.
\end{proof}

\begin{lemma}[The full norm weights]\label{h3:lem:full-norm-weights}
On $H^*_{\hthreects}(H,R)$ the full residual norm quotient has the two
characters in \eqref{h3:co:completed-input}.  The distribution target
$R(\chi^{-1})$ has no $H'$-cohomology, whereas
$R\Gamma_{\hthreects}(H',R)$ is $k$ by its constants map.
These assertions hold for every $p\geq5$, without choosing a subgroup
section of $H'\to\mu_{p-1}$.
\end{lemma}

\begin{proof}
\Cref{h3:co:decompletion} transports the actual residual actions in
\eqref{h3:co:completed-input} to $R$.  The quotient $H'/H$ has order
prime to $p$, so its invariants functor is exact; the continuous
Hochschild--Serre complex thus computes $H'$-cohomology by taking those
invariants on $H$-cohomology.  The untwisted characters are $0$ and
$2\epsilon$.  Only the first is trivial modulo $p-1$.
After twisting by $\chi^{-1}$ the exponents are
\[
 -s_{\mathrm{dist}},\qquad 2\epsilon-s_{\mathrm{dist}},\qquad
 \epsilon,s_{\mathrm{dist}}\in\{1,-1\}.
\]
They belong to $\{1,-1,3,-3\}$ and are nonzero modulo $p-1\geq4$.
Both target rows therefore vanish.  In the source the invariant row is
exactly the given constants row.

The distinction from the center can be seen at $p=7$.  A central
$a\in\mu_6\subset P_3$ has norm $a^3$, so pulling a norm-weight-two
character to that central subgroup would make it trivial.  Here the
normalized \'etale generator transforms by the full norm unit $u$;
the source weight is $u^{\pm2}$ and the twisted exponents are
$\pm1,\pm3$, nontrivial modulo six.  The reduced norm surjects on tame
units, as follows from the residue-field norm
$\F_{p^3}^\times\to\F_p^\times$.  Since $H$ acts trivially on the
translation group and by inner automorphisms trivially on its own
cohomology, all residual actions used here factor through the quotient
itself.  Neither the eigenparameter construction nor the invariants
calculation requires it to split inside $P_3$.
\end{proof}

\begin{proof}[Proof of \Cref{h3:thm:coheight-one-comparison}]
The distribution sequence \eqref{h3:co:distribution-sequence} is exact on
the specified continuous cochain complexes.  By
\Cref{h3:lem:full-norm-weights}, its last term has zero $H'$-cohomology,
and its middle term is computed by constants.  Thus its first term is
also computed by the actual constants map, proving
\eqref{h3:co:main-constants}.  The finite-stage $H'$-cohomology is finite
by \Cref{h3:prop:finite-stage-finiteness} and exact tame invariants.
\Cref{h3:co:parameters} proves the parameterized assertion.  Although a
choice of $D$ was used to prove the comparison, the resulting
quasi-isomorphism is the original constants inclusion.  Its naturality
therefore retains the connected-frame action and the remaining
height-three action.
\end{proof}

\subsection{Descent to base fields and base lattices}

The ordinary-stratum argument needs finiteness on the coheight-one
\emph{base lattices}, including all conormal and periodicity twists.
We give descent separately for fields and for lattices; the latter does
not divide by a possibly ramified Galois degree.

\begin{lemma}[Continuous field descent]\label{h3:co:field-descent}
For every profinite space $T$, the natural augmentation is a
quasi-isomorphism
\begin{equation}\label{h3:co:field-descent-map}
 C(T,K)\longrightarrow
 \colim_U C^\bullet_{\hthreects}\bigl(J,C(T,B_U)\bigr).
\end{equation}
The same statement holds for a coefficient line pulled back from $K$.
These comparisons are natural in $T$ and equivariant for the commuting
$P_3$-action.
\end{lemma}

\begin{proof}
For fixed normal open $U$, set $W_U=C(T,B_U)$ with its uniform topology.
The action on it factors through $J/U$.  Compare locally constant
$J$-cochains, viewing $W_U$ as discrete, with continuous $J$-cochains
for its given topology.  A fixed finite-projective
$\F_p[[J]]$-resolution from \Cref{h3:co:compact-resolution} computes both.
Their Hom complexes are the same underlying finite summands of powers
of $W_U$: every module map factors through the finite-dimensional
quotient of a resolution term by the augmentation ideal of $U$, and
is continuous with either topology.  The canonical inclusion of the
two bar-cochain complexes is consequently a quasi-isomorphism.

After taking the colimit over normal open subgroups, every locally
constant cochain, its finite clopen partition, and its finitely many
values are represented at one finite quotient.  Refining $U$ if
necessary gives an equality of complexes
\[
 \colim_U C^\bullet_{\mathrm{lc}}(J,W_U)
 =\colim_U C^\bullet(J/U,W_U).
\]
A finite topological $K$-basis of $B_U$ gives
$W_U=B_U\otimes_K C(T,K)$.  The normal-basis theorem for the finite
Galois \'etale algebra $B_U/K$ makes this a coinduced $J/U$-module.
The finite-group cochain complex of a coinduced module has zero
positive cohomology: its bar contraction is evaluation in the induced
function coordinate.  Its invariants are exactly $C(T,K)$.
Thus the canonical augmentation is a quasi-isomorphism at each finite
quotient.  Exactness of the filtered colimit proves
\eqref{h3:co:field-descent-map}.

For a pulled-back coefficient line, tensor the finite-basis and
normal-basis argument with that line over $K$.  All choices were used
only to prove exactness.  The comparison maps themselves are the
canonical augmentations, hence are natural and retain the commuting
action.  This proof makes no claim about cochains into the metric
union $L$.
\end{proof}

\begin{corollary}[Base lattice finiteness]\label{h3:co:base-lattices}
Let $\mathcal L$ be a coefficient line over $K$ whose pullback to $L$
has an $H$-invariant frame at a finite stage.  Suppose that
$\Lambda\subset\mathcal L$ is a $P_3$-stable order lattice over
$k[[y]]$.  Then
\[
 H^a_{\hthreects}(P_3,\Lambda)
 \quad\text{and}\quad H^a_{\hthreects}(H,\Lambda)
\]
are finite for every $a$.  This applies to every periodicity power,
every finite determinant character, their inverses, and their products
with conormal or canonical-line twists.  In particular each adjacent
$x$-pole strip used in \Cref{h3:sec:ordinary} has finite $P_3$-cohomology.
\end{corollary}

\begin{proof}
Choose a sufficiently deep good normal stage $B=B_U$ containing the
frame and put $Q=J/U$.  Let $\mathcal O_B$ be its integral closure
lattice and set
\[
 M=\mathcal O_B\otimes_{k[[y]]}\Lambda.
\]
The actions of $H$ and $Q$ commute and $M^Q=\Lambda$.
In the $H$-invariant frame, $M$ is an $H$-stable order lattice at the
stage $B$, so \Cref{h3:prop:finite-stage-finiteness} gives finite
$H$-cohomology for it.

For $b>0$, $H^b(Q,M)$ is a finite group.  To see this without averaging,
the finite-group cochain terms are finitely generated modules over
$k[[y]]$.  Their cohomology is finitely generated as well.  Inverting
$y$ commutes with that cohomology and gives zero in positive degrees
by normal-basis descent on the field algebra.  Hence $H^b(Q,M)$ has
finite length over $k[[y]]$, whose residue field is finite.  The finite
group boundaries are compact and closed, so these finite cohomology
groups have their discrete quotient topology and continuous $H$-action.

Use the bounded finite-projective $H$-resolution together with the
finite-$Q$ bar complex on $M$.  Taking $H$-cohomology first gives finite
groups $H^i(Q,H^j_{\hthreects}(H,M))$ in every bidegree, hence finite total
cohomology in every degree.  Taking $Q$-cohomology first gives
\[
 E_2^{a,b}=H^a_{\hthreects}(H,H^b(Q,M)).
\]
All rows with $b>0$ are finite by
\Cref{h3:co:compact-resolution}.  The bottom term
$E_2^{a,0}=H^a_{\hthreects}(H,\Lambda)$ has no outgoing differential and
only finitely many incoming differentials, each with finite image.
Its surviving quotient is finite because the total cohomology is
finite.  Thus it is finite.  These two spectral sequences compute the
same first-quadrant complex: finite projective Hom is exact on the
compact $Q$-cycle and boundary sequences, and the finite $Q$-cochain
functors are finite products.  No strictness assertion about the still
unknown base-lattice cohomology was used.

Finally $P_3/H=\Z_p^\times$ has finite cohomology on finite discrete
modules: take exact invariants for $\mu_{p-1}$ and the two-term
resolution for $1+p\Z_p$.  The continuous Hochschild--Serre sequence
for $H\subset P_3$ now proves finite $P_3$-cohomology of $\Lambda$.
The residual action on its finite $H$-cohomology is continuous, since
the coefficient lattice is compact and the finite-resolution
cohomology has the quotient topology.

The connected marking trivializes all periodicity powers and the full
norm kernel fixes determinant characters.  The canonical-line and
conormal identifications express the other indicated lines as products
of these.  An adjacent $x$-pole strip is a power-series lattice in $y$
with precisely one such conormal and coefficient-line twist.  It
therefore satisfies the hypotheses just proved.
\end{proof}

\begin{remark}[Finite constant fields]\label{h3:co:constant-descent}
Finite enlargement of $k$ does not lose either these finiteness
statements or the natural comparison maps.  Scalar extension of a
coefficient module, of each order lattice, or of each continuous
cochain complex is a finite direct sum in its given topology.
The continuous trace-one contraction of
\Cref{h3:lem:framework-semilinear} makes the associated semilinear Galois
complex acyclic in positive degrees; its invariants are the original
coefficient complex.  This remains true with profinite parameters,
because finite scalar extension commutes with continuous functions
and the contraction consists of finite sums.
For the original extended stabilizer one retains the finite coefficient
Galois action on the Honda stabilizer as well; the equivariant
augmentations give the corresponding semidirect-product descent.
Thus constants, units, line twists, and base lattice finiteness descend
to the original coefficient field.  No division by the constant-field
extension degree is involved.
\end{remark}

\section{Exact descent and the height-two map}\label{h3:sec:descent}

The coheight-one coefficient comparison first gives two actual maps at
height two: the unit and the determinant class.  We then use those maps
in the finiteness argument needed for the ordinary comparison.
Throughout these sections \(p\geq5\).  Put \(U=u^{-1}\), so that
\(|U|=2\) and \(v_1=xU^{p-1}\).

\subsection{Relative descent and exact functors}

The tests below are applied to already formed, completed Morava
cooperation spectra.  Their preservation of descent is a finite
categorical statement, which we establish first.  We then identify the
relative cooperation terms and their actual cofaces.
A commutative algebra is \emph{descendable} if the unit belongs to the
thick tensor ideal it generates.

\begin{lemma}[Exact Amitsur descent]\label{h3:lem:exact-amitsur}
Let $\mathcal C$ be a stable symmetric monoidal $\infty$-category with
limits and with tensor product exact in each variable.  Let $E$ be a
descendable commutative algebra, write $\mathbf1$ for the unit, and put
$I_E=\fib(\mathbf1\to E)$.  For its full Amitsur object
$C^q=E^{\otimes(q+1)}$, the partial totalization
\[
 \operatorname{Tot}_s C^\bullet
   =\lim_{[q]\in\Delta_{\leq s}}C^q
\]
is a finite cubical limit.  Its augmentation fiber, including its map
to the unit, is
\begin{equation}\label{h3:des:partial-amitsur-fiber}
 \bigl(\fib(\mathbf1\longrightarrow\operatorname{Tot}_s C^\bullet)
       \longrightarrow\mathbf1\bigr)
 \simeq \bigl(I_E^{\otimes(s+1)}\longrightarrow\mathbf1\bigr).
\end{equation}
The transition between successive fibers applies $I_E\to\mathbf1$
to the last factor.  There is an integer $N$ such that every composite
of $N$ successive maps between these augmentation fibers is null.
Every exact functor from
$\mathcal C$ to a stable category preserves the augmented
totalization $\mathbf1\simeq\operatorname{Tot}C^\bullet$ and this
same null-composite bound.
\end{lemma}

\begin{proof}
Here $\Delta_{\leq s}$ is the full simplex category on
$[0],\ldots,[s]$, so its limit includes the codegeneracies.
Let $\mathcal P_s^\times$ be the finite poset of nonempty subsets of
$\{0,\ldots,s\}$.  The functor
\[
 \theta_s:\mathcal P_s^\times\longrightarrow\Delta_{\leq s},
 \qquad J\longmapsto[|J|-1],
\]
sends an inclusion to the order-preserving injection between the
ordered subsets.  It is initial for limits.  To verify the cofinality
criterion, fix $[q]$, where $q\leq s$.  An object of
$\theta_s\downarrow[q]$ is a nonempty subset $J$ together with a
weakly increasing map $J\to\{0,\ldots,q\}$; its morphisms are
extensions of these partial maps.  Thus this comma category is the
nonempty face poset of the simplicial complex $P(s,q)$ whose simplices
are the lists
\[
 (i_0,j_0),\ldots,(i_b,j_b),\qquad
 0\leq i_0<\cdots<i_b\leq s,
 \quad 0\leq j_0\leq\cdots\leq j_b\leq q.
\]
Its nerve is the barycentric subdivision of $P(s,q)$.

The complex $P(s,q)$ is contractible whenever $s\geq q$.
The base case $P(0,0)$ is a point.  For the induction, the simplices
containing the vertex $(s,j)$, together with their faces, form the
cone on $P(s-1,j)$.
Start with the cone with apex $(s,q)$ on $P(s-1,q)$; this includes
all simplices with first coordinates less than $s$ and is
contractible even if its base is not.  Attach the cones with apices
$(s,j)$, for $j=q-1,\ldots,0$.  The intersection of each such cone
with the preceding union is exactly its base $P(s-1,j)$: a simplex
cannot contain two vertices with first coordinate $s$.
Every attaching base is contractible by induction, since
$j<q\leq s$ implies $j\leq s-1$.  These are inclusions of finite
simplicial subcomplexes, hence cofibrations; attaching a cone along a
contractible subcomplex preserves contractibility.  This proves the
claim, including the endpoint $q=s$, where no assertion about the
contractibility of the initial base $P(s-1,s)$ was used.  The comma
categories are therefore contractible, proving initiality of
$\theta_s$.

It follows that $\operatorname{Tot}_s C^\bullet$ is the limit of the
punctured finite $(s+1)$-cube
\[
 J\longmapsto\bigotimes_{i=0}^s
   \begin{cases}E,&i\in J,\\ \mathbf1,&i\notin J,\end{cases}
 \qquad \varnothing\ne J\subseteq\{0,\ldots,s\},
\]
whose maps insert units.  Successive fibers and exactness of tensor
product identify its augmentation fiber as the arrow in
\eqref{h3:des:partial-amitsur-fiber}.  Omitting the last cubical
direction gives the stated transition.  Thus the model and its
transition refer to full partial totalizations, not a coface-only
truncation.  An exact functor preserves this finite cubical limit.
The finite-cube formula appears in \cite[Proposition~2.14]{MNN2017};
the comparison with full partial totalizations and the augmentation
fiber have been proved here in the stated setting.

For a descendable $E$, the partial Amitsur tower represents the
constant pro-object with value the unit
\cite[Proposition~3.20]{Mathew}.  Its limit is preserved by every
finite-limit-preserving functor \cite[Proposition~3.10]{Mathew}.
Together with the finite cubical comparison, this proves the
asserted identification of the transformed totalization.  There is
also an integer $N$ such that every composite of $N$ maps which
become null after tensoring with $E$ is null
\cite[Proposition~3.27]{Mathew}.  One can read this bound from a finite
thick-ideal construction of the unit: it is one for $E$, is unchanged
by retracts and tensor products, and adds across a cofiber sequence.
The arrow $I_E\to\mathbf1$ becomes null after tensoring with $E$:
its tensor is the composite through the null map $I_E\to E$,
followed by multiplication.  Thus the error transitions in
\eqref{h3:des:partial-amitsur-fiber} become null after tensoring with
$E$, and their $N$-fold composites are null.  An exact functor
preserves those null composites, supplying the same convergence
bound after applying the functor.
\end{proof}

\begin{proposition}[Exact relative Morava descent]
\label{h3:prop:exact-descent}
Let \(n\geq1\), let \(k\) be a finite field containing \(F=\F_{p^n}\), and let
\(R_{n,k}=L_{K(n)}S_k\).  In the \(K(n)\)-local category of
\(S_k\)-modules, the algebra \(E_n(k)\) is descendable over the unit
\(R_{n,k}\).  Write \(\widehat\otimes\) for the tensor product in this
local category and put
\[
 \mathcal N_n^q=E_n(k)^{\widehat\otimes_{R_{n,k}}(q+1)}
 \qquad(q\geq0).
\]
Its augmented Amitsur object totalizes to \(R_{n,k}\).  Every exact
functor from this local category to a stable category preserves that
totalization and a uniform nilpotence bound for its error tower.
In particular this applies to the inclusion in spectra followed by
ordinary smash, a finite Moore quotient, a chromatic localization, or
a monochromatic fiber.

There are identifications of graded rings
\begin{equation}\label{h3:des:relative-cooperations}
 \pi_*\mathcal N_n^q
    \cong C_{\hthreects}(P_n^q,E_n(k)_*),
\end{equation}
with the maximal-ideal topology on the coefficients.  Under these
identifications the actual cofaces and codegeneracies induce those of
the continuous action cobar construction, and the augmented unit is
the constant unit.  Each \(\mathcal N_n^q\) is an ordinary
\(E_n(k)\)-module through its first factor.  This is the module
structure used when an ordinary exact functor is applied to a term.
\end{proposition}

\begin{proof}
We begin with the standard coefficient field, keeping the scalar
extension separate.  Set \(R_n=L_{K(n)}S\),
\(E_F=E_n(F)\), \(\Gamma_n=\Gal(F/\F_p)\), and
\(G_n^0=P_n\rtimes\Gamma_n\).  The Hopkins--Ravenel theorem in the
form of \cite[Theorem~4.18]{Mathew} says that \(E_F\) generates the unit
\(L_n\mathbb S_{(p)}\) as a thick tensor ideal in the \(E(n)\)-local
category.  Applying the symmetric monoidal localization \(L_{K(n)}\)
preserves this finite generation statement
\cite[Corollary~3.21]{Mathew}.  The map \(\mathbb S_{(p)}\to S\) is a
mod-\(p\) equivalence, hence a \(K(n)\)-equivalence for \(n\geq1\),
so the resulting unit is \(R_n\).  The same localized base case is
stated directly in \cite[Proposition~10.10]{Mathew} for Morava theory
attached to a perfect residue field; the relative assertion still needs
the following scalar step.  Now base change along
\(R_n\to R_{n,k}\) inside the \(K(n)\)-local category.  Thus
\[
 B=E_F\widehat\otimes_{R_n}R_{n,k}
\]
is descendable over \(R_{n,k}\).

We identify one factor of this base change and prove that it is
itself descendable.  A \(\Z_p\)-basis of \(W(k)\) gives an equivalence
\(R_{n,k}\simeq R_n^{\vee [k:\F_p]}\) of underlying \(R_n\)-modules.
More precisely, the natural map \(R_n\otimes_S S_k\to R_{n,k}\) is
an equivalence: its source is a finite sum of \(K(n)\)-local copies of
\(R_n\), and localizing the finite free decomposition of \(S_k\) gives
the same map.  The natural scalar map therefore identifies the graded coefficient
ring of \(B\) with \(E_{F,*}\otimes_{\Z_p}W(k)\); no infinite limit is
interchanged with scalar extension here.  The finite étale algebra
of constants splits as
\begin{equation}\label{h3:des:embedding-idempotents}
 W(F)\otimes_{\Z_p}W(k)
   \xrightarrow{\ \cong\ }
   \prod_{\sigma:F\hookrightarrow k}W(k),
 \qquad a\otimes b\longmapsto(\sigma(a)b)_\sigma.
\end{equation}
Let \(e_\sigma\) denote its embedding idempotents.  Idempotent
localization splits \(B\) into the finite product of
\(B_\sigma=B[e_\sigma^{-1}]\): on homotopy groups the map to this
product is the displayed orthogonal-idempotent decomposition, hence
is an equivalence.  Each factor is even periodic, with degree-zero
ring \(W(k)[[u_1,\ldots,u_{n-1}]]\).  Its orientation is the universal
Honda deformation after the indicated unramified base change.  We may
therefore take the factor for the chosen inclusion \(\iota:F\subseteq k\)
as the model \(E_n(k)\) used here.

For \(\gamma\in\Gamma_n\), the Galois automorphism of the standard
factor \(E_F\), base changed with the identity on \(R_{n,k}\), is an
\(R_{n,k}\)-algebra automorphism of \(B\).  It permutes the idempotents
by \(e_\sigma\mapsto e_{\sigma\gamma^{-1}}\).  In particular it
identifies all the factors \(B_\sigma\) as \(R_{n,k}\)-algebras while
the second scalar action stays fixed.  This is the factor equivalence
being used; coefficient Frobenius on the chosen \(E_n(k)\) alone is
semilinear over \(W(k)\).  The underlying module of \(B\) is consequently
a finite sum of modules equivalent to \(E_n(k)\).  The thick tensor
ideal generated by \(E_n(k)\) contains \(B\), and the ideal generated
by \(B\) contains the unit.  This proves descendability of the chosen
factor.

Apply \Cref{h3:lem:exact-amitsur} in the \(K(n)\)-local category of
\(S_k\)-modules, with its local tensor product and unit \(R_{n,k}\).
The inclusion into spectra is exact, as are the ordinary tests in the
statement.  The lemma therefore proves the asserted totalization and
uniform error-tower bound for the relative Amitsur object and for each
of those tests.  We now keep the same \(B\) and its embedding
idempotents to identify the terms and maps of that object.

It remains to prove \eqref{h3:des:relative-cooperations} with its
maps.  The completed-cooperations theorem of Devinatz--Hopkins applies
to the standard \(E_F\) and the standard extended group \(G_n^0\):
\cite[Proposition~2.2 and formulas~(2.5)--(2.7)]{DevinatzHopkins}
identifies
\[
 \pi_*\bigl(E_F^{\widehat\otimes_{R_n}(q+1)}\bigr)
     \cong C_{\hthreects}((G_n^0)^q,E_{F,*}).
\]
The tensor on the left is the \(K(n)\)-local tensor; with its local
unit \(R_n\) it is the completed smash appearing in that theorem.
Base change gives a canonical equivalence of actual terms
\[
 B^{\widehat\otimes_{R_{n,k}}(q+1)}
 \simeq
 E_F^{\widehat\otimes_{R_n}(q+1)}
       \widehat\otimes_{R_n}R_{n,k}.
\]
Because \(R_{n,k}\) is finite free over \(R_n\), the preceding
coefficient identification becomes
\begin{equation}\label{h3:des:basechanged-cooperations}
 \pi_*\bigl(B^{\widehat\otimes_{R_{n,k}}(q+1)}\bigr)
 \cong C_{\hthreects}((G_n^0)^q,E_{F,*})\otimes_{\Z_p}W(k)
 \cong C_{\hthreects}((G_n^0)^q,B_*).
\end{equation}
The last isomorphism uses a finite basis of \(W(k)\), so it only
commutes continuous functions with a finite direct sum.

We spell out the coordinate convention because it controls the first
coface and the field component.  Write the action of \(G_n^0\) on
\(B\) for its action on \(E_F\) tensored with the identity on
\(R_{n,k}\).  In inhomogeneous coordinates the evaluation of a class
in the left side of \eqref{h3:des:basechanged-cooperations} at
\((g_1,\ldots,g_q)\) is induced by the actual multiplication map
\begin{equation}\label{h3:des:inhomogeneous-evaluation}
 \mu\circ\bigl(1\widehat\otimes g_1\widehat\otimes
       (g_1g_2)\widehat\otimes\cdots\widehat\otimes
       (g_1\cdots g_q)\bigr):
 B^{\widehat\otimes_{R_{n,k}}(q+1)}\longrightarrow B.
\end{equation}
For \(q=0\) this is the identity.  For \(q\geq1\), the evaluation in
Devinatz--Hopkins applies
\(h_i^{-1}\) to the first \(q\) factors and leaves the last factor
fixed.  Put \(h_i=g_i\cdots g_q\) and then apply \(h_1\) to the value.
This continuous change of variables gives exactly
\eqref{h3:des:inhomogeneous-evaluation}.  Their natural evaluation
formulas apply to every homotopy class, so this argument does not assume
that products of factor coefficients generate the completed term.

The projection \(B\to E_n(k)=B_\iota\) is a unital algebra map.
In degree \(q\), the induced projection to \(\mathcal N_n^q\) is the
all-\(\iota\) idempotent summand of the tensor product.  Under
\eqref{h3:des:inhomogeneous-evaluation}, its projector acts on a function
by multiplication by
\[
 e_\iota\,g_1(e_\iota)\,(g_1g_2)(e_\iota)\cdots
                 (g_1\cdots g_q)(e_\iota).
\]
The action of \(\Gamma_n\) on the embedding idempotents is free and
transitive, and its kernel in \(G_n^0\) is \(P_n\).  This product is
\(e_\iota\) precisely when every \(g_i\) lies in \(P_n\), and is zero
otherwise.  Since \(P_n^q\) is clopen in \((G_n^0)^q\), the selected
summand is exactly \(C_{\hthreects}(P_n^q,(B_\iota)_*)\).  The
\(P_n\)-action on the selected coefficient ring fixes its \(W(k)\)
scalars.  This proves \eqref{h3:des:relative-cooperations}.

The same evaluation maps identify all cobar maps.  For a coefficient
function \(f\) in degree \(q\), their formulas are
\[
 \begin{aligned}
 (d^0f)(g_1,\ldots,g_{q+1})
     &=g_1\bigl(f(g_2,\ldots,g_{q+1})\bigr),\\
 (d^if)(g_1,\ldots,g_{q+1})
     &=f(g_1,\ldots,g_i g_{i+1},\ldots,g_{q+1})
       &&(1\leq i\leq q),\\
 (d^{q+1}f)(g_1,\ldots,g_{q+1})
     &=f(g_1,\ldots,g_q),\\
 (s^if)(g_1,\ldots,g_{q-1})
     &=f(g_1,\ldots,g_i,1,g_{i+1},\ldots,g_{q-1})
       &&(0\leq i<q).
 \end{aligned}
\]
They follow by inserting a unit or multiplying adjacent factors in
\eqref{h3:des:inhomogeneous-evaluation}.  The unital projection from
the \(B\)-Amitsur object to the \(E_n(k)\)-Amitsur object commutes with
these maps.  Thus its selected cofaces are obtained by projecting the
actual cofaces and then restricting the displayed formulas to \(P_n\).
Their higher compatibilities are those of that actual Amitsur object.
The augmented unit evaluates to the constant unit.  The first coface is
\(R_{n,k}\)-linear and has the displayed varying stabilizer action on
coefficients; the other cofaces and all codegeneracies are linear for
the first-factor \(E_n(k)\)-module structures.  This proves the claims
about maps and naturality.
\end{proof}

The notation \(C_{\hthreects}(P_n^q,E_n(k))\) below refers to the actual
term \(\mathcal N_n^q\), equipped with
\eqref{h3:des:relative-cooperations}; it does not infer an equivalence
of spectra from a coefficient isomorphism.  The inclusion of the
\(K(n)\)-local category in spectra is exact, but its tensor product is
still the completed one.  Only after forming a term do we regard it as
an ordinary module and apply an ordinary exact functor.

We also record the finite-field descent, including its unit.  For
\(\tau\in\Gal(k/\F_p)\), combine the standard automorphism
\(\tau|_F\) on \(E_F\) with \(\tau\) on the second factor
\(R_{n,k}\).  On constants, evaluation at \(\iota\) sends
\((\tau|_F)(a)\otimes\tau(b)\) to \(\tau(\iota(a)b)\).
Consequently this combined action preserves \(B_\iota\), is semilinear
over \(W(k)\), and conjugates \(P_n\) by the usual Honda action.
It gives the coefficient-field action on the relative object just
constructed and respects its augmentation and cofaces.

Let \(\Gamma=\Gal(k/\F_p)\).  A normal-basis element \(\bar a\in k\)
has nonzero trace: otherwise the sum of its conjugates would be a
nontrivial \(\F_p\)-linear dependence.  Lift it to \(a_0\in W(k)\).
The conjugates of \(a_0\) are a \(\Z_p\)-basis, since their determinant
is a unit modulo \(p\), and its trace is a unit.  Replacing \(a_0\) by
\(a=a_0/\operatorname{Tr}(a_0)\) gives a normal basis with trace one.
The map
\[
 \operatorname{Ind}_{1}^{\Gamma}S=\bigvee_{\gamma\in\Gamma}S
       \longrightarrow S_k,
 \qquad 1_\gamma\longmapsto\gamma(a),
\]
is a coherent \(\Gamma\)-equivariant equivalence of underlying
\(S\)-modules: it is the induced map adjoint to the class \(a\), and
it is an isomorphism on every homotopy group.  The diagonal unit maps
to \(\sum_\gamma\gamma(a)=1\).  Equivalently, under the
induction--coinduction adjunction its invariant unit is the specified
unit of \(S_k\).  Localizing gives the same statement for
\(R_n\to R_{n,k}\).

Let an exact functor be defined on the underlying \(K(n)\)-local
\(S\)-modules, as are the inclusion followed by the ordinary tests
used for spectral finiteness below.  It preserves this finite sum, its
action, and the diagonal map.  Finite induction equals finite
coinduction, so the homotopy fixed points of the transformed
\(R_{n,k}\) recover the transformed \(R_n\) and its specified unit.  This argument does not require that the functor
commute with an arbitrary homotopy limit.  On coefficient cochains the
trace-one averaging and contraction of
\Cref{h3:lem:framework-semilinear} likewise give exact semilinear descent,
even when \(p\) divides \(|\Gamma|\).  Finally, scalar extension is
faithful: its underlying finite free module of positive rank detects
a zero cofiber.

\subsection{The ordinary residue test}

\begin{lemma}\label{h3:des:flat-functions}
Put \(\widetilde A=W(k)[[x,y]]\), and let \(Q\) be profinite.
The module \(C_{\hthreects}(Q,\widetilde A)\), for the
\((p,x,y)\)-adic topology, is pro-free and flat over
\(\widetilde A\).  For every ordinary \(E_3(k)\)-module \(N_Q\) whose graded
coefficient module is \(C_{\hthreects}(Q,(E_3(k))_*)\), ordinary
smash with \(E_2(k)\) over
the \(p\)-local sphere \(\mathbb S_{(p)}\) has homotopy
\begin{equation}\label{h3:des:ordinary-kunneth}
 (E_2(k))_*E_3(k)\otimes_{(E_3(k))_*}
 C_{\hthreects}(Q,(E_3(k))_*).
\end{equation}
The first tensor factor in \eqref{h3:des:ordinary-kunneth} denotes
absolute ordinary cooperations over \(\mathbb S_{(p)}\).  This applies
in particular to the actual term \(N_Q=\mathcal N_3^q\) for
\(Q=P_3^q\), with its first-factor module structure.
\end{lemma}

\begin{proof}
For each power \(\mathfrak m^a\) of \(\mathfrak m=(p,x,y)\),
coefficient reduction gives
\[
 C_{\hthreects}(Q,\widetilde A)/\mathfrak m^a
 \cong C_{\mathrm{lc}}(Q,\widetilde A/\mathfrak m^a).
\]
Surjectivity follows by lifting on a finite clopen partition.
For the kernel, partition the monomials in \(p,x,y\) among the
finitely many monomial generators of \(\mathfrak m^a\); the resulting
division operations on coefficients are continuous and write a
kernel function as a sum of these generators times continuous
functions.  The right side is the filtered colimit of finite free
modules associated to finite clopen partitions.  It is flat over
the Artinian local ring \(\widetilde A/\mathfrak m^a\).

Choose a basis of the residue module over \(k\), lift the basis
elements, and let \(F\) be the free \(\widetilde A\)-module on those
lifts.  An element of \(\widehat F\) is a family of coefficients for which,
modulo each \(\mathfrak m^a\), only finitely many are nonzero.
Consequently \(\widehat F/\mathfrak m^a\) is the free
\(\widetilde A/\mathfrak m^a\)-module on the chosen basis; continuous
monomial division gives the kernel equality used here.  The map
\(\widehat F\to C_{\hthreects}(Q,\widetilde A)\) is an isomorphism
modulo \(\mathfrak m\).  It is an isomorphism modulo every
\(\mathfrak m^a\): both modules there are flat, and successive
nilpotent-ideal reduction proves injectivity and surjectivity.
Taking the inverse limit proves pro-freeness.  For the ordinary
flatness conclusion, \(\widetilde A=W(k)[[x,y]]\) is a regular complete
Noetherian local ring, and \(\mathfrak m=(p,x,y)\) is generated by the
minimal regular sequence \(p,x,y\).  The module \(F\) is free, hence
flat, and \(F/\mathfrak mF\) is free over \(\widetilde A/\mathfrak m\).
Thus \cite[Proposition~3.13]{BarthelStapleton2016} applies and makes the
ordinary \(\mathfrak m\)-adic completion \(\widehat F\) flat.  The
pro-free characterizations are also described in
\cite[Theorem~A.9 and Proposition~A.13]{HS99}.

Now use the ordinary module tensor identity
\[
 E_2(k)\wedge N_Q
 \simeq
 (E_2(k)\wedge E_3(k))\otimes_{E_3(k)}N_Q.
\]
The module Künneth spectral sequence has no positive Tor terms
by the flatness just proved.  Periodicity reduces graded modules over this coefficient ring to
parity-graded modules over the regular local ring
\(W(k)[[x,y]]\), of global dimension three.  A projective resolution
of length at most three therefore gives bounded convergence of this
Künneth calculation.
This proves \eqref{h3:des:ordinary-kunneth}.
\end{proof}

\begin{lemma}\label{h3:des:residue-cooperations}
Let
\[
 \overline E_2=E_2(k)/(p,u_1).
\]
Apply the ordinary relative tensor \(\overline E_2\otimes_{S_k}(-)\)
to the \(K(3)\)-completed Amitsur object for \(E_3(k)\).  Its homotopy is
even periodic.  After using the test periodicity generator, the
degree-zero coefficient in cosimplicial degree \(q\) is
\[
 \colim_B C_{\hthreects}(P_3^q,B),
\]
where \(B\) runs through the finite connected-marking algebras
over \(k((y))\) of \Cref{h3:sec:coheight-one}.  The cofaces give their
actual \(P_3\)-action, and the test of the unit is the constant
cochain.
\end{lemma}

\begin{proof}
Before imposing the residue ideal, both Morava theories are
Landweber exact over \(BP\): the sequence
\(p,u_1,\ldots,u_{n-1}\) is regular and the height-\(n\) coefficient
is a unit, and these facts are retained by finite unramified scalar
extension.  Landweber exactness gives
the homology base-change formula for all spectra, as recalled in the
introduction and Example~0.1(d) of \cite{HoveyStrickland}.  Applying it
first to \(E_2(k)_*(E_3(k))\), then to \(E_3(k)_*(BP)\) and using the
symmetry of ordinary smash, gives
\begin{equation}\label{h3:des:absolute-landweber-cooperations}
 (E_2(k))_*E_3(k)
 \cong (E_2(k))_*\otimes_{BP_*}BP_*BP
                 \otimes_{BP_*}(E_3(k))_*.
\end{equation}
The two tensor products use the two unit maps of \(BP_*BP\).
The strict \(p\)-typical formal-law groupoid represented by
\((BP_*,BP_*BP)\), including its units and composition, is recalled in
\cite[Section~2, Lemma~2.9 and the following discussion, p.~7]{BhattacharyaEgger2019}.
Naturality of the two multiplicative orientations identifies these
base-changed operations with those of the oriented formal isomorphisms.
This is an absolute ordinary cooperation calculation,
which is precisely the first factor in
\eqref{h3:des:ordinary-kunneth}.

We also record the needed flatness over each coefficient factor.
For a Landweber-exact \(BP_*\)-algebra \(B\),
\cite[Lemma~2.2]{HoveyStrickland} makes
\(B\otimes_{BP_*}BP_*BP\) flat for its other \(BP_*\)-action.
Hopf-algebroid conjugation gives the analogous assertion for
\(BP_*BP\otimes_{BP_*}B\).  Applying these two assertions with
\(B=(E_2(k))_*\) and \(B=(E_3(k))_*\), and then base changing the
remaining unit, proves that
\eqref{h3:des:absolute-landweber-cooperations} is flat over either
coefficient factor.  Hence the sequence \((p,u_1)\) on the height-two
factor acts regularly on these cooperations.
Tensoring over the height-three factor with the flat module in
\Cref{h3:des:flat-functions} preserves that regularity.  The two
residue cofibers therefore create no odd homotopy in the absolute
smash product.

We next pass from this absolute residue calculation to the
relative test in the statement.  Let
\[
 C=\cofib(\mathbb S_{(p)}\longrightarrow S).
\]
The completion map is a mod-\(p\) equivalence, so \(C/p=0\).
Thus multiplication by \(p\) is invertible on \(C\), and the
\(p\)-local spectrum \(C\) is rational.  The spectrum
\(\overline E_2\) is rationally acyclic, already after its
first residue cofiber.  Hence \(\overline E_2\wedge C=0\).
The unit map \(\overline E_2\to\overline E_2\wedge S\) is
an equivalence of spectra.  Its composite with the action map
\(\overline E_2\wedge S\to\overline E_2\) is the identity.
The action map is therefore an equivalence of right
\(S\)-modules.  Consequently, for every
\(S\)-module \(N\), there is a natural equivalence
\begin{equation}\label{h3:des:absolute-relative-residue}
 \overline E_2\wedge N
 \simeq(\overline E_2\wedge S)\otimes_S N
 \simeq\overline E_2\otimes_S N.
\end{equation}

For \(S_k\)-modules, the two scalar actions on the last
object give an action of
\[
 S_k\otimes_S S_k
       \simeq\prod_{\sigma\in\Gal(k/\F_p)}S_k.
\]
The multiplication map \(S_k\otimes_S S_k\to S_k\) is projection
to the identity factor.  Indeed, balancing the two scalar actions is
\[
 (\overline E_2\otimes_S N)
   \otimes_{S_k\otimes_S S_k}S_k
 \simeq \overline E_2\otimes_{S_k}N.
\]
Thus the identity idempotent selects
\(\overline E_2\otimes_{S_k}N\).  On the absolute
residue coefficient algebra already calculated, this is the
equal-embeddings summand of
\(k\otimes_{\F_p}k\simeq\prod_\sigma k\).
It is a retract and therefore remains even.  Applying this
to \(N=C_{\hthreects}(P_3^q,E_3(k))\) gives the required relative
term.  The comparison \eqref{h3:des:absolute-relative-residue}
and the identity idempotent are natural for \(S_k\)-linear
maps.  Every \(P_3\)-cobar coface is \(S_k\)-linear,
so these identifications commute with all cofaces and
with the constant unit.  Regularity was needed only for
the absolute calculation before the residue cut.

In this identity summand, the coefficients are described
by the formal-group-isomorphism Hopf algebroid over \(k\).
Under an isomorphism of formal groups the ideals
\((p,v_1)\) agree on the two sides.  Thus this cut is \(p=x=0\)
on the height-three factor.  The height-two \(v_2\) is a unit,
so \(y\) is a unit as well.  After fixing the test periodicity generator,
the conversion from graded to ungraded coordinates retains the invertible
leading coefficient of the isomorphism
\cite[(2.10)--(2.11) and (2.15)--(2.16), pp.~7--9]{BhattacharyaEgger2019}.
For the oriented isomorphism used here, write
\(aT^{p^2}\) and \(bT^{p^2}\) for the first nonzero terms of the
two \(p\)-series and \(rT\) for the linear term of the isomorphism.
The equation intertwining the \(p\)-series gives
\(ra=br^{p^2}\), or \(r^{p^2-1}=a/b\).  Both \(a\) and \(b\) are units,
and \(p^2-1\) is prime to \(p\), so this is the finite étale equation
identifying the nonzero height-two coefficients.  The remaining equations
are precisely the coefficients of an isomorphism from the specialized
law to the Honda law.

The finite connected-marking stages of
\Cref{h3:thm:coordinate-comparison,h3:prop:integral-frames} represent
these truncated equations.  Their union represents the full
isomorphism algebra; every polynomial in the cooperation
coordinates belongs to a finite such stage.  Thus the ordinary
algebraic tensor product gives this algebraic union, without its
valuation completion.

It remains to check the parameter topology.  Reduction of
continuous power-series functions modulo \((p,x)\) gives
continuous functions into \(k[[y]]\): coefficient lifting proves
surjectivity, and continuous monomial division by \(p,x\) identifies
the kernel with \((p,x)C_{\hthreects}(Q,\widetilde A)\).
There is an equality
\begin{equation}\label{h3:des:compact-laurent}
 C_{\hthreects}(Q,k[[y]])[y^{-1}]
   = C_{\hthreects}(Q,k((y))).
\end{equation}
Indeed, the image of compact \(Q\) in the local field is bounded,
so multiplication by one power of \(y\) puts that image in
\(k[[y]]\).  The reverse inclusion in
\eqref{h3:des:compact-laurent} is immediate.  A finite étale
algebra over \(k((y))\) is a finite product of finite-dimensional
complete vector spaces over it; a basis gives the identical
assertion at each finite marking stage.  Consequently the final
union is a colimit on continuous cochain complexes with a common
stage for every compact family.

Lastly, all the identifications were made in the isomorphism
Hopf algebroid.  Its composition law transports the tautological
marking and is exactly the coface action.  The identity arrow
gives the constant test-unit.  This proves the assertions about
the maps, not just the coefficient modules.
\end{proof}

\subsection{The determinant map and the height-two basis}

Choose once and for all a topological generator of
\(\Z_p^\times/\mu_{p-1}\), and let \(c_{\mathrm{det}}\) be the reduced-determinant
logarithm normalized to send it to \(1\).  On the standard extended
height-three stabilizer set \(c_{\mathrm{det}}(g,\sigma)=c_{\mathrm{det}}(g)\) for \(g\in P_3\).
The reduced norm is invariant under Galois conjugation, so this is a
surjective continuous homomorphism to \(\Z_p\).  Let \(T^1\) be the
Devinatz--Hopkins fixed-point spectrum for its kernel, and let
\(\gamma\) represent the chosen generator in the residual quotient.
Their Proposition~8.1 gives a fiber sequence with middle map
\(1-\gamma\) and first map the unit.  Negating the target of the middle
map gives the convention used here,
\[
 T\longrightarrow T^1\xrightarrow{\gamma-1}T^1
       \xrightarrow{\partial_+}\Sigma T.
\]
The first map stays the same; if \(\partial_-\) denotes the boundary in
the \(1-\gamma\) sequence, the resulting boundary is
\(\partial_+=-\partial_-\).  By
\cite[Theorem~6 and Propositions~8.1--8.2]{DevinatzHopkins}, the boundary
of the unit is an actual permanent degree-one class detected by
\(\pm c_{\mathrm{det}}\).  We determine its sign in our cobar convention locally.

We use the cochain differential \(\partial=d^0-d^1\) in degree zero,
so \((\partial b)(g)=g b-b\).  For the standard field \(F=\F_{p^3}\),
\cite[Theorem~2(i)--(ii)]{DevinatzHopkins} identifies the Morava module
of \(T^1\) with \(C_{\hthreects}(\Z_p,E_{F,*})\) and its Adams
spectral sequence with continuous cohomology.  Realize this module with
the \(T^1\) factor first and the test \(E_F\) factor second.  If
\(i:T^1\to E_F\) is the standard map and \(h\) represents the coset with
\(c_{\mathrm{det}}(h)=t\), its homogeneous evaluation is
\(\operatorname{ev}_t=\mu\circ(h^{-1}i\otimes1)\).
The \(\Z_p\)-valued submodule \(C_{\hthreects}(\Z_p,\Z_p)\) contains
the unit and the lift needed below.  Write \(t\) for the coordinate
with \(c_{\mathrm{det}}(\gamma)=1\).  The actual evaluation formula gives two commuting
actions: the residual action on the \(T^1\) factor is
\[
 (\gamma_L f)(t)=f(t-1),
\]
while the action induced by the other Morava factor is
\((g_R f)(t)=g(f(t+c_{\mathrm{det}}(g)))\).  On the \(\Z_p\)-valued submodule used
here this is simply
\[
 (g_R f)(t)=f(t+c_{\mathrm{det}}(g)).
\]
Indeed, equivariance of \(i\) gives
\(\operatorname{ev}_t\circ(\gamma\otimes1)=\operatorname{ev}_{t-1}\),
using \(\gamma^{-1}h\) for the new representative, while
\(\operatorname{ev}_t\circ(1\otimes g)
 =g\circ\operatorname{ev}_{t+c_{\mathrm{det}}(g)}\), using \(hg\) for that
representative.  These are identities of the actual multiplication
maps, and the \(\Z_p\) values have trivial coefficient action.
Both displayed Morava modules are even,
so completed Morava homology of the fiber sequence gives the exact
coefficient sequence
\[
 0\longrightarrow E_{F,*}\longrightarrow
 C_{\hthreects}(\Z_p,E_{F,*})
 \xrightarrow{D=\gamma_L-1}
 C_{\hthreects}(\Z_p,E_{F,*})\longrightarrow0,
\]
where the first map includes the constant functions.  We fix the sign of every descent edge used below on its coefficient
Postnikov layer.  For normalized cosimplicial degree \(s\) and internal
cochain degree \(r=-t\), use total degree \(s+r\) and differential
\[
 d_{\mathrm{Tot}}=d_{\mathrm{cos}}+(-1)^s d_{\mathrm{int}},
 \qquad d_{\mathrm{cos}}=\sum_i(-1)^i d^i.
\]
Thus the unit edge is the constant \(1\), and the internal-degree-zero
cochain differential is the one just specified.  Model suspension by
\(K[1]^n=K^{n+1}\) with differential \(-d_K\), with the strict inverse
shift for desuspension.  The comparison
\(\operatorname{Tot}(K[1])\to(\operatorname{Tot}K)[1]\) is
multiplication by \((-1)^s\) in column \(s\).

Here is the resulting boundary sign.  For a degreewise exact sequence
of normalized cochain complexes \(0\to A\xrightarrow{i}B\xrightarrow{D}C\to0\),
use
\[
 \operatorname{Cone}(i)^n=B^n\oplus A^{n+1},\qquad
 d(b,a)=(d_Bb+i(a),-d_Aa),
\]
with inclusion \(j(b)=(b,0)\) and projection \(p(b,a)=a\) to \(A[1]\).
This cone has the rotation convention of Section~2: the map
\(a\mapsto((0,a),-i(a))\) from \(A[1]\) to \(\operatorname{Cone}(j)\)
identifies the next cone projection with \(-i[1]\).
The quotient \(q:\operatorname{Cone}(i)\to C\), \(q(b,a)=D(b)\), is a
quasi-isomorphism under \(B\), so the boundary is the roof \(pq^{-1}\).
For a cocycle \(c_0=D(b)\), put \(a_1=i^{-1}d_Bb\).  The cone cocycle
lifting \(c_0\) is \((b,-a_1)\), and its boundary is \(-a_1\).
In particular this boundary is the negative of the short-exact-cochain
connecting cocycle \(+a_1\).

This is also the sign of the actual desuspended boundary at its first
Adams edge.  On the coefficient Postnikov layer, the comparison from
the totalization of a termwise cone to the cone of the totalized map is
\((b,a)\mapsto(b,(-1)^s a)\) in column \(s\).  It commutes with the cone
projection and the displayed suspension comparison.  Evenness gives
the displayed short exact coefficient sequence, and the unit's boundary
first occurs in bidegree \((1,0)\); hence this layer computes its leading
edge.  The strict inverse shift contributes no further desuspension sign.

For the constant unit, take \(b(t)=-t\).  Then
\[
 (Db)(t)=b(t-1)-b(t)=1,\qquad
 (d_{\mathrm{cos}}b)(g)(t)=b(t+c_{\mathrm{det}}(g))-b(t)=-c_{\mathrm{det}}(g).
\]
The boundary-unit map \(\Sigma^{-1}(\partial_+\circ1)\) is therefore
detected by \(-d_{\mathrm{cos}}b=+c_{\mathrm{det}}\).  Define
\[
 \zeta=\Sigma^{-1}(\partial_+\circ1)\in\pi_{-1}T.
\]
It is an actual sphere map detected by the inflated normalized logarithm
\(c_{\mathrm{det}}\).  The normalization also agrees with the two-term group-ring
resolution: its degree-one generator maps to the bar generator
\([\gamma]\), whose boundary is \(\gamma-1\), so the corresponding
one-cocycle takes value \(1\) on \(\gamma\).
This normalization is used for all later descent edges.  On the
internal-degree-zero rows the total-complex Alexander--Whitney product
has no interchange sign, so these edges retain the usual cup products
and their naturality under the Moore quotient maps.
The relative coface
calculation restricts this class to the same logarithm on \(P_3\), and
finite-field descent retains this map and the unit.

\begin{theorem}\label{h3:thm:height-two-test}
Set
\[
 W=L_2S\vee\Sigma^{-1}L_2S,\qquad
 w=(1,\zeta):W\longrightarrow X.
\]
The first component is the canonical unit, and \(L_{K(2)}w\)
is an equivalence.  More explicitly, the specified maps give
\[
 (1,\zeta):R_2\vee\Sigma^{-1}R_2
       \xrightarrow{\ \simeq\ }L_{K(2)}T.
\]
\end{theorem}

\begin{proof}
The unit and the integral class just constructed are maps out
of \(S\).  Thus they define \(w\) by applying \(L_2\); no
extension of a map defined only on \(R_2\) is being asserted.

Use \Cref{h3:prop:exact-descent,h3:des:residue-cooperations} over
\(S_k\).  By \Cref{h3:thm:coheight-one-comparison}, the actual
coefficient augmentation gives
\[
 R\Gamma_{\hthreects}(H',L)\simeq k.
\]
It retains compact parameters and the action of \(P_3/H'\).
This quotient is \(\Z_p\), with trivial action on \(k\).
The continuous group-ring resolution
\[
 0\longrightarrow k[[\Z_p]]
   \xrightarrow{\gamma-1}k[[\Z_p]]
   \longrightarrow k\longrightarrow0
\]
therefore shows that the residue-tested spectral sequence has
only columns zero and one.  They are generated by the constant
unit and the normalized logarithm.  There is no possible
differential, and the specified maps \(1,\zeta\) detect these
two generators by the coface identification in
\Cref{h3:des:residue-cooperations}.

Because \(L_2\) is smashing and \(\overline E_2\) is
\(E(2)\)-local, the ordinary relative tests satisfy
\[
 \overline E_2\otimes_{S_k}(W\otimes_S S_k)
       \simeq \overline E_2\vee\Sigma^{-1}\overline E_2,
 \qquad
 \overline E_2\otimes_{S_k}(X\otimes_S S_k)
       \simeq \overline E_2\otimes_{S_k}R_{3,k}.
\]
The preceding two detected classes therefore prove that the test of
\(w\otimes_S S_k\) is an equivalence.  By associativity this test is
\(\overline E_2\otimes_S w\), and
\eqref{h3:des:absolute-relative-residue} identifies it with the
absolute ordinary residue test of \(w\).
The residue theory \(\overline E_2\) has Bousfield class \(K(2)\);
its coefficient ring is a graded field after a finite field
extension.  Consequently \(w\) is a \(K(2)\)-equivalence.
Applying \(L_{K(2)}\) to \(w\) therefore gives an equivalence with
source \(R_2\vee\Sigma^{-1}R_2\) and target \(L_{K(2)}T\).
The scalar passage preserved the two original \(S\)-maps by the
naturality of the test and the unit-compatible finite-field descent.
Its components are consequently the specified unit and \(\zeta\).
\end{proof}

\section{Removing completion on the ordinary stratum}\label{h3:sec:ordinary}

The ordinary stratum has two enlargements to remove: adjoining all
Frobenius roots, and allowing unbounded negative powers of \(x\).
Equivariant root retractions give the first comparison.  For the second,
we first prove finiteness for every compact power-series lattice using a
Laurent-tail estimate and a compact stable image.  The intermediate
finiteness argument below uses those lattices and the height-two map to prove the separate bounded-pole
finiteness needed to remove unbounded tails.

\subsection{Coefficient spaces and their topology}

Put $G=P_3$, $A=k[[x,y]]$, and $B_0=A[x^{-1}]$.
For $a,b>0$ write
\[
 v_{a,b}\left(\sum_{i\in\Z,\,j\geq0}c_{ij}x^iy^j\right)
   =\inf_{c_{ij}\ne0}(ai+bj),\qquad v_{a,b}(0)=+\infty.
\]
The ring $B_{\hthreehol}$ consists of the integral-exponent series for which,
for every $a,b>0$ and every real $N$, only finitely many nonzero terms
have $ai+bj\leq N$.  This is precisely convergence on
$0<|x|<1$, $|y|<1$ over trivially valued $k$.
Rational positive $a,b$ give a countable defining family of norms
$\exp(-v_{a,b})$.  Coefficientwise limits of Cauchy sequences show
completeness; Laurent polynomials give a countable dense subset.
Thus $B_{\hthreehol}$ is a Polish additive group.
The algebra $B_{\hthreepk}$ introduced before
\Cref{h3:thm:completed-cohomology} is, equivalently, the completion for
these norms of the union of the finite root levels
$B_{\hthreehol}^{1/p^n}$.  This identification is the comparison of function
algebras in
\Cref{h3:thm:coordinate-comparison,h3:thm:completed-cohomology}.
The countable union of the root Laurent polynomials is dense, so this
complete space is Polish as well.  Frobenius and its inverse are
continuous on $B_{\hthreepk}$, with valuations multiplied and divided
by $p$, respectively.

In particular, \Cref{h3:thm:completed-cohomology} concerns the ordinary
continuous $G$-cochain complex of this space, including continuous
profinite parameters.  We use both its natural constants comparison and its
finiteness assertion for coefficient lines.

The lines needed below are tensor products of integral powers of the
periodicity line and finite determinant-character lines; their duals
are included.  The canonical line is among them by the
canonical-line calculation used in \Cref{h3:lem:invariant-differential}.
Write $\mathcal L$ for such a line, with its original free
$A$-lattice $\mathcal L_A$, and set
\[
 A_d(\mathcal L)=x^{-d}\mathcal L_A\quad(d\in\Z),
 \qquad B_{\hthreehol}(\mathcal L)=B_{\hthreehol}\otimes_A\mathcal L_A,
 \qquad B_{\hthreepk}(\mathcal L)=B_{\hthreepk}\otimes_A\mathcal L_A.
\]
Each $A_d(\mathcal L)$ is $G$-stable and compact, with its
$(x,y)$-adic topology.  This is also its subspace topology in
$B_{\hthreehol}(\mathcal L)$: on a fixed pole lattice, every $v_{a,b}$
is bounded below by a positive multiple of total-degree order minus a
constant, while $v_{1,1}$ recovers that order up to the fixed lattice
shift.  The group preserves $(x,y)$ and carries $x$ to $x$ times a
unit; its action on the chosen line lattice is by units.

For clarity, the notation for algebraic localization always means
\begin{equation}\label{h3:ord:cochain-colimit}
 R\Gamma_{\hthreects}(G,B_0(\mathcal L))
   :=\colim_d C^\bullet_{\hthreects}(G,A_d(\mathcal L)).
\end{equation}
An additional profinite parameter space $T$ means applying
$C(T,-)$ at each displayed finite stage before taking the colimit.
No topology on an unrestricted union of cochain values is being used.

\subsection{Natural root projections}

\begin{lemma}[Equivariant root retractions]\label{h3:ord:root-retractions}
There is a continuous $G$-equivariant $B_0$-linear retraction
$B_0^{1/p}\to B_0$.  Its iterates give compatible retractions
from every finite root level and extend to a continuous retraction
\[
 \rho:B_{\hthreepk}\longrightarrow B_{\hthreehol}
\]
of the natural inclusion.  For every line $\mathcal L$ above, the same
construction gives a continuous equivariant retraction from
$B_{\hthreepk}(\mathcal L)$ to $B_{\hthreehol}(\mathcal L)$.
There are also an integer $e>0$, $q=p^e$, and continuous
$G$-equivariant additive
operators $\mathrm{Fr}$ and $\mathcal P$ on its holomorphic
coefficients such that
\begin{equation}\label{h3:ord:retraction-identity}
 \mathcal P\mathrm{Fr}=1.
\end{equation}
Here $\mathrm{Fr}$ is the $q$-power map in a horizontal frame.
It is invertible on the perfected coefficients, preserves some
$A_{-r}(\mathcal L)$, and $\mathcal P$ preserves $A_D(\mathcal L)$
for every sufficiently large $D$.
All these maps are compatible with constant-field descent.
\end{lemma}

\begin{proof}
By \Cref{h3:lem:ordinary-torsor}, the finite scheme of splittings of the
ordinary connected--\'{e}tale sequence at level $p^\ell$ is a torsor
under
\[
 \underline{\Hom}(\mathcal E[p^\ell],\mathcal C[p^\ell]),
\]
a form of $\mu_{p^\ell}^2$.  Its algebra is
$B_0^{1/p^\ell}$ over $B_0$.
This includes the divisor $y=0$.  A unit $p$-basis on this open is
$x,1+y$; adjoining a root of $1+y$ also adjoins the corresponding
root of $y$.  Thus multiplicative torsor coordinates can be units
throughout the open.  The additive monomial basis $x^iy^j$ used below
only expresses the matrix of the resulting operator; its equivariance
comes from the splitting torsor.
We recall why the construction supplies a retraction without a choice
of a splitting.  After an \'{e}tale cover the group is diagonalizable,
and a comodule is graded by its character group.  Projection onto the
degree-zero summand is linear over the invariant algebra.  Every
comodule map preserves this summand.  Changing the \'{e}tale frame
permutes characters and fixes zero, so the projections descend by
faithfully flat descent.  The invariant subalgebra of a torsor is the
base algebra, and its inclusion is fixed by this projection.
This proves the required retraction.  It is an additive, base-linear
projection; no multiplicativity or preservation of the ideal $(y)$ is
used.  The construction commutes with automorphisms of the ordinary sequence,
hence with $G$ and with coefficient-field automorphisms.  It uses no division by the order of
$\mu_{p^\ell}^2$.

For the untwisted coefficients choose $q=p^e$ with $q$ acting
trivially on $k$, and let $\rho_1:B_0^{1/q}\to B_0$ be the
corresponding retraction.  Set
$\mathcal P(f)=\rho_1(f^{1/q})$ and $\mathrm{Fr}(f)=f^q$.
Then $\mathcal P(g^qf)=g\mathcal P(f)$ and
\eqref{h3:ord:retraction-identity} holds.  On the level
$B_0^{1/q^n}$ define
\begin{equation}\label{h3:ord:finite-root-projection}
 \rho_n(f^{1/q^n})=\mathcal P^n f.
\end{equation}
These maps are $B_0$-linear.  They agree on overlapping levels because
$\mathcal P^{n+1}(f^q)=\mathcal P^nf$.
Thus a single natural root retraction already supplies the whole
compatible system.

Here is also the construction for the actual line twists.  If
$\omega$ denotes the periodicity line, the invariant height-one
section $v_1$ trivializes $\omega^{p-1}$ on this open.
In a free frame it is $x$ times that frame to the power $p-1$,
with the exponent reversed if the dual convention for $\omega$ is
used.  Adjoin a $(p-1)$st root of this section.  This is a finite
\'{e}tale frame cover on the open, and the resulting horizontal frame
has finite tame transition characters.  Include the finite
determinant character and choose $q$ to fix all these characters and
the constants.  The $q$-power map on coefficients in that frame then
descends to the line, as does the base-changed root retraction.
Projection onto the required tame eigensummand is exact, since its
order is prime to $p$.
On the $n$th root level of a line, the intrinsic version of
\eqref{h3:ord:finite-root-projection} is
$\rho_n=\mathcal P^n\mathrm{Fr}^n$; thus its notation does not
require taking a root of a chosen line generator.

In the original rational frame the coefficients in a horizontal frame
have $x$-exponents in a fixed translate $\lambda+\Z$, with
$(q-1)\lambda\in\Z$.  More explicitly, write the horizontal
frame as $h=x^{-\lambda}e_0$, where $e_0$ is the original line
frame, and put $\kappa=(q-1)\lambda$.  Then
\[
 \mathrm{Fr}(fe_0)=f^q x^\kappa e_0.
\]
Thus Frobenius multiplies normalized order by $q$.
For line coefficients we use these normalized Gauss
valuations: relative to the original frame they add the fixed number
$a\lambda$ to $v_{a,b}$.  They define the same topology, and satisfy
$v_{a,b}(\mathrm{Fr}f)=qv_{a,b}(f)$ exactly.
Choosing $r+\lambda\geq0$ makes
$x^r\mathcal L_A$ Frobenius-stable.  This also shows explicitly that
passing to the tame frame introduces only a fixed shift of pole
orders.  The continuity, extension to completion, and assertion about
$A_D$ follow from the estimates proved next; these estimates apply to
the descended eigensummand as well.
\end{proof}

\begin{lemma}[Uniform loss bounds]\label{h3:ord:loss-bounds}
Fix the operators of \Cref{h3:ord:root-retractions} and a free rational
frame of $\mathcal L$, using the normalized Gauss valuations just
specified.  There are constants $c\geq0$ and
$c_{a,b}\geq0$ such that
\begin{align}
 \mathcal P(A_L(\mathcal L))
   &\subseteq A_{L/q+c}(\mathcal L),\label{h3:ord:one-pole-bound}\\
 v_{a,b}(\mathcal Pf)
   &\geq q^{-1}v_{a,b}(f)-c_{a,b}.
   \label{h3:ord:one-norm-bound}
\end{align}
For real indices the pole inequality defines the lattice, so
$A_t=A_{\lfloor t\rfloor}$ in the original integral frame.  With
$C=c/(1-q^{-1})$ and $C_{a,b}=c_{a,b}/(1-q^{-1})$, iteration gives
\begin{align}
 \mathcal P^n(A_L(\mathcal L))
   &\subseteq A_{q^{-n}L+C}(\mathcal L),\label{h3:ord:iterated-poles}\\
 v_{a,b}(\mathcal P^nf)
   &\geq q^{-n}v_{a,b}(f)-C_{a,b}.
   \label{h3:ord:iterated-norms}
\end{align}
The losses for $\rho_n$ in \eqref{h3:ord:finite-root-projection} are
therefore independent of the root level.
\end{lemma}

\begin{proof}
Let $e_0$ be the original free line frame, and temporarily write
$w_{a,b}$ for the unshifted scalar Gauss valuation.  Every scalar
coefficient has the unique decomposition
\[
 f=\sum_{0\leq i,j<q}x^iy^j f_{ij}^{\,q}.
\]
The inverse-semilinear operator is consequently given by a finite
matrix, here a single row:
\[
 \mathcal P(fe_0)=\left(\sum_{0\leq i,j<q} b_{ij}f_{ij}\right)e_0,
 \qquad b_{ij}\in B_0.
\]
The terms of the decomposition have disjoint exponent residue classes,
so
\[
 w_{a,b}(f)=\min_{i,j}\{ai+bj+qw_{a,b}(f_{ij})\}.
\]
Put
\[
 c^0_{a,b}=\max\left(0,
      \max_{i,j}\left\{\frac{ai+bj}{q}-w_{a,b}(b_{ij})\right\}\right).
\]
The scalar output has valuation at least
$q^{-1}w_{a,b}(f)-c^0_{a,b}$.
Since $v_{a,b}(fe_0)=w_{a,b}(f)+a\lambda$, its normalized
valuation is at least
\[
 q^{-1}v_{a,b}(fe_0)-c^0_{a,b}+a\lambda(1-q^{-1}).
\]
Thus \eqref{h3:ord:one-norm-bound} holds, for example, with
$c_{a,b}=c^0_{a,b}+a|\lambda|(1-q^{-1})$.
All $b_{ij}$ have one common finite pole bound.
For $f\in A_L$, its components have pole at most
$L/q+(q-1)/q$.  This proves \eqref{h3:ord:one-pole-bound} with a
constant independent of $L$.
Each iteration divides the previous loss by $q$, giving the displayed
geometric sums and proving \eqref{h3:ord:iterated-poles} and
\eqref{h3:ord:iterated-norms}.

The component decomposition and finite matrix are continuous on
$B_{\hthreehol}$: taking an exponent residue class and then a root
preserves convergence in every defining norm.  On a root level the
last inequality reads
\[
 v_{a,b}(\rho_n h)\geq v_{a,b}(h)-C_{a,b}.
\]
Thus the compatible maps extend continuously from the dense root
union to $B_{\hthreepk}$, with values in the complete space $B_{\hthreehol}$.
Their common restriction is the identity on $B_{\hthreehol}$.
Finally $D>C$ implies $D/q+c<D$, proving stability of $A_D$.
The proof on the tame frame cover uses its finitely many character
components and the same finite matrix calculation.  Descent therefore
retains both estimates for every specified twist.
\end{proof}

\begin{proposition}[Holomorphic comparison]\label{h3:prop:holomorphic-comparison}
For every specified coefficient line the natural map
\[
 R\Gamma_{\hthreects}(G,B_{\hthreehol}(\mathcal L))
   \longrightarrow R\Gamma_{\hthreects}(G,B_{\hthreepk}(\mathcal L))
\]
is a quasi-isomorphism.  Its cohomology groups are finite, and the
chosen Frobenius power acts invertibly on them.  For a profinite
parameter space $T$, the comparison and Frobenius invertibility persist,
and the cohomology is naturally
\[
 C\bigl(T,H^a_{\hthreects}(G,B_{\hthreehol}(\mathcal L))\bigr),
\]
where the finite point-value cohomology group has its discrete topology.
Finiteness is asserted for the cohomology without parameters.
The comparison is compatible with coefficient-field descent.
For the untwisted line it identifies
the natural constants comparison of \Cref{h3:prop:ordinary-constants}.
\end{proposition}

\begin{proof}
The retraction induces a split injection on cohomology into the finite
groups of \Cref{h3:thm:completed-cohomology}.  Hence holomorphic
cohomology is finite before any bounded-pole assertion is used.
The identity $\mathcal P\mathrm{Fr}=1$ makes Frobenius injective
on these finite groups, hence bijective.  Both cochain complexes are
Polish, so their finite cohomology groups are discrete by
\Cref{h3:co:strictness}.

For completeness, consider a completed class, rather than merely a
class in the algebraic root union.  On perfected coefficients define
\[
 R_\ell=\mathrm{Fr}^{-\ell}\rho\mathrm{Fr}^{\ell}:
 B_{\hthreepk}(\mathcal L)\longrightarrow
 B_{\hthreehol}(\mathcal L)^{1/q^\ell}.
\]
This is a projection onto the indicated root level and satisfies
\begin{equation}\label{h3:ord:root-approximation-bound}
 v_{a,b}(R_\ell f)\geq v_{a,b}(f)-q^{-\ell}C_{a,b}.
\end{equation}
It converges to the identity uniformly on compact subsets in every
defining norm.  Indeed, approximate a compact subset, at any prescribed
accuracy in finitely many norms, by finitely many elements at a common
finite root level.  For every later $\ell$, $R_\ell$ fixes these
approximants.  The bound \eqref{h3:ord:root-approximation-bound} controls
the errors uniformly; increasing the accuracy proves the assertion.

If $z$ is a continuous completed cocycle, its image is compact.
Thus $R_\ell z\to z$ in the cochain topology.  These are cocycles
because the projections are equivariant.  Discreteness of completed
cohomology gives $[R_\ell z]=[z]$ for all sufficiently large $\ell$.
Every completed class therefore comes from a finite root level.
Under Frobenius transport, the image of this level is the image of
holomorphic cohomology, because Frobenius is already invertible on
that image.  This proves surjectivity and hence the comparison.

The finite-cohomology lifting assertion of \Cref{h3:co:parameters}
now applies to these two Polish complexes.  It proves the comparison
with any profinite parameters.  All the maps used to define the
comparison are the natural inclusions; the retractions prove their
properties.  In particular the comparison preserves products and the
original constants maps.  Naturality of the torsor construction also
proves its coefficient-field equivariance.
\end{proof}

\subsection{Uniform Laurent tails and compact stable images}

Let $T_D$ discard the terms of pole order at most $D$, retaining only
$x^iy^j$ with $i<-D$, in a chosen rational frame.  It is a continuous
linear operator on $B_{\hthreehol}(\mathcal L)$; its complementary
projection has image $A_D(\mathcal L)$.  The frame shifts in
\Cref{h3:ord:root-retractions} can be absorbed in the constants of
\Cref{h3:ord:loss-bounds}.

\begin{lemma}[Tail contraction on compact families]\label{h3:ord:tail-contraction}
Choose an integer $D>C$ in \Cref{h3:ord:loss-bounds}.
For every compact subset $K\subset B_{\hthreehol}(\mathcal L)$,
\[
 T_D\mathcal P^nf\longrightarrow0
 \quad\hbox{uniformly for }f\in K
\]
in every defining Gauss norm.
\end{lemma}

\begin{proof}
Choose $0<\delta<D-C$ and set $L_n=\lfloor\delta q^n\rfloor$.
The pole estimate implies
\[
 \mathcal P^n(1-T_{L_n})f\in A_D(\mathcal L),\qquad
 T_D\mathcal P^nf=T_D\mathcal P^nT_{L_n}f.
\]
Fix $a,b>0$.  For any $R>0$, compactness gives a finite lower bound
\[
 m_R:=\inf_{f\in K}v_{a+R,b}(f)>-\infty.
\]
A monomial with pole $i>L$ has normalized $v_{a,b}$-value equal
to its $v_{a+R,b}$-value plus $R(i-\lambda)$.  Consequently
\[
 \inf_{f\in K}v_{a,b}(T_{L_n}f)\geq m_R+R(L_n-\lambda).
\]
Truncation cannot lower a Gauss valuation.  By the iterated norm bound,
\begin{equation}\label{h3:ord:quantitative-tail}
 \inf_{f\in K}v_{a,b}(T_D\mathcal P^nf)
   \geq q^{-n}(m_R-R\lambda)+Rq^{-n}L_n-C_{a,b}.
\end{equation}
The lower limit is at least $R\delta-C_{a,b}$.  Since $R$ can
be arbitrarily large, it is $+\infty$.  This includes $\lambda=0$
for the untwisted line.
\end{proof}

\begin{lemma}[The compact kernel and its stable image]\label{h3:ord:stable-image}
For $D$ as above and every $a\geq0$, put
\[
 M=H^a_{\hthreects}(G,A_D(\mathcal L)),\qquad
 \phi:M\longrightarrow H^a_{\hthreects}(G,B_{\hthreehol}(\mathcal L)).
\]
Then $M$ is compact Hausdorff, and
$\mathcal P^nz\to0$ for every $z\in\ker\phi$.
Moreover its stable image
\[
 M_\infty=\bigcap_{n\geq0}\mathcal P^nM
\]
is finite.
\end{lemma}

\begin{proof}
Compactness, Hausdorffness and the asserted cohomology topology follow
from the finite completed resolution in \Cref{h3:co:compact-resolution}.
The continuous map $\phi$ has finite discrete target by
\Cref{h3:prop:holomorphic-comparison}, so $K=\ker\phi$ is compact.
For $a=0$ it is zero, since the coefficient inclusion is injective.
For $a>0$, represent $z\in K$ by an $A_D$-valued cocycle $z_0$
and write $z_0=\partial b$ with a holomorphic continuous cochain $b$.
Set
\[
 b_n=(1-T_D)\mathcal P^nb,\qquad
 e_n=\mathcal P^nz_0-\partial b_n
      =\partial(T_D\mathcal P^nb).
\]
The first expression shows that $e_n$ is an $A_D$-valued cocycle.
The image of $b$ is compact, so \Cref{h3:ord:tail-contraction} makes
$T_D\mathcal P^nb$ tend uniformly to zero.
The group differential preserves this convergence in total-degree
order: every group element preserves $(x,y)$, sends $x$ to $x$ times
a unit and sends a line generator to a unit multiple.  Thus its
action does not decrease total-degree order in a fixed rational
frame.  The finite sum defining $\partial$ has the same property.
It follows that $e_n\to0$ in the compact lattice topology, and
$\mathcal P^nz=[e_n]\to0$ in $M$.

We give the uniformity step explicitly.  For an open additive subgroup
$U$ of $K$, let
\[
 K_N(U)=\bigcap_{n\geq N}\{z\in K:\mathcal P^nz\in U\}.
\]
These are increasing closed subgroups covering $K$ by the convergence
just proved.  The Baire theorem makes one of them have interior; as a
subgroup it is open and has finite index.  Choose representatives for
its finitely many cosets.  Each representative has all sufficiently
late iterates in $U$.  Taking the largest of these finitely many
bounds gives $\mathcal P^nK\subseteq U$ for all sufficiently
large $n$.  Hence
\begin{equation}\label{h3:ord:kernel-intersection}
 \bigcap_{n\geq0}\mathcal P^nK=0.
\end{equation}

The operator $\mathcal P$ is bijective on holomorphic cohomology:
it is the inverse of Frobenius there by finiteness and
\eqref{h3:ord:retraction-identity}.
If $z\in M_\infty\cap K$, choose, for each $n$, an $m_n\in M$
with $\mathcal P^nm_n=z$.  Applying $\phi$ and using this
bijectivity gives $m_n\in K$.  Equation
\eqref{h3:ord:kernel-intersection} therefore forces $z=0$.
Thus $M_\infty$ injects into the finite holomorphic cohomology group,
which proves its finiteness.
\end{proof}

\begin{proposition}[Finiteness of all ordinary lattices]
\label{h3:prop:ordinary-lattice-finiteness}
For every $a\geq0$, $d\in\Z$, and coefficient line $\mathcal L$
specified above, the group
\[
 H^a_{\hthreects}(P_3,A_d(\mathcal L))
\]
is finite.  This includes all periodicity, determinant, conormal, and
canonical-line twists and their inverses, with the original
$(x,y)$-adic lattice topology.  The conclusion uses no finiteness
or cohomology comparison for $B_0$.
\end{proposition}

\begin{proof}
Each adjacent quotient
$A_j(\mathcal L)/A_{j-1}(\mathcal L)$ is a free
$k[[y]]$-line.  Its action is the restriction of $\mathcal L$
tensored with the appropriate power of the conormal of $x=0$.
That conormal is a periodicity power.  Thus
\Cref{h3:co:base-lattices} says that every such quotient has finite
$G$-cohomology.  Finite strips of adjacent quotients have finite
cohomology as well, by their finite filtration.  The sequences are
exact on continuous cochains: the finite Laurent-coordinate
truncations give continuous nonequivariant sections.

Choose $A_{-r}(\mathcal L)\subseteq A_D(\mathcal L)$ with the
former Frobenius-stable and the latter $\mathcal P$-stable, as in
\Cref{h3:ord:root-retractions}.  If
$z\in H^a(G,A_{-r}(\mathcal L))$, then its image in $M$ satisfies
\[
 \iota(z)=\mathcal P^n\iota(\mathrm{Fr}^nz)
 \quad\hbox{for every }n.
\]
It therefore belongs to the finite subgroup $M_\infty$ of
\Cref{h3:ord:stable-image}.  The kernel of this map is a quotient of
$H^{a-1}(G,A_D/A_{-r})$, which is finite by the strip calculation
(and is zero for $a=0$).
Hence $H^a(G,A_{-r})$ is finite.  The same strip sequence then
proves finiteness for $A_D$; further finite strips reach every
$A_d$, in either direction.  All twists have been retained throughout.
\end{proof}

\subsection{Intermediate finiteness and bounded poles}
\label{h3:sec:intermediate-finiteness}

The compact lattice calculation now supplies a finite chromatic
abutment.  We use it to bound the algebraic bounded-pole cohomology
before invoking the ordinary comparison.

Write \(M_3\) for the fiber of \(L_3\to L_2\).

\begin{lemma}\label{h3:des:monochromatic-finiteness}
The groups
\[
 \pi_j(T/p),\qquad \pi_j(M_3S/p),\qquad
 \pi_j(X/p),\qquad \pi_j(L_1X/p)
\]
are finite for every integer \(j\).  This assertion uses
\Cref{h3:prop:ordinary-lattice-finiteness} but does not use
\Cref{h3:thm:ordinary-comparison}.
\end{lemma}

\begin{proof}
\emph{The height-three local term \(T/p\).}
Apply the finite Moore quotient to the relative object of
\Cref{h3:prop:exact-descent}.  Its abutment is
\(\pi_*(R_{3,k}/p)\), and its coefficient in internal degree
\(2j\) is \(A\otimes\omega^j\); odd coefficients vanish.
These compact modules have finite \(P_3\)-cohomology by
\Cref{h3:prop:ordinary-lattice-finiteness}.  They are complete and
separated with invariant open neighborhoods, so the reduced
\(\F_p[[P_3]]\)-resolution of \Cref{h3:co:compact-resolution},
which has length nine, computes their continuous cochains.  Hence only
columns \(0\leq s\leq9\) occur.  The bounded convergence from
\Cref{h3:prop:exact-descent} gives finitely many finite associated-graded
groups in each total degree, so \(\pi_*(R_{3,k}/p)\) is degreewise
finite.  The underlying equivalence
\(R_{3,k}/p\simeq(T/p)^{\vee [k:\F_p]}\) then gives finiteness for
\(T/p\) over \(S\).

\emph{The monochromatic term \(M_3S/p\).}
The chromatic fracture square gives natural equivalences
\[
 M_3S\simeq M_3T\simeq\fib(T\longrightarrow L_2T).
\]
The first identity is also \cite[Theorem~6.19]{HS99}.
Put $E=E_3(k)$, $\widetilde A=\pi_0E=W(k)[[x,y]]$ and
$I=(p,x,y)$. The module-spectrum form of chromatic localization is
needed here. For every ordinary $E$-module $N$, including a completed
continuous-function term, there is a natural fiber sequence
\[
 \operatorname{Kos}_E(I)\otimes_E N\longrightarrow N
       \longrightarrow L_2N,
 \qquad
 \operatorname{Kos}_E(I)=
 \bigotimes_{z\in\{p,x,y\}}\fib(E\longrightarrow E[z^{-1}]).
\]
This is \cite[Proposition~3.4(1) and Remark~3.5]{BarthelStapleton2016},
with its support functor described in Lemmas~2.3 and~2.5 of that paper.
The statement is about the underlying chromatic localization of
arbitrary module spectra; it imposes no flatness or finite-generation
condition. Since $E$ is $E(3)$-local and $L_3$ is smashing, every
$E$-module is $E(3)$-local. Thus the first term is $M_3N$.
In particular it applies to the actual completed Amitsur spectrum
$N_Q=\mathcal N_3^q$ for $Q=P_3^q$, equipped with
\eqref{h3:des:relative-cooperations} and viewed as an ordinary
$E$-module; no completed tensor product is replaced by an ordinary one.

The spectrum $N_Q/p$ has zero $p$-inversion, so its $p$-support
factor is the whole spectrum. The remaining two factors give the
finite \v Cech support complex for $(x,y)$. Its homotopy spectral
sequence is
\[
 H^s_{(x,y)}(\pi_t(N_Q/p))
       \Longrightarrow \pi_{t-s}(M_3(N_Q/p)),\qquad 0\leq s\leq2.
\]
By \Cref{h3:des:flat-functions}, the even coefficient modules of
$N_Q/p$ are $C_{\hthreects}(Q,A)\otimes\omega^{t/2}$,
where $A=k[[x,y]]$, and are flat over $A$; the odd ones vanish.
The regular sequence $x,y$ therefore makes the displayed support
cohomology vanish except for $s=2$. For $Q$ a point it is
\[
 H^2_{(x,y)}A
 =A[x^{-1},y^{-1}]/(A[x^{-1}]+A[y^{-1}]).
\]
Thus the shift is exactly two. The finite support complex introduces
no further convergence issue.

These identifications also respect the action-cobar maps.  All faces
other than the first, and all degeneracies, are $E$-linear; the first face uses the
continuous stabilizer action. The invariant ideal $I$ is preserved
even for this varying family: in
$C_{\hthreects}(P_3^{q+1},\widetilde A)$ one has
$(p,g_1(x),g_1(y))=(p,x,y)$. Indeed, the continuous monomial-division
operations used in \Cref{h3:des:flat-functions} express any element
of $I$ as a continuous sum of $p,x,y$ multiples. Apply them to
$g_1^{-1}(x)$ and $g_1^{-1}(y)$, then apply $g_1$; joint continuity
gives continuous coefficients expressing $x,y$ in the image ideal.
The reverse containment follows in the same way. For precision, let $I_q$ be the image of $I$ in $\pi_0N_Q$.
Base change of the displayed finite fiber construction from $E$ to
$N_Q$ identifies the ordinary $E$-support term with
$\operatorname{Kos}_{N_Q}(I_q)$.  Along an actual coface ring map
$N_Q\to N_{Q'}$, its further base change is the construction for the
image generators in $\pi_0N_{Q'}$.  The support functor on ordinary
$N_{Q'}$-modules is the canonical right adjoint for modules supported
at that finitely generated ideal, so it depends on the ideal, not on
the chosen generators.  The equality of ideals just proved identifies
the first-face image with the target support ideal.  Uniqueness of this
right adjoint makes the resulting comparisons compatible with the first face and with compositions.  No individual generator
is assumed fixed by the action. The calculation therefore applies to the
whole completed Amitsur object, and \Cref{h3:prop:exact-descent}
permits its termwise use before totalization.


For a continuous function term the same quotient is
\begin{equation}\label{h3:des:principal-parts-parameters}
 H^2_{(x,y)}C_{\hthreects}(Q,A)
     =C_{\mathrm{lc}}(Q,H^2_{(x,y)}A),
\end{equation}
where the target coefficient module is discrete.
To verify this, express \(H^2_{(x,y)}A\) as the colimit of
\(A/(x^a,y^a)\), with transition multiplication by \(xy\).
At each stage, quotienting continuous functions gives continuous
functions to the finite discrete quotient, using coefficient
truncation to lift them.  A locally constant map from compact
\(Q\) has finite image, so one common stage represents it.
This proves \eqref{h3:des:principal-parts-parameters} with the
required denominator and parameter bounds.

Residue pairing, followed by \(\hthreeTr_{k/\F_p}\), identifies the
continuous dual of
\(H^2_{(x,y)}(A\otimes\omega^j)\) with the compact module
\[
 M=A\otimes\omega^{-j}\otimes
       {\textstyle\bigwedge^2}\Omega^1_{A/k,\hthreects}.
\]
Concretely the pairing is the coefficient of \(x^{-1}y^{-1}\)
in a product with a two-form.  It is perfect on finite
principal-parts quotients and hence on their colimit against
the inverse-limit power-series module.  Its change-of-variables
formula proves equivariance.  The canonical-line identity of
\Cref{h3:lem:invariant-differential} makes this an admissible twist
in \Cref{h3:prop:ordinary-lattice-finiteness}.
The principal-parts module is discrete and \(M\) is compact.
The reversed compact/discrete duality
\eqref{h3:co:compact-duality}, applied to \(P_3\) in dimension nine,
therefore gives
\[
 H^s_{\hthreects}\bigl(P_3,H^2_{(x,y)}(A\otimes\omega^j)\bigr)
   \cong H^{9-s}_{\hthreects}(P_3,M)^\vee.
\]
The right side is finite by
\Cref{h3:prop:ordinary-lattice-finiteness}, so the principal-parts
cohomology is finite in every degree.
The relative descent spectral sequence for
\(M_3(R_{3,k}/p)\) is bounded in columns by the same reduced
length-nine resolution and has finite entries in every total degree.
It therefore has degreewise finite abutment.  Its underlying spectrum
is a sum of \([k:\F_p]\) copies of \(M_3(T/p)\), because the exact
functor \(M_3(-/p)\) preserves the finite free decomposition of
\(R_{3,k}\).  Using \(M_3T\simeq M_3S\) proves finiteness of
\(\pi_*(M_3S/p)\).

\emph{The overlap \(X/p\).}
The fiber sequence \(M_3S/p\to T/p\to X/p\) gives finiteness
for \(X/p\).

\emph{The height-one localization \(L_1X/p\).}
By \Cref{h3:thm:height-two-test},
\(L_{K(2)}X/p\) is the sum of two shifts of \(R_2/p\).
The completion map identifies \(S/p\) with the finite torsion Moore
spectrum \(\mathbb S_{(p)}/p\), and exactness of \(L_{K(2)}\) gives
\[
 R_2/p\simeq L_{K(2)}(S/p).
\]
For \(p\geq5\), \cite[Theorem~15.1]{HS99} applied at height two to this
finite torsion spectrum says that every integer-graded homotopy group is
finite; see also \cite[p.~49(c)]{BB}.

For the other height-two boundary put \(E=L_1R_2\).  The completion
map from the \(p\)-local sphere to \(S\) has rational cofiber, so it is
a \(K(2)\)-equivalence.  Thus \(R_2\) is the \(K(2)\)-local sphere in
the group calculation of \cite[Remark~7.8]{Behrens2012}.  We use the
following features of that formula.  Write \(d_v=|v_1|=2(p-1)\).
Besides finitely many shifted copies of \(\Z_p\), \(\Q_p\), and
\(\Q/\Z_{(p)}\), the formula contains finite cyclic \(p\)-groups with
generators \(1_{sp^n/(n+1)}\), where \(n\geq0\), \(p\nmid s\), and the
internal degree is \(d_vsp^n\).  Its exterior factor on the degree-minus-one
class denoted here by \(\zeta_{\mathrm B}\) adds only two shifts.  This
notation distinguishes that height-two class from the map \(\zeta\)
constructed above.

For a fixed nonzero internal degree, division by \(d_v\), when possible,
determines at most one integer of the form \(sp^n\) with \(p\nmid s\):
\(n\) is its \(p\)-adic valuation and \(s\) its remaining \(p\)-unit
integer.  The cyclic family therefore contributes at most one summand
in each such degree before the two exterior shifts.  Multiplication by
\(p\) has finite kernel and cokernel on each listed group: these pairs
are \((0,\F_p)\), \((0,0)\), and \((\F_p,0)\) on \(\Z_p\), \(\Q_p\),
and \(\Q/\Z_{(p)}\), respectively, and they are finite on a finite
cyclic \(p\)-group.  Hence both \(\pi_j(E)[p]\) and \(\pi_j(E)/p\) are
finite for every integer \(j\).  The Moore cofiber sequence gives
\[
 0\longrightarrow\pi_j(E)/p\longrightarrow\pi_j(E/p)
   \longrightarrow\pi_{j-1}(E)[p]\longrightarrow0.
\]
It follows that \(\pi_j(E/p)\) is finite in every degree.  Exactness
identifies \(E/p\simeq L_1(R_2/p)\), the boundary needed here.
The local height-two equivalence then identifies \(L_1L_{K(2)}(X/p)\)
with two shifts of \(E/p\), so that corner is degreewise finite.
Finally, apply the height-two fracture square to \(X/p\).
Its other three corners are degreewise finite, so its long
exact sequence proves finiteness of \(L_1X/p\).
No bounded-pole comparison has been used in this argument.
\end{proof}

\begin{proposition}\label{h3:prop:bounded-pole-finiteness}
For every \(a\geq0\), the group
\(H^a_{\hthreects}(P_3,B_0)\), with the prescribed bounded-pole
cochain convention, is finite.
\end{proposition}

\begin{proof}
Apply \(L_1(-/p)\) to the completed relative descent object.
We first identify this exact test on an actual term \(N_Q=\mathcal N_3^q\).
For odd \(p\), an Adams self-map of the finite Moore spectrum has degree
\(2(p-1)\) and acts nontrivially in \(K(1)\)-homology
\cite[Example~2.4.1(ii)]{RavenelNilpotence}.  Its action on
\(BP_*(S/p)=BP_*/p\) is multiplication by an element of degree
\(2(p-1)\), and that graded group is \(\F_p v_1\).  This even-degree
coefficient detects the entire \(BP\)-module map: applying maps into
\(BP/p\) to \(BP\xrightarrow p BP\to BP/p\) shows that the only
possible extra term is \(\pi_{2(p-1)+1}(BP/p)=0\).  If the coefficient
were zero, the \(BP\)-module map would be null, and so would its base
change to \(K(1)\), contrary to the chosen Adams map.  Its coefficient
is therefore a unit multiple of \(v_1\).  Rescaling the self-map by an
integer unit makes its \(BP\)-Hurewicz action exactly \(v_1\).
The height-three orientation sends this element to \(v_1=xU^{p-1}\)
modulo \(p\).  This unit rescaling does not change the telescope:
powers of the unit give an isomorphism between the two defining towers.

The height-one telescope theorem for this finite Moore spectrum is the
consequence of \cite[Theorem~4.11 and Corollary~4.12]{Miller1981}
identified as the \(n=1\) case in
\cite[Section~7.5, after Conjecture~7.5.5]{RavenelNilpotence}.
Since \(L_1\) is smashing, it gives, for the already formed ordinary
module \(N_Q\),
\[
 L_1(N_Q/p)\simeq N_Q\wedge L_1(S/p)
       \simeq N_Q\wedge\operatorname{Tel}(v_1:S/p\to\Sigma^{-2(p-1)}S/p).
\]
Thus the test is a termwise ordinary telescope of the same finite Moore
self-map.  Homotopy groups commute with that telescope.  In internal
degree zero the resulting coefficient is exactly
\[
 \colim_d x^{-d}C_{\hthreects}(P_3^q,A).
\]
Its other coefficients are the periodic translates.  The actual action
preserves every bounded-pole lattice modulo \(p\), because
\(g(x)/x\) is a unit in \(A\).  This is the prescribed cochain complex
of \(B_0\); no completion of the algebraic localization is taken.

The central tame units act trivially on \(A\) and by their
scalar powers on \(U\).  Averaging over \(\mu_{p-1}\) kills
internal degrees not divisible by \(2(p-1)\).
Modulo \(p\), \(v_1\) is invariant, so multiplication by
its powers identifies all the remaining coefficient complexes.
For each compact lattice \(x^{-d}A\), the reduced finite
\(\F_p[[P_3]]\)-resolution of \Cref{h3:co:compact-resolution}
computes its continuous cochains and has length nine.  Taking the exact
filtered colimit of these complexes proves the same column bound for
this prescribed bounded-pole union; it does not apply a compact-module
theorem directly to the union.  Thus the spectral sequence has
\[
 E_2^{a,\,2b(p-1)}
       =H^a_{\hthreects}(P_3,B_0)\,v_1^b,\qquad
 0\leq a\leq9,
\]
and no other nonzero internal degrees.  The abutment is the relative
one, \(\pi_*L_1(R_{3,k}/p)\).  If \(m=[k:\F_p]\), the underlying
finite free decomposition gives
\[
 L_1(R_{3,k}/p)\simeq L_1(T/p)^{\vee m}
       \simeq L_1(X/p)^{\vee m}.
\]
Here \(L_1L_2=L_1\), and all functors involved preserve finite sums.
Thus this relative abutment is degreewise finite by
\Cref{h3:des:monochromatic-finiteness}.  The error-tower bound in
\Cref{h3:prop:exact-descent} and the finite column range give the
bounded filtration whose associated graded is \(E_\infty\).

The differential has bidegree \((r,r-1)\), so a nonzero
differential requires \(2(p-1)\mid r-1\).
For \(p\geq7\), no such \(r\geq2\) fits the column range;
the finite abutment then gives the result directly.
For \(p=5\), the only possible differential is
\[
 d_9:E_9^{0,t}\longrightarrow E_9^{9,t+8}.
\]
Its source is finite.  Indeed, coefficient inclusion gives
\[
 H^0(P_3,B_0)\lhook\joinrel\longrightarrow
 H^0(P_3,B_{\hthreehol}),
\]
and the target is finite by
\Cref{h3:prop:holomorphic-comparison}.
All earlier differentials vanish by the displayed sparsity,
and hence the \(d_9\)-image is finite.  Column zero is already finite
by this invariants argument.  Columns one through eight are unchanged
and are finite from their \(E_\infty\) groups.  In column \(9\),
\(E_2^{9,t}=E_9^{9,t}\) is an extension of its finite
\(E_\infty\)-quotient by that finite image, and is finite.
This proves the proposition also at five, without asserting
that \(d_9\) vanishes.
\end{proof}

\subsection{The natural ordinary comparison}
\label{h3:sec:ordinary-comparison}

The preceding proposition supplies the finiteness needed for the second
decompletion.  We now remove unbounded Laurent tails by combining that
finiteness with the contraction already proved on every fixed quotient.

\begin{theorem}[Natural ordinary comparison]\label{h3:thm:ordinary-comparison}
The natural coefficient maps induce quasi-isomorphisms
\begin{equation}\label{h3:ord:natural-comparison}
 \colim_d C^\bullet_{\hthreects}(P_3,x^{-d}A)
 \longrightarrow C^\bullet_{\hthreects}(P_3,B_{\hthreehol})
 \longrightarrow C^\bullet_{\hthreects}(P_3,B_{\hthreepk}).
\end{equation}
They preserve products, constants, and coefficient-field descent, and
remain quasi-isomorphisms with any profinite space of continuous
parameters in the convention \eqref{h3:ord:cochain-colimit}.
Every Frobenius power, including the $p$th-power map, acts invertibly
on these untwisted cohomology groups.
Consequently there is the natural identification
\[
 H^*_{\hthreects}(P_3,B_0)
   \cong H^*_{\hthreects}(P_1\times\GL_2(\Z_p),k).
\]
\end{theorem}

\begin{proof}
By \Cref{h3:prop:bounded-pole-finiteness}, the cohomology of the
left-hand complex in \eqref{h3:ord:natural-comparison} is finite.
The lattice, height-two, and spectral arguments establishing this input
have already been completed.

First keep $D>C$ fixed.  The quotient
$Q_D=B_{\hthreehol}/A_D$ is Polish and has the continuous tail section
$T_D$.  Hence
\[
 0\longrightarrow C^\bullet_{\hthreects}(G,A_D)
 \longrightarrow C^\bullet_{\hthreects}(G,B_{\hthreehol})
 \longrightarrow C^\bullet_{\hthreects}(G,Q_D)\longrightarrow0
\]
is exact degreewise.  Its long exact sequence, together with
\Cref{h3:prop:holomorphic-comparison,h3:prop:ordinary-lattice-finiteness},
shows that $H^a(G,Q_D)$ is finite.  It is discrete by
\Cref{h3:co:strictness}.  Since $\mathcal P$ preserves $A_D$, it
acts on this fixed quotient.
Represent a quotient cocycle by its continuous holomorphic tail
section $f$.  Its iterates are represented by
$T_D\mathcal P^nf$, which converge to zero by
\Cref{h3:ord:tail-contraction}, uniformly on the compact cochain domain.
These are quotient cocycles.  Discreteness therefore implies that the
given class is killed by some power of $\mathcal P$.
Thus $\mathcal P$ is locally nilpotent on the cohomology of every
fixed $Q_D$.

Now form the quotient complex
\[
 Q^\bullet=
 C^\bullet_{\hthreects}(G,B_{\hthreehol})\big/
       \colim_d C^\bullet_{\hthreects}(G,A_d)
   =\colim_{D>C}C^\bullet_{\hthreects}(G,Q_D).
\]
The equality follows from the preceding degreewise exact sequences.
Filtered colimits of vector spaces are exact, so the fixed-quotient
argument makes $\mathcal P$ locally nilpotent on $H^a(Q^\bullet)$.
These groups are finite, using the long exact sequence for this
quotient, \Cref{h3:prop:holomorphic-comparison}, and the separately
proved \Cref{h3:prop:bounded-pole-finiteness}.
Both Frobenius and $\mathcal P$ preserve the bounded-pole cochain
union.  Their descended maps on $Q^\bullet$ satisfy
$\mathcal P\mathrm{Fr}=1$.  Thus $\mathcal P$ is surjective,
hence bijective, on its finite cohomology groups.  A bijective locally
nilpotent endomorphism can act only on the zero group.  This proves
that the first arrow of \eqref{h3:ord:natural-comparison} is a
quasi-isomorphism.  The second is
\Cref{h3:prop:holomorphic-comparison}.

Every $A_d$, $B_{\hthreehol}$, and $B_{\hthreepk}$ has finite point-value
cohomology and a Polish cochain complex.  The family-lifting result
\Cref{h3:co:parameters} identifies the cohomology with a profinite
parameter space $T$ at each such stage with continuous maps from $T$ to
its finite discrete cohomology.  In the filtered union, a continuous
map from compact $T$ to the discrete colimit has finite image, and those
finitely many classes occur at one common stage.  Thus this identification
commutes with the stagewise filtered colimit and gives the comparison
for every profinite $T$.  This argument does not assert that $B_0$, or
its quotient cochain complex, is Polish.
Finally the arrows being proved to be quasi-isomorphisms are the
coefficient inclusions, so they retain products and constants.
Write $\phi_p(f)=f^p$ for this untwisted coefficient Frobenius.
It is distinct from the fixed $q$-power operator $\mathrm{Fr}$ used
with the line retractions above.  The inclusions commute with $\phi_p$,
which is an automorphism on perfected cochains; hence $\phi_p$ is an
isomorphism on each of the compared cohomology groups.
\Cref{h3:prop:ordinary-constants} identifies their endpoint with the
displayed group cohomology.  All constructions are natural under the
semilinear coefficient-field action, so the comparison also descends
to the original constants.
\end{proof}

\begin{corollary}[The ordinary exterior classes]\label{h3:ord:exterior-classes}
For $p\geq5$ the last group in \Cref{h3:thm:ordinary-comparison} is
\[
 \Lambda_k(h,e,d),\qquad |h|=|e|=1,\quad |d|=3.
\]
Here $h$ is the logarithmic class of $P_1=\Z_p^\times$, $e$ is
the determinant logarithm of $\GL_2(\Z_p)$, and $d$ restricts to
the orientation class of $\hthreeSL_2(\Z_p)$.  All three classes have
integral lifts in continuous $\Z_p$-cohomology of these frame groups.
\end{corollary}

\begin{proof}
The group $\hthreeSL_2(\Z_p)$ has dimension three, no element of order
$p$, and trivial adjoint orientation.  Its first mod-$p$ cohomology
vanishes.  To see this directly, let $f:\hthreeSL_2(\Z_p)\to k$ be a
continuous homomorphism.  Choose $a\in\Z_p^\times$ with
$a^2-1$ a unit, possible for $p\geq5$.  Conjugation by
$\operatorname{diag}(a,a^{-1})$ gives
\[
 f\left(\begin{smallmatrix}1&a^2t\\0&1\end{smallmatrix}\right)
   =f\left(\begin{smallmatrix}1&t\\0&1\end{smallmatrix}\right).
\]
Additivity on this root subgroup and invertibility of $a^2-1$
force its restriction to vanish.  The lower root subgroup is treated
by the same displayed conjugation, with $a^{-2}$ in place of $a^2$.
These subgroups generate $\hthreeSL_2(\Z_p)$ by row reduction over the
local ring $\Z_p$, so $f=0$.
Compact group duality now gives $H^2=0$ and $H^3=k$, with the
orientation generator, while $H^0=k$.

The determinant quotient $\Z_p^\times$ has cohomology
$\Lambda_k(e)$: average its tame subgroup and use the two-term
resolution of its procyclic quotient.  Its action on $H^3$ of
$\hthreeSL_2$ is trivial, since conjugation on $\mathfrak{sl}_2$ has
determinant one.  The Hochschild--Serre sequence for
\[
 1\longrightarrow\hthreeSL_2(\Z_p)\longrightarrow\GL_2(\Z_p)
 \xrightarrow{\det}\Z_p^\times\longrightarrow1
\]
has only rows $0,3$ and columns $0,1$.  It has no possible
differential, and its edge map onto the orientation row is surjective.
Choose a lift $d$ of that row.  The product $ed$ is nonzero on
the associated graded, and odd squares vanish because $p$ is odd.
This proves the exterior algebra on $e,d$.  The finite-resolution
K\"unneth calculation with $P_1$ adds $h$.

Integrally the same finite resolutions have finitely generated
$\Z_p$-cohomology.  Integral duality gives the invariant
$\Z_p$-orientation line in degree three for $\hthreeSL_2$.
The preceding elementary-matrix argument also gives
$H^1(\hthreeSL_2,\Z_p)=0$.  Reduction embeds
$H^2(\hthreeSL_2,\Z_p)/p$ into $H^2(\hthreeSL_2,\F_p)=0$;
finite generation and Nakayama's lemma therefore give
$H^2(\hthreeSL_2,\Z_p)=0$.
For clarity, vanishing in the integral degree needed for this reduction
also follows from the finite coefficient models.  The length-three
reduced resolution and the finite filtration by powers of $p$ make
$H^a(\hthreeSL_2,\Z/p^\ell)$ finite and zero for $a>3$ at every
$\ell$.  Cochain reductions are surjective by continuous set-theoretic
sections, and the finite cohomology towers are Mittag--Leffler.  The
inverse-limit cochain sequence therefore gives
$H^4(\hthreeSL_2,\Z_p)=0$.  The coefficient reduction long exact
sequence now proves that the integral orientation surjects onto its
mod-$p$ class.  The integral determinant Hochschild--Serre sequence uses the same
two-term completed resolution for the procyclic quotient, after exact
tame averaging.  It again has no possible differential out of the
orientation row, so that orientation lifts to $\GL_2$.  Normalized continuous
logarithms give integral $e$ and $h$.  Reduction of these choices
gives the three specified mod-$p$ classes.
\end{proof}

\section{An integral basis at height one}\label{h3:sec:integral-basis}

The ordinary comparison now identifies the actual coefficient map and
its mod-\(p\) cup products.  We use it to construct four maps of
\(K(1)\)-local spectra.  The first two are the unit and the determinant
class already used at height two.  The other two will be an integral
degree-minus-three class \(\delta\) and its product with that same
determinant class.  The main point is to lift the ordinary orientation
through finite Witt coefficients while keeping the transition maps and
the maps from constants visible.

Throughout this section \(p\geq5\), \(T=L_{K(3)}S\), and
\(R_1=L_{K(1)}S\).  We use the local relative Morava descent result
\Cref{h3:prop:exact-descent}, including its finite coefficient-field
extension and its naturality for the ordinary tests.  The ordinary
cochain convention remains the colimit over stable compact pole
lattices from \Cref{h3:sec:framework}.

\subsection{The two logarithms and the orientation class}

Write
\[
 G_{\mathrm f}=P_1\times\GL_2(\Z_p).
\]
The ordering of these two factors is immaterial; they are the
connected and étale frame groups on the ordinary stratum.

\begin{lemma}\label{h3:des:frame-cohomology}
For every \(\ell\geq1\), there are compatible choices of generators
for which
\[
 H^*_{\hthreects}(G_{\mathrm f},\Z/p^\ell)
       =\Lambda_{\Z/p^\ell}(h,e,d),
       \qquad |h|=|e|=1,\quad |d|=3.
\]
The degree-one generators are the two normalized logarithms,
and the restriction of \(d\) to \(\hthreeSL_2(\Z_p)\) is its orientation
generator.  Coefficient extension to \(W_\ell(k)\) gives the
same exterior algebra over \(W_\ell(k)\).  These generators are
fixed by the changes of integral frames and by coefficient
Frobenius.
\end{lemma}

\begin{proof}
The group \(\hthreeSL_2(\Z_p)\) has no element of order \(p\):
a nontrivial \(p\)-th root of unity in a two-dimensional
\(\Q_p\)-representation has minimal polynomial of degree
\(p-1>2\).  It is an analytic duality group of dimension three.
Its orientation character is trivial, since conjugation on
its trace-zero Lie algebra has determinant one.
The continuous duality theorem, in the finite completed
resolution form used in \Cref{h3:sec:coheight-one}, gives perfect
pairings between degrees \(a\) and \(3-a\) with coefficients
\(\Z/p^\ell\).

Its first cohomology is zero by the elementary-matrix argument in
\Cref{h3:ord:exterior-classes}, now with coefficients \(\Z/p^\ell\).
That argument applies at every length: choose
\(b\in\Z_p^\times\) with \(b^2-1\in\Z_p^\times\);
conjugation invariance forces a homomorphism to vanish on the upper
root subgroup because multiplication by \(b^2-1\) is surjective
on its parameter group \(\Z_p\). The same holds for the lower
root subgroup using \(b^{-2}-1\), and these two subgroups generate
\(\hthreeSL_2(\Z_p)\).
Duality gives \(H^2=0\) and identifies \(H^3\) with
\(\Z/p^\ell\).  The orientation in the completed resolution
gives these generators compatibly in \(\ell\).

The determinant extension
\[
 1\longrightarrow\hthreeSL_2(\Z_p)\longrightarrow\GL_2(\Z_p)
   \longrightarrow\Z_p^\times\longrightarrow1
\]
has the section \(a\mapsto\operatorname{diag}(a,1)\).
Its action on the orientation generator is trivial, again
by determinant one on the adjoint Lie algebra.
The finite subgroup \(\mu_{p-1}\) has exact invariants,
and the remaining \(\Z_p\) has its two-term continuous
resolution.  The Hochschild--Serre sequence thus has only
quotient columns zero and one, and kernel rows zero and three.
It gives the logarithm \(e\), a lift \(d\) of the orientation,
and their nonzero product.  Their squares are zero:
\(e^2=0\) because \(2\) is invertible and \(e\) is odd,
and \(d^2=0\) also follows from the cohomological dimension.
This proves the exterior formula for \(\GL_2(\Z_p)\).
The factor \(P_1=\Z_p^\times\) contributes the second logarithm
\(h\), by its two-term resolution.  Tensoring these finite
resolution calculations gives the asserted exterior algebra.

The constructions use determinant and oriented Lie-algebra
classes, which are preserved under conjugation and transport
of a lattice.  They have coefficients in \(\Z/p^\ell\).
Consequently coefficient Frobenius fixes them.  Extension
to the finite free coefficient ring \(W_\ell(k)\) commutes
with the finite resolution and proves the last assertion.
\end{proof}

\begin{lemma}\label{h3:des:unit-directions}
Under the natural ordinary coefficient comparison, the map
from the height-one sphere sends its degree-minus-one
mod-\(p\) generator to the connected logarithm \(h\).
Under the same comparison, the mod-\(p\) reduction of the normalized
height-three logarithm detecting \(\zeta\) has image
\[
 a h+b e,\qquad a,b\in\F_p,\quad b\ne0.
\]
Thus, writing \(\zeta\) also for its image, the ordinary
mod-\(p\) coefficient algebra is
\[
 \F_p[v_1^{\pm1}]\otimes\Lambda_{\F_p}(h,\zeta,d).
\]
In this formula the degree-minus-one \(h\)-action is the
action of the actual unit \(R_1\to L_{K(1)}T\).
\end{lemma}

\begin{proof}
We first identify the unit direction, since an abstract
cohomology isomorphism would not specify it.  In the integral \(BP_*BP\) cobar construction use the
degree-zero differential \(d=\eta_R-\eta_L\), and write bars for
reduction modulo \(p\).  The invariance of \(\bar v_1\) makes
\[
 c_{v_1}=\frac{\eta_R(v_1)-\eta_L(v_1)}p
\]
integral; \(d^2=0\) and the \(p\)-torsion-free cobar terms make it a
cocycle.  For \(0\to BP_*\xrightarrow p BP_*\to BP_*/p\to0\),
the coefficient connecting map is
\[
 \partial_{\mathrm{coeff}}(\bar v_1)=[c_{v_1}]
 \in\Ext^{1,2(p-1)}_{BP_*BP}(BP_*,BP_*).
\]
Its mod-\(p\) reduction is represented by \(\bar c_{v_1}\).
On the mod-\(p\) height-one open, \(\bar v_1\) is an invariant
unit, so \(\bar c_{v_1}/\bar v_1\) is a cocycle of cohomological
degree one and internal degree zero.  Its naturality is
literal cobar naturality: for two composable formal-group
arrows, the difference of the two images of \(v_1\) is the
sum of the two corresponding differences after transport.
Changing a trivialization by one tautological arrow
therefore changes the pulled-back cocycle by a coboundary.

Pull the height-one specialized formal law to Honda form
on the connected-marking torsor of
\Cref{h3:prop:integral-frames}.  The preceding coboundary
identity identifies \(\bar c_{v_1}/\bar v_1\) with the Honda-height-one
class.  It is nonzero: on an automorphism \(a\in1+p\Z_p\)
the integral periodic frame changes \(v_1\) by
\(a^{p-1}\), and the cocycle is
\[
 \frac{a^{p-1}-1}{p}\pmod p.
\]
At \(a=1+p\) this is \(p-1\pmod p\), a unit.
It is a unit multiple of the normalized logarithm on \(P_1\).
We choose \(h\) with the corresponding normalization.
The leading coefficient of the formal-group arrow has
been retained here through the periodic frame \(U\);
discarding that frame would discard this calculation.

The height-one sphere calculation identifies its mod-\(p\)
groups with
\(\F_p[v_1^{\pm1}]\otimes\Lambda(h)\).
For the Adams self-map normalized in
\Cref{h3:prop:bounded-pole-finiteness}, the image of the Moore bottom
cell has \(BP\)-Hurewicz coefficient \(\bar v_1\).  Its Moore boundary
is a unit multiple of \(\alpha_1\).  The normalized cobar terms for
\(BP\) are even and \(p\)-torsion-free, so the comparison in
\Cref{h3:sec:descent} between the termwise cone and the cone of
totalization applies to this cofiber sequence.
For the lift \(v_1\), the coefficient connecting cocycle is
\(c_{v_1}\), while the actual desuspended boundary has leading
detector \(-[c_{v_1}]\).  The internal degree \(2(p-1)\) is
even, and strict desuspension contributes no further sign.
Choose the negative of this boundary class, map it to the Moore
quotient, and invert the same self-map.  The resulting source
generator in degree \(-1\) has detector
\(\bar c_{v_1}/\bar v_1\), the logarithm normalized above.
The odd-prime Adams-summand fiber is realized
with the unit as its first map in \cite[Section~4.2, pp.~30--32]{BB}.
Its logarithm, mod-\(p\) algebra, and unit-induced rational splitting
are given in \cite[Section~4.3, Remarks~4.16--4.17 and (4.18)]{BB}.
The maps of the preceding paragraph are the unit maps of
the formal-group cobar construction.  Hence their \(h\)
is precisely the \(R_1\)-action, rather than an independently
chosen degree-one direction.

For the second direction use the actual determinant
identity in \Cref{h3:prop:integral-frames}: the determinant
of the middle height-three frame is the product of the
norm of the connected height-one frame and the determinant
of the rank-two kernel frame.  It is an identity of
integral determinant lattices.  Applying the normalized
logarithm yields the sum of the two logarithms, with
the signs introduced by inverse-frame conventions.
In particular its coefficient on \(e\) is a unit modulo
\(p\).  This is the class \(\zeta\), by its determinant
construction.  The change of basis from \((h,e)\) to
\((h,\zeta)\) is invertible.

Finally, \Cref{h3:thm:ordinary-comparison,h3:des:frame-cohomology}
identify the actual coefficient algebra with this exterior
algebra.  The comparisons preserve products and the
two unit maps, proving all assertions.
\end{proof}

The coefficient calculation now has columns only from zero
to five.  Since the internal degrees are still multiples
of \(2(p-1)\), no mod-\(p\) differential is possible.
Cup products give the associated-graded multiplication.
For the actions used here one may use a Moore pairing
and the height-one Adams self-map; their cobar actions
are the same cup actions.  We next lift the degree-three
orientation before using these products to construct an
equivalence of spectra.

\subsection{Finite Witt length and ordinary integral coefficients}

Put
\[
 A_\ell=(W(k)/p^\ell)[[x,y]],\qquad
 D_\ell=A_\ell[x^{-1}].
\]
Write \(\phi_p(a)=a^p\) for coefficient \(p\)-Frobenius on the
characteristic-\(p\) rings \(B_0\) and \(B_{\hthreepk}\), and let
\(\sigma_k\) be the Witt automorphism of \(W(k)\) lifting
\(a\mapsto a^p\) on \(k\).  This notation is distinct from the
chosen \(q\)-power operator used in the ordinary root comparison.
The symbol \(D_\ell\) is separate from the distribution parameter
of \Cref{h3:sec:coheight-one}.
The action of \(P_3\) preserves \(A_\ell\) and satisfies
\(g(x)=xu_g+ph_g\), with \(u_g\) a unit and \(h_g\in A_\ell\).
It therefore extends to \(D_\ell\), since the second term
is nilpotent relative to the invertible first term.
An action on the uncompleted mixed-characteristic
localization \(W(k)[[x,y]][x^{-1}]\) is not needed.

To specify the cochain topologies at finite length, use
the following stable compact additive lattices, for \(d\geq0\):
\begin{equation}\label{h3:des:witt-pole-lattices}
 \mathcal L_{\ell,d}=[x]^{-d}W_\ell(A),\qquad
 \mathcal M_{\ell,d}
       =\sum_{j=0}^{\ell-1}p^j x^{-d-j}A_\ell.
\end{equation}
The first has its \(i\)-th Witt coordinate in
\(x^{-dp^i}A\).  These lattices exhaust \(W_\ell(B_0)\);
the bounds are cofinal with finite coordinate pole bounds.
Moreover,
\[
 x^{-d}A_\ell\subset\mathcal M_{\ell,d}
       \subset x^{-d-\ell+1}A_\ell,
\]
so the second family is cofinal with bounded poles in
\(D_\ell\).  Stability of the first follows from
\(g([x])=[x][\overline u_g]\).
For the second, expand \((xu_g+ph_g)^{-d-j}\)
modulo \(p^{\ell-j}\); a term with an extra \(p^a\)
has pole at most \(d+j+a\), as required by the
\((j+a)\)-th summand.  These actions are continuous
in the compact lattice topologies.
All cochains on either union mean the filtered colimit
of the continuous cochains on these stable lattices.
In particular, no coordinatewise pole-bound set that
fails to be additive is being used as a coefficient module.

\begin{lemma}\label{h3:lem:ghost-comparison}
For each \(\ell\geq1\), the last-ghost formula
\begin{equation}\label{h3:des:ghost-formula}
 \gamma_\ell(a_0,\ldots,a_{\ell-1})
   =\sum_{i=0}^{\ell-1}
       p^i\widetilde a_i^{\,p^{\ell-1-i}}\quad\hbox{in }D_\ell
\end{equation}
for arbitrary lifts \(\widetilde a_i\in D_\ell\) defines
a natural equivariant unital ring map
\[
 W_\ell(B_0)\xrightarrow{\gamma_\ell}D_\ell.
\]
Both this map and the natural coefficient map
\[
 \beta_\ell:W_\ell(B_0)\longrightarrow W_\ell(B_{\hthreepk})
\]
induce isomorphisms on \(P_3\)-cohomology with the prescribed
cochain conventions.  For \(\ell\geq2\), let
\(R_\ell:W_\ell(B_0)\to W_{\ell-1}(B_0)\) be Witt truncation and
let \(\operatorname{red}_{\ell,\ell-1}:D_\ell\to D_{\ell-1}\)
be coefficient reduction.  Then
\begin{equation}\label{h3:des:ghost-transition}
 \operatorname{red}_{\ell,\ell-1}\gamma_\ell
     =\gamma_{\ell-1}R_\ell W_\ell(\phi_p)
       \qquad(\ell\geq2).
\end{equation}
The map \(\gamma_\ell\) fixes the constant
\(\Z/p^\ell=W_\ell(\F_p)\).
\end{lemma}

\begin{proof}
The binomial theorem gives
\[
 a\equiv b\pmod p
 \quad\Longrightarrow\quad
 a^{p^j}\equiv b^{p^j}\pmod {p^{j+1}}
 \qquad(j\geq0).
\]
For the \(i\)-th summand of \eqref{h3:des:ghost-formula},
changing a lift consequently changes the result by a
multiple of \(p^\ell\).  This proves independence of all
lifts.
Lift two Witt vectors to \(W_\ell(D_\ell)\), and use their
Witt sum or product to lift the corresponding operation
over \(B_0\).  The last ghost component on
\(W_\ell(D_\ell)\) is a ring homomorphism.  Its reduction
and the lift independence prove the ring identities for
\(\gamma_\ell\).  The construction is natural for maps
of the reduction \(D_\ell\to B_0\), and hence equivariant
for the actual \(P_3\)-action on \(D_\ell\).
These are the elementary ghost identities of
\cite[Lemma~1 and Proposition~2]{Hesselholt}.

The map is continuous on each compact pole lattice.
Indeed, coefficientwise lifts and the finite polynomial
\eqref{h3:des:ghost-formula} prove continuity on \(W_\ell(A)\);
its ring property then gives
\[
 \gamma_\ell(\mathcal L_{\ell,d})
       \subset x^{-dp^{\ell-1}}A_\ell
       \subset\mathcal M_{\ell,dp^{\ell-1}}.
\]
It therefore acts on the specified cochain colimits.

Filter the source by \(V^iW_{\ell-i}(B_0)\) and the target
by \(p^iD_\ell\), for \(0\leq i\leq\ell\).
The source filtration iterates the one-coordinate
restriction--Verschiebung sequence in
\cite[proof of Lemma~9, p.~6]{Hesselholt}.
Their \(i\)-th graded pieces are additively \(B_0\);
the induced map is
\[
 a\longmapsto a^{p^{\ell-1-i}}.
\]
The quotient sequences remain exact on cochains.
For a Witt quotient the section is the coordinate
section \(a\mapsto V^i[a]\); for the corresponding
target quotient use \(a\mapsto p^i\widetilde a\) with
coefficientwise lifts.  These sections are continuous
and send a compact bounded-pole family into some
bounded-pole stage of \eqref{h3:des:witt-pole-lattices}.
They need not be additive or equivariant: their role
is to give degreewise surjectivity of continuous
cochains.  The kernels and differentials are the
equivariant ones.

By \Cref{h3:thm:ordinary-comparison}, the coefficient
\(p\)-Frobenius \(\phi_p\) is an isomorphism on
\(H^*(P_3,B_0)\).
The long exact cohomology sequences and induction on
the finite filtration now prove the assertion for
\(\gamma_\ell\).
For \(\beta_\ell\), use the two Verschiebung
filtrations.  Their graded maps are the natural
comparison \(B_0\to B_{\hthreepk}\), which is a cohomology
isomorphism by the same theorem.  Witt-coordinate
sections on the analytic side are continuous, so the
identical finite-filtration argument proves the
assertion for \(\beta_\ell\).

In characteristic \(p\), truncated Witt Frobenius is
\[
 R_\ell W_\ell(\phi_p)(a_0,\ldots,a_{\ell-1})
          =(a_0^p,\ldots,a_{\ell-2}^p);
\]
see \cite[Lemma~8]{Hesselholt}.
Substitution into \eqref{h3:des:ghost-formula} proves
\eqref{h3:des:ghost-transition}.  Since \(\gamma_\ell\)
is unital, it is the identity on \(\Z/p^\ell\).
Its restriction to \(W_\ell(k)\), under the identification
with \(W(k)/p^\ell\), is \(\sigma_k^{\ell-1}\);
we do not identify that restriction with the identity.
\end{proof}

\begin{lemma}\label{h3:des:witt-constants}
The actual ordinary constants map gives natural equivalences
\[
 \eta_\ell:
 R\Gamma_{\hthreects}(G_{\mathrm f},W_\ell(k))
    \xrightarrow{\ \simeq\ }
 R\Gamma_{\hthreects}(P_3,W_\ell(B_{\hthreepk})).
\]
They commute with truncation, coefficient Frobenius, and
the maps from \(\Z/p^\ell\) constants.  The maps
\(\beta_\ell,\gamma_\ell\) in \Cref{h3:lem:ghost-comparison}
preserve these latter constant maps as well.
\end{lemma}

\begin{proof}
On the ordinary quotient stack, the comparison of
\Cref{h3:prop:ordinary-constants} is induced by its map to
\(BG_{\mathrm f}\) and by the coefficient unit
\(k\to\mathcal O^\flat\).  Apply \(W_\ell\) to the
coefficient sheaves on this same diagram.  It gives
the displayed map before taking cohomology, and all
the asserted naturalities follow from the natural
Witt operations.  Its Verschiebung filtration has
the length-one constants comparison on each graded
piece.  The continuous coordinate sections make
the associated short exact sequences exact on the
computing cochains, including their profinite
parameters.  Finite-filtration induction from
\Cref{h3:prop:ordinary-constants} proves the equivalence.

The coefficient map \(\beta_\ell\) preserves constants
by definition.  For \(\gamma_\ell\), the composite
\[
 C^*_{\hthreects}(P_3,\Z/p^\ell)
   \longrightarrow C^*_{\hthreects}(P_3,W_\ell(B_0))
   \longrightarrow C^*_{\hthreects}(P_3,D_\ell)
\]
is exactly the ordinary coefficient inclusion, since
\(\gamma_\ell\) is the identity on \(\Z/p^\ell\).
Thus these are statements about the specified maps
from constants, not about an abstract isomorphism
of their images in cohomology.
\end{proof}

\begin{lemma}\label{h3:des:finite-length-model}
The ordinary spectrum \(L_1(T/p^\ell)\) is the totalization
of the termwise \(L_1(-/p^\ell)\) Morava descent object.
Its terms are even, and their internal degree-zero
cochain complex is the \(D_\ell\) complex above.
The corresponding spectral sequence is strongly
convergent, has columns \(0\leq s\leq5\), and has
no differentials.  All its abutment groups are finite
in each stem.  In particular there is a natural
identification
\begin{equation}\label{h3:des:minus-three-edge}
 \pi_{-3}L_1(T/p^\ell)
       \cong H^3_{\hthreects}(G_3(k),D_\ell),
\end{equation}
where the right side includes the coefficient
Galois descent just described.
\end{lemma}

\begin{proof}
The totalization assertion is
\Cref{h3:prop:exact-descent}.
Put \(M_\ell=\mathbb S_{(p)}/p^\ell\) and \(E=E_3(k)\).
The completion map identifies \(M_\ell\) with \(S/p^\ell\),
naturally in the ordinary Moore quotient transitions: its cofiber
has zero mod-\(p\) quotient, so multiplication by \(p\) is an
equivalence there and every mod-\(p^\ell\) quotient vanishes.
Each finite \(M_\ell\) has type one, since its rational homology
vanishes and its \(K(1)\)-homology is nonzero.
Let \(\psi:\Sigma^{2(p-1)}M_1\to M_1\) be the normalized Adams
\(v_1\)-map chosen in the proof of
\Cref{h3:prop:bounded-pole-finiteness}; its action on \(E/p\)
is multiplication by \(xU^{p-1}\).
The periodicity theorem
\cite[Theorem~1.5.4(i)--(ii), pp.~9--10]{RavenelNilpotence}
supplies a \(v_1\)-map \(f_\ell:\Sigma^{d_\ell}M_\ell\to M_\ell\)
and, for the ordinary reduction \(r_\ell:M_\ell\to M_1\),
positive integers \(i_\ell,j_\ell\) such that
\(d_\ell i_\ell=2j_\ell(p-1)\) and
\[
 r_\ell\circ f_\ell^{i_\ell}
   \simeq \psi^{j_\ell}\circ
          \Sigma^{2j_\ell(p-1)}r_\ell.
\]
Here the theorem is applied after a common suspension to finite
CW-complexes and then desuspended to spectra.  Multiplication of
the Adams map by the integer unit used in its normalization
preserves its \(v_1\)-map property.  Fix this iterate
\(\varphi_\ell=f_\ell^{i_\ell}\), of degree
\(e_\ell=2j_\ell(p-1)\).

The cofiber sequence \(E\xrightarrow{p^\ell}E\to E/p^\ell\)
and evenness of \(E/p^\ell\) give an isomorphism
\[
 [\Sigma^{e_\ell}E/p^\ell,E/p^\ell]_E
       \xrightarrow{\ \simeq\ }\pi_{e_\ell}(E/p^\ell)
\]
by evaluation on the unit: the possible kernel is a quotient of
\(\pi_{e_\ell+1}(E/p^\ell)=0\), and multiplication by
\(p^\ell\) vanishes on the group on the right.
Thus the \(E\)-module map induced by \(\varphi_\ell\) is determined
by a scalar \(\theta_\ell\) in this group.  Applying \(E\wedge-\)
to the chosen reduction square proves
\[
 \theta_\ell\bmod p=(xU^{p-1})^{j_\ell},\qquad
 \theta_\ell=U^{j_\ell(p-1)}c_\ell,\qquad
 c_\ell=x^{j_\ell}+p b_\ell\quad(b_\ell\in A_\ell).
\]
A lift of \(\theta_\ell\) to \(\pi_{e_\ell}E\) gives the same
\(E\)-module map by the evaluation isomorphism.  Tensoring with an
ordinary \(E\)-module \(N_Q\) therefore makes
\(\varphi_\ell\) act on \(\pi_*(N_Q/p^\ell)\) by multiplication
by this scalar.

The two localizations of \(A_\ell\) at \(c_\ell\) and at \(x\)
agree canonically.  Indeed, in \(A_\ell[x^{-1}]\) the factorization
\(c_\ell=x^{j_\ell}(1+p b_\ell x^{-j_\ell})\) makes \(c_\ell\)
invertible.  In \(A_\ell[c_\ell^{-1}]\), the factorization
\(x^{j_\ell}=c_\ell(1-p b_\ell c_\ell^{-1})\) makes \(x\)
invertible.  Both parenthesized factors have finite geometric-series
inverses because \(p^\ell=0\).
Miller's odd-prime result for \(M_1\)
\cite[Theorem~4.11 and Corollary~4.12]{Miller1981}, together with
the thick-subcategory argument in
\cite[Section~7.5, after Conjecture~7.5.5, p.~85]{RavenelNilpotence},
identifies the telescope of a \(v_1\)-map with \(L_1M_\ell\)
for every finite type-one \(M_\ell\).  Taking the chosen iterate
does not change the telescope, since it selects a cofinal subsequence.
As \(L_1\) is smashing, this is an equivalence under \(N_Q/p^\ell\)
between \(L_1(N_Q/p^\ell)\) and the telescope of its
\(\varphi_\ell\)-action.

For the actual term \(N_Q=\mathcal N_3^q\), with \(Q=P_3^q\),
\Cref{h3:des:flat-functions} gives even, \(p\)-torsion-free
coefficients.  Coefficientwise continuous lifts identify
\(\pi_0(N_Q/p^\ell)=C_{\hthreects}(Q,A_\ell)\).
Homotopy groups commute with the telescope, so the scalar calculation
and the two localizations above give
\[
 \pi_0L_1(N_Q/p^\ell)
   \cong C_{\hthreects}(Q,A_\ell)[x^{-1}]
   =\colim_d x^{-d}C_{\hthreects}(Q,A_\ell).
\]
The odd groups vanish, and the other even groups are its periodic
translates.  This is the algebraic bounded-pole localization;
the stable cofinal lattices \(\mathcal M_{\ell,d}\) specify it in
\eqref{h3:des:witt-pole-lattices}.

These coefficient identifications preserve the actual transition
maps.  Each is the unique extension of the map on coefficients
induced by the \(L_1\)-localization unit, by the universal property
of localization of a module at \(x\).  The ordinary Moore quotient
transition is induced by the cofiber square with vertical maps \(p\)
and \(\hthreeid\); it reduces coefficients and preserves the fixed
coordinate \(x\).  Naturality of the localization unit therefore
identifies its localized coefficient map with ordinary reduction
\(D_\ell\to D_{\ell-1}\).  The same argument applies to every
coface and codegeneracy of the actual descent object: coefficient
actions send \(x\) to \(xu_g+ph_g\), which is invertible in
\(D_\ell\), and precomposition on continuous functions preserves
the bounded-pole colimit.  Hence the degree-zero coefficient
complex is the stated \(D_\ell\) complex, with its actual cochain
maps, ordinary Moore transitions, and maps from constants.

The tame central units kill internal degrees
not divisible by \(2(p-1)\).
For all remaining degrees the finite \(p\)-filtration
has associated graded the mod-\(p\) coefficient
complex, trivialized by powers of \(v_1\).
By \Cref{h3:thm:ordinary-comparison,h3:des:frame-cohomology}
that complex has finite cohomology in columns
zero through five, and vanishes above five.
The long exact sequences of the finite
\(p\)-filtration give those same finiteness and
vanishing conclusions in every length.
There is no assertion that \(v_1\) itself is an
integral invariant in length \(\ell\); only the
associated-graded trivialization is used.

A differential has \(2(p-1)\mid r-1\), so the first
candidate has \(r=2p-1\geq9\).  It cannot fit the
column range \(0\leq s\leq5\).
Thus the spectral sequence collapses, and its
finite column range proves degreewise finiteness
of the abutment.  Strong convergence comes from
\Cref{h3:prop:exact-descent}.

Finally, in stem \(-3\) one must have \(t-s=-3\)
with \(0\leq s\leq5\) and \(t\in2(p-1)\Z\).
The unique solution is \((s,t)=(3,0)\).
There is consequently no additive extension or
choice of filtration quotient in
\eqref{h3:des:minus-three-edge}.  The construction is
natural in the Moore quotient maps and in
coefficient constants.  This proves its stated
naturality as well.
\end{proof}

The finite coefficient calculation has now isolated one degree-three
line at every length, with ordinary reduction as the spectral
transition.  To turn a compatible family on these lines into a map of
spectra, we spell out the completion step.

\begin{lemma}[Completion inside height one]\label{h3:lem:height-one-completion}
Let \(M\) be an \(E(1)\)-local \(S\)-module.  The natural map
\[
 M\longrightarrow \widehat M_p:=
       \operatorname*{holim}_{\ell\geq1} M/p^\ell
\]
is its \(K(1)\)-localization in \(S\)-modules.  Here
\(M/p^\ell\) is the cofiber of multiplication by \(p^\ell\),
with the ordinary Moore quotient transitions.
\end{lemma}

\begin{proof}
Write \(S/p^\infty=\operatorname*{colim}_\ell S/p^\ell\), with
the usual inclusions of Moore spectra.  Duality of the finite Moore
\(S\)-modules identifies
\(M/p^\ell\simeq F_S(\Sigma^{-1}S/p^\ell,M)\).  The inclusion
\(S/p^\ell\to S/p^{\ell+1}\) is induced by the square with vertical
maps \(\hthreeid\) and \(p\); dualizing gives the ordinary quotient
transition induced by \(p\) and \(\hthreeid\).  Thus these
identifications respect the tower.  Applying the internal function functor to
\[
 \Sigma^{-1}S/p^\infty\longrightarrow S\longrightarrow S[1/p]
\]
therefore gives a natural fiber sequence
\[
 F_S(S[1/p],M)\longrightarrow M\longrightarrow\widehat M_p.
\]
Multiplication by \(p\) is invertible on the first term.  Moreover
\[
 F_S(S[1/p],\widehat M_p)
 \simeq F_S(S[1/p]\otimes_S\Sigma^{-1}S/p^\infty,M)=0.
\]
Thus \(\widehat M_p\) is local with respect to the Moore quotient:
maps from any \(p\)-invertible \(S\)-module vanish, since such a
module is an \(S[1/p]\)-module and the displayed function object is
zero.  Each \(M/p^\ell\) is \(E(1)\)-local, and local objects are
closed under limits, so \(\widehat M_p\) is also \(E(1)\)-local.

In the \(E(1)\)-local category the Moore-acyclic objects are exactly
the rational objects: the Moore condition makes \(p\) invertible,
and conversely rationalization kills every Moore spectrum.  The
height-one fracture square identifies these with the \(K(1)\)-acyclic
\(E(1)\)-local objects.  Hence Moore localization and
\(K(1)\)-localization have the same acyclic objects and the same
local objects in this category.  The \(E(1)\)-localization unit is
itself a \(K(1)\)-equivalence, so passing first to \(E(1)\)-local
objects does not change \(K(1)\)-localization.  The displayed natural
map is therefore the asserted localization in \(S\)-modules.
\end{proof}

\begin{proposition}\label{h3:prop:integral-delta}
There is a class
\[
 \delta\in\pi_{-3}L_{K(1)}T
\]
whose reduction is the orientation generator \(d\)
in \Cref{h3:des:unit-directions}.
Moreover,
\[
 \pi_{-3}L_{K(1)}T\cong\Z_p,
\]
with \(\delta\) a generator after the chosen orientation
normalization.  The finite-length identifications used
to construct it preserve the given maps from
\(P_3\)-cochains with integral constant coefficients.
\end{proposition}

\begin{proof}
Choose the compatible classes
\(o_\ell\in H^3(G_{\mathrm f},\Z/p^\ell)\)
of \Cref{h3:des:frame-cohomology}.
Write \(\rho_\ell\) for the cohomology isomorphism
induced by \(\beta_\ell\), and set
\[
 \nu_\ell=\rho_\ell^{-1}\eta_\ell(o_\ell),\qquad
 z_\ell=\gamma_\ell(\nu_\ell)
       \in H^3(P_3,D_\ell).
\]
Here the constant coefficient extension to \(W_\ell(k)\)
is understood.  The naturalities in
\Cref{h3:des:witt-constants} give
\(R_\ell\nu_\ell=\nu_{\ell-1}\).
Coefficient Frobenius fixes \(o_\ell\), because it
fixes \(\Z/p^\ell\); injectivity of \(\rho_\ell\)
then gives
\[
 W_\ell(\phi_p)(\nu_\ell)=\nu_\ell.
\]
Using \eqref{h3:des:ghost-transition} we obtain
\[
 \operatorname{red}_{\ell,\ell-1}(z_\ell)
   =\gamma_{\ell-1}R_\ell
          W_\ell(\phi_p)(\nu_\ell)
   =z_{\ell-1}.
\]
Thus the Frobenius twist in the ghost transition
does not obstruct compatibility of these actual
integral orientation classes.

By \Cref{h3:des:frame-cohomology,h3:des:witt-constants,h3:lem:ghost-comparison}, the degree-three cohomology
over \(k\) is one free copy of \(W_\ell(k)\).
The determinant and orientation constructions are
preserved under coefficient Galois action and
transport of the integral frames.  Normal-basis
descent therefore identifies its descended group
with \(\Z/p^\ell\).  The elements \(z_\ell\) are
generators and their transitions are ordinary
reduction.  This proves the corresponding assertions
for \(\pi_{-3}L_1(T/p^\ell)\) by
\eqref{h3:des:minus-three-edge}.

Apply \Cref{h3:lem:height-one-completion} to \(L_1T\).  Exactness of
\(L_1\) identifies its Moore quotients with \(L_1(T/p^\ell)\).
Consequently
\[
 L_{K(1)}T\simeq
       \operatorname*{holim}_\ell L_1(T/p^\ell).
\]
All finite-length homotopy groups are finite by
\Cref{h3:des:finite-length-model}; their inverse systems
are Mittag--Leffler, so the Milnor
\(\lim^1\)-term vanishes in every stem.
The compatible generators \(z_\ell\) therefore
give an actual class \(\delta\), and their inverse
limit is \(\Z_p\).
The last assertion follows from the equality of
the constant cochain maps in
\Cref{h3:des:witt-constants} and from the natural
edge identification in
\Cref{h3:des:finite-length-model}.
\end{proof}

\begin{remark}\label{h3:des:constants-order-bridge}
The preceding argument does not identify the usual
full Witt tower on the uncompleted \(B_0\) with the
ordinary \(D_\ell\) tower.  For later use, the exact
statement needed for any class from integral
\(P_3\)-constants is simpler.
Its image in \(W_\ell(B_0)\) is fixed by coefficient
Frobenius, so \eqref{h3:des:ghost-transition} gives an
ordinary reduction-compatible family in \(D_\ell\).
The two cohomology isomorphisms in
\Cref{h3:lem:ghost-comparison} preserve the order of
each finite-length element.
Hence unbounded \(p\)-power order of the corresponding
analytic Witt images implies unbounded order of
these specified ordinary images.  Combined with
\Cref{h3:prop:integral-delta}, this is the bridge used
by the rational detector in \Cref{h3:sec:rational}.
\end{remark}

\subsection{The actual height-one equivalence}

\begin{theorem}\label{h3:thm:height-one-maps}
The integral classes \(\zeta\) and \(\delta\) above
and the canonical unit define an equivalence
\[
 (1,\zeta,\delta,\zeta\delta):
 R_1\vee\Sigma^{-1}R_1\vee
       \Sigma^{-3}R_1\vee\Sigma^{-4}R_1
      \xrightarrow{\ \simeq\ }L_{K(1)}T.
\]
The maps use the \(R_1\)-action and the multiplication
induced from \(T\).  In particular the first two
components are the \(K(1)\)-localizations of the
same maps from \(S\) used in
\Cref{h3:thm:height-two-test}.
\end{theorem}

\begin{proof}
The target is a \(K(1)\)-local algebra and hence a
module over the local unit \(R_1\).  Each homotopy
class displayed defines an actual map from the
indicated suspension of \(R_1\).  The product
\(\zeta\delta\) is formed using this algebra
multiplication, so no independent choice of a
degree-minus-four line is being made.

Modulo \(p\), the source groups are the direct sum
of the four shifts of
\[
 \pi_*(R_1/p)
     =\F_p[v_1^{\pm1}]\otimes\Lambda(h).
\]
By \Cref{h3:des:unit-directions}, the images of
\(1,\zeta,d,\zeta d\), together with their
\(h\)-multiples, are exactly the eight exterior
basis elements in the target coefficient
calculation.  The class \(\delta\) reduces to \(d\)
by \Cref{h3:prop:integral-delta}.
Products are detected by the cup action, so these
statements specify the leading filtered class of
each of the eight actual homotopy images.
The spectral sequence has collapsed with finite filtration.  In each
stem, filter the source homotopy group by assigning these eight
periodic basis elements the cohomological degrees of their displayed
leading classes.  The map is filtered, and its associated-graded map
is an isomorphism in every degree.  Induction on the finite filtration,
using the short exact sequences for successive filtration quotients,
shows that the original map of homotopy groups is an isomorphism.
This argument does not require an a priori splitting of the target
filtration.  The displayed map is therefore an isomorphism on
mod-\(p\) homotopy in every stem.

Its cofiber is \(K(1)\)-local, hence derived
\(p\)-complete by \Cref{h3:lem:height-one-completion}.  The vanishing of its mod-\(p\)
quotient makes multiplication by \(p\) an
equivalence, so all its quotients by \(p^\ell\)
vanish.  Derived completeness then makes the
cofiber zero.  This proves the equivalence,
including the asserted compatibility of its
first two components.
\end{proof}

\section{The rational primitive and its localization}\label{h3:sec:rational}

The integral basis gives a distinguished line
$\Z_p\delta\subset\pi_{-3}L_{K(1)}T$.  We now construct a rational
class in $\pi_{-3}L_0T$ whose image is nonzero in its rational span
$\Q_p\delta$.

Fix $C=\widehat{\overline{\Q}_p}$, with tilted base $C^\flat$ and
untilt divisor $\infty$.  Let $\mathcal M$ classify $P_3$-structured
forms $E$ of $V_3$ equipped with a quotient line at $\infty$.
The lower modification $E^-$ is a slope-zero rank-three bundle,
hence a rational local system; its determinant inherits an integral
lattice from the $P_3$-structure.  The two frame descriptions give
\[
 \begin{aligned}
 \mathcal M&\simeq[\mathbf P^2_C/P_3]\simeq[\Omega_C^2/J_3],\\
 J_3&=\{g\in\GL_3(\Q_p):v_p(\det g)=0\},
 \end{aligned}
\]
where $\Omega_C^2$ is the projective plane minus its
$\Q_p$-rational hyperplanes.  The first presentation records quotient
lines of a framed $E$; the second records upper modifications of a
rationally framed $E^-$ with the prescribed determinant lattice.
These identifications, including their common arithmetic action, are
proved in \Cref{h3:rat:modification-presentations}.

The projective presentation singles out the reduced-trace line from
$P_3$ as the arithmetic-invariant part of rational degree-three cohomology.
The Drinfeld presentation places a nonzero matrix-trace class in that
same line.  Pulling their proportionality back along the ordinary
sequence then restricts the rank-three matrix trace to the trace of
its original integral rank-two kernel.  We retain these actual
classifying maps through finite Witt coefficients and exact Morava
descent to identify the nonzero image in $\Q_p\delta$ under localization.

\subsection{Integral coefficients and continuous geometric descent}

Put $\Lambda_\ell=\Z/p^\ell$.  For a small v-stack $Z$ in this section, set
\[
 R\Gamma_{\mathrm{int}}(Z)
   =R\varprojlim_{\ell}R\Gamma_v(Z,\Lambda_\ell),\qquad
 H^i_{\mathrm{int}}(Z)=H^i(R\Gamma_{\mathrm{int}}(Z)),
 \qquad H^i_{\mathrm b}(Z)=H^i_{\mathrm{int}}(Z)[1/p].
\]
For a locally profinite group $G$ we use a separate complex in $D(\Z)$,
defined independently of a geometric base:
\[
 R\Gamma_{\mathrm{int}}(G)
    =R\varprojlim_\ell C^\bullet_{\hthreects}(G,\Lambda_\ell),
 \qquad H^i_{\mathrm{int}}(G)=H^i(R\Gamma_{\mathrm{int}}(G)),
 \qquad H^i_{\mathrm b}(G)=H^i_{\mathrm{int}}(G)[1/p].
\]
The coefficients here are finite discrete modules with trivial action,
and the cochains are the inhomogeneous cochains of
\Cref{h3:sec:framework}.  At each finite level the cochain complex is the
external derived cohomology of the condensed classifying topos
$\mathrm B_{\mathrm{cond}}G$, the topos of condensed sets with $G$-action.
Indeed, the standard resolution there has terms indexed by $G^a$.
Finite discrete coefficients are solid, and their higher condensed
cohomology on a locally profinite set vanishes.  Its degree-zero
sections are $C_{\hthreects}(G^a,\Lambda_\ell)$, so the standard
resolution gives the natural identification in the derived category
\[
 R\Gamma_{\mathrm{Cond}}
   (\mathrm B_{\mathrm{cond}}G,\underline{\Lambda_\ell})
       \simeq C^\bullet_{\hthreects}(G,\Lambda_\ell)
\]
by \cite[Section~2, standard resolution and Lemmas~2.1--2.2,
pp.~5--7]{AnschuetzSolidGroupCohomology}.  This is an external complex
of ordinary modules.  We reserve $R\Gamma(G,K)$ for internal condensed
derived invariants of a condensed coefficient object $K$ with its
actual $G$-action; evaluation at the point, denoted $(*)$, gives its
external complex.

Twists are made at finite level.  For compact $G$, reduction of finite
cochains is surjective: lift the finitely many values on their clopen
fibers.  Thus the derived coefficient limit is represented by
$C^\bullet_{\hthreects}(G,\Z_p)$.  Every continuous function from the
compact space $G^a$ to $\Q_p$ is bounded, so exact inversion of $p$
identifies this complex with $C^\bullet_{\hthreects}(G,\Q_p)$.
For the noncompact groups below we retain the displayed definition
using finite coefficients before the derived limit and inversion of $p$.

We now specify how a group class enters relative geometric cohomology.
For a small v-base $B$, let $\underline G_B=B\times\underline G$ be
the constant locally profinite v-group and put
$\mathrm B_B G=[B/\underline G_B]$, with trivial action on $B$.
Products with group parameters below are over this base.

\begin{lemma}[Finite classifying maps]\label{h3:rat:finite-classifying}
For a $G$-space $X$ over $B$ there is a natural finite coefficient map
\begin{equation}\label{h3:eq:rat-finite-classifying}
 \mathsf u_{X,G,\ell}:
 C^\bullet_{\hthreects}(G,\Lambda_\ell)
 \longrightarrow
 \operatorname{Tot}_{[a]\in\Delta}
 R\Gamma_v(X\times\underline{G^a},\Lambda_\ell)
 \simeq R\Gamma_v([X/G],\Lambda_\ell).
\end{equation}
For a $G$-torsor $T\to Z$ over $B$ with relative classifying map
$c_T:Z\to\mathrm B_B G$, this gives
\[
 \operatorname{cl}_{c_T,\ell}
   :=c_T^*\mathsf u_{B,G,\ell}:
 C^\bullet_{\hthreects}(G,\Lambda_\ell)
       \longrightarrow R\Gamma_v(Z,\Lambda_\ell).
\]
It equals \eqref{h3:eq:rat-finite-classifying} for $X=T$ under
$[T/G]\simeq Z$.  These maps commute with continuous group
homomorphisms and extension of torsor structure, equivariant maps,
base change, base automorphisms with compatible descent, coefficient
reductions, and cup products.  They also commute with the finite
coefficient units $\Lambda_\ell=W_\ell(\F_p)\to W_\ell\mathcal O^\flat$.
Their derived coefficient limits and their rationalizations define
maps denoted $\operatorname{cl}_{c_T,\mathrm{int}}$ and
$\operatorname{cl}_{c_T,\mathrm b}$ on the corresponding cohomology.

For an algebraically closed perfectoid field $C$ of characteristic zero, put
$B_C=\operatorname{Spa}(C,\mathcal O_C)^\diamond$.  For every
profinite $Q$ and every finite constant coefficient module $A$ on
$B_C$, the actual constants map is an equivalence
\begin{equation}\label{h3:eq:rat-geometric-point-constants}
 C_{\hthreects}(Q,A)[0]
       \xrightarrow{\ \simeq\ }
 R\Gamma_v(B_C\times\underline Q,A).
\end{equation}
This is natural in $Q$, in $A$, and in base automorphisms with their
actions on $A$.  In particular it includes the finite Tate fibers
$\Lambda_\ell(-j)$.  For compact $G$, it makes
$\mathsf u_{B_C,G,\ell}$ an equivalence.
\end{lemma}

\begin{proof}
For a locally profinite parameter $Q$, regard
$A_Q=C_{\hthreects}(Q,\Lambda_\ell)$ as a discrete ring.  Evaluation
gives a map of ring sheaves on $X\times\underline Q$,
\[
 \underline{A_Q}\longrightarrow\underline{\Lambda_\ell},
 \qquad (f,q)\longmapsto f(q).
\]
It is defined on the finitely many clopen fibers of each $f$; this
uses no compactness of $Q$.  The unit of the derived
constant-sheaf/global-sections adjunction, followed by derived
sections of evaluation, gives
\[
 C_{\hthreects}(Q,\Lambda_\ell)[0]
       \longrightarrow R\Gamma_v(X\times\underline Q,\Lambda_\ell).
\]
Equivalently, this sends a continuous $\Lambda_\ell$-valued function
on $Q$ to its ordinary section and then to derived sections.  With $Q=G^a$, the
projections of the action-groupoid nerve commute with all bar faces
and codegeneracies, including the action face.  These maps therefore
form a cosimplicial map.  Its totalization is
\eqref{h3:eq:rat-finite-classifying} by quotient \v{C}ech descent.

The map of nerves that forgets $T$ over $B$ proves the stated equality
with $c_T^*\mathsf u_{B,G,\ell}$.  For a continuous homomorphism
$h:H\to G$, let $S\to Z$ be an $H$-torsor over $B$ with a specified
identification $S\times^H G\simeq T$ over $Z$.  The resulting map of
torsor nerves acts in degree $a$ by $h^a$ on the transition elements.  Evaluation commutes with these maps, with equivariant
maps of $X$, and with every base change.  A base automorphism whose
descent commutes with the constant group acts on a torsor nerve by
$(t,g_1,\ldots,g_a)\mapsto(\gamma(t),g_1,\ldots,g_a)$, which proves
the asserted arithmetic naturality.  The construction uses maps of
coefficient ring sheaves and the lax monoidal derived sections functor,
so it preserves cup products.  Composing the same ring-sheaf maps
with $\Lambda_\ell\to\Lambda_{\ell-1}$ or with
$\Lambda_\ell=W_\ell(\F_p)\to W_\ell\mathcal O^\flat$ proves the
coefficient naturalities term by term.  In particular these are maps
of the derived coefficient towers before cohomology is taken.

For \eqref{h3:eq:rat-geometric-point-constants}, use all finite clopen
quotients $Q_i$ of $Q$.  The example after Definition~7.8 and
Proposition~7.16 of \cite{ScholzeDiamonds} identify
$B_C\times\underline Q$ with the inverse limit of $B_C\times Q_i$
and make $B_C$ a strictly totally disconnected geometric point.
Its \'etale covers split, so its finite constant cohomology is $A[0]$.
Continuity and comparison with v-cohomology in
\cite[Propositions~14.9--14.10]{ScholzeDiamonds} now give
\[
 R\Gamma_v(B_C\times\underline Q,A)
   \simeq\colim_i\hthreeMap(Q_i,A)[0]
   =C_{\hthreects}(Q,A)[0].
\]
Every map here is induced by the coefficient unit and is natural in
the finite partitions, coefficient fiber, and base.  This proves the
stated naturality, including the arithmetic action on a Tate fiber.
Apply the equivalence to $Q=G^a$ for compact $G$ and totalize to obtain
the last assertion.  Over the finite-field base used below we use the
finite classifying map itself; this point comparison is only asserted
over $B_C$.
\end{proof}

We retain profinite parameters throughout geometric descent.  Write
\[
 \mathscr C_{Z,\ell}(Q)
   =R\Gamma_v(Z\times\underline Q,\Lambda_\ell),\qquad
 \mathscr C_Z=R\varprojlim_\ell\mathscr C_{Z,\ell},
\]
where $Q$ is profinite; these are derived sections on $Q$.

\subsection{Trace classes and a parabolic subgroup}

The invariant form
\[
 (X,Y,Z)\longmapsto\hthreeTr(X[Y,Z])
\]
uses reduced trace for the division algebra and matrix trace for the
split groups.  This is the classical degree-three cocycle attached to
an invariant symmetric form, as in Chevalley--Eilenberg
\cite[Section~21]{ChevalleyEilenberg1948} and Koszul
\cite[Section~11]{Koszul1950}.
The next lemma constructs the corresponding group-cohomology classes
and identifies the restriction needed for the ordinary kernel.

\begin{lemma}\label{h3:rat:trace-groups}
The groups $P_3$, $J_3$, and $\GL_2(\Z_p)$ have rational cohomology
$H^1_{\mathrm b}=\Q_p$, $H^2_{\mathrm b}=0$, and
$H^3_{\mathrm b}=\Q_p$.  They admit nonzero classes
$e_3$, $e_J$, and $e_2$, respectively, whose rational Lie-algebra
images are nonzero multiples of the displayed trace form.
For every compact open $K\subset J_3$, restriction followed by
the compact Lie-algebra comparison identifies $H^*_{\mathrm b}(J_3)$
with $H^*(\mathfrak{gl}_3,\Q_p)$; the identifications for different
$K$ agree under further restriction.  For
\[
 D=\left\{\begin{pmatrix}A&v\\0&a\end{pmatrix}:
 A\in\GL_2(\Z_p),\ a\in\Z_p^\times,\ v\in\Q_p^2\right\},
\]
write $\iota:D\to J_3$ and $\rho:D\to\GL_2(\Z_p)$ for
inclusion and projection.  Then
\begin{equation}\label{h3:eq:rat-parabolic-trace}
 \iota^*e_J=b\rho^*e_2\quad\text{for some }b\in\Q_p^\times.
\end{equation}
\end{lemma}

\begin{proof}
We first make the continuous comparison on compact groups precise.
Each compact group $K$ under consideration has an open normal uniform
pro-$p$ subgroup $U$.  For $P_3$ and $\GL_2(\Z_p)$ one may use a
sufficiently deep principal congruence subgroup.  A compact open
$K\subset J_3$ preserves a lattice in $\Q_p^3$: the orbit of any
lattice is finite, and its sum is stable.  A sufficiently deep
principal congruence subgroup for this lattice lies in $K$ and is
normal there.  If the depth is at least one, the corresponding Lie
lattice $\mathfrak u=p^m\mathcal O$, for the relevant endomorphism
order $\mathcal O$, satisfies
$[\mathfrak u,\mathfrak u]\subset p\mathfrak u$.  At odd $p$,
logarithm and exponential converge on these lattices, identify the
$p$th-power map with multiplication by $p$, and show that the group
is finitely generated, powerful, and torsion free.  It is therefore
uniform.

Give $U$ its lower-$p$-series valuation.  For $p\geq5$ it is saturated
and equi-$p$-valued, with value denominator $e=1$ and graded generators
in degree one.  The trivial finite free coefficient module $\Z_p$
has action in $1+p\operatorname{End}_{\Z_p}(\Z_p)$.  The continuous
comparison of \cite[Theorem~3.3.3]{HuberKingsNaumann} therefore gives
a cup-compatible isomorphism
\[
 H^*_{\hthreects}(U,\Z_p)\simeq H^*(\mathfrak u,\Z_p),
\]
where $\mathfrak u$ is the integral Lie algebra of $U$.
Coefficient naturality and rational agreement in
\cite[Theorem~3.1.1(1)--(2)]{HuberKingsNaumann}, together with the
boundedness of compact rational cochains, identify its rationalization
with Lazard's comparison for $\operatorname{Lie}(K)=\mathfrak u[1/p]$.
Continuous group homomorphisms preserve lower $p$-series, so
Theorem~3.1.1(3) of the same source gives naturality for restrictions
and for the block maps after passing to such small opens.

The rational Hochschild--Serre sequence for $U\triangleleft K$
collapses to finite quotient invariants: positive cohomology of the
finite group $K/U$ vanishes by averaging in $\Q_p$.  Here
$\operatorname{Lie}(K)$ is the Lie algebra of a $\GL_n$-form, with $n=2$ or $3$.
The exterior calculation and triviality of its adjoint action are
given by \cite[Lemma~3.8.2]{BSSW}.  They give generators in degrees
$1,3,\ldots,2n-1$, without choosing explicit cocycles for all of them.
Thus the finite quotient fixes the cohomology, and in particular the
low-degree dimensions are $1,0,1$ in degrees $1,2,3$.

We identify the particular degree-three generator needed here.
For the split Lie algebra put $q(X,Y)=\operatorname{Tr}(XY)$;
for the inner form use reduced trace.  Cyclicity gives
$q([X,Y],Z)=\operatorname{Tr}(X[Y,Z])$.  After an algebraically
closed splitting-field extension the restriction of $q$ to
$\mathfrak{sl}_n$, for $n=2,3$, is nonzero and nondegenerate.
The degree-three invariant-form calculation of
\cite[Section~11, Theorems~11.1--11.2]{Koszul1950}
therefore makes its class nonzero.  The finite Chevalley--Eilenberg
cochain complex commutes with field extension, so the class already
defined on the form over $\Q_p$ is nonzero.  Comparing restrictions through a
common small uniform subgroup shows that the group identifications
are independent of the chosen open subgroup.  The upper-left block
map restricts the displayed degree-three matrix trace form to the
identical form on the $2\times2$ block, and hence restricts it
nontrivially.

For $P_3$, choose $e_3$ to have exactly this reduced-trace Lie-algebra
class under the natural compact comparison.  The coefficient-field
group $\Gamma=\Gal(k/\F_p)$ acts by automorphisms of the Honda division
algebra and preserves reduced trace.  Naturality of that comparison
therefore makes this normalized rational class $\Gamma$-invariant.

For the noncompact group $J_3$, use the reduced affine building in
\Cref{h3:lem:rank-three-building} with $F=\Q_p$.  That lemma supplies
a contractible two-dimensional complex on which $J_3$ preserves types,
has compact open face stabilizers fixing their faces pointwise, and has
a closed chamber as a fundamental domain.  The cellular resolution and continuous
Shapiro give, first with finite coefficients and then with the
derived integral limit,
\[
 E_1^{a,s}=\bigoplus_{\substack{\sigma\subset\Delta\cr
                         \dim\sigma=a}}
 H^s_{\mathrm{int}}(K_\sigma)
       \ \Longrightarrow\ H^{a+s}_{\mathrm{int}}(J_3).
\]
There are finitely many summands and cellular degrees.  After
inversion of $p$, every restriction in a fixed row identifies with
the identity on $H^s(\mathfrak{gl}_3,\Q_p)$; changing a face
representative only introduces an inner conjugation.  That row
is therefore the cellular cochain complex of the simplex $\Delta$
with these constant coefficients.  It has cohomology only in
cellular degree zero.  The edge map, and then further restriction
to an arbitrary compact open subgroup using intersections, gives
the required assertion for $J_3$.

It remains to account for the unipotent radical of $D$.
Fix $\ell$ and write $\Q_p^2=\bigcup_{r\geq0}p^{-r}\Z_p^2$.
The two-variable Koszul resolution computes
\[
 H^q(p^{-r}\Z_p^2,\Z/p^\ell)
       =\bigwedge^q(\Z/p^\ell)^2.
\]
Restriction from stage $r+1$ to $r$ multiplies degree $q$ by
$p^q$.  These systems are pro-zero for $q>0$ and constant for
$q=0$.  The restrictions of cochains themselves are surjective:
the smaller subgroup is clopen, so a cochain extends by zero.
Thus their inverse limit computes the derived limit, and the
Milnor sequence gives
$R\Gamma(\Q_p^2,\Z/p^\ell)=\Z/p^\ell[0]$.
The calculation remains valid with profinite parameters, by the
same finite Koszul complexes and extension of cochains.
Thus this equality holds for internal derived condensed invariants.
The constants map is equivariant for the Levi action, and its
underlying equivalence is therefore an equivariant equivalence.
Hochschild--Serre consequently makes inflation from the Levi subgroup
an equivalence at each finite level and then integrally.
The rational K\"unneth calculation for
$\GL_2(\Z_p)\times\Z_p^\times$ has degree three equal to the
degree-three group of its first factor, since the second factor
has cohomological dimension one and the first has $H^2=0$.
Restriction along $A\mapsto\operatorname{diag}(A,1)$ detects the
nonzero trace form.  This proves \eqref{h3:eq:rat-parabolic-trace}.
\end{proof}

\subsection{Two presentations of a modification}

Write
\[
 B_k=\Spd k,\qquad
 B_C=\Spd C^\flat\simeq\operatorname{Spa}(C,\mathcal O_C)^\diamond.
\]
The chosen embedding $F=W(k)[1/p]\subset C$ gives the embedding
$k\subset C^\flat$ and hence the base map $B_C\to B_k$.
All quotient presentations in this subsection are over $B_C$.
A determinant lattice on a slope-zero rank-three bundle is a
$\Z_p$-lattice in its determinant rational local system.  Its frame
group is $J_3$.  Fix a lattice in the one-dimensional rational
local system associated to
\[
 \Delta_3=\det(V_3)(-\infty).
\]
The action of $D_3^\times$ on it is reduced norm.  Its stabilizer
is therefore exactly $P_3$.

\begin{lemma}\label{h3:rat:modification-presentations}
Let $\mathcal M$ classify a form of $V_3$ with $P_3$-structure
and a quotient line at $\infty$.  There are equivalences
\begin{equation}\label{h3:eq:rat-two-presentations}
 \mathcal M\simeq[\mathbf P^2_C/P_3]
                 \simeq[\Omega_C^2/J_3].
\end{equation}
The relative classifying maps
\[
 c_{3,\mathcal M}:\mathcal M\longrightarrow\mathrm B_{B_C}P_3,
 \qquad
 c_{J,\mathcal M}:\mathcal M\longrightarrow\mathrm B_{B_C}J_3
\]
record the positive bundle and its lower modification, respectively.
For the prescribed finite unramified field $F=W(k)[1/p]$, these
presentations carry a common semilinear arithmetic action of $G_F$
over its action on $B_C$, commuting with both constant groups and
inducing the usual action on $\Omega_C^2$.
\end{lemma}

\begin{proof}
The kernel of a length-one quotient of a form $E$ of $V_3$ is
a rank-three bundle $E^-$ of degree zero.  If it had a positive
Harder--Narasimhan subbundle of rank $b<3$ and integral degree
$d\geq1$, its inclusion in $E$ would contradict semistability:
$d/b\geq1/b>1/3$.  Thus $E^-$ has slope zero.  By the relative
slope-zero correspondence it is a rational local system
\cite[Corollary~8.7.10]{KedlayaLiu}; see \Cref{h3:geom:curve-inputs}.

Conversely an upper length-one modification of $\mathcal O^3$
has, at a geometric point with local parameter $t$ at $\infty$,
lattice
\[
 (B_{\mathrm{dR}}^+)^3+
 B_{\mathrm{dR}}^+t^{-1}\widetilde\ell
\]
for a direction line $\ell$.  If $\ell$ lies in a proper
rational subspace $W$, modifying $W\otimes\mathcal O$ gives a
degree-one subbundle of rank less than three, which destabilizes.
In the other direction, let $F'$ be a saturated destabilizing
subbundle of the upper modification, of rank $b<3$ and degree
$d\geq1$.  Its intersection $F''$ with $\mathcal O^3$ has
degree $d-\epsilon$, where $0\leq\epsilon\leq1$.  Semistability
of $\mathcal O^3$ forces $d=\epsilon=1$ and $\deg F''=0$.
The saturation of $F''$ in $\mathcal O^3$ has degree at least zero,
while semistability forces its degree to be at most zero.  Its torsion
quotient therefore has degree zero, so $F''$ is already saturated.
Thus $F''$ is a
slope-zero subbundle of $\mathcal O^3$, hence equals
$W\otimes\mathcal O$ for a rational subspace $W$; maps between
trivial bundles are matrices over $H^0(\mathcal O)=\Q_p$.
The equality $\epsilon=1$ says that $\ell\subset W_C$.
This proves precisely the Drinfeld condition.

Let $\mathcal Z$ be the sheaf of injections
$\mathcal O^3\hookrightarrow V_3$ with length-one cokernel at
$\infty$ preserving the chosen determinant lattices.  The
commuting actions are postcomposition by $P_3$ and inverse
precomposition by $J_3$.  The preceding slope calculation and
the relative trivial/basic-stratum theorems
\cite[Theorems~III.2.4 and~III.4.5]{FarguesScholze} give actual sheaves of
frames in families.  They identify
$\mathcal Z/J_3=\mathbf P^2_C$.  In the other direction they
identify $\mathcal Z/P_3=\Omega_C^2$: the full frame group of
the positive bundle is $D_3^\times$, and its surjective norm
valuation lets one match the determinant lattice locally.
Taking the remaining quotient proves
\eqref{h3:eq:rat-two-presentations}.

The arithmetic field is $F=W(k)[1/p]$; the coefficient field
of the curve remains $\Q_p$.  Let $N$ be the fixed Honda
isocrystal over $k$, whose period-ring construction gives
$V_3$.  All Honda endomorphisms are defined over $k$.
On the period coefficients tensored over $F$ with $N$, define
the action of $\gamma\in G_F$ by $\gamma\otimes\hthreeid_N$.
It commutes with Frobenius and so descends to the bundle.
Every element of $D_3$ is an $F$-linear endomorphism of $N$
commuting with its Frobenius; its matrix entries are fixed
by $G_F$.  Thus this specified descent commutes with the
entire $D_3$-action.  We have not trivialized a torsor of
frames of an arbitrary descended basic bundle.

The untilt ideal $\ker\theta$ is functorial, so $\Delta_3$
inherits descent.  The relative slope-zero correspondence
makes its rational section space a continuous one-dimensional
representation of $G_F$.  This representation need not be
trivial.  Its valuation character has compact subgroup image
in $\Z$, hence zero, so every lattice in it is stable;
no generator is asserted to be invariant.
Define the arithmetic action on an injection $u$ by pullback
and the specified descent on source and target.  The standard
rational coordinate vectors of $\mathcal O^3$ are fixed, so
\[
 \gamma(guj^{-1})=g\gamma(u)j^{-1}
       \qquad(g\in P_3,\ j\in J_3).
\]
The determinant condition is preserved by lattice stability.
On the upper side, directions lie in
$\Q_p^3\otimes(\mathcal O(\infty)/\mathcal O)|_\infty$.
The semilinear scalar in the second factor is common to all
coordinates and cancels in projective space.  This is the usual
arithmetic action on $\Omega_C^2$, commuting with $J_3$.
No arithmetic invariant choice of a frame is required.
\end{proof}

\begin{lemma}\label{h3:rat:fixed-projective}
The first presentation in \eqref{h3:eq:rat-two-presentations} has
a natural arithmetic-equivariant splitting
\begin{equation}\label{h3:eq:rat-projective-h3}
 H^3_{\mathrm b}(\mathcal M)
   =H^3_{\mathrm b}(P_3)\oplus H^1_{\mathrm b}(P_3)(-1),
\end{equation}
whose first inclusion is the classifying map
$\operatorname{cl}_{c_{3,\mathcal M},\mathrm b}$ of
\Cref{h3:rat:finite-classifying}.
\end{lemma}

\begin{proof}
For a profinite set
$Q=\varprojlim_iQ_i$ with finite $Q_i$, continuity on spatial
diamonds and comparison of finite coefficients give
\[
 R\Gamma_v(\mathbf P^2_C\times\underline Q,\Lambda_\ell)
  \simeq\colim_i\hthreeMap\bigl(Q_i,
                 R\Gamma_{\acute et}(\mathbf P^2_C,\Lambda_\ell)\bigr).
\]
Here one applies \cite[Proposition~14.9]{ScholzeDiamonds} to
$\mathbf P^2_C\times Q_i$ and then
\cite[Proposition~14.10 and Lemma~15.6]{ScholzeDiamonds}; these statements
allow the coefficient ring $\Lambda_\ell$ also when its torsion
prime is $p$.  The finite-coefficient cohomology of the analytic
projective space agrees with its algebraic cohomology: the
proper comparison of \cite[Theorem~3.12]{ScholzeSurvey},
restating \cite[Theorem~3.7.2]{Huber}, applies to
$\mathbf P^2_C\to\Spec C$ with $\F_p$ coefficients, and
the coefficient exact sequences extend it to $\Lambda_\ell$.
Thus the projective-space calculation is valid as the actual
condensed geometric complex.

The anticanonical bundle has canonical projective equivariance,
including arithmetic descent, and restricts to $\mathcal O(3)$.
Since $p\geq5$, the class
$h_\ell=3^{-1}c_1(\omega_{\mathbf P^2}^{-1})$
on the quotient restricts to the hyperplane class.  Its powers
give an actual morphism of equivariant derived coefficient objects
\[
 \bigoplus_{j=0}^2\Lambda_\ell(-j)[-2j]
       \longrightarrow\mathscr C_{\mathbf P^2_C,\ell}.
\]
The preceding calculation shows that this is an equivalence on
every profinite parameter.  Apply it to $Q=P_3^a$ in the quotient
nerve and totalize.  The finite direct sum commutes with
totalization, giving an equivalence
\[
 \bigoplus_{j=0}^2
 C^\bullet_{\hthreects}(P_3,\Lambda_\ell(-j))[-2j]
       \xrightarrow{\ \simeq\ }
 R\Gamma_v(\mathcal M,\Lambda_\ell).
\]
In its $j=0$ component, each bar term sends a continuous constant
function to its coefficient-unit section.  By
\Cref{h3:rat:finite-classifying}, this component is exactly
$\operatorname{cl}_{c_{3,\mathcal M},\ell}$.
The Kummer classes commute with coefficient reduction.  Take the
derived limit in $\ell$ and then invert $p$; the finite direct sum
commutes with both operations.  Only $j=0,1$ contribute to degree
three, giving \eqref{h3:eq:rat-projective-h3}.  Naturality of the
Kummer classes gives its arithmetic compatibility and the
displayed twists.
\end{proof}

The projective splitting identifies the classifying image of the
$P_3$ trace.  We next need to show that the $J_3$ trace has a
nonzero classifying image on the same modification stack.  The
Drinfeld cohomology theorem gives arithmetic scalar actions on
point-valued cohomology; the following finite resolution justifies
their use on the compact-group descent rows.

\begin{lemma}[Finite resolutions for geometric coefficients]
\label{h3:rat:solid-resolution}
The complexes $\mathscr C_Z$ are bounded below, derived solid, and
derived $p$-complete.  Let $G$ be a compact $p$-adic analytic group
without $p$-torsion.  There is a finite projective resolution
$P_\bullet\to\Z_p$ over $\Lambda_G=\Z_p[[G]]$, with finitely
generated terms and length $\dim G$, such that, for every derived
solid $\underline{\Z_p}$-module complex $K$ with
$\underline{\Z_p}$-linear continuous $G$-action,
\begin{equation}\label{h3:eq:rat-solid-hom}
 R\Gamma(G,K)(*)\simeq\Hom_{\Lambda_G}(P_\bullet,K).
\end{equation}
Here $\Hom$ is the external complex of morphisms in solid
$\underline{\Lambda_G}$-modules.  A continuous action means an
actual module over $\underline{\Z_p}[\underline G]$ in the
derived condensed category, rather than an action on $K(*)$.
If $K$ is bounded below, and an operator $\gamma$ commuting with $G$
acts by the scalar $\lambda_j\in\Q_p$ on $H^j(K)(*)[1/p]$, it
acts by $\lambda_j$ on
$H^s(G,H^j(K))(*)[1/p]$ for every $s$.  For $K=\mathscr C_Z$ the
resulting spectral sequence is
\begin{equation}\label{h3:eq:rat-solid-descent}
 H^s(G,H^j(\mathscr C_Z))(*)[1/p]
       \ \Longrightarrow\ H^{s+j}_{\mathrm b}([Z/G]),
\end{equation}
with a finite filtration in each total degree.
\end{lemma}

\begin{proof}
Let $V=\operatorname{Spa}(R,R^+)$ be affinoid perfectoid, and write
$R^\flat$ for its tilt (equal to $R$ in characteristic $p$).
Its product with $\underline Q$ is affinoid perfectoid with ring
$C_{\hthreects}(Q,R)$ and tilt $C_{\hthreects}(Q,R^\flat)$.
Affinoid acyclicity gives
\[
 R\Gamma_v(V\times\underline Q,\mathcal O^\flat)
       =C_{\hthreects}(Q,R^\flat)[0]
\]
by \cite[Proposition~8.8]{ScholzeDiamonds}.  Choose a decreasing countable basis of
open additive subgroups $R^\flat_a$ of the complete additive group $R^\flat$.
Then $\underline{R^\flat}=\varprojlim_a(R^\flat/R^\flat_a)_{\mathrm{disc}}$ is solid.
This limit is also the derived limit: continuous maps from a
profinite set to each discrete quotient have finite image, and
lifts of their finitely many values give surjectivity at each
successive quotient.  A v-hypercover by disjoint unions of affinoid
perfectoids now expresses the $\mathcal O^\flat$ complex, with all
parameters, as a limit of products of solid modules.  Derived solid
modules are closed under limits and extensions
\cite[Theorem~4.4(ii)]{Tang}.  Artin--Schreier gives the assertion
for $\F_p$, the coefficient exact sequences give it for $\Z/p^\ell$,
and the derived limit gives it for $\mathscr C_Z$.  These complexes
are bounded below.  The finite coefficient complexes are derived
$p$-complete, so their limit is too.  Evaluation at the point is
exact and preserves limits; in particular
\[
 H^j(\mathscr C_Z)(*)=H^j(R\Gamma_{\mathrm{int}}(Z)).
\]

The compact-resolution argument in \Cref{h3:co:compact-resolution} gives
Noetherianity of $\Lambda_G$ and projective dimension $\dim G$ for
its trivial profinite module, using
\cite[Sections~1.1--1.2, pp.~275--276]{Venjakob} together with
\cite[Section~1, Corollary~(1)]{Serre1965CohomologicalDimension}.
Successive finite free presentations therefore give the stated
finite resolution.  We give the derived comparison needed to use
it.  Let $\mathcal C=D(\operatorname{Cond}_{\underline{\Z_p}})$
and $\mathcal D=D(\operatorname{Solid}_{\underline{\Z_p}})$.
Solidification is a symmetric monoidal localization
$L:\mathcal C\to\mathcal D$ with fully faithful right adjoint
$I$ \cite[Theorem~4.4(ii)]{Tang}.
For the free condensed group ring $A=\underline{\Z_p}[\underline G]$,
its solidification $LA$ is $\underline{\Lambda_G}$ in degree zero,
including multiplication \cite[Section~3.6 and the proof of
Theorem~4.4]{Tang}.
The underived assertion also holds for a free module on any
profinite set $Q$: choose an extremally disconnected hypercover
$Q_\bullet\to Q$ and solidify its free resolution.  The resulting
augmented complex is $\Z_p[[Q_\bullet]]\to\Z_p[[Q]]$.
Its Pontryagin dual is the augmented complex of continuous
$\Q_p/\Z_p$-valued functions.  Discrete coefficients are acyclic
on a profinite set, by passage to finite clopen partitions, so
that dual complex is exact.  Pontryagin duality and exact
condensation prove the assertion.  Functoriality with respect
to multiplication on $G$ gives the ring identification.

For any symmetric monoidal localization there is an induced
adjunction
\[
 L_A:\hthreeMod_A(\mathcal C)\ \rightleftarrows\
       \hthreeMod_{LA}(\mathcal D):J,
\]
where $J$ is inclusion followed by restriction along
$A\to I(LA)$.  Its unit has underlying map $M\to ILM$,
and its counit has underlying map $LIN\to N$.
The counit is an equivalence; hence $J$ is fully faithful,
with essential image exactly the modules whose underlying
complex is derived solid.  This proves the assertion in
unbounded degrees.  To identify the target, apply this same
localization to the already solid ring
$B=\underline{\Lambda_G}$.  Its essential image inside
$D(\operatorname{Cond}_B)$ consists of precisely the
complexes with derived solid underlying coefficients.
By \cite[Theorem~4.4]{Tang}, the unbounded category
$D(\operatorname{Solid}_B)$ has that same essential image.
This identifies the two derived module categories; it is
stronger than an equivalence of their hearts.  Equivalently,
the comparison is computed on free modules $B[Q]$ for
extremally disconnected profinite $Q$ and their solidifications,
which give the projective resolutions in these categories.
Thus the actual condensed action on $K$ belongs to this
category, and full faithfulness identifies its external
derived invariants with those computed there.  For geometric
$K$, the action and its coherent group law come from the
maps on $Z\times\underline Q$ with all group parameters
retained; the preceding solidity argument applies to its
underlying complex.

Condensation preserves the exact profinite resolution
\cite[Theorem~3.2]{Tang}.  Every $P_a$ is a retract of a finite
free $\Lambda_G$-module, and
$\Hom_{\Lambda_G}(\Lambda_G,M)=M(*)$ is exact.  Thus these terms
are projective against arbitrary solid coefficient modules and
prove \eqref{h3:eq:rat-solid-hom}.  The complex
$\Hom_{\Lambda_G}(P_\bullet,H^j(K))$ consists of retracts of
finite sums of $H^j(K)(*)$.  Its retracts commute with $\gamma$.
After exact inversion of $p$, the entire complex therefore has
scalar action $\lambda_j$, proving the assertion about its
cohomology.  No assertion about a condensed sheaf is inferred from
its point value.

Finally the nerve of $Z\to[Z/G]$ has terms
$Z\times\underline G^a$.  With parameters retained, its
descent complex is precisely the continuous bar construction on
$\mathscr C_Z$.  At finite level the equivariant coefficient unit
$\underline{\Lambda_\ell}\to\mathscr C_{Z,\ell}$ is the evaluation
map of \Cref{h3:rat:finite-classifying} on every profinite parameter.
Its map on $G$-invariants is therefore exactly
$\mathsf u_{Z,G,\ell}$, since their degree-$a$ bar maps agree.
This identity also holds after the derived coefficient limit and
inversion of $p$.  Its source after that limit is
$R\Gamma(G,\underline{\Z_p})(*)\simeq R\Gamma_{\mathrm{int}}(G)$:
the finite reductions of constants are surjective on every profinite
parameter, and internal invariants and point evaluation preserve
limits.  The resolution bounds the group direction by
$\dim G$, giving \eqref{h3:eq:rat-solid-descent} and the asserted
finite filtration before and after rationalization.
\end{proof}

\begin{lemma}\label{h3:rat:two-traces}
On $\mathcal M$, abbreviate the classifying images
$\operatorname{cl}_{c_{3,\mathcal M},\mathrm b}(e_3)$ and
$\operatorname{cl}_{c_{J,\mathcal M},\mathrm b}(e_J)$ by $e_3$ and
$e_J$.  Both are nonzero and satisfy $e_3=\kappa e_J$ for some
$\kappa\in\Q_p^\times$.  We retain this abbreviation for subsequent
pullbacks from $\mathcal M$.
\end{lemma}

\begin{proof}
By \eqref{h3:eq:rat-projective-h3}, the class $e_3$ is nonzero.
Arithmetic acts trivially on its first summand and by
$\chi_{\mathrm{cyc}}^{-1}$ on the second.  The cyclotomic
character of $G_F$ has open image, so the invariant subspace
is precisely $\Q_p e_3$.

Choose a compact uniform open subgroup $U\subset J_3$.  The
restriction $\operatorname{res}_U(e_J)$ of the source group class
$e_J\in H^3_{\mathrm b}(J_3)$ is nonzero by \Cref{h3:rat:trace-groups}.
Let $q_U:[\Omega_C^2/U]\to\mathcal M$ be the quotient map and
$c_U:[\Omega_C^2/U]\to\mathrm B_{B_C}U$ its relative classifying map.
Extension of torsor structure gives
$c_{J,\mathcal M}q_U\simeq(\mathrm B_{B_C}(U\hookrightarrow J_3))c_U$.
Thus \Cref{h3:rat:finite-classifying} identifies the restriction of
the geometric $e_J$ with
$\operatorname{cl}_{c_U,\mathrm b}(\operatorname{res}_U(e_J))$.
Apply \eqref{h3:eq:rat-solid-descent} to $[\Omega_C^2/U]$.
The integral Drinfeld computation
\cite[Theorems~1.1 and~5.1]{CDN} identifies
$H^j(\mathscr C_{\Omega_C^2})(*)[1/p]$ as a representation
on which arithmetic acts by the scalar
$\chi_{\mathrm{cyc}}^{-j}$, with $j=0,1,2$ and no higher
rows.  Its integral convention is the derived limit of finite
$p$-power coefficients followed by inversion of $p$, exactly
the convention here: the finite reduction tower defines
$\widehat{\Z}_p$ on the pro-etale site, and derived global
sections preserve its derived inverse limit.  The calculation is for hyperplanes
rational over the ground field, here $\Q_p$; we merely restrict
its arithmetic action to $G_F$.

In degree zero the same computation gives a one-dimensional rational
point value.  The finite coefficient units form a map
$\eta:\underline{\Z_p}\to\mathscr C_{\Omega_C^2}$ after their derived
limit: on every profinite parameter the finite reductions of constants
are surjective, so their limit is $\underline{\Z_p}$.
Pullback to any geometric point sends the unit to $1$.  Hence
\[
 H^0(\eta)(*)[1/p]:\Q_p
       \xrightarrow{\ \simeq\ }
 H^0(\mathscr C_{\Omega_C^2})(*)[1/p].
\]
For each finite projective term $P_a$ in the $U$-resolution of
\Cref{h3:rat:solid-resolution}, applying $\Hom_{\Lambda_U}(P_a,-)$
gives a retract of a finite sum of these point maps, and the unit
commutes with its defining idempotent.  It is therefore an isomorphism
after inversion of $p$ in every term.  Naturality of the descent
spectral sequence for $\eta$ maps its single constant row to the
bottom row of geometric descent.  By the compact bar identity, the
$E_2^{3,0}$ bottom-row map of
the spectral-sequence morphism induced by the finite units defining
$\operatorname{cl}_{c_U,\mathrm b}$ is the resulting isomorphism
\[
 H^3_{\mathrm b}(U)
       \xrightarrow{\ \simeq\ }
 H^3(U,H^0(\mathscr C_{\Omega_C^2}))(*)[1/p].
\]
It sends $\operatorname{res}_U(e_J)$ to the
nonzero bottom-row class.  This argument uses the rational point
value and the natural finite projective complex.

By \Cref{h3:rat:solid-resolution}, an element
$\gamma\in G_F$ with cyclotomic character of infinite order
acts on the actual $j$th descent row by this scalar.  The only
possible incoming differentials to $(3,0)$ are from $(1,1)$
and $(0,2)$.  Their source scalars differ from $1$, so both
differentials vanish.  There is no outgoing differential from
the bottom row.  Thus the nonzero bottom-row image of
$\operatorname{res}_U(e_J)$ survives in
$H^3_{\mathrm b}([\Omega_C^2/U])$, proving that the geometric class
$e_J$ on $\mathcal M$ is nonzero.  The map to $\mathrm B_{B_C}J_3$ is
arithmetic equivariant, so this class is invariant and belongs
to the line $\Q_p e_3$.  This proves the assertion.
\end{proof}

\subsection{The actual rank-two kernel on the ordinary stratum}

Let $\mathcal Y_1$ be the ordinary stratum stack of
\Cref{h3:prop:integral-frames}, over $B_k$.  Put
$G_o=\GL_2(\Z_p)\times P_1$.  Its two frame presentations give
relative classifying maps
\[
 c_{3,k}:\mathcal Y_1\longrightarrow\mathrm B_{B_k}P_3,
 \qquad c_{o,k}:\mathcal Y_1\longrightarrow\mathrm B_{B_k}G_o.
\]
On this stack, $e_3$ means $\operatorname{cl}_{c_{3,k},\mathrm b}(e_3)$,
and $e_2$ means
$\operatorname{cl}_{c_{o,k},\mathrm b}(\operatorname{pr}_1^*e_2)$,
where $\operatorname{pr}_1:G_o\to\GL_2(\Z_p)$ is the first projection.
Set $\mathcal Y_{1,C}=\mathcal Y_1\times_{B_k}B_C$ and write
$c_{3,C},c_{o,C}$ for the base-changed classifying maps.  The finite
base-change identity in \Cref{h3:rat:finite-classifying} identifies
the corresponding classes on $\mathcal Y_{1,C}$ with the base changes
of these same classes over $k$.

\begin{proposition}[Transport along the ordinary kernel]
\label{h3:prop:trace-transport}
There is $a\in\Q_p^\times$ such that
\begin{equation}\label{h3:eq:rat-ordinary-traces}
 e_3=a e_2\quad\text{in }H^3_{\mathrm b}(\mathcal Y_{1,C}).
\end{equation}
The class $e_2$ is pulled back from the actual integral frame
group of the rank-two kernel.
\end{proposition}

\begin{proof}
The universal sequence on this stack is
\[
 0\longrightarrow\mathcal L_2\otimes_{\Z_p}\mathcal O
 \longrightarrow\mathcal V_3\longrightarrow\mathcal V_1
 \longrightarrow0,
\]
where $\mathcal L_2$ is the integral rank-two local system of
\Cref{h3:prop:integral-frames}.  The canonical quotient
$\mathcal V_1\to i_{\infty*}i_\infty^*\mathcal V_1$ defines
$\mathcal V_1^- =\mathcal V_1(-\infty)$ and, by pullback,
a lower modification $\mathcal V_3^-$ of $\mathcal V_3$.
It gives a map $m:\mathcal Y_{1,C}\to\mathcal M$ and the exact
sequence of slope-zero bundles
\[
 0\longrightarrow\mathcal L_2\otimes\mathcal O
 \longrightarrow\mathcal V_3^-
 \longrightarrow\mathcal V_1^-\longrightarrow0.
\]
Apply the relative slope-zero equivalence on the pro-etale
site of a perfectoid test object over $\mathcal Y_{1,C}$.
To verify exactness without assuming it, take a common
pro-etale cover trivializing the three rational local systems.
Full faithfulness identifies the maps with matrices of
sections of $\underline{\Q_p}$.  Write the quotient map as
$(b_1,b_2,b_3)$.  The open loci $b_i\ne0$ cover, since its
rank is one on every geometric fiber.  On such a locus,
$1\mapsto b_i^{-1}e_i$ is a continuous section.  The kernel
has basis $e_j-(b_j/b_i)e_i$ for $j\ne i$.  The given map
from $\mathcal L_2\otimes\Q_p$ to this kernel is a square
matrix invertible on every geometric fiber; its inverse is
continuous by the determinant formula.  Thus the sequence
of rational local systems is exact and locally split.

The $P_3$-structure gives an integral determinant lattice
on the middle local system.  Its tensor quotient by
$\det\mathcal L_2$ is an integral lattice on the quotient
line.  Refine the cover to choose an integral basis of the
original $\mathcal L_2$ and an integral generator of this
quotient lattice, and lift the generator by the displayed
splitting.  The resulting middle frame has determinant
generating its prescribed determinant lattice.  Its possible
changes are exactly
$\bigl(\begin{smallmatrix}A&v\\0&a\end{smallmatrix}\bigr)$
with $A\in\GL_2(\Z_p)$, $a\in\Z_p^\times$, and
$v\in\Q_p^2$.  Its determinant is a unit, so these adapted
frames form exactly the $D$-torsor of \Cref{h3:rat:trace-groups}.
Forgetting the quotient generator and its lift gives the
original integral Tate-frame torsor under $D\to\GL_2(\Z_p)$.
These constructions commute with all perfectoid base changes
and descend on the v-site, giving the actual relative map
$c_D:\mathcal Y_{1,C}\to\mathrm B_{B_C}D$.  The identifications of
the associated torsors give homotopies of relative classifying maps
\[
 \begin{aligned}
 c_{3,\mathcal M}m&\simeq c_{3,C},\\
 (\mathrm B_{B_C}\iota)c_D&\simeq c_{J,\mathcal M}m,\\
 (\mathrm B_{B_C}\rho)c_D
     &\simeq(\mathrm B_{B_C}\operatorname{pr}_1)c_{o,C}.
 \end{aligned}
\]
The first remembers the same structured positive bundle.  The second
forgets the adapted filtration and retains the middle determinant
lattice induced from the same chosen $\Delta_3$ lattice.  The third
forgets the quotient generator and
its lift, retaining the original integral Tate frame.  These
forgetful identifications can be checked on the local frames just
constructed and then descend; they commute with base change and with
the specified semilinear arithmetic descent.  Applying the finite
classifying naturality and then its rational limit to
$e_3=\kappa e_J$ and \eqref{h3:eq:rat-parabolic-trace} therefore gives
$e_3=\kappa b e_2$.  Both factors are nonzero.
\end{proof}

\subsection{Witt constants and the localization map}

Define $H_W^i(Z)=H^i(R\varprojlim_\ell
R\Gamma_v(Z,W_\ell\mathcal O^\flat))$.  For $A=k$ and $A=C^\flat$,
write $\mathcal Y_{1,A}$ for $\mathcal Y_1$ and $\mathcal Y_{1,C}$,
respectively, and use $c_{o,A}$ for the corresponding ordinary
classifying map.  Give $W_\ell(A)$ its finite Witt coordinate
topology; it is discrete when $A=k$.  At each length there is a
commutative square of external derived complexes
\begin{equation}\label{h3:eq:rat-finite-witt-square}
\begin{tikzcd}[column sep=large]
 C^\bullet_{\hthreects}(G_o,\Lambda_\ell)
   \arrow[r,"\operatorname{cl}_{c_{o,A},\ell}"] \arrow[d]
 & R\Gamma_v(\mathcal Y_{1,A},\Lambda_\ell) \arrow[d]\\
 R\Gamma(G_o,\underline{W_\ell(A)})(*) \arrow[r,"\mathrm{constants}"]
 & R\Gamma_v(\mathcal Y_{1,A},W_\ell\mathcal O^\flat).
\end{tikzcd}
\end{equation}
The vertical arrows are induced by
$\Lambda_\ell=W_\ell(\F_p)\to W_\ell(A)$ and by the Witt coefficient
unit.  The bottom arrow is the actual ordinary constants map: over
$B_C$ it comes from \Cref{h3:prop:ordinary-constants}, and over $B_k$
from its base-field Frobenius descent in
\Cref{h3:thm:completed-cohomology}.  Its extension to Witt
coefficients is proved below.
On the torsor nerve both composites send a finite constant function
to the same section through
$W_\ell(\F_p)\to W_\ell(A)\to W_\ell\mathcal O^\flat$.
Thus \Cref{h3:rat:finite-classifying} proves the square term by term,
including its base change from $k$ to $C^\flat$ and its reductions in
$\ell$.  Over $B_k$, the same identity for the $P_3$ presentation
identifies the Witt image of $\operatorname{cl}_{c_{3,k},\ell}$ with
the map from $C^\bullet_{\hthreects}(P_3,\Lambda_\ell)$ to the
ordinary perfected Witt cochains in \Cref{h3:des:witt-constants}.
Here the canonical interchange of the two frame factors identifies
$G_o$ with the group $G_{\mathrm f}$ used there; the earlier geometric
maps written $BG_{\mathrm f}$ are these relative classifying maps on
their displayed bases.
Those proofs retain profinite parameters through the identification
of this quotient nerve with the computing cochains.

Taking the derived coefficient limits of the ordinary constants maps
gives the following square; their finite equivalences are justified
in the proof below:
\begin{equation}\label{h3:eq:rat-witt-square}
\begin{tikzcd}[column sep=large]
 H^i(G_o,W(k)) \arrow[r,"\sim"] \arrow[d]
     & H_W^i(\mathcal Y_1) \arrow[d]\\
 H^i(G_o,W(C^\flat)) \arrow[r,"\sim"]
     & H_W^i(\mathcal Y_{1,C}).
\end{tikzcd}
\end{equation}

\begin{lemma}\label{h3:rat:witt-nonvanishing}
The right vertical map of \eqref{h3:eq:rat-witt-square} is injective
after inversion of $p$.  The Witt coefficient image of $e_3$
on $\mathcal Y_1$ is nonzero.  Its corresponding image under
the actual ordinary deformation-coefficient map is nonzero.
\end{lemma}

\begin{proof}
For completeness, the Witt extension of the constants calculation
uses the natural additive exact sequences
\[
 0\longrightarrow R\xrightarrow{V^{\ell-1}}W_\ell(R)
    \longrightarrow W_{\ell-1}(R)\longrightarrow0.
\]
They are exact for the solid coefficient objects and v-sheaves
in question: in Witt coordinates, restriction has the continuous
set-theoretic section that appends zero.  Apply the sequence to
both sides of the characteristic-$p$ constants equivalence.
Induction using the resulting morphisms of exact triangles
proves the equivalence at every length.  Take derived inverse
limits and then the residual group invariants.  All maps in this
construction are induced by constants, so it proves the square
with its asserted arrows and their compatibility with cup products.

The group $G_o$ has no $p$-torsion.  Choose its resolution in
\Cref{h3:rat:solid-resolution} and put
$Q_\bullet=\Z_p\otimes_{\Z_p[[G_o]]}P_\bullet$.
Each $Q_j$ is finite projective, hence finite free over $\Z_p$.
For a trivial coefficient module $M$, the computing complex is
$\Hom_{\Z_p}(Q_\bullet,M)$, a fixed bounded complex of finite
matrices over $\Z_p$ tensored with $M$.
Since $W(C^\flat)$ is torsion free over the discrete valuation
ring $W(k)$, it is flat, and its fraction field extension is
faithfully flat.  Consequently
\[
 H^i(G_o,W(k))[1/p]\otimes_{W(k)[1/p]}W(C^\flat)[1/p]
       \simeq H^i(G_o,W(C^\flat))[1/p].
\]
This proves the injection in the square.  It also proves that
the image of $e_2$ is nonzero: it is the scalar extension of
the nonzero trace class of the first factor of $G_o$.

By \eqref{h3:eq:rat-finite-witt-square} and its $P_3$ counterpart,
the two Witt images of the group classes on $\mathcal Y_{1,C}$ are
the coefficient-unit images of the exact classifying images in
\eqref{h3:eq:rat-ordinary-traces}.  Their maps over $k$ base change
to these maps before the derived limit is taken.  Applying that
limit and inversion of $p$ to the geometric equality, and then
using the injection in \eqref{h3:eq:rat-witt-square}, proves the
same equality over $k$ in Witt cohomology.  Its right side is
nonzero, as just shown.

It remains to follow this nonzero Witt class through the actual
deformation-coefficient map.  Put
\[
 \mathsf C^W_\ell=C^\bullet_{\hthreects}(P_3,W_\ell B_0),\qquad
 \mathsf C^D_\ell=C^\bullet_{\hthreects}(P_3,D_\ell),\qquad
 D_\ell=(W(k)/p^\ell)[[x,y]][x^{-1}].
\]
By \Cref{h3:lem:ghost-comparison}, the two maps from $W_\ell B_0$
to $D_\ell$ and to $W_\ell B_{\hthreepk}$ induce cohomology
isomorphisms.  The graded maps for $\gamma_\ell$ are powers of
$\phi_p$, while those for $\beta_\ell$ are the natural maps
$B_0\to B_{\hthreepk}$.  Both are cohomological equivalences by
\Cref{h3:thm:ordinary-comparison}.  This statement is used at
each finite length; the two inverse systems have different
transition maps.

Choose an integral constant class
$\beta\in H^3_{\mathrm{int}}(P_3)$ with
$\beta[1/p]=p^N e_3$ for some $N\geq0$, where $e_3$ denotes the
source class in $H^3_{\mathrm b}(P_3)$.  The normalized rational
trace is fixed by $\Gamma=\Gal(k/\F_p)$, as proved in
\Cref{h3:rat:trace-groups}.  Thus each difference
$\gamma\beta-\beta$ is killed by a power of $p$.  Since $\Gamma$
is finite, one common power kills all these differences.  Replacing
$\beta$ by that multiple and increasing $N$ makes the integral
cohomology class $\Gamma$-invariant while retaining the same nonzero
rational trace line.  Let $\bar\beta_\ell$ be its canonical projection
to $H^3(C^\bullet_{\hthreects}(P_3,\Lambda_\ell))$.  These finite
cohomology classes are reduction-compatible and $\Gamma$-invariant.
Let
$a\in H^3(R\varprojlim_{R_\ell}\mathsf C^W_\ell)$ be its coefficient image,
where $R_\ell:W_\ell B_0\to W_{\ell-1}B_0$ is ordinary Witt restriction.
The maps to perfected and geometric Witt coefficients commute with
these restrictions.  The nonvanishing already proved therefore implies
$a[1/p]\ne0$.  The groups $H^2(\mathsf C^W_\ell)$ are finite, by
the finite-length additive Witt filtration and
\Cref{h3:prop:bounded-pole-finiteness,h3:thm:ordinary-comparison}.
Their tower is Mittag--Leffler, so the Milnor sequence gives
\[
 H^3(R\varprojlim_{R_\ell}\mathsf C^W_\ell)
      =\varprojlim_{R_\ell} H^3(\mathsf C^W_\ell).
\]
Consequently the component classes $a_\ell$ have unbounded
$p$-power orders.  By naturality of the projections, $a_\ell$ is the
coefficient image of $\bar\beta_\ell$ in $W_\ell B_0$ cohomology.

Write $d_\ell$ for their ghost images in $H^3(\mathsf C^D_\ell)$.
These have exactly the same orders.  The ghost maps fix
$\Z/p^\ell$, so $d_\ell$ is the direct coefficient image
of $\bar\beta_\ell$ and is compatible with coefficient reduction.
More explicitly the transition identity is
\[
 \operatorname{red}_{\ell,\ell-1}\gamma_\ell
       =\gamma_{\ell-1}R_\ell W_\ell(\phi_p),
\]
where $\phi_p(b)=b^p$ is the characteristic-$p$ Frobenius defined
in \Cref{h3:lem:ghost-comparison}.  The class $a_\ell$ is fixed by
$W_\ell(\phi_p)$ because it comes from $W_\ell(\F_p)$.
The naturalities of the finite units also make each $d_\ell$
$\Gamma$-invariant.  Exact semilinear descent in
\Cref{h3:lem:framework-semilinear,h3:prop:exact-descent} identifies, by restriction,
\[
 H^3_{\hthreects}(G_3(k),D_\ell)
       \simeq H^3_{\hthreects}(P_3,D_\ell)^\Gamma.
\]
Here descent is applied after extension to the semilinear $D_\ell$
coefficients.
The descended class restricts to $d_\ell$, so it has the same order;
these identifications also preserve ordinary reduction.  Thus the
specified descended constant classes have unbounded $p$-power order.
The natural
spectral comparison in \Cref{h3:des:minus-three-edge,h3:prop:integral-delta}, with its
unique bidegree $(s,t)=(3,0)$ in stem $-3$, identifies this family
with the finite projections of a class in $\pi_{-3}L_{K(1)}T$.
Its unbounded finite orders make that class nonzero after
rationalization, and the identification preserves the actual maps
from constants.
\end{proof}

\begin{lemma}[Origin in the rational sphere]\label{h3:lem:constants-origin}
The reduced-trace class in $H^3_{\mathrm b}(P_3)$, after the
constant coefficient map and coefficient-field descent, is the
degree-three class of an element
$\alpha_3\in\pi_{-3}L_0T$.  This element is nonzero and its
image in $\pi_{-3}L_0L_{K(1)}T$ is nonzero.
\end{lemma}

\begin{proof}
Reduction followed by the ordinary coefficient localization
$E_{3,0}(k)\to D_\ell$ sends the integral constant class chosen in
\Cref{h3:rat:witt-nonvanishing} to its specified finite classes
$d_\ell$.  These have unbounded $p$-power order and are compatible
with ordinary reduction.  Hence the coefficient image of the
rational trace in $H^3(P_3,E_{3,0}(k))[1/p]$ is nonzero.
The trace construction and its coefficient map commute with the
semilinear finite Galois action; their class is invariant and
descends by \Cref{h3:lem:framework-semilinear}.

We explain why this cohomology class is an actual rational
homotopy class.  A central element $z=1+p$ acts trivially on
$E_{3,0}(k)$ and by the nontrivial scalar $z^{-q}$ on the
nonzero periodicity twist $E_{3,2q}(k)$, for $q\ne0$.
This is the action on $\omega^{\otimes q}$ in the
contravariant deformation convention of
\Cref{h3:lem:invariant-differential}.  A central group element acts as the
identity on group cohomology when its coefficient action is
combined with its trivial inner conjugation.  Equivalently,
the natural transformation furnished by that element gives
a homotopy to the identity on the bar resolution.  Thus
$z^{-q}-1$ annihilates
$H^s(P_3,E_{3,2q}(k))$.  It is a nonzero element of $\Z_p$,
so every such row vanishes after inversion of $p$.

The descent of \Cref{h3:prop:exact-descent} is exact and strongly
convergent.  The finite stabilizer resolutions and exact semilinear
descent of \Cref{h3:rat:solid-resolution,h3:lem:framework-semilinear}
bound its cohomological amplitude.
Rationalization therefore preserves its finite filtrations.
Only internal degree zero remains.  After finite coefficient-field
descent, the actual coefficient maps fit into the commutative square
\[
\begin{tikzcd}[column sep=2em]
 H^3_{\hthreects}(G_3(k),E_{3,0}(k))[1/p]
   \arrow[r] \arrow[d,"\simeq"]
 & \bigl(\varprojlim_\ell
       H^3_{\hthreects}(G_3(k),D_\ell)\bigr)[1/p]
   \arrow[d,"\simeq"] \\
 \pi_{-3}L_0T \arrow[r]
 & \pi_{-3}L_0\operatorname*{holim}_\ell L_1(T/p^\ell).
\end{tikzcd}
\]
The upper arrow is induced by the coefficient reductions and
localizations just described.  The left isomorphism is the
single-row rational descent calculation.  The right isomorphism
uses the natural, unique $(s,t)=(3,0)$ edge of
\Cref{h3:des:finite-length-model} and the vanishing of the Milnor
$\lim^1$ term in \Cref{h3:prop:integral-delta}.  The finite-length
model identifies those coefficient maps with the actual Moore
quotient and localization maps on every descent term.
For commutativity, first choose an integral representative of a
$p$-power multiple of a rational source cohomology class.
Finite source filtration and rational collapse allow a further
$p$-power multiple to survive to homotopy.  Naturality in
\Cref{h3:prop:exact-descent} identifies its finite target edges
with that same multiple of the coefficient images, for every
$\ell$.  This multiplier is chosen in the source, independently
of $\ell$.  Taking the coefficient limit and then inverting $p$
gives the displayed square.

By \Cref{h3:lem:height-one-completion}, the bottom map is the
rationalization of $T\to L_{K(1)}T$.  The nonzero rational trace
image in the upper left therefore determines
$\alpha_3\in\pi_{-3}L_0T$, and its image in the upper right is the
nonzero rational class of the compatible finite constants.
The square proves the asserted nonvanishing of its actual
localization without an additional arithmetic coefficient splitting.
\end{proof}

\begin{theorem}[Rational compatibility]\label{h3:thm:rational-compatibility}
The generators in \eqref{h3:eq:framework-rational-boundary} may be
chosen so that $\alpha_1=\zeta$ and the actual localization map
$T\to L_{K(1)}T$ satisfies
\begin{equation}\label{h3:eq:rat-actual-images}
 \alpha_3\longmapsto c\delta,\qquad
 \zeta\alpha_3\longmapsto c\zeta\delta
       \qquad(c\in\Q_p^\times)
\end{equation}
after rationalization.  In particular the rational homotopy of
$T$ has basis
\[
 1,\ \zeta,\ \alpha_3,\ \zeta\alpha_3,\
 \alpha_5,\ \zeta\alpha_5,\ \alpha_3\alpha_5,\
 \zeta\alpha_3\alpha_5
\]
in degrees $0,-1,-3,-4,-5,-6,-8,-9$, respectively.
\end{theorem}

\begin{proof}
The rational exterior calculation
\eqref{h3:eq:framework-rational-boundary} makes degree $-3$
one-dimensional.  Choose its generator to be the nonzero
constant-origin class of \Cref{h3:lem:constants-origin}.  The
determinant class is the degree-one group primitive by its
construction in \Cref{h3:sec:descent}, and its actual image is
nonzero rationally by \Cref{h3:thm:height-one-maps}.  Thus it
may be used as $\alpha_1$.  Choose any remaining exterior generator in
degree $-5$.

By \Cref{h3:prop:integral-delta}, the rational target in degree
$-3$ is precisely $\Q_p\delta$.  Nonvanishing of the actual
image gives $\alpha_3\mapsto c\delta$ with $c\ne0$.
All coefficient comparisons and localizations used above
respect products, and the same actual determinant class
occurs on source and target.  Multiplication by it gives
the second formula with the same scalar.  The exterior
calculation then supplies the displayed basis.
\end{proof}

\section{Applying fracture to the height-three overlap}\label{sec:application}

The coefficient arguments have now produced the maps required by the
finite construction of \Cref{prop:fracture}.  We apply it with one
choice of local maps and compatibility homotopy, obtaining both the
filtration and its attachment identification.  The local basis results
used here were proved without the rational-obstruction theorem, which
enters only in the final subsection.

\subsection{Checking the fracture data}

\begin{proof}[Proof of \Cref{thm:main} and the height-three instance of
\Cref{thm:attachment-identification}]
Put \(T=L_{K(3)}S\), \(X=L_2T\), and \(R_1=L_{K(1)}S\).
Recall \(Q=L_0S\simeq H\Qp\), and set
\[
 W=W_1\vee W_2=L_2S\vee\Sigma^{-1}L_2S.
\]
\Cref{h3:thm:height-two-test} supplies the actual unit and determinant
map \(w=(1,\zeta):W\to X\), a \(K(2)\)-equivalence.  This verifies
\textup{(i)} of \Cref{prop:fracture}.  Put \(Y=\cofib(w)\) and
write \(q:X\to Y\).

For the middle block, set
\[
 U_{\mathrm{mid}}=U_3\vee U_4=\Sigma^{-3}L_1S\vee\Sigma^{-4}L_1S.
\]
\Cref{h3:thm:height-one-maps} supplies the equivalence of actual
\(R_1\)-module maps
\begin{equation}\label{eq:height-one-basis}
 (1,\zeta,\delta,\zeta\delta):
 R_1\vee\Sigma^{-1}R_1\vee\Sigma^{-3}R_1\vee\Sigma^{-4}R_1
 \xrightarrow{\simeq}L_{K(1)}T.
\end{equation}
Its first two components are the localizations of the same maps from
\(S\) used in \(w\).  Their restrictions are the same unit and
determinant maps, and the module localization adjunction
\eqref{eq:module-localization} identifies their extensions to \(R_1\),
including the required homotopies.  The fourth component is the specified
product \(\zeta\delta\).  Localizing the cofiber sequence at \(K(1)\)
therefore identifies
\(L_{K(1)}Y\) with the cofiber of the first two components of
\eqref{eq:height-one-basis} using those specified maps.
Consequently the last two components of
\eqref{eq:height-one-basis}, followed by the quotient, give
the required equivalence
\begin{equation}\label{eq:kappa-application}
 \kappa:L_{K(1)}U_{\mathrm{mid}}\xrightarrow{\simeq}L_{K(1)}Y.
\end{equation}
They are represented by the projections of \(\delta\) and
\(\zeta\delta\).

For the rational part of the middle map,
\Cref{h3:thm:rational-compatibility} gives
\[
 \pi_*L_0T=\Lambda_{\Qp}(\zeta,\alpha_3,\alpha_5),
 \qquad |\zeta|=-1,\quad|\alpha_3|=-3,\quad|\alpha_5|=-5,
\]
and identifies the actual rationalized localization map by
\begin{equation}\label{eq:rational-images}
 \alpha_3\longmapsto c\delta,\qquad
 \zeta\alpha_3\longmapsto c\zeta\delta
 \quad\text{for one }c\in\Qp^\times.
\end{equation}
Set
\[
 \beta_3=c^{-1}\alpha_3.
\]
This normalization is made only in the rational source.
No condition \(c\in\Zp^\times\) is required.
The projections of \(\beta_3\) and \(\zeta\beta_3\) to
\(L_0Y\) define a \(Q\)-module map
\[
 r:L_0U_{\mathrm{mid}}=\Sigma^{-3}Q\vee\Sigma^{-4}Q\longrightarrow L_0Y.
\]
By \eqref{eq:rational-images}, their images in
\(L_0L_{K(1)}Y\) are the projections of
\(\delta\) and \(\zeta\delta\).  These are also the images,
under \(L_0\kappa\), of the two unit generators of \(L_0U_{\mathrm{mid}}\).
The maps in \eqref{eq:compatibility} therefore agree on the
generators of this finite free \(Q\)-module.  An identifying
homotopy of \(Q\)-module maps exists; choose one as the
compatibility datum in \textup{(ii)} of \Cref{prop:fracture}.

It remains to check the final quotient.  Since \(L_0L_iS=Q\)
for \(i=1,2\), rationalization of \(w\) is the map from
\(Q\vee\Sigma^{-1}Q\) represented by \(1,\zeta\).
These are two distinct members of the exterior basis in
\Cref{h3:thm:rational-compatibility}, so this map is injective
on rational homotopy.  Its cofiber \(L_0Y\) has one copy of
\(\Qp\) in each degree
\[
 -3,\ -4,\ -5,\ -6,\ -8,\ -9.
\]
There are no suspended kernel terms.  The map \(r\) is likewise
injective on homotopy and includes the first two lines, represented
by \(\beta_3,\zeta\beta_3\).

Put
\[
\begin{aligned}
 V&=V_5\vee V_6\vee V_7\vee V_8\\
  &=\Sigma^{-5}Q\vee\Sigma^{-6}Q
       \vee\Sigma^{-8}Q\vee\Sigma^{-9}Q.
\end{aligned}
\]
Define \(b:V\to L_0X\simeq L_0T\) by the four ordered classes
\begin{equation}\label{eq:rational-basis-application}
 \alpha_5,\qquad \zeta\alpha_5,\qquad
 \beta_3\alpha_5,\qquad\zeta\beta_3\alpha_5.
\end{equation}
They represent the remaining quotient lines in degrees
\(-5,-6,-8,-9\).  The composite \eqref{eq:quotient-basis}
is therefore an isomorphism on all homotopy groups, and hence
an equivalence of rational modules.  This verifies
\textup{(iii)}.

Carry out \Cref{prop:fracture} with this single choice of
\(w,\kappa,r,b\) and compatibility homotopy, choosing coherent
inverse data for \(L_{K(2)}w\) and for \(L_Ja\) at the steps where
those equivalences are established.  This gives the same
\(h,F_i,a,Z,t,e\) used by \Cref{thm:attachment-identification},
with exactly the ordered cofibers in \Cref{thm:main}.
The terminal projection \(e:F_8\to X\) is an equivalence.  Its
restriction to \(F_1=L_2S\) is the first component of \(w\),
namely the canonical map \(L_2(S\to T)\), proving \Cref{thm:main}.
For these same choices, \Cref{thm:attachment-identification}
identifies all eight connecting maps by its two signed roofs.
This proves its height-three instance without replacing the
filtration or choosing new local maps.
\end{proof}

\begin{remark}
The exact compatibility in \eqref{eq:rational-images} is needed
for this application.  Rational dimensions alone would not give
the homotopy in \eqref{eq:compatibility}; independently chosen
degree-three and degree-four maps would not ensure that the same
normalization works for the product with \(\zeta\).
\end{remark}

\subsection{The first middle attachment}\label{sec:attachment}

The filtration has been constructed without a nonvanishing
assumption.  We now ask whether its first height-one layer attaches
nontrivially.  The answer depends on the actual canonical
\(K(2)\)-localization map, not on the rational dimensions of
the layers.

First note why rationalizing the filtration does not answer this
question.  Its rationalized cofibers are single \(Q\)-modules
in the strictly decreasing degrees
\(0,-1,-3,-4,-5,-6,-8,-9\).
Inductively, if the previous rationalized boundaries vanish,
\(L_0F_{i-1}\) is the sum of the earlier shifts of \(Q\).
The degree of the new cofiber is smaller than every earlier
degree, and hence cannot be a homotopy degree of
\(\Sigma L_0F_{i-1}\).  A \(Q\)-module map from the new shift
to that suspension is zero.  This proves by induction that
every rationalized boundary vanishes.  A spectral boundary can
nevertheless be nonzero.

\subsubsection{A conditional lifting criterion}

Let
\[
 u_X:L_0X\longrightarrow L_0L_{K(2)}X
\]
be induced by the canonical localization unit.  Retain the
specific \(h:U_{\mathrm{mid}}\to Y\), \(\beta_3\), and
\(d_3=\partial_wh|_{U_3}\) constructed in
\Cref{sec:application}, where \(U_3=\Sigma^{-3}L_1S\).

\begin{proposition}[Canonical-map criterion]\label{prop:attachment-criterion}
If \(u_X\) does not kill
\(\beta_3\in\pi_{-3}L_0X\), then
\[
 d_3:U_3=\Sigma^{-3}L_1S\longrightarrow\Sigma W
\]
is nonzero as a map of underlying spectra.
\end{proposition}

\begin{proof}
Suppose \(d_3\) were null after forgetting the module structure.
Applying spectral mapping spaces from \(U_3\) to the
cofiber sequence
\[
 W\longrightarrow X\xrightarrow qY\xrightarrow{\partial_w}\Sigma W
\]
shows that the chosen nullhomotopy lifts \(h|_{U_3}\) to a map
\(\widetilde h_3:U_3\to X\).
Rationally, \(L_0W\) has homotopy only in degrees \(0\) and
\(-1\).  Its groups in degrees \(-3\) and \(-4\) are zero,
so the cofiber sequence makes
\(\pi_{-3}L_0X\to\pi_{-3}L_0Y\) an isomorphism.
The construction of \(h\) sends the degree-\(-3\) generator
of \(L_0U_3\) to the projection of \(\beta_3\).
Thus \(L_0\widetilde h_3\) sends it to \(\beta_3\) itself.
This uses only the image of that underlying homotopy element.

But \(L_{K(2)}U_3=0\), since \(U_3\) is \(E(1)\)-local.
By localization adjunction, every map from \(U_3\) to the
\(K(2)\)-local target \(L_{K(2)}X\) is null.  In particular,
the composite of \(\widetilde h_3\) with the unit
\(X\to L_{K(2)}X\) is null.  Rationalizing and using naturality
would make \(u_X\) kill \(\beta_3\), contrary to the hypothesis.
\end{proof}

This criterion is independent of the existence argument:
its antecedent was not used to construct \(h\), the stages, or
their terminal equivalence.

\subsubsection{The companion canonical-map theorem}

For every prime \(p\geq5\), the canonical map for the uncompleted
sphere \(\mathbb S\),
\begin{equation}\label{eq:companion-map}
 L_0L_{K(3)}\mathbb S
 \longrightarrow L_0L_{K(2)}L_{K(3)}\mathbb S,
\end{equation}
is nonzero on \(\pi_{-3}\).
This is \cite[Theorem~1.1]{CompanionHeightThreeObstruction}.

\begin{corollary}[Nonzero first middle attachment]\label{cor:nonzero-attachment}
For every prime \(p\geq5\), retain the common height-three choice
in \Cref{sec:application}, which realizes \Cref{thm:main} and the
height-three instance of \Cref{thm:attachment-identification}.
Then
\[
 d_3:\Sigma^{-3}L_1S\longrightarrow
       \Sigma\bigl(L_2S\vee\Sigma^{-1}L_2S\bigr)
\]
is nonzero as a map of underlying spectra.
\end{corollary}

\begin{proof}
Set \(T_0=L_{K(3)}\mathbb S\).  The completion map
\(\mathbb S\to S\) is a mod-\(p\) equivalence.
After \(p\)-localization its cofiber has invertible
multiplication by \(p\), hence is rational and \(K(3)\)-acyclic.
Completion therefore induces an equivalence
\(e_c:T_0\xrightarrow{\simeq}T\).
Naturality of the \(K(2)\)-localization unit gives
\begin{equation}\label{eq:completion-square}
\begin{tikzcd}[column sep=large]
 L_0T_0 \arrow[r] \arrow[d,"L_0e_c"',"\simeq"]
   & L_0L_{K(2)}T_0
       \arrow[d,"L_0L_{K(2)}e_c","\simeq"']\\
 L_0T \arrow[r] & L_0L_{K(2)}T.
\end{tikzcd}
\end{equation}
The top map is nonzero on \(\pi_{-3}\) by
\cite[Theorem~1.1]{CompanionHeightThreeObstruction}.  Hence so is the bottom map.

Let \(\ell:T\to X=L_2T\) be the localization map.
Its cofiber is \(E(2)\)-acyclic.  Since
\(\langle E(2)\rangle=\langle K(0)\vee K(1)\vee K(2)\rangle\),
that cofiber is both rationally and \(K(2)\)-acyclic.
The vertical maps in the next naturality square are therefore
equivalences:
\begin{equation}\label{eq:localization-square}
\begin{tikzcd}[column sep=large]
 L_0T \arrow[r] \arrow[d,"L_0\ell"',"\simeq"]
   & L_0L_{K(2)}T
       \arrow[d,"L_0L_{K(2)}\ell","\simeq"']\\
 L_0X \arrow[r,"u_X"] & L_0L_{K(2)}X.
\end{tikzcd}
\end{equation}
Thus \(u_X\) is nonzero on \(\pi_{-3}\), and it is the same
canonical map under these identifications.  An arbitrary nonzero
map between the same two objects would not suffice.

The square \eqref{eq:localization-square} is a square of
\(S\)-modules by \eqref{eq:module-localization}.
Rationalization is extension to \(Q=H\Qp\), so its homotopy
maps are \(\Qp\)-linear.  Equivalently, the action of
\(\pi_0S=\Zp\) on rational homotopy extends uniquely across
inversion of \(p\).  By \Cref{h3:thm:rational-compatibility},
\[
 \pi_{-3}L_0X=\Qp\alpha_3.
\]
A nonzero \(\Qp\)-linear map from this one-dimensional space
cannot kill \(\alpha_3\) or its nonzero scalar multiple
\(\beta_3=c^{-1}\alpha_3\).  This conclusion requires no
numerical identification of a primitive chosen in the companion
proof with the chosen \(\alpha_3\).

The hypothesis of \Cref{prop:attachment-criterion} is now
satisfied, so \(d_3\ne0\).
\end{proof}

Thus the companion theorem enters only in verifying the antecedent of
the conditional criterion.  The stages and local bases in
\Cref{thm:main} were constructed independently.

\appendix
\section{The rational boundary}\label{h3:app:rational-boundary}

We give the height-three argument of
\cite[Sections~2.5--2.6 and 3.9--6.3]{BSSW} for the rational algebra
used in \Cref{h3:thm:rational-compatibility}.  The separate trace
calculation in \Cref{h3:sec:rational} identifies the localization of
its degree-three generator.

\begin{proposition}[The rational height-three sphere]\label{h3:prop:rational-boundary}
For every prime \(p\geq5\), there is an isomorphism of graded algebras
\[
 \pi_*L_0T\simeq\Lambda_{\Q_p}(\alpha_1,\alpha_3,\alpha_5),
 \qquad |\alpha_i|=-i.
\]
\end{proposition}

The proof reduces this homotopy calculation to a comparison with
constant coefficients.  Put
\[
 \begin{aligned}
 \overline W&=W(\overline{\F}_p),&
 \breve K&=\overline W[1/p],\\
 \overline A&=\overline W[[u_1,u_2]],&
 \overline G&=P_3\rtimes\Gal(\overline{\F}_p/\F_p).
 \end{aligned}
\]
The map to be proved an isomorphism is the constants inclusion
\begin{equation}\label{h3:eq:rational-boundary-constants}
 H^*(\overline G,\overline W)[1/p]
 \longrightarrow H^*(\overline G,\overline A)[1/p].
\end{equation}
The compact Lie-algebra calculation identifies the source with the
exterior algebra in cohomological degrees $1,3,5$.  Comparisons of
the Lubin--Tate and Drinfeld towers identify the target with the
same graded vector space after adjoining to both sides an exterior
class $\epsilon$ of degree one.  A continuous equivariant additive
retraction onto $\overline W$ makes the source a direct summand.
Comparing dimensions, including the common $\epsilon$ factor,
then forces the complementary cohomology to vanish.  Since the
original constants inclusion is multiplicative, this proves the
algebra comparison.  Rational Morava descent converts it to the
asserted homotopy algebra.

The tower comparisons must be made integrally: the arithmetic
estimates give torsion bounds uniform over the charts and radii,
so that the derived limits and group invariants can be taken before
inverting $p$.  We first establish the Drinfeld model's actual
scalar equivalence.  Its building calculation also supplies the
properties used for the group $J_3$ in \Cref{h3:rat:trace-groups}.

\subsection{The Drinfeld model and its building}

\begin{lemma}[The reduced rank-three building]\label{h3:lem:rank-three-building}
Let $F$ be a non-Archimedean local field with normalized valuation $\nu$,
valuation ring $\mathcal O$, uniformizer $\pi$, and residue field $\kappa$
of cardinality $q$.  Let $\mathcal B$ be the reduced affine building of
$\mathrm{GL}_3(F)$, equivalently the building of $\mathrm{PGL}_3(F)$, and put
\[
 J_3(F)=\{g\in\mathrm{GL}_3(F):\nu(\det g)=0\}.
\]
Then $\mathcal B$ is a contractible, locally finite, two-dimensional
simplicial complex.  Each vertex has $2(q^2+q+1)$ adjacent vertices.
The group $J_3(F)$ preserves vertex types.  The setwise stabilizer of every
nonempty face is compact open in $J_3(F)$ and fixes that face pointwise.
Any closed chamber is a fundamental domain, including all its faces.
\end{lemma}

\begin{proof}
For $\mathrm{GL}_3$, the semisimple quotient used to define the reduced
building is the split group $\mathrm{PGL}_3$ of rank two.
The geometry statement in
\cite[Section~1.1.2, p.~6]{MeyerSolleveldBuildings}
applies over the local field $F$: the reduced building is locally finite,
has dimension the rank of this quotient, and is equivariantly contractible
for every compact subgroup of the quotient.  Taking the trivial compact
subgroup gives contractibility.

In the lattice-chain model of
\cite[Section~5.6, pp.~952--953]{DeligneFlicker2013}, vertices are homothety
classes $[L]$ of $\mathcal O$-lattices in $F^3$, and a chamber can be
represented by a strict cyclic chain
\[
 L_0\supset L_1\supset L_2\supset\pi L_0.
\]
Modulo $\pi L_0$ this is a complete flag in $\kappa^3$.  Thus a chamber
has three vertices and is a two-simplex.  An adjacent vertex to $[L]$ has
a unique representative $M$ with $\pi L\subsetneq M\subsetneq L$.
These representatives correspond to the nonzero proper subspaces of
$L/\pi L\cong\kappa^3$.  There are $q^2+q+1$ lines and the same number
of planes, giving the stated valency and a direct local-finiteness check.
Since the complex is locally finite, its metric polyhedral and weak
simplicial topologies agree locally on finite unions of simplices.  Its
realization is therefore a contractible CW complex.  The augmented
cellular chain complex over $\Z$ is exact, with free direct-sum terms
in degrees $0,1,2$.

Fix the standard lattice $\mathcal O^3$ and define
\[
 \operatorname{type}([g\mathcal O^3])=\nu(\det g)\pmod 3.
\]
Changing an integral basis changes the determinant by a unit; scalar
homothety changes its valuation by a multiple of three.  This type is
therefore well defined, and $J_3(F)$ preserves it.  The three vertices of
a chamber have all three types, since successive flag steps change the
lattice determinant valuation by one.

To send the standard chamber to a given chamber, start at the latter's
unique type-zero vertex.  Write its lattice as $M=g\mathcal O^3$ with
$\nu(\det g)=3m$, and replace $M$ by $M_0=\pi^{-m}M$.  Reanchor the
cyclic chain at $M_0$.  A basis $f_1,f_2,f_3$ of $M_0$ adapted to its
residue flag gives
\[
 M_i=\pi\mathcal O f_1+\cdots+\pi\mathcal O f_i
       +\mathcal O f_{i+1}+\cdots+\mathcal O f_3
 \quad(i=1,2).
\]
The map from the standard basis to $(f_1,f_2,f_3)$ has determinant
valuation zero, because its image lattice is $M_0$.  It lies in $J_3(F)$
and sends the standard chamber to the given chamber.

If $h\in J_3(F)$ stabilizes $[L]$, then $hL=aL$ for some $a\in F^\times$.
Taking determinant valuations relative to $L$ gives
$0=\nu(\det h)=3\nu(a)$, so $a$ is a unit and $hL=L$.  Hence the vertex
stabilizer is exactly $\operatorname{Aut}_{\mathcal O}(L)$, a conjugate
of the compact open group $\mathrm{GL}_3(\mathcal O)$.  A setwise face
stabilizer fixes each vertex because their types are distinct.  It is
therefore the finite intersection of these vertex stabilizers, hence is
compact open and fixes the face pointwise.

Every face lies in a chamber, and chamber transitivity sends it into the
standard chamber.  Its set of vertex types picks out a unique face there.
Type preservation also preserves the barycentric coordinates labeled by
those types, so no distinct points of the closed chamber are identified.
This proves the fundamental-domain assertion, including its faces.
\end{proof}

\begin{lemma}[The structure sheaf of the height-three Drinfeld model]
\label{h3:lem:drinfeld-structure-sheaf}
Let $\mathfrak H$ be the $p$-adic completion of the standard separated
strictly semistable weak formal model of $\Omega^2_{\Q_p}$ over $\Z_p$.
The canonical scalar map
\[
 \underline{\Z_p}\longrightarrow R\Gamma(\mathfrak H,\mathcal O)
\]
is an equivalence of condensed complexes, and is equivariant for the
action of $\GL_3(\Z_p)$.  Here the right side retains the derived
sections on $\mathfrak H\times\underline Q$ for every profinite set $Q$.
\end{lemma}

\begin{proof}
The imported model facts are given in
\cite[Sections~3.6,~3.8,~6.1--6.2,~6.4 and~7.1, and the proof of Proposition~6.5]{GKModel}.
After ordinary $p$-adic completion the local strictly semistable
equations and the special fiber are unchanged.  The model is separated:
each of the increasing weak-formal opens lies in a separated blowup
stage, and any two lie in one common such stage; the local closed-diagonal
criterion is preserved by completion.  The components and their
nonempty intersections are products of successive blowups of
projective spaces along rational linear subspaces.  They are indexed
by the vertices and simplices of the locally finite two-dimensional
Bruhat--Tits building $\mathcal B$ of \Cref{h3:lem:rank-three-building}.  Only the factors in dimensions
$0,1,2$ occur here.

We compute the structure-sheaf cohomology of exactly these factors.
For $\mathbf P^n_k$, $0\leq n\leq2$, over a field $k$, use the standard
affine cover.  On an intersection indexed by a nonempty set
$I\subset\{0,\ldots,n\}$, its degree-zero homogeneous Laurent
monomials have exponent vector $(a_0,\ldots,a_n)$ with
$\sum a_i=0$ and $a_i\geq0$ when $i\notin I$.  The Čech complex
decomposes by these monomials.  The constant monomial gives the
cochain complex of the full simplex and hence $k$ in degree zero.
For any other monomial its negative support
$N=\{i:a_i<0\}$ is nonempty and proper, and the monomial occurs
exactly for $I\supset N$.  Adding and deleting a fixed index outside
$N$, with the usual alternating sign, contracts this summand.
It follows that
\[
 R\Gamma(\mathbf P^n_k,\mathcal O)=k[0].
\]

The only nontrivial blowups needed for a surface factor are blowups
at the specified smooth rational points.  The calculation for one
such point is local on the surface and can be made in étale
coordinates $(u,v)$.  For the blowup of $\mathbf A^2_k$ at the
origin, the two standard charts and their overlap have rings
\[
 B_u=k[u,t],\qquad B_v=k[v,s],\qquad
 B_{uv}=k[u,t,t^{-1}],\qquad v=ut,\quad s=t^{-1}.
\]
In the common Laurent ring a monomial $u^a t^b$, with $a\geq0$ and
$b\in\Z$, belongs to $B_u$ when $b\geq0$ and to $B_v$ when $b\leq a$.
These two conditions cover all monomials.  Their intersection is
spanned by $0\leq b\leq a$, where
$u^a t^b=u^{a-b}v^b$.  Thus the affine Čech sequence
\[
 0\longrightarrow k[u,v]\longrightarrow B_u\oplus B_v
   \longrightarrow B_{uv}\longrightarrow0
\]
is exact.  The same sequence remains exact after localization and
flat étale base change, which are the local coordinate changes in
question.  The affine chart calculation therefore proves
$R\pi_*\mathcal O=\mathcal O$ for each required smooth-point blowup
$\pi$.  Blowing up an invertible ideal is an isomorphism.  This covers
the later strict transforms of rational lines on the surface and
the point centers on a projective line, all of which are Cartier
divisors.  Iterating the result, and using the finite affine-cover
double Čech complex for products over $k$, proves
\[
 R\Gamma(Y,\mathcal O_Y)=k[0]
\]
for every stratum $Y$ occurring in this height-three model.  This
argument makes no assertion about other rational varieties.
For a profinite parameter $Q$ the same calculations give
$C_{\hthreects}(Q,k)[0]$: continuous maps to each discrete section
module factor through a finite clopen partition of $Q$, and the
filtered union of the resulting finite powers preserves the exact
Čech sequences.

We next justify the resolution by intersections of components.
On a standard special-fiber chart its augmented form is the
alternating complex
\[
 k[t_0,\ldots,t_d]/(t_0\cdots t_r)
 \longrightarrow \bigoplus_i k[t_0,\ldots,t_d]/(t_i)
 \longrightarrow \bigoplus_{i<j} k[t_0,\ldots,t_d]/(t_i,t_j)
 \longrightarrow\cdots .
\]
For a nonzero monomial, let $Z$ be the nonempty set of indices
$i\leq r$ at which its exponent is zero.  Its coefficient complex
is the augmented cochain complex of the simplex on $Z$, so it is
exact.  This proves exactness monomial by monomial.  The calculation
is unchanged by the torus variables, localizations, and flat étale
base changes in the semistable charts.  Local finiteness of the
components now gives the global resolution by the products of their
intersection sheaves.  Taking derived global sections with a
profinite parameter and using the preceding stratum calculation
computes $R\Gamma(\mathfrak H_{\F_p}\times\underline Q,\mathcal O)$
by the three product terms below.  The arrow from
$C_{\hthreects}(Q,\F_p)$ is the separate canonical constants
augmentation:
\[
 C_{\hthreects}(Q,\F_p)\longrightarrow
 \prod_{\sigma\in\mathcal B_0}C_{\hthreects}(Q,\F_p)
 \longrightarrow
 \prod_{\sigma\in\mathcal B_1}C_{\hthreects}(Q,\F_p)
 \longrightarrow
 \prod_{\sigma\in\mathcal B_2}C_{\hthreects}(Q,\F_p),
\]
with the cellular incidence signs.  Products occur because the
strata are indexed by all cells of the building.

Here is a contraction valid for these product cochains.  By
\Cref{h3:lem:rank-three-building}, the augmented bounded cellular chain complex
$C_\bullet(\mathcal B;\Z)\to\Z$ is exact.  Its terms are free abelian,
and each cycle subgroup is free because it is a subgroup of a free
abelian group.  The short exact sequences of cycles therefore split.
Choosing splittings gives a chain contraction of the augmented
complex.  This contraction is not asserted to be equivariant.
Each basis cell maps to a finite cellular chain, since a cellular
chain group is a direct sum of its cells.  Dualizing into
$C_{\hthreects}(Q,\F_p)$ consequently gives a contraction of the
displayed product cochains: every output coordinate is a finite
linear combination of input coordinates and is continuous for the
product topology.  The contraction is natural in $Q$.

The canonical constants map into cellular cochains is equivariant
for the building action.  The contraction proves that its underlying
condensed map is an equivalence; the forgetful functor from equivariant
derived objects is conservative, so this same canonical map is an
equivariant equivalence.  In particular
\[
 \underline{\F_p}\xrightarrow{\ \simeq\ }
 R\Gamma(\mathfrak H_{\F_p},\mathcal O).
\]
Finally the formal structure sheaf is derived $p$-complete, and
derived global sections preserve its inverse limit of reductions.
The cone of the actual scalar map
$\underline{\Z_p}\to R\Gamma(\mathfrak H,\mathcal O)$ is therefore
derived $p$-complete.  Flatness of the semistable model identifies
its reduction modulo $p$ with the zero cone just computed.  Every
successive reduction modulo $p^\ell$ is then zero, and derived
$p$-completeness makes the cone itself zero.  This proves the
claimed equivariant equivalence.
\end{proof}

\subsection{The rational boundary calculation}

\begin{proof}[Proof of \Cref{h3:prop:rational-boundary}]
All cohomology complexes below retain their condensed parameters;
all inverse limits and group invariants are derived and taken
before inverting $p$.

\medskip\noindent\emph{1. Arithmetic bounds.} Let
$K$ be complete discretely valued of mixed characteristic $(0,p)$
with perfect residue field, and let $C$ be a completed algebraic
closure.  The scalar class obtained from the cyclotomic logarithm
gives a natural map
\[
 \mathcal O_K[\epsilon]\longrightarrow
 R\Gamma(G_K,\mathcal O_C),\qquad |\epsilon|=1,
 \quad \epsilon^2=0.
\]
Its cone is cohomologically nonnegative, and in each degree its
cohomology is killed by a power of $p$.  The required bounds can
be chosen uniformly after finite tame extensions of $K$.
For the finitely many nonzero twists $j=1,2$, the same assertion
of bounded torsion holds for $H^i(G_K,\mathcal O_C(-j))$, with
$H^0=0$.  Here is the part of the proof that gives uniformity,
\cite[Sections~4.2--4.4]{BSSW}.  Choose a finite cyclotomic stage $B/K$ far enough out that the
remaining $\Z_p$-tower is totally ramified and sufficiently ramified
in the sense of \cite[Definition~4.2.5 and Lemma~4.2.6]{BSSW}, generated
by $\sigma$, with $\chi(\sigma)$ a generator of $1+p^s\Z_p$,
$s\geq2$.  Write $M=\mathcal O_{\widehat K_\infty}$ and
$V=\ker(t:\widehat K_\infty\to B)$ for normalized trace.
The estimates in \cite[(4.2.9), (4.2.15), Lemmas~4.2.12 and
4.2.17, and the proof of Proposition~4.2.18]{BSSW} are
\[
 \|t\|\leq |p|^{-1/(p-1)},\qquad
 \|(\sigma-1)^{-1}\|_V
       \leq |p|^{-1-1/[p(p-1)]}.
\]
For existence of the inverse, at a finite cyclic level
$\sigma-1$ has kernel $B$ and image the trace-zero subspace.
Its inverse there is bounded by the second estimate, so the
inverses extend compatibly to the completion.  The estimate follows
by induction from normalized layer trace
$t_n=p^{-1}\operatorname{Tr}$ and
\[
 x-t_nx=p^{-1}\sum_{a=1}^{p-1}
       (1-\sigma^{a p^{n-1}})x,\qquad
 \|t_n\|\leq |p|^{-p^{-n}}.
\]
Set $M_0=\mathcal O_B\oplus(V\cap M)\subset M$.
The first norm bound gives $pM\subset M_0$, so $M/M_0$ is killed
by $p$.  For $j\in\{-2,-1,1,2\}$, put
$\lambda=\chi(\sigma)^{-j}$;
then $v_p(\lambda-1)=s$, since $p\geq5$.
The inequality
$|\lambda-1|\,\|(\sigma-1)^{-1}\|_V<1$ follows from
$s\geq2>1+1/[p(p-1)]$.  A Neumann series therefore inverts
$\sigma-\lambda$ on $V$ with the same bound
\cite[Lemma~4.4.1]{BSSW}, so its cokernel on $V\cap M$ is killed by $p^2$.
On $\mathcal O_B$ its cokernel is killed by $p^s$.
Thus $H^1(\langle\sigma\rangle,M_0(j))$ is killed by $p^s$,
while the invariants vanish.  The sequence $M_0\to M\to M/M_0$
gives a $p^{s+1}$ bound for $H^1(\langle\sigma\rangle,M(j))$.
The rational decomposition $B\oplus V$ also shows directly that
$M(j)$ has no invariants under the tail.
Restriction from $\Gal(K_\infty/K)$ to this open tail injects
$H^1$, because the tail has no invariants.  The cyclotomic fixed-field calculation and affinoid perfectoid
integral acyclicity compare this with $\mathcal O_C(j)$
\cite[Lemmas~4.3.2 and~4.3.5]{BSSW}.
Before cyclotomic invariants, the perfectoid-tail degree-zero
comparison has almost-zero cohomological error.  The maximal
ideal of the completed valuation ring is flat and idempotent,
so this error is annihilated by $p$ in the equivariant derived
category.  In the resulting triangle after invariants, the error
$Y$ has $H^0(Y)=0$ and $pH^i(Y)=0$ for $i>0$
\cite[Lemma~4.3.5 and Theorem~4.4.3]{BSSW}.
This gives the conservative bound $p^{s+2}$ in degree one;
the remaining degrees are bounded as in the same triangle.
The untwisted calculation retains the constant summand and its
cyclotomic-logarithm class, giving the displayed map
$\mathcal O_K[\epsilon]$ and its bounded-torsion cone.
Sufficient ramification persists under tame base change.  For a
finite tame extension $L/K$, let $e=e(LB/B)$ be the effective
tame ramification index after the chosen initial cyclotomic stage
$B$.  If $d_n$ is the different exponent of a $p$-layer above $B$,
then the base-changed exponent is
\[
 d'_n=e d_n-(e-1)(p-1).
\]
Consequently the same tail, $s$, and bounds work for all the
tame extensions in question.
Only existence of these tame-uniform bounds is used; the sharp
numerical constants in the preprint are unnecessary.

\medskip\noindent\emph{2. Semistable comparison.}
\par\noindent First apply this to a small affine semistable formal $\mathcal O_K$-scheme
$\mathfrak U=\operatorname{Spf}(R_0)$ of relative dimension $d\leq2$,
where small means admitting an \'etale morphism to a standard semistable chart.
Write $U$ for its generic fiber, $\mathfrak U_C$ for its completed base
change to $\mathcal O_C$, and $U_C$ for the geometric generic fiber. Use the divisorial log
structure and the valuation log structure on the base. The
semistable logarithmic comparison of
\cite[Sections~1.5--1.6 and Theorem~4.11]{CK}, extending the smooth
comparison of \cite[Theorem~8.3]{BMS}, applies on \'etale charts
\[
 \mathcal O_C\langle T_0,\ldots,T_r,
             T_{r+1}^{\pm1},\ldots,T_d^{\pm1}\rangle/
             (T_0\cdots T_r-a),\qquad 0\ne a\in\mathfrak m_C.
\]
They identify the cohomology of
$L\eta_{\zeta_p-1}R\nu_*\widehat{\mathcal O}^{+}$ with
$\Omega^j_{\mathfrak U_C,\log}\{-j\}$ for $0\leq j\leq d$.
No properness is required for these local statements.
The logarithmic differentials on $\mathfrak U$ are coherent and flat
over $\mathcal O_K$, and their higher cohomology on this affine
vanishes. Their completed base changes are the displayed
differentials on $\mathfrak U_C$.
The needed completed projection formula can be checked modulo
$p$: a flat module modulo $p$ is free over the artinian local
ring $\mathcal O_K/p$, and continuous cochains on a compact group
commute with these direct sums, since a map to a discrete direct
sum has finite image.  Derived $p$-completeness then proves the
projection formula before reduction.  For $j>0$, the normalized
map from the Breuil--Kisin twist to the Tate twist has cokernel
killed by $p$, so the preceding arithmetic bounds apply to each
of the finitely many positive graded pieces.

For the comparison with the original complex, put
$C_U=R\Gamma(U_{C,\mathrm{pro\acute et}},\widehat{\mathcal O}^{+})$,
$I=(\zeta_p-1)\subset\mathcal O_C$, and $A_U=L\eta_I C_U$.
The ideal $I$ is $G_K$-stable, although its displayed generator need not
be fixed.  On the semistable charts above, the perfectoid-cover edge map
of \cite[Section~3.3, Equation~(3.3.1)]{CK} has almost-zero cone.
Its source is the continuous cochain complex for
$\Delta\simeq\Z_p^d$, represented by a $d$-variable Koszul complex
\cite[Lemma~3.7]{CK}.  Thus $H^q(C_U)$ is almost zero for $q>d$.
Moreover, \cite[Theorems~3.9 and 4.11]{CK} identifies $A_U$ with a
complex concentrated in degrees $[0,d]$.  The degree-zero cohomology
of $C_U$ is $I$-torsion-free; indeed, it is the formal structure ring
\cite[Proposition~4.4]{CK}.

Set $B_U=\tau^{\leq d}C_U$.  Lemma~6.9 of \cite{BMS}, together
with compatibility of $L\eta_I$ with truncation, gives maps
\[
 a:A_U\longrightarrow B_U,
 \qquad b:I^{\otimes d}\otimes B_U\longrightarrow A_U
\]
whose composites are the natural ideal inclusions.  The invariant
element $p^d\in I^d$ gives an equivariant map
$\mathcal O_C\to I^{\otimes d}$; composing it with $b$ gives
$b_p:B_U\to A_U$ with $ab_p=p^d$ and $b_pa=p^d$.
Hence the kernel and cokernel of each $H^q(a)$ are killed by $p^d$,
and the long exact sequence gives
$p^{2d}H^q(\operatorname{cone}(a))=0$.
The triangle
\[
 \operatorname{cone}(a)\longrightarrow
 \operatorname{cone}(A_U\to C_U)\longrightarrow
 \tau^{>d}C_U
\]
then shows that $p^{2d+1}$ kills every cohomology group of the full
cone: the last term has almost-zero cohomology, which is killed by $p$.
The full cone is connective.  This argument concerns its cohomology
and does not assert that it carries a strict
$\Z/p^{2d+1}$-module structure.

After $G_K$-invariants, the first-quadrant spectral sequence has
terms $H^s(G_K,H^t(\operatorname{cone}(A_U\to C_U)))$.
In total degree $k$ its filtration has at most $k+1$ pieces, each
killed by $p^{2d+1}$.  Therefore $p^{(2d+1)(k+1)}$ kills the
cohomology in that degree.  These conservative bounds depend on
$d$ and $k$, and are independent of the chart and its tame extension.
The positive logarithmic graded pieces retain the arithmetic
bounds proved in Step~1. Tame descent preserves these bounds, by
unramified semilinear descent and averaging over tame inertia.
We obtain the natural comparison
\[
 \mathcal O(\mathfrak U)[\epsilon]\longrightarrow
 R\Gamma(U_{\mathrm{pro\acute et}},\widehat{\mathcal O}^{+})
\]
with nonnegative cone and degreewise bounds independent of the
chosen charts and their tame extensions. For a general semistable
affine, choose an \'etale hypercover refining the \v Cech nerve of
a small-chart cover, with each term a finite disjoint union of
small affine charts. Such finite covers of the matching objects
exist because the base is quasi-compact and quasi-separated.
Apply \'etale hyperdescent to the actual natural comparison map.
Its levelwise cones are connective and have the uniform degreewise
bounds just proved. In a fixed total degree the descent spectral
sequence has only finitely many cosimplicial and cohomological
positions, so the same bounded-torsion conclusion follows. Gluing next
by a separated affine formal cover likewise preserves a bound in
every total degree, since the relevant \v Cech diagonal is finite.  This proves the part of
\cite[Sections~5.3--5.6]{BSSW} needed here.

\medskip\noindent\emph{3. The two global calculations.} The Lubin--Tate space
$\mathrm{LT}_{\breve K}$ is the open two-dimensional ball.
Exhaust it by closed balls of radii $|p|^{1/\ell}$ for primes
$\ell\ne p$; each has a smooth model after a tame extension.
The two-index system of integral functions modulo $p^j$, indexed
by radius and $j$, is Mittag--Leffler.  This coefficient limit is
taken over the fixed base $\breve K$, whose uniformizer is $p$;
the tame extensions were used for the local comparisons.  More
generally, with uniformizer $\varpi$, for a fixed inner radius $r$
and reduction exponent $j$ take
$M=\lceil(j+1)/\log_{|\varpi|}(r)\rceil$.  Restriction from any
$s\geq r^{1/(j+1)}$, with source reduced modulo any
$\varpi^{j'}$ for $j'\geq j$, kills modulo $\varpi^j$ all terms of
total degree at least $M$, and the surviving coefficients are
integral.  Conversely every monomial of degree less than $M$ is
already integral on the larger ball.  The image is therefore the
fixed module spanned by those monomials.  Continuous functions from a
profinite parameter preserve this image statement by finite clopen
partitions.  A cofinal diagonal reduces the two indices to a
countable Mittag--Leffler tower.  Thus the derived two-index limit is
$\overline A$ in degree zero, as in
\cite[Lemmas~6.1.2--6.1.3]{BSSW}.  The preceding local comparisons
and the countable inverse-limit spectral sequence give
\[
 \overline A[\epsilon]\longrightarrow
 R\Gamma(\mathrm{LT}_{\breve K,\mathrm{pro\acute et}},
          \widehat{\mathcal O}^{+})
\]
with nonnegative, degreewise bounded-torsion cone: in any degree
the limit spectral sequence uses only a limit in that degree and
a $\lim^1$ from the preceding degree, both with bounds uniform
over the radii.  No uniform bound over all cohomological degrees is needed.

The other space is $\Omega^2_{\Q_p}$, with the formal model
$\mathfrak H$ of \Cref{h3:lem:drinfeld-structure-sheaf}.  That lemma
gives the actual equivariant scalar equivalence
$\underline{\Z_p}\simeq R\Gamma(\mathfrak H,\mathcal O)$, including
all profinite parameters.  It supplies the structure-sheaf assertion
needed from \cite[Theorem~6.2.1]{BSSW}.
The semistable comparison now gives
\[
 \Z_p[\epsilon]\longrightarrow
 R\Gamma(\Omega^2_{\Q_p,\mathrm{pro\acute et}},
          \widehat{\mathcal O}^{+})
\]
with the same degreewise bounded-torsion conclusion.

\medskip\noindent\emph{4. The tower comparison.} The bridge between these calculations is an equivalence of
\emph{families}, not merely a bijection of geometric points.
The relative Rapoport--Zink/shtuka comparison
\cite[Theorem~24.2.5 and its proof, pp.~227--229]{ScholzeWeinsteinBerkeley},
using the family comparison of its Theorem~22.3.1, identifies the
universal integral
Tate local system with that in a modification
\[
 0\longrightarrow \mathcal T\otimes_{\Z_p}\mathcal O
   \longrightarrow\mathcal E\longrightarrow i_*\mathcal L
   \longrightarrow0,
\]
where $\mathcal T=T_p(\mathcal X)$ is the actual integral Tate local
system of the universal $p$-divisible group $\mathcal X$, of rank
three, and $\mathcal L$ is a line at the
untilt divisor, and $\mathcal E$ is fiberwise $V_3$.
The relative basic-stratum frame torsors
\cite[Theorems~III.2.4 and III.4.5]{FarguesScholze} identify the two quotient
presentations of this moduli problem.  Honda endomorphisms are
already defined over $\F_{p^3}$; their canonical action and the
basic-stratum equivalence descend from the geometric residue
field as in \Cref{h3:geom:curve-inputs}.  After the height-component
normalization on the Lubin--Tate side, this gives
\[
 [\mathrm{LT}_{\breve K}^{\diamond}/\overline G]
 \simeq[\Omega^{2,\diamond}_{\Q_p}/\GL_3(\Z_p)].
\]
Pulling back the two atlases gives a common space with commuting
pro-\'etale torsor actions.  Their \v Cech nerves therefore compute
the same integral condensed cohomology object; this is the
actual descent argument behind \cite[Theorem~3.9.1 and (3.9.3)]{BSSW}.
Taking invariants of the comparison cones remains harmless:
they are nonnegative, so each diagonal of the group-cohomology
spectral sequence contains finitely many bounded-torsion rows.
After inverting $p$ we obtain
\[
 H^*(\overline G,\overline A)[1/p]\otimes\Q_p[\epsilon]
 \simeq H^*(\GL_3(\Z_p),\Q_p)\otimes\Q_p[\epsilon].
\]
The extended residue-field action is retained here.
Witt Artin--Schreier gives the continuous unramified calculation
over $\overline W$ in \cite[Lemma~3.8.5(1)]{BSSW};
its consequence for the full extended stabilizer is
\cite[Lemma~3.8.5(2)]{BSSW}.  Here $P_3$ fixes $\overline W$ and the
residue-field group acts by Witt Frobenius, exactly as in the chosen
Honda semilinear action.  The coefficient conventions also agree:
for every profinite $Q$, cochains on $Q\times\overline G^a$ with
values in the discrete $W_\ell(\overline{\F}_p)$ have finite image.
Lifting those finitely many values makes cochain reduction surjective.  Their derived inverse limit is therefore
the continuous $\overline W$ cochain complex.  Every continuous
cochain on this compact domain with values in $\overline W[1/p]$
has bounded image, so inverting $p$ gives the rational cochains
used in part~(2).

To pass to the prescribed finite field $k$, put
$\Delta_k=\Gal(\overline{\F}_p/k)$.  Its action fixes $u_1,u_2$
and acts by coefficient Frobenius.  Each quotient
$\overline A/(p,u_1,u_2)^a$ is a finite direct sum of truncated
Witt coefficient modules.  Witt Artin--Schreier therefore identifies
its $\Delta_k$-invariants with
$W(k)[[u_1,u_2]]/(p,u_1,u_2)^a$ and makes its positive
$\Delta_k$-cohomology vanish.  The coefficient reductions, including
cochains with profinite parameters, are surjective by finite-value
lifting.  Taking the derived inverse limit gives
\[
 R\Gamma(\Delta_k,\overline A)\simeq W(k)[[u_1,u_2]],
\]
equivariantly for $P_3$, since $k$ contains the Honda endomorphism
field.  Hochschild--Serre now identifies the algebraically closed
residue-field calculation with the finite-field one.  The trace-one
contraction of \Cref{h3:lem:framework-semilinear} supplies the finite
semilinear descent used by \Cref{h3:prop:exact-descent}.

\medskip\noindent\emph{5. Constants and rational homotopy.} The additive splitting in
\cite[Proposition~2.5.1]{BSSW} has a direct construction.
Compose each power operation with the $K(3)$-local transfer and
write the resulting operations as $\beta_m$.  Transfer transitivity
and the power-operation addition formula give
$\beta_m(a+b)=\sum_{i+j=m}\beta_i(a)\beta_j(b)$.
Thus $a\mapsto\sum_m\beta_m(a)t^m$, followed by reduction to
$\overline\F_p$ and projection from big to $p$-typical Witt vectors,
is an equivariant additive map
$\gamma:\overline A\to\overline W$.
For continuity, write the localized unit as the filtered colimit of its
compact Moore stages $M_k$. The spectrum
$L_{K(3)}\Sigma^\infty_+B\Sigma_m$ is dualizable for the finite
group $\Sigma_m$ \cite[Corollary~8.7]{HS99}; its tensor product with
the compact object $M_k$ is therefore compact. Compactness makes the map from
$M_k\otimes L_{K(3)}\Sigma^\infty_+ B\Sigma_m$ factor through some
finite tensor stage $(M_j^{\otimes m})_{h\Sigma_m}$; hence each
power operation is continuous for the Moore topology
\cite[Lemma~2.5.4]{BSSW}.
Transfers are maps of $K(3)$-local spectra and are continuous
for the same topology \cite[Proposition~11.1(b)]{HS99}.
For $\overline A=\pi_0\overline E$, where $\overline E$ is the
Morava theory used here, the Moore topology is the
$\mathfrak m=(p,u_1,u_2)$-adic topology by the following local check.
Choose the unit-compatible generalized Moore spectra $\mathbb S/J_j$
of \cite[Proposition~4.22]{HS99}.  Their $K(3)$-local duals
$M_j=D_{K(3)}L_{K(3)}(\mathbb S/J_j)$ are compact stages for the unit
by \cite[Corollary~7.11 and Theorem~8.5]{HS99}.  Duality carries restriction
along the dual unit exactly to the map induced by
$\mathbb S\to\mathbb S/J_j$.  The latter spectrum is finite, so
$\overline E\otimes(\mathbb S/J_j)$ remains $K(3)$-local.
The unit-normalized Landweber formula
\cite[Introduction and Example~0.1(d)]{HoveyStrickland} and the Moore
coefficient calculation \cite[Definition~4.12]{HS99} identify the restriction as
\[
 \overline A\longrightarrow[M_j,\overline E]
   \cong\pi_0(\overline E\otimes(\mathbb S/J_j))
   \cong\overline A/\overline J_j.
\]
Here $\overline J_j$ is the image of the Moore ideal after
removing invertible periodicity factors.  The normalized images of
$p,v_1,v_2$ form a regular sequence of length three in the
three-dimensional regular local ring $\overline A$.  Their ideal
is therefore $\mathfrak m$-primary.  The pure-power ideals
$\overline J_j$, whose exponents tend to infinity, are cofinal with
the powers of $\mathfrak m$,
which proves the claimed equality of topologies from the restriction
kernels.  Thus each composite $\beta_m:\overline A\to\overline A$
is $\mathfrak m$-adically continuous.  The coefficientwise
Witt-vector construction is therefore continuous.
Its restriction $f$ to $\overline W$ is the identity modulo $p$;
continuity and additivity make it $\Z_p$-linear.  Write
$f=\hthreeid+p h$ with a continuous $\Z_p$-linear endomorphism $h$.
The convergent series $\sum_{n\geq0}(-p h)^n$ is its inverse.
Hence $f^{-1}\gamma$ is an equivariant retraction and
$\overline A=\overline W\oplus\overline A^c$ additively.

The compact comparison used to compute the constants
has the following integral hypotheses.  Choose a deep normal uniform
open in $P_3$ and in $\GL_3(\Z_p)$.  At $p\geq5$ its lower-$p$-series
valuation is saturated and equi-$p$-valued with denominator $e=1$
and generators in degree one.  The trivial finite free module $\Z_p$
satisfies the hypotheses of \cite[Theorem~3.3.3]{HuberKingsNaumann}; its continuous
comparison is cup compatible.  The coefficient and restriction
naturality in \cite[Theorem~3.1.1(1)--(3)]{HuberKingsNaumann}
identifies its rationalization with Lazard's comparison
\cite[V.2.4.9--V.2.4.10]{Lazard} and makes
it equivariant for conjugation.  Rational Hochschild--Serre then
takes invariants under the finite quotient, whose higher cohomology
vanishes by averaging.  The full exterior calculation for the compact
matrix and stabilizer groups, together with the Lie-algebra
calculation and trivial adjoint action, is given by
\cite[Proposition~3.8.1 and Lemma~3.8.2]{BSSW}.
At height three its generator degrees are $1,3,5$.
The unramified and finite semilinear descent from Step~4, with
\cite[Lemma~3.8.5(2)]{BSSW} for the full extended stabilizer, therefore identifies both
\[
 H^*(\overline G,\overline W)[1/p]
 \quad\text{and}\quad H^*(\GL_3(\Z_p),\Q_p)
\]
with that exterior algebra over $\Q_p$.

The displayed tower comparison makes the cohomology of
$\overline A$ finite-dimensional in every degree.  Its additive
decomposition into constants and $\overline A^c$ is a decomposition
of continuous coefficient complexes.  If
$d_i=\dim_{\Q_p}H^i(\overline G,\overline A^c)[1/p]$, comparison
with the identical constant algebra, including the common factor
$\Q_p[\epsilon]$, gives $d_i+d_{i-1}=0$ (with $d_{-1}=0$).
Thus every $d_i$ is zero.
The constants inclusion is a ring map, so its resulting
cohomology isomorphism is multiplicative.  This proves
\eqref{h3:eq:rational-boundary-constants}.

Apply the local relative Morava descent of
\Cref{h3:prop:exact-descent} to $T$.  The central procyclic subgroup
generated by $1+p$ acts on internal degree $2t$ through
$(1+p)^{-t}$.  Its continuous two-term resolution has differential
$(1+p)^{-t}-1$ on that coefficient module.  For $t\ne0$ this scalar
is nonzero in $\Q_p$, so the rationalized two-term complex is
acyclic; Hochschild--Serre makes every nonzero internal row vanish.
The locally proved relative spectral sequence is strongly convergent
with finite cohomological amplitude.  Exact rationalization preserves
its bounded convergence.  Only the internal-degree-zero row remains, with
one filtration degree in each total degree.  Multiplicativity
therefore yields
\[
 \pi_*L_0T=\Lambda_{\Q_p}(\alpha_1,\alpha_3,\alpha_5),
 \qquad |\alpha_i|=-i.
\]
This establishes the rational boundary used in
\Cref{h3:thm:rational-compatibility}; the image of its chosen
primitive is determined by the separate map calculation in
\Cref{h3:sec:rational}.
\end{proof}

\bibliographystyle{plain}
\bibliography{references}
\end{document}
